% !TeX program = xelatex
% !TeX encoding = UTF-8
% 《D-模代数理论讲义》中文译本 —— Dragan Miličić 著
% 主文件：编译命令 xelatex -shell-escape main.tex （需编译 2~3 次以生成目录与交叉引用）
\documentclass[UTF8, a4paper, 12pt, openany]{ctexbook}

% ---------- 宏包 ----------
\usepackage{amsmath, amssymb, amsthm, mathtools}
\usepackage{fontspec}
\usepackage{geometry}
\usepackage{fancyhdr}
\usepackage{booktabs}
\usepackage{array}
\usepackage{tikz}
\usepackage{tikz-cd}
\usepackage{enumitem}
\usepackage{hyperref}
\usepackage{xcolor}

% ---------- 字体 ----------
% 本机只装了 Noto Serif SC / Noto Sans SC 的可变字体，没有独立的 Bold/Italic 字面，
% 故用 AutoFakeBold / AutoFakeSlant 合成粗斜体（否则 fontspec 报 Font shape ... undefined）。
\setCJKmainfont{Noto Serif SC}[AutoFakeBold=1.5, AutoFakeSlant=0.2]
\setCJKsansfont{Noto Sans SC}[AutoFakeBold=1.5]
\setCJKmonofont{Noto Sans SC}
\setmonofont{Consolas}[Scale=0.9]

% ---------- 页面版式 ----------
\geometry{
  left=2.8cm, right=2.6cm, top=2.8cm, bottom=2.6cm,
  headsep=0.6cm, footskip=0.9cm
}
\pagestyle{fancy}
\fancyhf{}
\renewcommand{\headrulewidth}{0.4pt}
\fancyhead[LE]{\nouppercase{\leftmark}}
\fancyhead[RO]{\nouppercase{\rightmark}}
\fancyfoot[C]{\thepage}
\fancypagestyle{plain}{%
  \fancyhf{}\renewcommand{\headrulewidth}{0pt}\fancyfoot[C]{\thepage}}

% ---------- 断行与容差优化（减少长行内公式引起的 overfull hbox） ----------
\tolerance=1500
\emergencystretch=3em
\hfuzz=2pt

% ---------- 定理类环境（按节编号：节.序号，与原书一致） ----------
\theoremstyle{plain}
\newtheorem{theorem}{定理}[section]
\newtheorem{lemma}[theorem]{引理}
\newtheorem{proposition}[theorem]{命题}
\newtheorem{corollary}[theorem]{推论}
\theoremstyle{definition}
\newtheorem{definition}[theorem]{定义}
\newtheorem{example}[theorem]{例}
\theoremstyle{remark}
\newtheorem{remark}[theorem]{注}

% 编号显示为「节.序号」（原书体系：第 I 章 §10 第 2 条显示为 10.2），
% 而 \thesection 仍保留「章.节」形式用于节标题。
\renewcommand{\thetheorem}{\arabic{section}.\arabic{theorem}}

% ---------- 校注（批注）环境：正文中插入的解释性小字块 ----------
\newenvironment{annote}{%
  \par\vspace{0.4em}\begingroup\small
  \noindent\textcolor{black!50}{\rule[-0.25em]{2.2pt}{1.1em}}\hspace{0.5em}%
  \textsf{\textbf{校注}}\hspace{0.4em}%
}{\par\endgroup\vspace{0.4em}}

% ---------- 常用记号 ----------
\newcommand{\Diff}{\operatorname{Diff}}
\newcommand{\Der}{\operatorname{Der}}
\newcommand{\Symb}{\operatorname{Symb}}
\newcommand{\boxtimesext}{\mathbin{\boxtimes}}
\newcommand{\ord}{\operatorname{ord}}
\newcommand{\Gr}{\operatorname{Gr}}
\newcommand{\Hom}{\operatorname{Hom}}
\newcommand{\End}{\operatorname{End}}
\newcommand{\Ext}{\operatorname{Ext}}
\newcommand{\Tor}{\operatorname{Tor}}
\newcommand{\im}{\operatorname{im}}
\newcommand{\Id}{\operatorname{Id}}
\newcommand{\Spec}{\operatorname{Spec}}
\newcommand{\Supp}{\operatorname{Supp}}
\newcommand{\rank}{\operatorname{rank}}
\newcommand{\ch}{\operatorname{ch}}
\newcommand{\N}{\mathbb{N}}
\newcommand{\Z}{\mathbb{Z}}
\newcommand{\Q}{\mathbb{Q}}
\newcommand{\R}{\mathbb{R}}
\newcommand{\C}{\mathbb{C}}
\newcommand{\D}{\mathcal{D}}
\newcommand{\calO}{\mathcal{O}}

% ---------- 超链接 ----------
\hypersetup{
  colorlinks=true, linkcolor=blue!70!black, citecolor=green!50!black,
  urlcolor=blue!70!black, bookmarks=true, bookmarksnumbered=true,
  pdftitle={D-模代数理论讲义（中文译本）}, pdfauthor={Dragan Miličić}
}

% ---------- 目录深度 ----------
\setcounter{tocdepth}{2}

\begin{document}

\frontmatter
\begin{titlepage}
\centering
\vspace*{1.4cm}

{\Huge\bfseries D-模代数理论讲义}\\[1.1cm]

{\Large Lectures on Algebraic Theory of $D$-Modules}\\[1.6cm]

{\large Dragan Miličić\ 著}\\[0.5cm]

{\large 中文校订批注本}\\[2cm]

{\normalsize 犹他大学数学系}\\[0.35cm]

{\normalsize 1999 年修订版}\\[1.6cm]

{\small 译者注：本译本依据作者在互联网上公开发行的讲义原文翻译而成，
力求忠实于原文的叙述与全部数学内容。}\\[1.2cm]

\begin{minipage}{0.86\textwidth}
\small
\textbf{校订说明}　本版为 2026 年 11 月的校订批注本。校订以英文原文逐页对照，逐条核对全部编号条目的陈述、证明与公式，统一术语（如 holonomic 一律作“完整”、direct image 作“直像”），并把编号恢复为原书的“节.序号”体系；初版第 I 章第 10 节 10.5--10.10 整段漏译，本版据原文补译。书中的灰色小字段为校注，用于补充讲义未展开的背景，每条末尾注明依据（Stacks Project 标签、nLab 条目及其修订日期，或原书页码）。
\end{minipage}

\vfill
\end{titlepage}

\tableofcontents
\chapter*{前言}
\addcontentsline{toc}{chapter}{前言}

本讲义是对 $D$-模代数理论的一个简要介绍。最初的版本写于 1986 年，当时作者在讲授一门为期一年的相关课程。后来在讲授类似课程时又做过一些小的修订。1999 年进行了大规模的重新组织。

作者感谢 Dan Barbasch 与 Pavle Pandžić 指出手稿以前版本中的若干错误。


\mainmatter
\chapter{多项式系数微分算子环上的模}

\section{Hilbert 多项式}

设 $A=\bigoplus_{n\in\mathbb{Z}} A^n$ 是一个分次诺特交换环，带有单位元 $1$，且 $1$ 含于 $A^0$ 中。于是 $A^0$ 是带单位元 $1$ 的交换环。假设当 $n<0$ 时 $A^n=0$。

\begin{lemma}\ 
\begin{enumerate}[label=(\roman*)]
  \item $A^0$ 是诺特环；
  \item $A$ 是有限生成的 $A^0$-代数。
\end{enumerate}
\end{lemma}

\begin{proof}
(i) 令 $A_+=\bigoplus_{n=1}^{\infty}A^n$。则 $A_+$ 是 $A$ 的理想且 $A^0=A/A_+$。

(ii) $A_+$ 是有限生成的。设 $x_1,x_2,\dots,x_s$ 是 $A_+$ 的一组齐次生成元，记 $d_i=\deg x_i$，$1\le i\le s$。令 $B$ 为 $x_1,x_2,\dots,x_s$ 生成的 $A^0$-子代数。我们断言对所有 $n\in\mathbb{Z}_+$ 有 $A^n\subseteq B$。显然 $A^0\subseteq B$。设 $n>0$ 且 $y\in A^n$。则 $y\in A_+$，故 $y=\sum_{i=1}^s y_i x_i$，其中 $y_i\in A^{n-d_i}$。由归纳假设可用于 $y_i$，$1\le i\le s$。这说明 $y\in B$。
\end{proof}

1.1 的逆命题由 Hilbert 定理得出，该定理指出：若环 $A^0$ 是诺特的，则多项式环 $A^0[X_1,X_2,\dots,X_n]$ 是诺特的。

设 $M=\bigoplus_{n\in\mathbb{Z}} M^n$ 是一个有限生成的分次 $A$-模。则每个 $M^n$（$n\in\mathbb{Z}$）都是 $A^0$-模。并且当 $n$ 充分小时 $M^n=0$。

\begin{lemma}
$A^0$-模 $M^n$（$n\in\mathbb{Z}$）是有限生成的。
\end{lemma}

\begin{proof}
设 $m_i$（$1\le i\le k$）是 $M$ 的齐次生成元，$\deg m_i=r_i$。对 $j\in\mathbb{Z}_+$，记 $z_i(j)$，$1\le i\le \ell(j)$，为 $x_1,x_2,\dots,x_s$ 中次数为 $j$ 的所有齐次单项式。设 $m\in M^n$。则 $m=\sum_{i=1}^k y_i m_i$，其中 $y_i\in A^{n-r_i}$。由 1.1，$y_i=\sum_j a_{ij} z_j(n-r_i)$，其中 $a_{ij}\in A^0$。这意味着 $m=\sum_{i,j}a_{ij}z_j(n-r_i)m_i$，故 $M^n$ 由 $(z_j(n-r_i)m_i;\ 1\le j\le \ell(n-r_i),\ 1\le i\le k)$ 生成。
\end{proof}

令 $\mathcal{M}_{fg}(A^0)$ 为有限生成 $A^0$-模的范畴。设 $\lambda$ 是 $\mathcal{M}_{fg}(A^0)$ 上取值于 $\mathbb{Z}$ 的函数。若对任意短正合列
$$0\longrightarrow M'\longrightarrow M\longrightarrow M''\longrightarrow 0$$
都有
$$\lambda(M)=\lambda(M')+\lambda(M''),$$
则称 $\lambda$ 是\emph{可加的}。显然，可加性蕴含 $\lambda(0)=0$。

\begin{lemma}
设
$$0\to M_0\to M_1\to M_2\to\cdots\to M_n\to 0$$
是 $\mathcal{M}_{fg}(A^0)$ 中的正合列。则
$$\sum_{i=0}^n (-1)^i\lambda(M_i)=0.$$
\end{lemma}

\begin{proof}
显然。
\end{proof}

设 $\mathbb{Z}[[t]]$ 为系数在 $\mathbb{Z}$ 中的形式幂级数环。记 $\mathbb{Z}((t))$ 为 $\mathbb{Z}[[t]]$ 关于乘法系 $\{t^n\mid n\in\mathbb{Z}_+\}$ 的局部化。

设 $M$ 是有限生成分次 $A$-模。则 $M$ 关于 $\lambda$ 的 Poincaré 级数为
$$P(M,t)=\sum_{n\in\mathbb{Z}}\lambda(M^n)\,t^n\ \in \mathbb{Z}((t)).$$

例如，设 $A=k[X_1,X_2,\dots,X_s]$ 是 $s$ 个变量、系数在域 $k$ 中、按全次数分次的多项式代数。则 $A^0=k$，且对每个有限生成分次 $A$-模 $M$ 有 $\dim_k M^n<\infty$。因此可对 $\lambda=\dim_k$ 定义 Poincaré 级数。特别地，对 $A$ 本身有
$$P(A,t)=\sum_{n\in\mathbb{Z}}\dim_k A^n t^n=\sum_{n=0}^{\infty}\binom{s+n-1}{s-1}t^n=\frac{1}{(1-t)^s}.$$

下一个结果表明一般情形的 Poincaré 级数具有类似形式。

\begin{theorem}[Hilbert, Serre]
对任意有限生成分次 $A$-模 $M$ 有
$$P(M,t)=\frac{f(t)}{\prod_{i=1}^s (1-t^{d_i})},$$
其中 $f(t)\in\mathbb{Z}[t,t^{-1}]$。
\end{theorem}

\begin{proof}
对 $s$ 归纳证明。若 $s=0$，则 $A=A^0$，$M$ 是有限生成 $A^0$-模。这意味着当 $n$ 充分大时 $M^n=0$。因此除有限多个 $n\in\mathbb{Z}$ 外 $\lambda(M^n)=0$，且 $P(M,t)\in\mathbb{Z}[t,t^{-1}]$。

现设 $s>0$。乘以 $x_s$ 定义了 $M$ 的一个 $A$-模自同态 $f$。令 $K=\ker f$，$I=\operatorname{im} f$，$L=M/I$。则 $K,I,L$ 是分次 $A$-模，且有正合列
$$0\longrightarrow K\longrightarrow M\xrightarrow{\ f\ } M\longrightarrow L\longrightarrow 0.$$
这意味着
$$0\longrightarrow K^n\longrightarrow M^n\xrightarrow{\ x_s\ } M^{n+d_s}\longrightarrow L^{n+d_s}\longrightarrow 0$$
对所有 $n\in\mathbb{Z}$ 都是 $A^0$-模的正合列。特别地，由 1.3，
$$\lambda(K^n)-\lambda(M^n)+\lambda(M^{n+d_s})-\lambda(L^{n+d_s})=0$$
对所有 $n\in\mathbb{Z}$ 成立。这意味着
\begin{align*}
(1-t^{d_s})P(M,t)&=\sum_{n\in\mathbb{Z}}\lambda(M^n)t^n-\sum_{n\in\mathbb{Z}}\lambda(M^n)t^{n+d_s}\\
&=\sum_{n\in\mathbb{Z}}\bigl(\lambda(M^{n+d_s})-\lambda(M^n)\bigr)t^{n+d_s}\\
&=\sum_{n\in\mathbb{Z}}\bigl(\lambda(L^{n+d_s})-\lambda(K^n)\bigr)t^{n+d_s}\\
&=P(L,t)-P(K,t)\,t^{d_s},
\end{align*}
即
$$(1-t^{d_s})P(M,t)=P(L,t)-t^{d_s}P(K,t).$$
由构造可知 $x_s$ 在 $L$ 和 $K$ 上的作用都是乘以 $0$，即可以把它们看作 $A/(x_s)$-模。因此归纳假设适用，这就立即得到结论。
\end{proof}

由于 Poincaré 级数 $P(M,t)$ 是有理函数，可以讨论它在一点的极点阶数。记 $d_\lambda(M)$ 为 $P(M,t)$ 在 $1$ 处极点的阶数。

由定理，$f(t)=\sum_{k\in\mathbb{Z}}a_k t^k$，其中 $a_k\in\mathbb{Z}$ 且除有限多个 $k$ 外 $a_k=0$。设 $p$ 为 $f$ 在 $1$ 处零点的阶数。假设 $p>0$。则 $f(t)=(1-t)g(t)$，其中 $g(t)=\sum_{k\in\mathbb{Z}}b_k t^k$，$b_k\in\mathbb{Q}$ 且除有限多个 $k$ 外 $b_k=0$。此外对所有 $k\in\mathbb{Z}$ 有 $a_k=b_k-b_{k-1}$。对 $k$ 归纳可得 $b_k\in\mathbb{Z}$。必要时重复此过程，得到 $f(t)=(1-t)^p g(t)$，其中 $g(t)=\sum_{k\in\mathbb{Z}}b_k t^k$，$b_k\in\mathbb{Z}$ 且除有限多个 $k$ 外 $b_k=0$，并且 $g(1)\neq 0$。

\begin{corollary}
若对 $1\le i\le s$ 有 $d_i=1$，则函数 $n\longmapsto \lambda(M^n)$ 对充分大的 $n\in\mathbb{Z}$ 等于一个有理系数多项式，次数为 $d_\lambda(M)-1$。
\end{corollary}

\begin{proof}
设 $p$ 为 $f$ 在 $1$ 处零点的阶数。则可写 $f(t)=(1-t)^p g(t)$，其中 $g(1)\neq 0$。此外令 $d=d_\lambda(M)=s-p$，于是
$$P(M,t)=\frac{g(t)}{(1-t)^d}.$$
另一方面
$$(1-t)^{-d}=\sum_{k=0}^{\infty}\frac{d(d+1)\cdots(d+k-1)}{k!}t^k=\sum_{k=0}^{\infty}\binom{d+k-1}{d-1}t^k,$$
若写 $g(t)=\sum_{k=-N}^{N}a_k t^k$，则对所有 $n\ge N$，
$$\lambda(M^n)=\sum_{k=-N}^{N}a_k\binom{d+n-k-1}{d-1}.$$
这等于
$$\sum_{k=-N}^{N}a_k\frac{(d+n-k-1)!}{(d-1)!(n-k)!}
=\sum_{k=-N}^{N}a_k\frac{(n-k+1)(n-k+2)\cdots(n-k+d-1)}{(d-1)!},$$
故 $\lambda(M^n)$ 是 $n$ 的多项式，首项为
$$\Bigl(\sum_{k=-N}^{N}a_k\Bigr)\frac{n^{d-1}}{(d-1)!}=g(1)\frac{n^{d-1}}{(d-1)!}\neq 0.$$
\end{proof}

我们把对充分大的 $n\in\mathbb{Z}$ 给出 $\lambda(M^n)$ 的多项式称为 $M$ 关于 $\lambda$ 的 \emph{Hilbert 多项式}。由证明可见 $M$ 的 Hilbert 多项式的首项系数为 $\dfrac{g(1)}{(d-1)!}$。

回到 $A=k[X_1,X_2,\dots,X_s]$ 的例子，我们看到
$$\dim_k A^n=\binom{n+s-1}{s-1}=\frac{n^{s-1}}{(s-1)!}+\cdots,$$
故对 $A=k[X_1,X_2,\dots,X_s]$，Hilbert 多项式的次数等于 $s-1$。

现在证明关于（系数在特征 $0$ 的域中、对充分大正整数取整数值的）多项式的一个刻画。首先注意，对任意 $s\in\mathbb{Z}_+$ 和 $q\ge s$ 有
$$q^s=s!\binom{q}{s}+Q(q),$$
其中 $Q$ 是次数为 $s-1$ 的多项式。因此任意 $d$ 次多项式 $P$ 对充分大的 $q$ 可唯一写成
$$P(q)=c_0\binom{q}{d}+c_1\binom{q}{d-1}+\cdots+c_{d-1}\binom{q}{1}+c_d,$$
其中系数 $c_i$（$0\le i\le d$）适当选取。由于二项式系数是整数，若所有 $c_i$（$0\le i\le d$）是整数，则对整数 $n\ge d$，$P$ 取整数值。下一个结果是这一观察的逆命题。

\begin{lemma}
若多项式
$$q\longmapsto P(q)=c_0\binom{q}{d}+c_1\binom{q}{d-1}+\cdots+c_{d-1}\binom{q}{1}+c_d$$
对充分大的 $n\in\mathbb{Z}$ 取整数值 $P(n)$，则所有系数 $c_i$（$0\le i\le d$）都是整数。
\end{lemma}

\begin{proof}
对 $d$ 归纳证明。若 $d=0$ 结论显然。此外
\begin{align*}
P(q+1)-P(q)&=\sum_{i=0}^d c_i\binom{q+1}{d-i}-\sum_{i=0}^d c_i\binom{q}{d-i}\\
&=\sum_{i=0}^d c_i\left(\binom{q+1}{d-i}-\binom{q}{d-i}\right)=\sum_{i=0}^{d-1}c_i\binom{q}{d-i-1},
\end{align*}
这里用到恒等式（对 $q\ge s\ge 1$）
$$\binom{q+1}{s}=\binom{q}{s}+\binom{q}{s-1}.$$
因此 $q\longmapsto P(q+1)-P(q)$ 是系数为 $c_0,c_1,\dots,c_{d-1}$ 的多项式，且对充分大的 $n\in\mathbb{Z}$ 有 $P(n)\in\mathbb{Z}$。由归纳假设，所有 $c_i$（$0\le i\le d-1$）是整数。这立即蕴含 $c_d$ 也是整数。
\end{proof}

还需要另一个相关的注记。若 $F$ 是首项系数为 $a_0$ 的 $d$ 次多项式，则
\begin{align*}
G(n)&=F(n)-F(n-1)\\
&=(a_0 n^d+a_1 n^{d-1}+\cdots)-(a_0(n-1)^d+a_1(n-1)^{d-1}+\cdots)\\
&=a_0 d\, n^{d-1}+\cdots
\end{align*}
是 $n$ 的 $d-1$ 次多项式，首项系数为 $d a_0$。下一个结果是这一事实的逆命题。

\begin{lemma}
设 $F$ 是 $\mathbb{Z}$ 上的函数，使得
$$G(n)=F(n)-F(n-1)$$
对充分大的 $n\in\mathbb{Z}$ 等于 $n$ 的 $d-1$ 次多项式。则 $F$ 对充分大的 $n\in\mathbb{Z}$ 等于 $n$ 的 $d$ 次多项式。
\end{lemma}

\begin{proof}
假设对 $n\ge N\ge d$ 有 $G(n)=P(n-1)$，其中 $P$ 是 $n$ 的 $d-1$ 次多项式。则由 1.6 有
$$P(n)=\sum_{i=0}^{d-1}c_i\binom{n}{d-i-1}.$$
于是对 $n\ge N+1$，
$$F(n)=\sum_{k=N+1}^{n}\bigl(F(k)-F(k-1)\bigr)+F(N)=\sum_{k=N+1}^{n}G(k)+F(N)=\sum_{k=d}^{n}P(k-1)+C,$$
其中 $C$ 是常数。此外，由上一定理中用到的恒等式，对 $q>s\ge 1$，
$$\binom{q}{s}=\sum_{j=s+1}^{q}\left(\binom{j}{s}-\binom{j-1}{s}\right)+1=\sum_{j=s+1}^{q}\binom{j-1}{s-1}+1=\sum_{j=s}^{q}\binom{j-1}{s-1}.$$
这意味着
\begin{align*}
\sum_{k=d}^{n}P(k-1)&=\sum_{k=d}^{n}\sum_{i=0}^{d-1}c_i\binom{k-1}{d-i-1}=\sum_{i=0}^{d-1}c_i\sum_{k=d}^{n}\binom{k-1}{d-i-1}\\
&=\sum_{i=0}^{d-1}c_i\sum_{k=d-i}^{n}\binom{k-1}{d-i-1}-\sum_{i=1}^{d-1}c_i\sum_{k=d-i}^{d-i-1}\binom{k-1}{d-i-1}\\
&=\sum_{i=0}^{d-1}c_i\binom{n}{d-i}+C_0
\end{align*}
对某个常数 $C_0$ 成立。
\end{proof}

特别地，由此可知 $\sum_{n\le N}\lambda(M^n)$ 对充分大的 $N\in\mathbb{Z}$ 等于次数为 $d_\lambda(M)$ 的多项式。此外，若对充分大的 $N\in\mathbb{Z}$ 写
$$\sum_{n\le N}\lambda(M^n)=a_0 N^d+a_1 N^{d-1}+\cdots+a_{d-1}N+a_d,$$
则 $d!\,a_0$ 是整数。

例如，若 $A=k[X_1,X_2,\dots,X_s]$，则次数 $\le N$ 的所有多项式空间的维数等于
$$\sum_{n=0}^{N}\dim_k(A^n)=\sum_{n=0}^{N}\binom{s+n-1}{s-1}=\binom{s+N}{s}=\frac{N^s}{s!}+\cdots.$$
\begin{annote}
讲义用 Hilbert 多项式 $P(M,t)$ 的次数与首项系数来定义 $\dim M$ 与重数 $e(M)$。这与交换代数里按素理想链定义的 Krull 维数一致：若 $R$ 是诺特局部环、$M$ 有限生成，则 $\dim M=\dim R/\mathrm{Ann}(M)$。文中反复使用的“$\dim M=\deg P(M,t)$”正是 Hilbert–Serre 定理的推论，它把“模的大小”换算成一条多项式的次数，后文特征簇的维数比较全靠这一点。[来源：Stacks Project Tag 00KE（Definition 10.60.2, Krull dimension）；原书 §1, p1–6]
\end{annote}



\section{局部环上模的维数}

\begin{lemma}[Nakayama]
设 $A$ 是以 $\mathfrak{m}$ 为极大理想的局部环。设 $V$ 是有限生成 $A$-模且满足 $\mathfrak{m}V=V$。则 $V=0$。
\end{lemma}

\begin{proof}
假设 $V\neq 0$。则可取 $V$ 作为 $A$-模的一组最小生成元 $v_1,\dots,v_s$。由假设，$v_s=\sum_{i=1}^s m_i v_i$，其中 $m_i\in\mathfrak{m}$。因此 $(1-m_s)v_s=\sum_{i=1}^{s-1}m_i v_i$。由于 $1-m_s$ 可逆，这说明 $v_1,\dots,v_{s-1}$ 生成 $V$，与假设矛盾。
\end{proof}

以下假设 $A$ 是诺特局部环，$\mathfrak{m}$ 为其极大理想，$k=A/\mathfrak{m}$ 为 $A$ 的剩余域。

\begin{lemma}
$\dim_k(\mathfrak{m}/\mathfrak{m}^2)<+\infty$。
\end{lemma}

\begin{proof}
由诺特性假设 $\mathfrak{m}$ 有限生成。若 $a_1,\dots,a_p$ 是 $\mathfrak{m}$ 的生成元，则它们在 $\mathfrak{m}/\mathfrak{m}^2$ 中的像 $\bar a_1,\dots,\bar a_p$ 张成它作为 $k$-向量空间。
\end{proof}

令 $s=\dim_k(\mathfrak{m}/\mathfrak{m}^2)$。则可找到 $a_1,\dots,a_s\in\mathfrak{m}$，使得 $\bar a_1,\dots,\bar a_s$ 构成 $\mathfrak{m}/\mathfrak{m}^2$ 的一组基。我们断言它们生成 $\mathfrak{m}$。设 $I$ 为由 $a_1,\dots,a_s$ 生成的理想。则 $I+\mathfrak{m}^2=\mathfrak{m}$ 且 $\mathfrak{m}(\mathfrak{m}/I)=\mathfrak{m}/I$。故由 2.1 得 $\mathfrak{m}/I=0$。因此我们证明了：

\begin{lemma}
正整数 $\dim_k(\mathfrak{m}/\mathfrak{m}^2)$ 等于 $\mathfrak{m}$ 的最小生成元个数。
\end{lemma}

$\mathfrak{m}$ 中满足 $(\bar a_1,\dots,\bar a_s)$ 构成 $\mathfrak{m}/\mathfrak{m}^2$ 的一组基的任意 $s$ 元组 $(a_1,\dots,a_s)$ 称为 $A$ 中的一个\emph{坐标系}。

显然，$(\mathfrak{m}^p;\ p\in\mathbb{Z}_+)$ 是 $A$ 的一个递减滤过。因此可以构成
$$\Gr A=\bigoplus_{p=0}^{\infty}\mathfrak{m}^p/\mathfrak{m}^{p+1}.$$
我们断言 $\Gr A$ 是 $k$ 上的有限生成代数，因而是诺特分次环。事实上，映射 $X_i\longmapsto \bar a_i\in\mathfrak{m}/\mathfrak{m}^2\subset\Gr A$ 扩充为 $k[X_1,\dots,X_s]$ 到 $\Gr A$ 上的满射。

设 $M$ 是有限生成 $A$-模。则可用 $\mathfrak{m}^p M$（$p\in\mathbb{Z}_+$）定义 $M$ 的递减滤过，并考虑分次 $\Gr A$-模
$$\Gr M=\bigoplus_{p=0}^{\infty}\mathfrak{m}^p M/\mathfrak{m}^{p+1}M.$$

\begin{lemma}
若 $M$ 是有限生成 $A$-模，则 $\Gr M$ 是有限生成 $\Gr A$-模。
\end{lemma}

\begin{proof}
由 $\Gr M$ 的定义可见对所有 $p\in\mathbb{Z}_+$ 有 $\mathfrak{m}\cdot \Gr_p M=\Gr_{p+1}M$，故 $\Gr_0 M=M/\mathfrak{m}M$ 生成 $\Gr M$。另一方面，$M/\mathfrak{m}M$ 是 $k$ 上的有限维线性空间。
\end{proof}

这意味着（由 1.2）$\dim_k(\mathfrak{m}^p M/\mathfrak{m}^{p+1}M)<+\infty$，特别地 $A$-模 $\mathfrak{m}^p M/\mathfrak{m}^{p+1}M$ 具有有限长度。由于长度显然是可加函数，由 1.5 可见 $p\longmapsto \operatorname{length}_A(\mathfrak{m}^p M/\mathfrak{m}^{p+1}M)=\dim_k(\mathfrak{m}^p M/\mathfrak{m}^{p+1}M)$ 对充分大的 $p\in\mathbb{Z}_+$ 等于有理系数多项式。此外，函数
$$p\longmapsto \operatorname{length}_A(M/\mathfrak{m}^p M)=\sum_{q=0}^{p-1}\operatorname{length}_A(\mathfrak{m}^q M/\mathfrak{m}^{q+1}M)$$
对充分大的 $p\in\mathbb{Z}_+$ 等于有理系数多项式，且其首项形如 $e\dfrac{p^d}{d!}$，其中 $e,d\in\mathbb{Z}_+$。我们记 $d(M)=d$，$e(M)=e$，并称这些数为 $M$ 的\emph{维数}和\emph{重数}。

现在讨论函数 $M\longmapsto d(M)$ 的一些性质。控制 $A$-模滤过的关键结果是 Artin--Rees 引理。

\begin{theorem}[Artin, Rees]
设 $M$ 是有限生成 $A$-模，$N$ 是其子模。则存在 $m_0\in\mathbb{Z}_+$ 使得对所有 $p\in\mathbb{Z}_+$，
$$\mathfrak{m}^{p+m_0}M\cap N=\mathfrak{m}^p(\mathfrak{m}^{m_0}M\cap N).$$
\end{theorem}

\begin{proof}
令 $A^*=\bigoplus_{n=0}^{\infty}\mathfrak{m}^n$。则 $A^*$ 具有天然的分次环结构。设 $(a_1,\dots,a_s)$ 是 $A$ 中的坐标系，则有自然的满射
$$A[a_1,\dots,a_s]\longrightarrow A^*,$$
且 $A^*$ 是分次诺特环。令 $M^*=\bigoplus_{n=0}^{\infty}\mathfrak{m}^n M$。则 $M^*$ 是分次 $A^*$-模。它显然作为 $A^*$-模由 $M^*_0=M$ 生成。由于 $M$ 是有限生成 $A$-模，我们断定 $M^*$ 是有限生成 $A^*$-模。

此外，令 $N^*=\bigoplus_{n=0}^{\infty}(N\cap \mathfrak{m}^n M)\subset M^*$。则由
$$\mathfrak{m}^p(N\cap \mathfrak{m}^n M)\subset \mathfrak{m}^p N\cap \mathfrak{m}^{n+p}M\subset N\cap \mathfrak{m}^{n+p}M$$
可知 $N^*$ 是 $M^*$ 的 $A^*$-子模。由于 $A^*$ 是诺特环，$N^*$ 有限生成。存在 $m_0\in\mathbb{Z}_+$ 使得 $\bigoplus_{n=0}^{m_0}(N\cap\mathfrak{m}^n M)$ 生成 $N^*$。于是对任意 $p\in\mathbb{Z}_+$，
$$N\cap\mathfrak{m}^{p+m_0}M=\sum_{s=0}^{m_0}\mathfrak{m}^{p+m_0-s}(N\cap\mathfrak{m}^s M)\subset \mathfrak{m}^p(N\cap\mathfrak{m}^{m_0}M)\subset N\cap\mathfrak{m}^{p+m_0}M.$$
\end{proof}

这一结果有如下推论——Krull 交定理。

\begin{theorem}[Krull]
设 $M$ 是有限生成 $A$-模。则
$$\bigcap_{p=0}^{\infty}\mathfrak{m}^p M=\{0\}.$$
\end{theorem}

\begin{proof}
令 $E=\bigcap_{p=0}^{\infty}\mathfrak{m}^p M$。则由 2.5，
$$E=\mathfrak{m}^{p+m_0}M\cap E=\mathfrak{m}^p(\mathfrak{m}^{m_0}M\cap E)=\mathfrak{m}^p E,$$
特别地 $\mathfrak{m}E=E$，由 Nakayama 引理得 $E=0$。
\end{proof}

\begin{lemma}
设
$$0\longrightarrow M'\longrightarrow M\longrightarrow M''\longrightarrow 0$$
是有限生成 $A$-模的正合列。则
\begin{enumerate}[label=(\roman*)]
  \item $d(M)=\max(d(M'),d(M''))$；
  \item 若 $d(M)=d(M')=d(M'')$，则 $e(M)=e(M')+e(M'')$。
\end{enumerate}
\end{lemma}

\begin{proof}
可把 $M'$ 看作 $M$ 的子模。若给 $M$ 配置滤过 $\mathfrak{m}^p M$（$p\in\mathbb{Z}_+$），给 $M'$ 和 $M''$ 配置诱导滤过 $M'\cap\mathfrak{m}^p M$（$p\in\mathbb{Z}_+$）和 $\mathfrak{m}^p M''$（$p\in\mathbb{Z}_+$），则得到正合列
$$0\longrightarrow \Gr M'\longrightarrow \Gr M\longrightarrow \Gr M''\longrightarrow 0.$$
这意味着对任意 $p\in\mathbb{Z}_+$，
$$\operatorname{length}_A(\mathfrak{m}^p M/\mathfrak{m}^{p+1}M)
=\operatorname{length}_A\bigl((M'\cap\mathfrak{m}^p M)/(M'\cap\mathfrak{m}^{p+1}M)\bigr)+\operatorname{length}_A(\mathfrak{m}^p M''/\mathfrak{m}^{p+1}M'')$$
并且求和得
$$\operatorname{length}_A(M/\mathfrak{m}^p M)=\operatorname{length}_A(M'/(M'\cap\mathfrak{m}^p M))+\operatorname{length}_A(M''/\mathfrak{m}^p M'').$$
因此函数 $p\longmapsto \operatorname{length}_A(M'/(M'\cap\mathfrak{m}^p M))$ 对充分大的 $p\in\mathbb{Z}_+$ 等于 $p$ 的多项式。另一方面，由 2.5，
$$\mathfrak{m}^{p+m_0}M'\subset \mathfrak{m}^{p+m_0}M\cap M'\subset \mathfrak{m}^p M';$$
故对充分大的 $p\in\mathbb{Z}_+$，函数 $p\longmapsto \operatorname{length}_A(M'/(M'\cap\mathfrak{m}^p M))$ 与 $p\longmapsto \operatorname{length}_A(M'/\mathfrak{m}^p M')$ 由首项相同的多项式给出。
\end{proof}

\begin{corollary}
设 $A$ 是诺特局部环，$s=\dim_k(\mathfrak{m}/\mathfrak{m}^2)$。则对任意有限生成 $A$-模 $M$ 有 $d(M)\le s$。
\end{corollary}

\begin{proof}
由 2.7，只需证明 $d(A)\le s$。这由 $k[X_1,\dots,X_s]$ 到 $\Gr A$ 的满同态的存在性，以及 $s$ 个变量中次数 $\le n$ 的多项式空间的维数是 $n$ 的 $s$ 次多项式这一事实立即得出。
\end{proof}

若诺特局部环满足 $d(A)=\dim_k(\mathfrak{m}/\mathfrak{m}^2)$，则称为\emph{正则的}。

\begin{theorem}
设 $A$ 是诺特局部环，$(a_1,a_2,\dots,a_s)$ 是 $A$ 中的坐标系。则下列条件等价：
\begin{enumerate}[label=(\roman*)]
  \item $A$ 是正则局部环；
  \item 由 $X_i\longmapsto \bar a_i$（$1\le i\le s$）定义的 $k[X_1,X_2,\dots,X_s]$ 到 $\Gr A$ 的典范态射是同构。
\end{enumerate}
\end{theorem}

\begin{proof}
由定义，$k[X_1,\dots,X_s]$ 到 $\Gr A$ 的典范态射是满射。设 $I$ 是 $k[X_1,\dots,X_s]$ 到 $\Gr A$ 的自然满射的核，它是分次理想。若 $I\neq 0$，它含有次数 $d>0$ 的齐次多项式 $P$。设 $J$ 是 $k[X_1,X_2,\dots,X_s]$ 中由 $P$ 生成的理想。则其 Poincaré 级数为 $P(J,t)=\dfrac{t^d}{(1-t)^s}$。显然，
$$P(k[X_1,\dots,X_s]/J,t)=P(k[X_1,\dots,X_s],t)-P(J,t)=\frac{1-t^d}{(1-t)^s}=\frac{1+t+\cdots+t^{d-1}}{(1-t)^{s-1}}.$$
Poincaré 级数 $P(k[X_1,\dots,X_s]/J,t)$ 在 $1$ 处的极点阶数为 $s-1$，由 1.5 函数 $\dim_k(k[X_1,\dots,X_s]/J)_n$ 对充分大的 $n\in\mathbb{Z}_+$ 由 $n$ 的 $s-2$ 次多项式给出。由此可知 $\dim_k(k[X_1,\dots,X_s]/I)_n=\dim_k\Gr_n A$ 对充分大的 $n\in\mathbb{Z}_+$ 由次数 $\le s-2$ 的多项式给出。这意味着 $d(A)\le s-1$。因此 $I=0$ 当且仅当 $d(A)=s$。
\end{proof}

\begin{theorem}
设 $A$ 是正则局部环。则 $A$ 是整环。
\end{theorem}

\begin{proof}
设 $a,b\in A$ 且 $a\neq 0$，$b\neq 0$。则由 2.6，可找到 $p,q\in\mathbb{Z}_+$ 使得 $a\in\mathfrak{m}^p$，$a\notin\mathfrak{m}^{p+1}$，且 $b\in\mathfrak{m}^q$，$b\notin\mathfrak{m}^{q+1}$。则它们在 $\Gr_p A$ 和 $\Gr_q A$ 中的像 $\bar a$ 与 $\bar b$ 均非零，且由于 $\Gr A$ 由 2.9 是整环，得 $\bar a\bar b\neq 0$。因此 $ab\neq 0$。
\end{proof}

最后讨论一个以后将起重要作用例子。设 $k$ 是域，$A=k[X_1,X_2,\dots,X_n]$ 是系数在 $k$ 中的 $n$ 元多项式环，$\widehat A=k[[X_1,X_2,\dots,X_n]]$ 是系数在 $k$ 中的 $n$ 元形式幂级数环。容易验证 $\widehat A$ 是以 $X_1,X_2,\dots,X_n$ 生成的极大理想 $\widehat{\mathfrak{m}}$ 的局部环。此外，$k[X_1,X_2,\dots,X_n]$ 到 $\Gr\widehat A$ 的典范态射显然是同构。

对任意 $x\in k^n$，记 $\mathfrak{m}_x$ 为 $A$ 中由 $X_i-x_i$（$1\le i\le n$）生成的极大理想。则它在 $A$ 中的补集是 $A$ 中的乘法系，记 $A_x$ 为 $A$ 相应的局部化。它同构于 $k^n$ 上在 $x$ 处正则的所有有理函数环。这显然是诺特局部环。$\mathfrak{m}_x$ 的局部化是极大理想 $\mathfrak{n}_x=(\mathfrak{m}_x)_x$，即所有在 $x$ 处消失的有理函数。由 $X_i\longmapsto X_i-x_i$（$1\le i\le n$）定义的 $A$ 的自同构给出对任意 $x\in k^n$，$A^0$ 与 $A_x$ 的一个同构。另一方面，$A$ 到 $\widehat A$ 的自然同态扩充为 $A^0$ 到 $\widehat A$ 的单射同态。此同态保持这些局部环上的滤过，并诱导 $\Gr A^0$ 到 $\Gr\widehat A$ 的典范同构。因此有如下结果。

\begin{proposition}
环 $A_x$（$x\in k^n$）是 $n$ 维正则局部环。
\end{proposition}
\begin{annote}
讲义 §2 证的正是 Artin–Rees 引理（Lemma 10.51.2）：$R$ 诺特、$I\subset R$ 理想、$N\subset M$ 有限生成，则存在 $c>0$ 使 $I^nM\cap N=I^{n-c}(I^cM\cap N)$ 对 $n\ge c$ 成立。文中引入的 $A^*=\bigoplus_{n\ge 0}\mathfrak m^n$ 就是标准证明里的 Rees 环 $S=R\oplus I\oplus I^2\oplus\cdots$，同一构造还给出 Krull 交定理 $\bigcap_{n\ge0}I^nM=0$（$R$ 诺特局部、$I$ 真理想）。这两条是 §3、§4 滤过论证的唯一外部输入。[来源：Stacks Project Tag 00IN（Lemma 10.51.2, Artin–Rees）、Tag 00IP（Lemma 10.51.4, Krull 交定理）]
\end{annote}



\section{滤过环上模的维数}

设 $D$ 是带单位的环，$(D_n;\ n\in\mathbb{Z})$ 是 $D$ 关于加法子群的递增滤过，满足
\begin{enumerate}[label=(\roman*)]
  \item $D_n=\{0\}$ 当 $n<0$；
  \item $\bigcup_{n\in\mathbb{Z}}D_n=D$；
  \item $1\in D_0$；
  \item 对任意 $n,m\in\mathbb{Z}$，$D_n\cdot D_m\subset D_{n+m}$；
  \item 对任意 $n,m\in\mathbb{Z}$，$[D_n,D_m]\subset D_{n+m-1}$。
\end{enumerate}
则 $\Gr D=\bigoplus_{n\in\mathbb{Z}}\Gr_n D=\bigoplus_{n\in\mathbb{Z}}D_n/D_{n-1}$ 是带单位的分次环。性质 (v) 蕴含它是交换的。特别地 $D_0=\Gr_0 D$ 是带单位的交换环。因此可以把 $\Gr D$ 看作 $D_0$-代数。进一步假设 $D$ 满足
\begin{enumerate}[label=(\roman*),resume]
  \item $\Gr D$ 是诺特环；
  \item $\Gr_1 D$ 作为 $D_0$-代数生成 $\Gr D$。
\end{enumerate}
则由 1.1，$D_0$ 是诺特环。此外，由 (vi)、(vii) 和 1.2 可知可取有限多个元素 $x_1,x_2,\dots,x_s\in\Gr_1 D$，使得 $\Gr D$ 作为 $D_0$-代数由它们生成。显然由 (vii)，对 $n\in\mathbb{Z}_+$ 有
$$\Gr_{n+1}D=\Gr_1 D\cdot\Gr_n D,$$
因此
$$D_{n+1}=D_n\cdot D_1\quad\text{对 }n\in\mathbb{Z}_+.$$

设 $D^\circ$ 是 $D$ 的反环。则滤过 $(D_n;\ n\in\mathbb{Z})$ 关于 $D^\circ$ 的乘法具有相同性质。此外，恒等映射 $D\longrightarrow D^\circ$ 诱导 $\Gr D$ 与 $\Gr D^\circ$ 的分次环同构。

设 $M$ 是 $D$-模。$M$ 关于加法子群的递增滤过 $F M=(F_n M;\ n\in\mathbb{Z})$ 称为 $D$-模滤过，若对 $n,m\in\mathbb{Z}$ 有 $D_n\cdot F_m M\subset F_{m+n}M$。特别地 $F_n M$ 是 $D_0$-模。

$D$-模滤过 $F M$ 称为\emph{Hausdorff 的}，若 $\bigcap_{n\in\mathbb{Z}}F_n M=\{0\}$。它是\emph{穷尽的}，若 $\bigcup_{n\in\mathbb{Z}}F_n M=M$。若存在 $m_0\in\mathbb{Z}$ 使得对所有 $n\in\mathbb{Z}_+$ 和 $m\ge m_0$ 有 $D_n\cdot F_m M=F_{m+n}M$，则称它是\emph{稳定的}。

$D$-模滤过称为\emph{好的}，若
\begin{enumerate}[label=(\roman*)]
  \item 对充分负的 $n\in\mathbb{Z}$，$F_n M=\{0\}$；
  \item 滤过 $F M$ 是穷尽的；
  \item $F_n M$（$n\in\mathbb{Z}$）是有限生成 $D_0$-模；
  \item 滤过 $F M$ 是稳定的。
\end{enumerate}
特别地，好滤过是 Hausdorff 的。

\begin{lemma}
设 $F M$ 是 $M$ 的穷尽 Hausdorff $D$-模滤过。则下列陈述等价：
\begin{enumerate}[label=(\roman*)]
  \item $F M$ 是好滤过；
  \item $\Gr D$-模 $\Gr M$ 是有限生成的。
\end{enumerate}
\end{lemma}

\begin{proof}
先证 (i)$\Rightarrow$(ii)。存在 $m_0\in\mathbb{Z}$ 使得对所有 $n\in\mathbb{Z}_+$ 有 $D_n\cdot F_{m_0}M=F_{n+m_0}M$。因此对 $n\in\mathbb{Z}_+$ 有 $\Gr_n D\cdot\Gr_{m_0}M=\Gr_{n+m_0}M$。由此可知 $\bigoplus_{n\le m_0}\Gr_n M$ 作为 $\Gr D$-模生成 $\Gr M$。由于 $F_n M$ 是有限生成 $D_0$-模，$\Gr_n M$ 也是有限生成 $D_0$-模。这说明（由于对充分负的 $n$，$F_n M=\{0\}$）$\bigoplus_{n\le m_0}\Gr_n M$ 是有限生成 $D_0$-模。

(ii)$\Rightarrow$(i)。显然对充分负的 $n\in\mathbb{Z}$ 有 $\Gr_n M=\{0\}$。此外由 1.2，所有 $\Gr_n M$ 是有限生成 $D_0$-模。正合列
$$0\longrightarrow F_{n-1}M\longrightarrow F_n M\longrightarrow \Gr_n M\longrightarrow 0$$
蕴含 $\bigcap_{n\in\mathbb{Z}}F_n M=F_{n_0}M$（对充分负的 $n$ 有 $F_n M=F_{n-1}M$，故存在 $n_0\in\mathbb{Z}$）。由于滤过 $F M$ 是 Hausdorff 的，$F_{n_0}M=\{0\}$。这由 $n$ 上的归纳蕴含所有 $F_n M$ 是有限生成 $D_0$-模。

设 $m_0\in\mathbb{Z}$ 使得 $\bigoplus_{n\le m_0}\Gr_n M$ 作为 $\Gr D$-模生成 $\Gr M$。设 $m\ge m_0$。则
\begin{align*}
\Gr_{m+1}M&=\bigoplus_{k\le m_0}\Gr_{m+1-k}D\cdot\Gr_k M
=\bigoplus_{k\le m_0}\Gr_1 D\cdot\Gr_{m-k}D\cdot\Gr_k M\\
&\subset\Gr_1 D\cdot\Gr_m M\subset\Gr_{m+1}M,
\end{align*}
即 $\Gr_1 D\cdot\Gr_m M=\Gr_{m+1}M$。这意味着
$$F_{m+1}M=D_1\cdot F_m M+F_m M=D_1\cdot F_m M,$$
并由 $n$ 上的归纳，
$$F_{m+n}M=D_1\cdot D_1\cdots D_1\cdot F_m M\subset D_n\cdot F_m M\subset F_{m+n}M.$$
因此对所有 $n\in\mathbb{Z}_+$ 有 $F_{m+n}M=D_n\cdot F_m M$。故 $F M$ 是好滤过。
\end{proof}

特别地，$(D_n;\ n\in\mathbb{Z})$ 是把 $D$ 看作左乘 $D$-模时的好滤过。

\begin{remark}
由证明可知，好滤过定义中的稳定性条件可替换为表面上较弱的条件：

(iv)$'$ 存在 $m_0\in\mathbb{Z}$ 使得对所有 $n\in\mathbb{Z}_+$ 有 $D_n\cdot F_{m_0}M=F_{m_0+n}M$。
\end{remark}

\begin{lemma}
设 $M$ 是带好滤过 $F M$ 的 $D$-模。则 $M$ 是有限生成的。
\end{lemma}

\begin{proof}
由定义 $\bigcup_{n\in\mathbb{Z}}F_n M=M$，且对某充分大的 $m_0\in\mathbb{Z}$ 和 $n\in\mathbb{Z}_+$ 有 $F_{n+m_0}M=D_n\cdot F_{m_0}M$。因此 $F_{m_0}M$ 作为 $D$-模生成 $M$。由于 $F_{m_0}M$ 是有限生成 $D_0$-模，结论成立。
\end{proof}

\begin{lemma}
设 $M$ 是有限生成 $D$-模。则 $M$ 有好滤过。
\end{lemma}

\begin{proof}
设 $U$ 是作为 $D$-模生成 $M$ 的有限生成 $D_0$-模。令 $n<0$ 时 $F_n M=0$，$n\ge 0$ 时 $F_n M=D_n\cdot U$。则 $U=\Gr_0 M$，且
$$\Gr_n M=F_n M/F_{n-1}M=(D_n\cdot U)/(D_{n-1}\cdot U)\subset\Gr_n D\cdot\Gr_0 M\subset\Gr_n M,$$
即对所有 $n\in\mathbb{Z}_+$ 有 $\Gr_n M=\Gr_n D\cdot\Gr_0 M$。故 $\Gr M$ 作为 $\Gr D$-模有限生成。结论由 3.1 得出。
\end{proof}

引理 3.1 和 3.3 蕴含：具有好滤过的 $D$-模恰好是有限生成 $D$-模。

\begin{proposition}
环 $D$ 是左诺特且右诺特的。
\end{proposition}

\begin{proof}
设 $L$ 是 $D$ 的左理想。$D$ 的自然滤过在 $L$ 上诱导滤过 $(L_n=L\cap D_n;\ n\in\mathbb{Z})$。这显然是 $D$-模滤过。分次模 $\Gr L$ 自然是 $\Gr D$ 的理想，且由于 $\Gr D$ 是诺特环，它作为 $\Gr D$-模有限生成。因此由 3.1 滤过 $(L_n;\ n\in\mathbb{Z})$ 是好滤过，由 3.3 $L$ 是有限生成的。这证明 $D$ 是左诺特的。

要得到右诺特性质，只需把 $D$ 换成其反环 $D^\circ$。
\end{proof}

若 $D$-模 $M$ 有两个滤过 $F M$ 和 $F'M$，若存在 $k\in\mathbb{Z}_+$ 使得对所有 $n\in\mathbb{Z}$ 有 $F_n M\subset F'_{n+k}M$，则称 $F M$ 比 $F'M$\emph{更细}。若 $F M$ 比 $F'M$ 更细且 $F'M$ 比 $F M$ 更细，则称它们\emph{等价}。

\begin{lemma}
设 $F M$ 是有限生成 $D$-模 $M$ 上的好滤过。则 $F M$ 比 $M$ 上任何其他穷尽 $D$-模滤过更细。
\end{lemma}

\begin{proof}
固定 $m_0\in\mathbb{Z}_+$ 使得对所有 $n\in\mathbb{Z}_+$ 有 $D_n\cdot F_{m_0}M=F_{n+m_0}M$。设 $F'M$ 是 $M$ 上另一个穷尽 $D$-模滤过。则 $F_{m_0}M$ 作为 $D_0$-模有限生成。由于 $F'M$ 是穷尽的，存在 $p\in\mathbb{Z}$ 使得 $F_{m_0}M\subset F'_p M$。由于 $F M$ 是好滤过，存在 $n_0$ 使得 $F_{n_0}M=\{0\}$。令 $k=p+|n_0|$。显然对 $m\le n_0$ 有 $F_m M=0\subset F'_{m+k}M$。对 $n_0<m\le m_0$，有 $-|n_0|\le n_0<m$ 且 $p=-|n_0|+k<m+k$。这给出
$$F_m M\subset F_{m_0}M\subset F'_p M\subset F'_{m+k}M.$$
最后对 $m>m_0$，由于 $m_0$ 为正有 $m-m_0\le m$，且 $p\le k$。于是
$$F_m M=D_{m-m_0}\cdot F_{m_0}M\subset D_m\cdot F'_p M\subset F'_{m+p}M\subset F'_{m+k}M.$$
\end{proof}

\begin{corollary}
有限生成 $D$-模上任意两个好滤过等价。
\end{corollary}

设 $M$ 是有限生成 $D$-模，$F M$ 是其上好滤过。则 $\Gr M$ 是有限生成 $\Gr D$-模，故可应用 §1 关于 Hilbert 多项式的结果。设 $\lambda$ 是有限生成 $D_0$-模上的可加函数。再假设 $\lambda$ 在有限生成 $D_0$-模范畴 $\mathcal{M}_{fg}(D_0)$ 的对象上仅取非负值。则由 1.5，
$$\lambda(F_n M)-\lambda(F_{n-1}M)=\lambda(\Gr_n M)$$
对充分大的 $n\in\mathbb{Z}_+$ 等于 $n$ 的多项式。由 1.7 这意味着 $\lambda(F_n M)$ 对充分大的 $n\in\mathbb{Z}_+$ 等于 $n$ 的多项式。若 $F'M$ 是 $M$ 上另一个好滤过，由 3.7 知 $F M$ 与 $F'M$ 等价，即存在 $k\in\mathbb{Z}_+$ 使得
$$F_n M\subset F'_{n+k}M\subset F_{n+2k}M$$
对所有 $n\in\mathbb{Z}$ 成立。由于 $\lambda$ 可加且仅取非负值，我们断定
$$\lambda(F_n M)\le\lambda(F'_{n+k}M)\le\lambda(F_{n+2k}M)$$
对所有 $n\in\mathbb{Z}$ 成立。这意味着对充分大的 $n$，表示 $\lambda(F_n M)$ 与 $\lambda(F'_n M)$ 的多项式有相等的首项。我们记这些多项式的公共次数为 $d_\lambda(M)$，并称之为 $D$-模 $M$ 关于 $\lambda$ 的\emph{维数}。由 1.6，这些多项式的首项系数形如 $e_\lambda(M)/d_\lambda(M)!$，其中 $e_\lambda(M)\in\mathbb{N}$。我们称 $e_\lambda(M)$ 为 $D$-模 $M$ 关于 $\lambda$ 的\emph{重数}。

设
$$0\longrightarrow M'\xrightarrow{\ f\ }M\xrightarrow{\ g\ }M''\longrightarrow 0$$
是 $D$-模的正合列。若给 $M$ 配置 $D$-模滤过 $F M$，它在 $M'$ 上诱导滤过 $F M'=(f^{-1}(f(M')\cap F_n M);\ n\in\mathbb{Z})$，在 $M''$ 上诱导滤过 $F M''=(g(F_n M);\ n\in\mathbb{Z})$。显然这些滤过是 $D$-模滤过。此外，序列
$$0\longrightarrow \Gr M'\xrightarrow{\ \Gr f\ }\Gr M\xrightarrow{\ \Gr g\ }\Gr M''\longrightarrow 0$$
是正合的。若滤过 $F M$ 是好的，则 $\Gr M$ 是有限生成 $\Gr D$-模，故 $\Gr M'$ 和 $\Gr M''$ 都是有限生成 $\Gr D$-模。由 3.1，$F M'$ 和 $F M''$ 是好滤过。因此我们证明了如下结果。

\begin{lemma}
设
$$0\longrightarrow M'\longrightarrow M\longrightarrow M''\longrightarrow 0$$
是 $D$-模的正合列。若 $F M$ 是 $M$ 上的好滤过，则诱导滤过 $F M'$ 和 $F M''$ 是好的。
\end{lemma}

由上述讨论
$$\lambda(\Gr_n M)=\lambda(\Gr_n M')+\lambda(\Gr_n M'')$$
对所有 $n\in\mathbb{Z}$ 成立。这由 $n$ 上的归纳蕴含
$$\lambda(F_n M)=\lambda(F_n M')+\lambda(F_n M'')$$
对所有 $n\in\mathbb{Z}$ 成立。这导出如下结果。

\begin{proposition}
设
$$0\longrightarrow M'\longrightarrow M\longrightarrow M''\longrightarrow 0$$
是有限生成 $D$-模的正合列。则
\begin{enumerate}[label=(\roman*)]
  \item $d_\lambda(M)=\max(d_\lambda(M'),d_\lambda(M''))$；
  \item 若 $d_\lambda(M)=d_\lambda(M')=d_\lambda(M'')$，则 $e_\lambda(M)=e_\lambda(M')+e_\lambda(M'')$。
\end{enumerate}
\end{proposition}

最后，设 $\varphi$ 是环 $D$ 的自同构且 $\varphi(D_0)=D_0$。可定义从 $D$-模范畴 $\mathcal{M}(D)$ 到自身的函子 $\tilde\varphi$，它把 $D$-模 $M$ 映为 $D$-模 $\tilde\varphi(M)$，其加法群结构不变，而 $D$ 的作用由 $(T,m)\longmapsto\varphi(T)m$（$T\in D$，$m\in M$）给出。显然 $\tilde\varphi$ 是范畴 $\mathcal{M}(D)$ 的自同构，且保持有限生成 $D$-模。

\begin{proposition}
设 $M$ 是有限生成 $D$-模。则
$$d_\lambda(\tilde\varphi(M))=d_\lambda(M).$$
\end{proposition}

\begin{proof}
设 $T_1,T_2,\dots,T_s$ 是 $D_1$ 中代表生成 $\Gr D$ 作为 $D_0$-代数的 $\Gr_1 D$ 中那些类的代表元。则存在 $d\in\mathbb{N}$ 使得对 $1\le i\le s$ 有 $\varphi(T_i)\in D_d$。由于 $T_1,T_2,\dots,T_s$ 和 $1$ 作为 $D_0$-模生成 $D_1$，我们断定 $\varphi(D_1)\subset D_d$。

设 $F M$ 是 $M$ 的好滤过。定义滤过 $F\tilde\varphi(M)$ 为
$$F_p\tilde\varphi(M)=F_{dp}M\quad (p\in\mathbb{Z}).$$
显然 $F\tilde\varphi(M)$ 是由有限生成 $D_0$-子模构成的 $\tilde\varphi(M)$ 的递增滤过。此外，对 $m\in\mathbb{Z}$，
$$D_1\cdot F_m\tilde\varphi(M)=\varphi(D_1)F_{dm}M\subset D_dF_{dm}M\subset F_{d(m+1)}M=F_{m+1}\tilde\varphi(M).$$
故由归纳，对所有 $n,m\in\mathbb{Z}$，
$$D_n\cdot F_m\tilde\varphi(M)=D_1\cdot D^{n-1}\cdot F_m\tilde\varphi(M)\subset D_1F_{m+n-1}\tilde\varphi(M)\subset F_{m+n}\tilde\varphi(M),$$
即 $F\tilde\varphi(M)$ 是 $D$-模滤过。由 3.6，存在比此滤过更细的好滤过 $F'\tilde\varphi(M)$，即存在 $k\in\mathbb{Z}_+$ 使得
$$F'_n\tilde\varphi(M)\subset F_{n+k}\tilde\varphi(M)=F_{d(n+k)}M$$
对所有 $n\in\mathbb{Z}$ 成立。因此
$$\lambda(F'_n\tilde\varphi(M))\le\lambda(F_{d(n+k)}M)$$
对 $n\in\mathbb{Z}$ 成立。对充分大的 $n\in\mathbb{Z}$，$\lambda(F_{d(n+k)}M)$ 等于 $n$ 的多项式，其首项为
$$\frac{e_\lambda(M)d^{d_\lambda(M)}}{d_\lambda(M)!}\,n^{d_\lambda(M)}.$$
由于对充分大的 $n\in\mathbb{Z}$，$\lambda(F'_n\tilde\varphi(M))$ 也由次数为 $d_\lambda(\tilde\varphi(M))$ 的多项式给出，我们断定 $d_\lambda(\tilde\varphi(M))\le d_\lambda(M)$。对 $\varphi^{-1}$ 应用同样推理也得到
$$d_\lambda(M)=d_\lambda(\tilde\varphi^{-1}(\tilde\varphi(M)))\le d_\lambda(\tilde\varphi(M)).$$
\end{proof}
\begin{annote}
“好滤过”者，指滤过模 $M$ 的分次 $\mathrm{gr}M$ 在 $\mathrm{gr}A$ 上有限生成；等价说法是 $M$ 为有限个 $A_{\le d}$ 元素的 $A$-线性组合。好滤过的选择不影响 $\mathrm{gr}M$ 的 Hilbert 函数的主项，这正是文中可放心用 $\mathrm{gr}M$ 替代 $M$ 计算维数的原因：两条好滤过生成的 $\mathrm{gr}$ 只差有限维向量空间。[来源：原书 §3, p15–22；译者按]
\end{annote}



\section{多项式环上模的维数}

设 $A=k[X_1,\dots,X_n]$，其中 $k$ 是代数闭域。可按多项式的次数滤过 $A$，即令
$$A_m=\Bigl\{\sum_I c_I x^I\ \Big|\ c_I\in k,\ |I|\le m\Bigr\}.$$
则 $\Gr A=k[X_1,\dots,X_n]$，故 $A$ 满足前一节的性质 (i)--(vii)。

由于 $A^0=k$，可取函数 $\dim_k$ 作为可加函数 $\lambda$。这导出有限生成 $A$-模 $M$ 的维数 $d(M)$ 和重数 $e(M)$ 的概念。已知对任意 $p\in\mathbb{Z}_+$，
$$\dim_k A_p=\binom{n+p}{n}=\frac{p^n}{n!}+\text{（关于 $p$ 的低阶项）},$$
即 $d(A)=n$，$e(A)=1$。此外，对任意有限生成 $A$-模 $M$ 有正合列
$$0\longrightarrow K\longrightarrow A^p\longrightarrow M\longrightarrow 0,$$
故由 3.9 有 $d(M)\le n$。以后将给出 $d(M)$ 的几何解释。

设 $x\in k^n$，记 $\mathfrak{m}_x$ 为 $k[X_1,\dots,X_n]$ 中所有在 $x$ 处消失的多项式构成的极大理想。记 $A_x$ 为 $A$ 在 $x$ 处的局部化，即 $k^n$ 上在 $x$ 处正则的所有有理 $k$-值函数环。如 2.11 中所见，$A_x$ 是以 $\mathfrak{n}_x=(\mathfrak{m}_x)_x$ 为极大理想的 $n$ 维正则局部环，$\mathfrak{n}_x$ 由 $k^n$ 上在 $x$ 处消失的所有有理 $k$-值函数组成。设 $M$ 是 $A$-模。它在 $x$ 处的局部化 $M_x$ 是 $A_x$-模。定义 $M$ 的\emph{支集}为
$$\operatorname{supp}(M)=\{x\in k^n\mid M_x\neq 0\}.$$

\begin{lemma}
设
$$0\longrightarrow M'\longrightarrow M\longrightarrow M''\longrightarrow 0$$
是 $A$-模的正合列。则
$$\operatorname{supp}(M)=\operatorname{supp}(M')\cup\operatorname{supp}(M'').$$
\end{lemma}

\begin{proof}
由局部化的正合性可知
$$0\longrightarrow M'_x\longrightarrow M_x\longrightarrow M''_x\longrightarrow 0$$
是 $A_x$-模的正合列。这立即蕴含我们的结论。
\end{proof}

对理想 $I\subset k[X_1,\dots,X_n]$，记 $V(I)=\{x\in k^n\mid f(x)=0\ \text{对 }f\in I\}$。

\begin{proposition}
设 $M$ 是有限生成 $A$-模，$I$ 是它在 $A$ 中的零化子。则 $\operatorname{supp}(M)=V(I)$。
\end{proposition}

\begin{proof}
对 $M$ 的生成元个数归纳证明。

先设 $M$ 有一个生成元，即 $M=A/I$。则 $M_x=(A/I)_x=A_x/I_x$。设 $x\in V(I)$。则 $I\subset\mathfrak{m}_x$，$I_x\subset\mathfrak{n}_x$，故 $I_x\neq A_x$。由此可知 $(A/I)_x\neq 0$，即 $x\in\operatorname{supp}(M)$。反之，若 $x\notin V(I)$，存在 $f\in I$ 使得 $f(x)\neq 0$，即 $f\notin\mathfrak{m}_x$。因此 $f$ 在局部环 $A_x$ 中可逆，而 $f\in I_x$ 蕴含 $I_x=A_x$。故 $(A/I)_x=0$，$x\notin\operatorname{supp}(A/I)$。因此 $\operatorname{supp}(A/I)=V(I)$。

现考虑一般情形。设 $m_1,\dots,m_p$ 是 $M$ 的一组生成元。记 $M'$ 为由 $m_1,\dots,m_{p-1}$ 生成的子模。则有正合列
$$0\longrightarrow M'\longrightarrow M\longrightarrow M''\longrightarrow 0$$
且 $M''$ 是循环模。此外由 4.1，$\operatorname{supp}(M)=\operatorname{supp}(M')\cup\operatorname{supp}(M'')$。故由归纳假设，$\operatorname{supp}(M)=V(I')\cup V(I'')$，其中 $I'$ 和 $I''$ 分别是 $M'$ 和 $M''$ 的零化子。

显然 $I'\cdot I''$ 含于 $M$ 的零化子 $I$ 中。另一方面，$I$ 零化 $M'$ 和 $M''$，故 $I\subset I'\cap I''$。于是
$$I'\cdot I''\subset I\subset I'\cap I''.$$
这意味着
$$V(I')\cup V(I'')\subset V(I'\cap I'')\subset V(I)\subset V(I'\cdot I'').$$
设 $x\notin V(I')\cup V(I'')$。则存在 $f\in I'$ 和 $g\in I''$ 使得 $f(x)\neq 0$ 且 $g(x)\neq 0$。由此可知 $(fg)(x)=f(x)g(x)\neq 0$，$x\notin V(I'\cdot I'')$。故 $V(I'\cdot I'')\subset V(I')\cup V(I'')$，上述所有包含均为等式。因此 $V(I)=V(I')\cup V(I'')$ 且 $\operatorname{supp}(M)=V(I)$。
\end{proof}

这立即蕴含如下推论。

\begin{corollary}
设 $M$ 是有限生成 $A$-模。则支集 $\operatorname{supp}(M)$ 是 $k^n$ 中的 Zariski 闭子集。
\end{corollary}

下一个引理在某些归约论证中很有用。

\begin{lemma}
设 $B$ 是诺特交换环，$M\neq 0$ 是有限生成 $B$-模。则存在 $M$ 由 $B$-子模构成的滤过
$$0=M_0\subset M_1\subset\cdots\subset M_{n-1}\subset M_n=M$$
以及 $B$ 的素理想 $\mathfrak{J}_i$，使得对 $1\le i\le n$ 有 $M_i/M_{i-1}\cong B/\mathfrak{J}_i$。
\end{lemma}

\begin{proof}
对 $x\in M$ 令 $\operatorname{Ann}(x)=\{a\in B\mid ax=0\}$。设 $\mathcal{A}$ 是所有这种理想 $\operatorname{Ann}(x)$（$x\in M$，$x\neq 0$）的族。由于 $B$ 是诺特环，$\mathcal{A}$ 有极大元。设 $I$ 是 $\mathcal{A}$ 中的极大元。我们断言 $I$ 是素的。设 $x\in M$ 使得 $I=\operatorname{Ann}(x)$。则 $ab\in I$ 蕴含 $abx=0$。假设 $b\notin I$，即 $bx\neq 0$。则 $I\subset\operatorname{Ann}(bx)$ 且 $a\in\operatorname{Ann}(bx)$。由 $I$ 的极大性，$a\in\operatorname{Ann}(bx)=I$，故 $I$ 是素的。因此存在 $x\in M$ 使得 $\mathfrak{J}_1=\operatorname{Ann}(x)$ 是素的。若令 $M_1=Bx$，则 $M_1\cong B/\mathfrak{J}_1$。

现记 $\mathcal{F}$ 为 $M$ 的所有具有滤过 $0=N_0\subset N_1\subset\cdots\subset N_k=N$ 且 $N_i/N_{i-1}\cong B/\mathfrak{J}_i$（对某素理想 $\mathfrak{J}_i$）的 $B$-子模 $N$ 的族。由于 $M$ 是诺特模，$\mathcal{F}$ 含有一个极大元 $L$。假设 $L\neq M$。则会有正合列
$$0\longrightarrow L\longrightarrow M\longrightarrow L'\longrightarrow 0,$$
由证明的第一部分，$L'$ 会有形如 $B/\mathfrak{J}'$（对某素理想 $\mathfrak{J}'$）的子模 $N'$，与 $L$ 的极大性矛盾。故 $L=M$。这证明了具有所需性质的滤过的存在性。
\end{proof}

\begin{theorem}
设 $M$ 是有限生成 $A$-模，$\operatorname{supp}(M)$ 是其支集。则
$$d(M)=\dim\operatorname{supp}(M).$$
\end{theorem}

这一结果有如下局部版本。$A$ 在 $x\in k^n$ 处的局部化 $A_x$ 是诺特局部环。此外其极大理想 $\mathfrak{n}_x$ 由多项式 $X_i-x_i$（$1\le i\le n$）生成，且它们在 $\mathfrak{n}_x/\mathfrak{n}_x^2$ 中的像作为 $k$-向量空间张成它。因此 $X_i-x_i$（$1\le i\le n$）构成 $A_x$ 中的坐标系。对任意有限生成 $A$-模 $M$，其在 $x$ 处的局部化 $M_x$ 是有限生成 $A_x$-模，故可考虑其维数 $d(M_x)$。

对 $k$ 上任意代数簇 $V$ 和 $x\in V$，记 $\dim_x V$ 为 $V$ 在 $x$ 处的局部维数。

\begin{theorem}
设 $M$ 是有限生成 $A$-模，$x\in\operatorname{supp}(M)$。则
$$d(M_x)=\dim_x(\operatorname{supp}(M)).$$
\end{theorem}

同时证明 4.5 和 4.6。首先注意，若有 $A$-模的正合列
$$0\longrightarrow M'\longrightarrow M\longrightarrow M''\longrightarrow 0$$
且 4.5、4.6 对 $M'$ 和 $M''$ 成立，则由 3.9 和 4.1，
\begin{align*}
d(M)&=\max(d(M'),d(M''))=\max(\dim\operatorname{supp}(M'),\dim\operatorname{supp}(M''))\\
&=\dim(\operatorname{supp}(M')\cup\operatorname{supp}(M''))=\dim\operatorname{supp}(M).
\end{align*}
此外，对任意 $x\in\operatorname{supp}(M)$，由局部化的正合性有正合列
$$0\longrightarrow M'_x\longrightarrow M_x\longrightarrow M''_x\longrightarrow 0;$$
故由 2.7 和 4.1，
\begin{align*}
d(M_x)&=\max(d(M'_x),d(M''_x))=\max(\dim_x\operatorname{supp}(M'),\dim_x\operatorname{supp}(M''))\\
&=\dim_x(\operatorname{supp}(M')\cup\operatorname{supp}(M''))=\dim_x\operatorname{supp}(M).
\end{align*}

假设 4.5 和 4.6 对所有 $M=A/J$（$J$ 为素理想）成立。则前面的注记、4.4 以及滤过长度上的归纳将证明一般情形的陈述。

故可设 $M=A/J$，$J$ 为素理想。先设 $J$ 使得 $A/J$ 是 $k$ 上的有限维向量空间。则 $A/J$ 是整环且在 $k$ 上是整的，故它是 $k$ 的代数扩张的域。由于 $k$ 代数闭，$A/J=k$，$J$ 是极大理想。这时由 Hilbert 零点定理，$\operatorname{supp}(M)=V(J)$ 是 $k^n$ 中一点 $x$，即 $\dim\operatorname{supp}(M)=0$。另一方面，由于 $M_x$ 是一维线性空间，$d(M_x)=0$，结论显然。于是可设 $J$ 在 $A$ 中不是有限余维的，特别地不是极大理想。设 $J_1\supsetneq J$ 是不同于 $J$ 的素理想。则存在 $f\in J_1$ 使得 $f\notin J$。由此可知 $J\subset (f)+J\subset J_1$ 且 $J\neq(f)+J$。因此 $A/J_1$ 是 $A/((f)+J)$ 的商，而 $A/((f)+J)$ 是 $A/J$ 的商。此外 $A/((f)+J)=M/fM$。考虑 $M$ 上由乘以 $f$ 给出的自同态。若 $g+J$ 在此映射的核中，则 $0=f(g+J)=fg+J$，$fg\in J$。由于 $J$ 是素理想且 $f\notin J$，得 $g\in J$，$g+J=0$，此映射是单射。因此有 $A$-模的正合列
$$0\longrightarrow M\xrightarrow{\ f\ }M\longrightarrow M/fM\longrightarrow 0.$$
由 3.9 这意味着 $d(M/fM)\le d(M)$。若 $d(M/fM)=d(M)$，还会有 $e(M)=e(M)+e(M/fM)$，故 $e(M/fM)=0$。这只有在 $d(M/fM)=0$ 时才可能，而这时还会蕴含 $d(M)=0$、$M$ 有限维，与假设矛盾。因此 $d(M/fM)<d(M)$。由于 $A/J_1$ 是 $M/fM$ 的商，这意味着 $d(A/J_1)<d(A/J)$。

设 $x\in V(J_1)$。则由局部化得到 $A_x$-模的正合列
$$0\longrightarrow M_x\xrightarrow{\ f\ }M_x\longrightarrow M_x/fM_x\longrightarrow 0.$$
由 2.7 这意味着 $d(M_x/fM_x)\le d(M_x)$。若 $d(M_x/fM_x)=d(M_x)$，还会有 $e(M_x)=e(M_x)+e(M_x/fM_x)$，故 $e(M_x/fM_x)=0$。这只有在 $d(M_x/fM_x)=0$ 时才可能，而这时会蕴含 $\mathfrak{m}_x(M_x/fM_x)=M_x/fM_x$，由 Nakayama 引理 $M_x/fM_x=0$。这将推出乘以 $f$ 在 $M_x$ 上是满射，且由于 $f\in\mathfrak{m}_x$，由 Nakayama 引理这会蕴含 $M_x=0$，与假设矛盾。因此 $d(M_x/fM_x)<d(M_x)$。由于 $A/J_1$ 是 $M/fM$ 的商，这意味着 $d((A/J_1)_x)<d((A/J)_x)$。

设
$$Z_0=\{x\}\subset Z_1\subset\cdots\subset Z_{n-1}\subset Z_n=k^n$$
是 $k^n$ 中非空不可约闭子集的极大链。则
$$I(Z_0)=\mathfrak{m}_x\supset I(Z_1)\supset\cdots\supset I(Z_{n-1})\supset I(Z_n)=\{0\}$$
是 $A$ 中素理想的极大链。由上述论证有以下严格不等式链
$$0\le d(A/I(Z_0))<d(A/I(Z_1))<\cdots<d(A/I(Z_n))=d(A)=n,$$
以及
$$0\le d((A/I(Z_0))_x)<d((A/I(Z_1))_x)<\cdots<d((A/I(Z_n))_x)=d(A_x)=n,$$
（由 2.11）。由此可知
$$d((A/I(Z_j))_x)=d(A/I(Z_j))=j=\dim Z_j\quad\text{对 }0\le j\le n.$$
由于每个闭不可约子集 $Z$ 都可置于极大链中，可知对任意闭不可约子集 $Z\subset k^n$ 和任意 $x\in Z$ 有 $d((A/I(Z))_x)=d(A/I(Z))=\dim Z$。另一方面，这意味着对 $A$ 中任意素理想 $J$ 和 $x\in V(J)$ 有 $d((A/J)_x)=d(A/J)=\dim V(J)$。由 4.2，这完成了 4.5 和 4.6 的证明。

下一个结果由 4.5 和 4.6 立即得出。

\begin{corollary}
设 $M$ 是有限生成 $A$-模。则
$$d(M)=\sup_{x\in\operatorname{supp}(M)} d(M_x).$$
\end{corollary}

最后证明一个以后需要的结果。

\begin{lemma}
设 $I$ 是 $A$ 中的理想。则 $\dim V(I)=\dim V(\Gr I)$。
\end{lemma}

\begin{proof}
$A$-模的短正合列
$$0\longrightarrow I\longrightarrow A\longrightarrow A/I\longrightarrow 0$$
（其中各模配置由 $A$ 的自然滤过诱导的滤过）导出分次 $A$-模的短正合列
$$0\longrightarrow\Gr I\longrightarrow A\longrightarrow\Gr(A/I)\longrightarrow 0.$$
因此
\begin{align*}
\dim_k F_p(A/I)&=\sum_{q=0}^{p}\bigl(\dim_k F_q(A/I)-\dim_k F_{q-1}(A/I)\bigr)\\
&=\sum_{q=0}^{p}\dim_k\Gr_q(A/I)=\sum_{q=0}^{p}\bigl(\dim_k\Gr_q A-\dim\Gr_q I\bigr)\\
&=\dim_k F_p A-\dim_k F_p\Gr I=\dim_k F_p(A/\Gr I).
\end{align*}
因此 $d(A/I)=d(A/\Gr I)$。结论由 4.4 和 4.5 得出。
\end{proof}

\section{多项式系数微分算子环}

设 $k$ 是特征 $0$ 的域，$A$ 是 $k$ 上的交换代数。设 $\End_k(A)$ 是 $A$ 的所有 $k$-线性自同态构成的代数。以换位子 $[S,T]=ST-TS$（对 $S,T\in\End_k(A)$）它是一个 Lie 代数。显然 $\End_k(A)$ 包含 $A$-线性自同态的集合 $\End_A(A)$ 作为子代数。对任意 $a\in A$，可附着由 $b\longmapsto ab$（$b\in A$）定义的 $A$-线性自同态。由于此自同态在 $1$ 处取值 $a$，该映射显然是代数的单射。

另一方面，若 $T\in\End_A(A)$，则对任意 $b\in A$ 有
$$T(b)=bT(1)=T(1)b,$$
即 $T$ 由乘以 $T(1)$ 给出。这导出如下结果。

\begin{lemma}
把 $a\in A$ 映为 $A$-线性自同态 $b\longmapsto ab$（$b\in A$）的代数同态是 $A$ 到 $\End_A(A)$ 上的同构。
\end{lemma}

以下把 $A$ 与 $\End_k(A)$ 的子代数 $\End_A(A)$ 等同。

$A$ 的一个 $k$-导子是满足
$$T(ab)=T(a)b+aT(b)\quad\text{对任意 }a,b\in A$$
的 $T\in\End_k(A)$。特别地 $[T,a](b)=T(ab)-aT(b)=T(a)b$，即对任意 $a\in A$ 有 $[T,a]=T(a)\in A$。这意味着对任意 $a_0,a_1\in A$ 有 $[[T,a_0],a_1]=0$。

这引出如下定义。设 $n\in\mathbb{Z}_+$。若对任意 $a_0,a_1,\dots,a_n\in A$ 有
$$[\cdots[[T,a_0],a_1],\dots,a_n]=0,$$
则称 $T\in\End_k(A)$ 是 $A$ 上阶数 $\le n$ 的（$k$-线性）微分算子。记 $\Diff_k(A)$ 为 $A$ 上所有微分算子的空间。

\begin{lemma}
设 $T,S$ 分别是阶数 $\le n$、$\le m$ 的两个微分算子。则 $T\circ S$ 是阶数 $\le n+m$ 的微分算子。
\end{lemma}

\begin{proof}
对 $n+m$ 归纳证明。若 $n=m=0$，则 $T,S\in\End_A(A)$，故 $T\circ S\in\End_A(A)$，是阶数 $0$ 的微分算子。

现设 $n+m\le p$。则
$$[T\circ S,a]=TSa-aTS=T[S,a]+[T,a]S,$$
且 $[T,a]$ 和 $[S,a]$ 分别是阶数 $\le n-1$ 和 $\le m-1$ 的微分算子。由归纳假设，此微分算子阶数 $\le n+m-1$。因此 $T\circ S$ 阶数 $\le n+m$。
\end{proof}

因此 $\Diff_k(A)$ 是 $\End_k(A)$ 的子代数。我们称它为 $A$ 上所有 $k$-线性微分算子的代数。此外，令 $n<0$ 时 $F_n\Diff_k(A)=\{0\}$，
$$F_n\Diff_k(A)=\{T\in\Diff_k(A)\mid \operatorname{order}(T)\le n\}\quad(n\ge 0).$$
显然这是 $\Diff_k(A)$ 关于 $k$-向量子空间的递增穷尽滤过。此滤过与 $\Diff_k(A)$ 的环结构相容，即对任意 $n,m\in\mathbb{Z}$，
$$F_n\Diff_k(A)\circ F_m\Diff_k(A)\subset F_{n+m}\Diff_k(A).$$

\begin{lemma}\ 
\begin{enumerate}[label=(\roman*)]
  \item $F_0\Diff_k(A)=A$；
  \item $F_1\Diff_k(A)=\Der_k(A)\oplus A$；
  \item 对任意 $n,m\in\mathbb{Z}_+$，$[F_n\Diff_k(A),F_m\Diff_k(A)]\subset F_{n+m-1}\Diff_k(A)$。
\end{enumerate}
\end{lemma}

\begin{proof}
(i) 显然。

(ii) 如前所述，$\Der_k(A)\subset F_1\Diff_k(A)$。此外，对任意 $T\in\Der_k(A)$ 有 $T(1)=T(1\cdot 1)=2T(1)$，故 $T(1)=0$。这意味着 $\Der_k(A)\cap A=0$。

设 $S\in F_1\Diff_k(A)$，$T=S-S(1)$。则 $T(1)=0$，故 $T(a)=[T,a](1)$，且
$$T(ab)=[T,ab](1)=([T,a]b)(1)+(a[T,b])(1)=(b[T,a])(1)+(a[T,b])(1)=T(a)b+aT(b),$$
即 $T\in\Der_k(A)$。

(iii) 设 $T,S$ 分别是阶数 $\le n$、$\le m$ 的算子。我们断言 $[T,S]$ 阶数 $\le n+m-1$。对 $n+m$ 归纳证明。若 $n=m=0$，无需证明。一般地，由 Jacobi 恒等式，
$$[[T,S],a]=[[T,a],S]+[T,[S,a]],$$
其中 $[T,a]$ 和 $[S,a]$ 分别是阶数 $\le n-1$ 和 $\le m-1$ 的算子。故由归纳假设，$[[T,S],a]$ 阶数 $\le n+m-2$，$[T,S]$ 阶数 $\le n+m-1$。
\end{proof}

这意味着分次环 $\Gr\Diff_k(A)$ 是交换 $A$-代数。此外，$\Diff_k(A)$ 满足 §3 的性质 (i)--(v)。

设 $n\ge 1$。设 $T$ 是 $A$ 上阶数 $\le n$ 的微分算子。则可定义从 $A^n$ 到 $\Diff_k(A)$ 的映射
$$\sigma_n(T)(a_1,a_2,\dots,a_{n-1},a_n)=[\cdots[[T,a_1],a_2],\dots,a_{n-1}],a_n].$$
由于 $\sigma_n(T)(a_1,\dots,a_n)$ 阶数 $\le 0$，可把此映射看作从 $A^n$ 到 $A$ 的映射。

\begin{lemma}
设 $T$ 是 $A$ 上阶数 $\le n$ 的微分算子。则：
\begin{enumerate}[label=(\roman*)]
  \item 映射 $\sigma_n(T):A^n\to A$ 是对称 $k$-多重线性映射；
  \item 算子 $T$ 阶数 $\le n-1$ 当且仅当 $\sigma_n(T)=0$。
\end{enumerate}
\end{lemma}

\begin{proof}
(i) 只需检验对称性。为此注意，由 Jacobi 恒等式，对任意 $S\in\Diff_k(A)$ 和 $a,b\in A$ 有
$$[[S,a],b]=[[S,b],a].$$
这意味着
\begin{align*}
&\sigma_n(T)(a_1,\dots,a_i,a_{i+1},\dots,a_n)\\
&=[\cdots[[\cdots[[T,a_1],a_2],\dots,a_i],a_{i+1}]\dots,a_{n-1}],a_n]\\
&=[\cdots[[\cdots[[T,a_1],a_2],\dots,a_{i+1}],a_i]\dots,a_{n-1}],a_n]\\
&=\sigma_n(T)(a_1,\dots,a_{i+1},a_i,\dots,a_n),
\end{align*}
故 $\sigma(T)$ 是对称的。

(ii) 显然。
\end{proof}

现在讨论一个特例。设 $A=k[X_1,X_2,\dots,X_n]$。则记 $D(n)=\Diff_k(A)$，称 $D(n)$ 为 $k^n$ 上所有微分算子的代数。设 $\partial_1,\partial_2,\dots,\partial_n$ 是 $k[X_1,X_2,\dots,X_n]$ 的标准导子。对 $I,J\in\mathbb{Z}_+^n$ 令
$$X^I=X_1^{i_1}X_2^{i_2}\cdots X_n^{i_n},\qquad \partial^J=\partial_1^{j_1}\partial_2^{j_2}\cdots\partial_n^{j_n}.$$
则 $X^I\partial^J\in D(n)$，且它是阶数 $\le |J|=j_1+j_2+\cdots+j_n$ 的微分算子。此外，若 $T$ 是由
$$T=\sum_{|I|\le p}P_I(X_1,X_2,\dots,X_n)\partial^I$$
给出的微分算子，其中 $P_I\in k[X_1,X_2,\dots,X_n]$，则 $T$ 阶数 $\le p$。

\begin{lemma}
导子 $(\partial_i;\ 1\le i\le n)$ 构成自由 $k[X_1,\dots,X_n]$-模 $\Der_k(k[X_1,X_2,\dots,X_n])$ 的一组基。
\end{lemma}

\begin{proof}
设 $T\in\Der_k(k[X_1,X_2,\dots,X_n])$。令 $P_i=T(X_i)$（$1\le i\le n$），定义 $S=\sum_{i=1}^n P_i\partial_i$。显然对所有 $1\le i\le n$，
$$S(X_i)=\sum_{j=1}^n P_j\partial_j(X_i)=P_i=T(X_i).$$
由于 $X_1,X_2,\dots,X_n$ 作为 $k$-代数生成 $k[X_1,X_2,\dots,X_n]$，可得 $T=S$。因此 $(\partial_i;\ 1\le i\le n)$ 生成 $k[X_1,X_2,\dots,X_n]$-模 $\Der_k(k[X_1,X_2,\dots,X_n])$。若 $\sum_{i=1}^n Q_i\partial_i=0$，其中 $Q_i\in k[X_1,X_2,\dots,X_n]$，则对所有 $1\le i\le n$ 有 $0=(\sum_j Q_j\partial_j)(X_i)=Q_i$。这说明 $\partial_i$（$1\le i\le n$）是自由生成元。
\end{proof}

设 $T$ 是 $k[X_1,X_2,\dots,X_n]$ 上阶数 $\le p$ 的微分算子。若 $p<0$，$T=0$，令 $\Symb_p(T)=0$。若 $p=0$，$T\in A$，令 $\Symb_0(T)=T$。对 $p\ge 1$，按如下方式定义 $k[X_1,\dots,X_n,\xi_1,\dots,\xi_n]$ 中的多项式 $\Symb_p(T)$。设 $(\xi_1,\dots,\xi_n)\in k^n$。定义线性多项式 $\ell_\xi=\sum_{i=1}^n\xi_i X_i\in k[X_1,X_2,\dots,X_n]$，以及函数
$$(\xi_1,\xi_2,\dots,\xi_n)\longmapsto \frac{1}{p!}\sigma_p(T)(\ell_\xi,\ell_\xi,\dots,\ell_\xi)$$
在 $k^n$ 上取值于 $k[X_1,X_2,\dots,X_n]$。显然可把此函数看作关于 $X_1,\dots,X_n,\xi_1,\dots,\xi_n$ 的多项式，关于 $\xi_1,\dots,\xi_n$ 齐次 $p$ 次，记为 $\Symb_p(T)$。多项式 $\Symb_p(T)$ 称为微分算子 $T$ 的 \emph{$p$-主符号}。由定义，当 $T$ 阶数 $<p$ 时 $\Symb_p(T)$ 为零。因此对 $p\ge 0$，它诱导 $\Gr_p D(n)$ 到 $k[X_1,\dots,X_n,\xi_1,\dots,\xi_n]$ 的 $k$-线性映射。记相应的 $\Gr D(n)$ 到 $k[X_1,\dots,X_n,\xi_1,\dots,\xi_n]$ 的 $k$-线性映射为 $\Symb$。

\begin{theorem}
映射 $\Symb:\Gr D(n)\longrightarrow k[X_1,X_2,\dots,X_n,\xi_1,\xi_2,\dots,\xi_n]$ 是 $k$-代数同构。
\end{theorem}

此结果的证明分几步。先证明符号映射是代数同态。

\begin{lemma}
设 $T,S\in D(n)$ 分别是阶数 $\le p$ 和 $\le q$ 的算子。则
$$\Symb_{p+q}(TS)=\Symb_p(T)\,\Symb_q(S).$$
\end{lemma}

\begin{proof}
设 $\xi\in k^n$，定义映射 $\tau_\xi:D(n)\to D(n)$ 为 $\tau_\xi(T)=[T,\ell_\xi]$。则
$$\tau_\xi(TS)=[TS,\ell_\xi]=TS\ell_\xi-\ell_\xi TS=[T,\ell_\xi]S+T[S,\ell_\xi]=\tau_\xi(T)S+T\tau_\xi(S).$$
因此对任意 $k\in\mathbb{Z}_+$，
$$\tau_\xi^k(TS)=\sum_{i=0}^k\binom{k}{i}\tau_\xi^{k-i}(T)\,\tau_\xi^i(S).$$
这意味着
\begin{align*}
\Symb_{p+q}(TS)&=\frac{1}{(p+q)!}\sigma_{p+q}(T)(\ell_\xi,\ell_\xi,\dots,\ell_\xi)=\frac{1}{(p+q)!}\tau_\xi^{p+q}(TS)\\
&=\frac{1}{p!\,q!}\tau_\xi^p(T)\,\tau_\xi^q(S)=\Symb_p(T)\,\Symb_q(S).
\end{align*}
\end{proof}

由于 $\Symb_0(X_i)=X_i$，$\Symb_1(\partial_i)=\xi_i$（$1\le i\le n$），可见对 $X^I\partial^J$（其中 $p=|J|$）有
$$\Symb_p(X^I\partial^J)=X^I\xi^J.$$
特别地，对
$$T=\sum_{|J|\le p}P_I(X_1,X_2,\dots,X_n)\partial^I$$
（其中 $P_I\in k[X_1,X_2,\dots,X_n]$），有
$$\Symb_p(T)=\sum_{|I|=p}P_I(X_1,X_2,\dots,X_n)\xi^I.$$
因此符号同态是满射。还需证明符号映射是单射。

\begin{lemma}
设 $T\in F_p D(n)$。则 $\Symb_p(T)=0$ 当且仅当 $T$ 阶数 $\le p-1$。
\end{lemma}

\begin{proof}
对 $p$ 归纳证明。$p=0$ 时显然。故设 $p>0$。设 $\xi\in k^n$，定义映射 $\tau_\xi:D(n)\to D(n)$ 为 $\tau_\xi(T)=[T,\ell_\xi]$。则对任意 $\lambda\in k$ 和 $\eta\in k^n$，
$$\tau_{\xi+\lambda\eta}(T)=[T,\ell_{\xi+\lambda\eta}]=[T,\ell_\xi]+\lambda[T,\ell_\eta]=\tau_\xi(T)+\lambda\tau_\eta(T).$$
由于 $\tau_\xi$ 和 $\tau_\eta$ 可交换，可见对任意 $k\in\mathbb{Z}_+$，
$$\tau_{\xi+\lambda\eta}^k(T)=\sum_{i=0}^k\binom{k}{i}\lambda^i\tau_\xi^{k-i}(\tau_\eta^i(T)).$$
由假设，对任意 $\lambda\in k$ 有 $\tau_{\xi+\lambda\eta}^p(T)=0$。因此由于域 $k$ 是无限的，对 $0\le i\le p$ 有 $\tau_\xi^{p-i}(\tau_\eta^i(T))=0$。特别地，对任意 $\xi,\eta\in k^n$ 有 $\tau_\xi^{p-1}(\tau_\eta(T))=0$。这意味着对任意 $\eta\in k^n$ 有 $\Symb_{p-1}([T,\ell_\eta])=0$，特别地
$$\Symb_{p-1}([T,X_i])=0$$
对 $1\le i\le n$ 成立。由归纳假设，$[T,X_i]$（$1\le i\le n$）阶数 $\le p-2$。设 $P,Q\in k[X_1,X_2,\dots,X_n]$。则
$$[T,PQ]=TPQ-PQT=[T,P]Q+P[T,Q],$$
故 $[T,PQ]$ 的阶数不超过 $[T,P]$ 和 $[T,Q]$ 阶数的最大值。由于 $X_i$（$1\le i\le n$）生成 $k[X_1,X_2,\dots,X_n]$，我们断定对任意多项式 $P$，$[T,P]$ 的阶数 $\le p-2$。这意味着 $T$ 的阶数 $\le p-1$。
\end{proof}

这也结束了 5.6 的证明。特别地，我们看到 $D(n)$ 满足 §3 的性质 (i)--(vii)。由 3.5 立即推出如下结果。

\begin{theorem}
环 $D(n)$ 是右诺特且左诺特的。
\end{theorem}

\begin{corollary}
$(X^I\partial^J;\ I,J\in\mathbb{Z}_+^n)$ 构成 $D(n)$ 作为 $k$-向量空间的一组基。
\end{corollary}

\begin{proof}
若 $|J|=p$，$X^I\partial^J$ 的 $p$-主符号等于 $X^I\xi^J$，而 $(X^I\xi^J;\ I,J\in\mathbb{Z}_+^n)$ 构成 $k[X_1,\dots,X_n,\xi_1,\dots,\xi_n]$ 作为 $k$-向量空间的一组基。
\end{proof}

$D(n)$ 的如下刻画经常有用。

\begin{theorem}
$k$-代数 $D(n)$ 是由 $X_1,X_2,\dots,X_n$ 和 $\partial_1,\partial_2,\dots,\partial_n$ 生成、满足定义关系 $[X_i,X_j]=0$、$[\partial_i,\partial_j]=0$ 以及 $[\partial_i,X_j]=\delta_{ij}$（对所有 $1\le i,j\le n$）的 $k$-代数。
\end{theorem}

\begin{proof}
设 $B$ 是由 $X_1,X_2,\dots,X_n$ 和 $\partial_1,\partial_2,\dots,\partial_n$ 生成、满足定义关系 $[X_i,X_j]=0$、$[\partial_i,\partial_j]=0$ 以及 $[\partial_i,X_j]=\delta_{ij}$（对所有 $1\le i,j\le n$）的 $k$-代数。由于这些关系在 $D(n)$ 中成立，且 $D(n)$ 由 $X_1,X_2,\dots,X_n$ 和 $\partial_1,\partial_2,\dots,\partial_n$ 生成，我们断定存在唯一的从 $B$ 到 $D(n)$ 上的满射，把生成元映为相应的生成元。显然 $B$ 由 $(X^I\partial^J;\ I,J\in\mathbb{Z}_+^n)$ 张成。因此由 5.10，此同态也是单射。
\end{proof}

\begin{proposition}
$D(n)$ 的中心等于 $k\cdot 1$。
\end{proposition}

\begin{proof}
设 $T$ 是 $D(n)$ 的中心元。则对任意多项式 $P$ 有 $[T,P]=0$，且 $T$ 阶数 $\le 0$。因此由 5.3，$T\in k[X_1,X_2,\dots,X_n]$。另一方面，$0=[\partial_i,T]=\partial_i(T)$（$1\le i\le n$）。这意味着 $T$ 是常数多项式。
\end{proof}

设 $D(n)^\circ$ 是 $D(n)$ 的反代数。则由 5.11，存在唯一的同构 $\varphi:D(n)^\circ\to D(n)$，定义为 $\varphi(X_i)=X_i$，$\varphi(\partial_i)=-\partial_i$（$1\le i\le n$）。此同态 $\varphi$ 称为 $D(n)$ 的\emph{主反自同构}。这证明了如下结果。

\begin{proposition}
代数 $D(n)^\circ$ 同构于 $D(n)$。
\end{proposition}

此外，由 5.11 可定义 $D(n)$ 的自同构 $F$：$F(X_i)=\partial_i$，$F(\partial_i)=-X_i$（$1\le i\le n$）。此自同构称为 $D(n)$ 的\emph{Fourier 自同构}。$F$ 的平方 $F^2$ 是 $D(n)$ 的自同构 $\iota$，作用为 $\iota(X_i)=-X_i$，$\iota(\partial_i)=-\partial_i$（$1\le i\le n$）。显然 $\iota^2=1$。

与按微分算子阶数的滤过不同，$D(n)$ 还有另一个与其环结构相容的滤过，它不能在更一般的微分算子环上定义。对 $p\in\mathbb{Z}$ 令
$$D_p(n)=\Bigl\{\sum_{|I|+|J|\le p}a_{IJ}X^I\partial^J\Bigr\}.$$
显然 $(D_p(n)\mid p\in\mathbb{Z})$ 是 $D(n)$ 关于 $k$ 上有限维向量空间的递增穷尽滤过。

\begin{lemma}
对任意 $p,q\in\mathbb{Z}$ 有
\begin{enumerate}[label=(\roman*)]
  \item $D_p(n)\circ D_q(n)\subset D_{p+q}(n)$；
  \item $[D_p(n),D_q(n)]\subset D_{p+q-2}(n)$。
\end{enumerate}
\end{lemma}

\begin{proof}
由 5.10 和滤过 $(D_p(n);\ p\in\mathbb{Z})$ 的定义，只需检验
$$[\partial^I,X^J]\in D_{|I|+|J|-2}(n).$$
对 $|I|$ 归纳证明。若 $|I|=1$，则对某 $1\le i\le n$ 有 $\partial^I=\partial_i$，且 $[\partial_i,X^J]=\partial_i(X^J)\in D_{|J|-1}(n)$。若 $|I|>1$，可写 $\partial^I=\partial^{I'}\partial_i$，其中 $I'\in\mathbb{Z}_+^n$，$1\le i\le n$。这导出
\begin{align*}
[\partial^I,X^J]&=[\partial^{I'}\partial_i,X^J]=\partial^{I'}\partial_i X^J-X^J\partial^{I'}\partial_i\\
&=\partial^{I'}[\partial_i,X^J]+[\partial^{I'},X^J]\partial_i\\
&=[\partial^{I'},[\partial_i,X^J]]+[\partial_i,X^J]\partial^{I'}+[\partial^{I'},X^J]\partial_i,
\end{align*}
故由归纳假设 $[\partial^I,X^J]\in D_{|I|+|J|-2}(n)$。
\end{proof}

这意味着 $(D_p(n);\ p\in\mathbb{Z})$ 是与 $D(n)$ 上环结构相容的滤过。此外，分次环 $\Gr D(n)$ 是交换 $k$-代数。若定义线性映射 $\Psi_p$（从 $D_p(n)$ 到 $k[X_1,\dots,X_n,\xi_1,\dots,\xi_n]$）为
$$\Psi_p\Bigl(\sum_{|I|+|J|\le p}a_{IJ}X^I\partial^J\Bigr)=\sum_{|I|+|J|=p}a_{IJ}X^I\xi^J,$$
则它是 $\Gr_p D(n)$ 到 $p$ 次齐次多项式的线性同构。因此它扩充为线性同构
$$\Psi:\Gr D(n)\longrightarrow k[X_1,X_2,\dots,X_n,\xi_1,\xi_2,\dots,\xi_n].$$
由 5.14 可见此映射是 $k$-代数同构。因此，配置滤过 $(D_p(n);\ p\in\mathbb{Z})$ 的环 $D(n)$ 满足 §3 的性质 (i)--(vii)。滤过 $(D_p(n);\ p\in\mathbb{Z})$ 称为 $D(n)$ 的\emph{Bernstein 滤过}。

显然，$D(n)$ 的主反自同构和 Fourier 自同构保持 Bernstein 滤过。
\begin{annote}
文中的 $\mathcal D(n)=k[x_1,\dots,x_n]\langle\partial_1,\dots,\partial_n\rangle$ 即 $n$ 阶 Weyl 代数 $A_n(k)$，也是 $A^n$ 上的微分算子环。一般簇 $X$ 上 $\mathcal D_X$ 由 $\mathcal O_X$ 与 $\mathcal T_X$ 生成、满足 Leibniz 关系，其滤过由“微分阶”给出；光滑点处 $\mathrm{gr}\mathcal D_X\cong\mathrm{Sym}\,\mathcal T_X$，所以 §7 的特征簇天然落在余切丛 $T^*X$ 里。正特征下 $\mathcal D_X$ 不再诺特，$\mathcal O_X$-凝聚也不等于 $\mathcal D_X$-凝聚，故全书默认 $k$ 特征零。[来源：nLab 条目 D-module（Last revised 2024-07-27）；原书 §5, p23–32]
\end{annote}



\section{多项式系数微分算子环上的模}

本节研究多项式系数微分算子环 $D(n)$ 上的模范畴。记 $\mathcal{M}_L(D(n))$ 和 $\mathcal{M}_R(D(n))$ 分别为左、右 $D(n)$-模的范畴。它们是 Abel 范畴。$D(n)$ 的主反自同构 $\varphi$ 定义一个从右模范畴 $\mathcal{M}_R(D(n))$ 到左模范畴 $\mathcal{M}_L(D(n))$ 的正合函子，它把模 $M$ 映为其\emph{转置} $M^t$——作为加法群等于 $M$，而 $D(n)$ 的作用由 $(T,m)\longmapsto\varphi(T)m$（$T\in D(n)$，$m\in M$）给出。类似地定义从 $\mathcal{M}_L(D(n))$ 到 $\mathcal{M}_R(D(n))$ 的函子。显然这些函子互为范畴的同构逆。若记 $\mathcal{M}_{fg}^L(D(n))$ 和 $\mathcal{M}_{fg}^R(D(n))$ 为相应的有限生成模全子范畴，则这些函子也诱导它们的等价。因此在以下讨论中可仅限于左模，并在记号中省略上标 $L$（除非想强调处理的是右模）。由于 $D(n)$ 是诺特环，全子范畴 $\mathcal{M}_{fg}(D(n))$ 在取子模、商模和扩张下封闭。

首先把 $D(n)$ 看作配置 Bernstein 滤过的环。由于此时 $D_0(n)=k$，可以用 $k$ 上有限维向量空间范畴上的可加函数 $\dim_k$ 定义 $\mathcal{M}_{fg}^L(D(n))$ 和 $\mathcal{M}_{fg}^R(D(n))$ 中模的维数。模 $M$ 的这个维数 $d(M)$ 及相应重数 $e(M)$ 分别称为 \emph{Bernstein 维数}和\emph{Bernstein 重数}。由于主反自同构保持 Bernstein 滤过，对任意有限生成 $D(n)$-模 $M$ 有 $d(M)=d(M^t)$。

对任意有限生成 $D(n)$-模 $M$ 有正合列 $D(n)^p\longrightarrow M\longrightarrow 0$，故 $d(M)\le d(D(n))$。此外，由 5.6 得出如下结果。

\begin{lemma}
对任意有限生成 $D(n)$-模 $M$ 有 $d(M)\le 2n$。
\end{lemma}

\begin{example}
考虑单变量的多项式微分算子代数 $D(1)$。设 $M$ 是非零有限生成 $D(1)$-模。则其 Bernstein 维数 $d(M)$ 可为 $0$、$1$ 或 $2$。显然 $d(M)=0$ 会蕴含：对 $M$ 的任意好滤过 $F M$，函数 $p\longmapsto\dim F_p M$ 对充分大的 $p\in\mathbb{Z}$ 为常数。由于 $F M$ 是穷尽的，这意味着 $M$ 有限维。记 $\pi(x)$ 和 $\pi(\partial)$ 为由 $x$ 和 $\partial$ 的作用诱导的 $M$ 上的线性变换。则 $[\pi(x),\pi(\partial)]=1_M$。对等式两边取迹可得 $\dim_k M=0$，与 $M\neq 0$ 的假设矛盾。故 $d(M)$ 只能是 $1$ 或 $2$。
\end{example}

$D(n)$ 维数理论的主要结果是把上述例子推广的下述陈述。

\begin{theorem}[Bernstein]
设 $M$ 是有限生成 $D(n)$-模且 $M\neq 0$。则 $d(M)\ge n$。
\end{theorem}

\begin{proof}
由于 $M$ 是有限生成 $D(n)$-模，由 3.4 可给它配置好滤过。此外，通过平移指标，显然可设 $n<0$ 时 $F_n M=0$ 且 $F_0 M\neq 0$。

对任意 $p\in\mathbb{Z}_+$，考虑线性映射 $D_p(n)\longrightarrow\Hom_k(F_p M,F_{2p}M)$，它把 $T\in D_p(n)$ 映为线性映射 $m\longmapsto Tm$。我们断言此映射是单射。$p\le 0$ 时显然。设它对 $p-1$ 成立，且 $T\in D_p(n)$ 满足对所有 $m\in F_p M$ 有 $Tm=0$。则对任意 $v\in F_{p-1}M$ 和 $1\le i\le n$，有 $X_i v\in F_p M$ 和 $\partial_i v\in F_p M$，故
$$[X_i,T]v=X_iTv-TX_iv=0$$
以及
$$[\partial_i,T]v=\partial_i T v-T\partial_i v=0,$$
且由 5.14 有 $[X_i,T],[\partial_i,T]\in D_{p-1}(n)$。由归纳假设这意味着对 $1\le i\le n$ 有 $[X_i,T]=0$、$[\partial_i,T]=0$，即 $T$ 在 $D(n)$ 的中心中。由于 $D(n)$ 的中心由 5.12 等于 $k$，我们断定 $T=0$。因此对任意 $p\in\mathbb{Z}$，
$$\dim_k\bigl(D_p(n)\bigr)\le\dim_k\bigl(\Hom_k(F_p M,F_{2p}M)\bigr)=\dim_k(F_p M)\cdot\dim_k(F_{2p}M).$$
另一方面，对充分大的 $p\in\mathbb{Z}_+$，左边等于 $p$ 的 $2n$ 次多项式（首项系数为正），右边等于 $p$ 的 $2d(M)$ 次多项式（首项系数为正）。这只有在 $d(M)\ge n$ 时才可能。
\end{proof}

下一节将给出 Bernstein 维数的几何解释。

最后，若 $M$ 是 $D(n)$-模，可定义其 Fourier 变换 $F(M)$ 为这样的模：作为加法群等于 $M$，$D(n)$ 的作用由 $(T,m)\longmapsto F(T)m$（$T\in D(n)$，$m\in M$）给出。显然 Fourier 变换是范畴 $\mathcal{M}(D(n))$ 的自同构。它也诱导 $\mathcal{M}_{fg}(D(n))$ 的自同构。由 Fourier 自同构 $F$ 保持 Bernstein 滤过（或 3.9）这一事实得出如下结果。

\begin{lemma}
设 $M$ 是有限生成 $D(n)$-模。则 $d(F(M))=d(M)$。
\end{lemma}

\section{特征簇}

现在研究有限生成 $D(n)$-模的一个更具几何味道的不变量。特别地，它将用 $D(n)$ 按微分算子次数的滤过 $F D(n)$ 来构造，而非 Bernstein 滤过。与 Bernstein 滤过不同，次数滤过对任意光滑仿射代数簇上的微分算子环都有意义。

首先，由于任何 $D(n)$-模 $M$ 都可看作 $k[X_1,X_2,\dots,X_n]$-模，可考虑其支集 $\operatorname{supp}(M)\subset k^n$。

\begin{proposition}
设 $M$ 是有限生成 $D(n)$-模。则 $\operatorname{supp}(M)$ 是 $k^n$ 的闭子簇。
\end{proposition}

\begin{proof}
固定 $M$ 上的好滤过 $F M$。则对 $x\in k^n$，$M_x=0$ 等价于对所有 $p\in\mathbb{Z}$ 有 $(F_p M)_x=0$。因此由局部化的正合性，等价于 $(\Gr M)_x=0$。设 $I_p$ 是 $k[X_1,X_2,\dots,X_n]$-模 $\Gr_p M$（$p\in\mathbb{Z}$）的零化子。由于 $\Gr_p M$ 是有限生成 $k[X_1,X_2,\dots,X_n]$-模，由 4.2 其支集 $\operatorname{supp}(\Gr_p M)=V(I_p)$。这意味着 $\operatorname{supp}(M)=\bigcup_{p\in\mathbb{Z}}V(I_p)$。设 $m_1,m_2,\dots,m_s$ 是 $\Gr D(n)$-模 $\Gr M$ 的一组齐次生成元。则 $m_1,m_2,\dots,m_s$ 在 $k[X_1,X_2,\dots,X_n]$ 中的零化子 $I$ 零化整个 $\Gr M$。因此存在 $\mathbb{Z}$ 的有限子集 $S$ 使得 $\bigcap_{p\in S}I_p=I\subset I_q$ 对所有 $q\in\mathbb{Z}$ 成立。这意味着 $\bigcup_{p\in S}V(I_p)=V(I)\supset V(I_q)$ 对所有 $q\in\mathbb{Z}$ 成立，且 $\operatorname{supp}(M)=V(I)$。
\end{proof}

设 $D$ 是配置满足 §3 开头性质 (i)--(vii) 的滤过 $F D$ 的滤过环。设 $M$ 是有限生成 $D$-模，$F M$ 是 $M$ 的好滤过。则 $\Gr M$ 是分次 $\Gr D$-模。设 $I$ 是 $\Gr M$ 在 $\Gr D$ 中的零化子。这显然是 $\Gr D$ 中的分次理想，故其根 $r(I)$ 也是分次理想。一般地 $I$ 依赖于 $M$ 上好滤过的选择，但也有如下结果。

\begin{lemma}
设 $M$ 是有限生成 $D$-模，$F M$ 和 $F'M$ 是 $M$ 上两个好滤过。设 $I$、$I'$ 分别是相应分次 $\Gr D$-模 $\Gr M$、$\Gr' M$ 的零化子。则 $r(I)=r(I')$。
\end{lemma}

\begin{proof}
设 $T\in r(I)\cap\Gr_p D$。则存在 $s\in\mathbb{Z}_+$ 使得 $T^s\in I$。若取 $Y\in F_p D$ 使得 $Y+F_{p-1}D=T$，则对所有 $q\in\mathbb{Z}$ 有 $Y^s F_q M\subset F_{q+sp-1}M$。故由归纳得
$$Y^{ms}F_q M\subset F_{q+msp-m}M$$
对所有 $m\in\mathbb{N}$ 和 $q\in\mathbb{Z}$ 成立。另一方面，由 3.7 知 $F M$ 与 $F'M$ 等价。故存在 $l\in\mathbb{Z}_+$ 使得对所有 $q\in\mathbb{Z}$ 有 $F_q M\subset F'_{q+l}M\subset F_{q+2l}M$。这导出
$$Y^{ms}F'_q M\subset Y^{ms}F_{q+l}M\subset F_{q+l+msp-m}M\subset F'_{q+2l+msp-m}M$$
对所有 $q\in\mathbb{Z}$ 和 $m\in\mathbb{N}$ 成立。若取 $m>2l$，则对所有 $q\in\mathbb{Z}$ 有 $Y^{ms}F'_q M\subset F'_{q+msp-1}M$，即 $T^{ms}\in I'$。因此 $T\in r(I')$，有 $r(I)\subset r(I')$。由于 $I$ 和 $I'$ 的角色对称，断定 $r(I)=r(I')$。
\end{proof}

因此 $\Gr M$ 零化子的根与 $\Gr M$ 上好滤过的选择无关。我们称之为 $M$ 的\emph{特征理想}，记为 $J(M)$。

现在可把此构造应用于 $D(n)$。由于由 5.6 有 $\Gr D(n)=k[X_1,\dots,X_n,\xi_1,\dots,\xi_n]$，可定义闭代数集
$$\operatorname{Ch}(M)=V(J(M))\subset k^{2n},$$
称之为 $M$ 的\emph{特征簇}。

由于 $J(M)$ 关于后 $n$ 个变量是齐次理想，立即得到如下结果。

\begin{lemma}
有限生成 $D(n)$-模 $M$ 的特征簇 $\operatorname{Ch}(M)$ 具有如下性质：若 $(x,\xi)\in\operatorname{Ch}(M)$，则对任意 $\lambda\in k$ 有 $(x,\lambda\xi)\in\operatorname{Ch}(M)$。
\end{lemma}

我们称 $\operatorname{Ch}(M)$ 是\emph{锥簇}。

\begin{proposition}
设
$$0\longrightarrow M'\longrightarrow M\longrightarrow M''\longrightarrow 0$$
是有限生成 $D(n)$-模的正合列。则
$$\operatorname{Ch}(M)=\operatorname{Ch}(M')\cup\operatorname{Ch}(M'').$$
\end{proposition}

\begin{proof}
设 $F M$ 是 $M$ 上的好滤过。则它在 $M'$ 上诱导滤过 $F M'$，在 $M''$ 上诱导滤过 $F M''$。由 3.8 知这些滤过也是好的。此外有有限生成 $k[X_1,\dots,X_n,\xi_1,\dots,\xi_n]$-模的正合列
$$0\longrightarrow\Gr M'\longrightarrow\Gr M\longrightarrow\Gr M''\longrightarrow 0,$$
由 4.2 它们的支集分别是 $D(n)$-模 $M$、$M'$、$M''$ 的特征簇。因此结论由 4.1 得出。
\end{proof}

接下来两个结果阐明特征簇与有限生成 $D(n)$-模支集之间的关系。设 $\pi:k^{2n}\longrightarrow k^n$ 是由 $\pi(x,\xi)=x$（对 $x,\xi\in k^n$）定义的映射。

\begin{proposition}
设 $M$ 是有限生成 $D(n)$-模。则 $\operatorname{supp}(M)=\pi(\operatorname{Ch}(M))$。
\end{proposition}

\begin{proof}
记 $m_1,m_2,\dots,m_s$ 为 $\Gr M$ 的一组齐次生成元。则如 7.1 的证明，$m_1,m_2,\dots,m_s$ 在 $k[X_1,X_2,\dots,X_n]$ 中的零化子 $I$ 满足 $\operatorname{supp}(M)=V(I)$。另一方面，若 $J$ 是 $m_1,m_2,\dots,m_s$ 在 $k[X_1,X_2,\dots,X_n,\xi_1,\xi_2,\dots,\xi_n]$ 中的零化子，它是关于 $\xi_1,\xi_2,\dots,\xi_n$ 的齐次理想，满足 $I=k[X_1,X_2,\dots,X_n]\cap J$，且 $\operatorname{Ch}(M)=V(J)$。这意味着 $x\in V(I)=\operatorname{supp}(M)$ 等价于 $(x,0)\in V(J)=\operatorname{Ch}(M)$。由于 $\operatorname{Ch}(M)$ 是锥的，这蕴含结论。
\end{proof}

设 $M$ 是有限生成 $D(n)$-模。定义 $M$ 的\emph{奇异支集}为
$$\operatorname{sing\,supp}(M)=\{x\in k^n\mid (x,\xi)\in\operatorname{Ch}(M)\ \text{对某个 }\xi\neq 0\}.$$
显然有 $\operatorname{sing\,supp}(M)\subset\operatorname{supp}(M)$。

\begin{lemma}
设 $M$ 是有限生成 $D(n)$-模。则 $\operatorname{sing\,supp}(M)$ 是 $\operatorname{supp}(M)$ 的闭子簇。
\end{lemma}

\begin{proof}
设 $p:k^n-\{0\}\longrightarrow P^{n-1}(k)$ 为自然投射。则
$$1\times p:k^n\times(k^n-\{0\})\longrightarrow k^n\times P^{n-1}(k)$$
把 $\operatorname{Ch}(M)-(k^n\times\{0\})$ 投射到 $k^n\times P^{n-1}(k)$ 中对应于理想 $J(M)$（它关于 $\xi_1,\xi_2,\dots,\xi_n$ 齐次）的闭子簇上。最后，到第一个因子的投射 $k^n\times P^{n-1}(k)\longrightarrow k^n$ 把它映到 $\operatorname{sing\,supp}(M)$ 上。由于 $P^{n-1}(k)$ 是完全簇，投射 $k^n\times P^{n-1}(k)\longrightarrow k^n$ 是闭映射。因此 $\operatorname{sing\,supp}(M)$ 是闭的。
\end{proof}

关于特征簇的基本结果如下定理，它也给出 Bernstein 维数的几何描述。

\begin{theorem}
设 $M$ 是有限生成 $D(n)$-模。则
$$\dim\operatorname{Ch}(M)=d(M).$$
\end{theorem}

证明此定理需要一些准备。先考虑 $D(n)$-模 $M=D(n)/L$，其中 $L$ 是 $D(n)$ 中的左理想。$D(n)$-模的正合列
$$0\longrightarrow L\longrightarrow D(n)\longrightarrow M\longrightarrow 0$$
（对 $D(n)$ 配置由微分算子次数滤过诱导的滤过）导出 $A$-模的正合列
$$0\longrightarrow\Gr L\longrightarrow\Gr D(n)\longrightarrow\Gr M\longrightarrow 0,$$
其中 $A=\Gr D(n)=k[X_1,\dots,X_n,\xi_1,\dots,\xi_n]$。因此 $\Gr M$ 是商 $A/\Gr L$，$\Gr M$ 的零化子等于 $\Gr L$。故由定义，$V(\Gr L)$ 是 $M$ 的特征簇。因此在这种情况下须建立 $\dim V(\Gr L)=d(D(n)/L)$。

为证明这一点，先作一些一般性注记。设 $A=k[X_1,X_2,\dots,X_n]$。设 $\mathbf{t}=(t_1,\dots,t_n)\in\mathbb{N}^n$。通过令 $\Gr^{\mathbf{t}}_m A$ 为满足 $\sum_{j=1}^n t_j i_j=m$ 的 $X^I$ 的线性张成，定义 $\Gr^{\mathbf{t}} A$ 的分次。显然这样 $A$ 成为分次环。此外可定义相应滤过 $F^{\mathbf{t}}A$ 为 $F^{\mathbf{t}}_p A=\sum_{m\le p}\Gr^{\mathbf{t}}_m A$。显然，若以 $F A$ 记 $A$ 按多项式次数的自然滤过，并令 $t=\max_{1\le i\le n}t_i$，则对任意 $p\in\mathbb{Z}$ 有
$$F^{\mathbf{t}}_p A\subset F_p A,\qquad F_p A\subset F^{\mathbf{t}}_{tp}A.$$

设 $I$ 是 $A$ 中的理想。则可考虑配置由 $A$ 上滤过诱导的滤过的 $A$-模正合列
$$0\longrightarrow I\longrightarrow A\longrightarrow A/I\longrightarrow 0.$$
则对任意 $p\in\mathbb{Z}$ 有
$$F^{\mathbf{t}}_p(A/I)\subset F_p(A/I),\qquad F_p(A/I)\subset F^{\mathbf{t}}_{tp}(A/I).$$
这进而蕴含如下引理。

\begin{lemma}
对任意 $p\in\mathbb{Z}$，有
$$\dim_k F^{\mathbf{t}}_p(A/I)\le\dim_k F_p(A/I),\qquad \dim_k F_p(A/I)\le\dim_k F^{\mathbf{t}}_{tp}(A/I).$$
\end{lemma}

设 $s\in\mathbb{N}$。则定义代数 $D(n)$ 的滤过 $F^{(s)}D(n)$ 为
$$F^{(s)}_m D(n)=\Bigl\{T\in D(n)\ \Big|\ T=\sum_{|I|+s|J|\le m}c_{I,J}X^I\partial^J,\ c_{I,J}\in k\Bigr\}.$$
显然 $F^{(1)}D(n)$ 是 $D(n)$ 的 Bernstein 滤过。滤过 $F^{(s)}D(n)$ 具有 §3 中考虑的环滤过的性质 (i)--(iii)。此外，$T\in F^{(s)}_m D(n)$ 当且仅当 $T\in F^{(1)}_p D(n)$ 且 $T$ 的阶数 $\le q$，其中 $p,q$ 满足 $m=p+(s-1)q$。因此，若 $T\in F^{(s)}_m D(n)$、$S\in F^{(s)}_{m'}D(n)$，则存在 $p,p'$ 和 $q,q'$ 使得 $m=p+(s-1)q$、$m'=p'+(s-1)q'$，$T\in F^{(1)}_p D(n)$，$S\in F^{(1)}_{p'}D(n)$，且 $T$、$S$ 的阶数分别 $\le q$、$\le q'$。这意味着 $TS$ 的阶数 $\le q+q'$ 且 $TS\in F^{(1)}_{p+p'}D(n)$。由此可知 $TS\in F^{(s)}_{m+m'}D(n)$。故滤过 $F^{(s)}D(n)$ 也满足 (iv)，即是环滤过。同样可检验 (v) 成立，即分次环 $\Gr^{(s)}D(n)$ 是交换的。此外，分次环 $\Gr^{(s)}D(n)$ 同构于分次环 $A=k[X_1,\dots,X_n,\xi_1,\dots,\xi_n]$，其分次结构对应于 $\mathbf{t}=(1,\dots,1,s,\dots,s)$。我们记该分次模为 $\Gr^{(s)}A$，其相应滤过为 $F^{(s)}A$。

此外，对任意 $p\in\mathbb{Z}$ 有
$$F^{(s)}_p D(n)\subset F^{(1)}_p D(n),\qquad F^{(1)}_p D(n)\subset F^{(s)}_{sp}D(n).$$

再考虑配置由 $D(n)$ 上滤过诱导的滤过的 $D(n)$-模正合列
$$0\longrightarrow L\longrightarrow D(n)\longrightarrow D(n)/L\longrightarrow 0.$$
则对任意 $p\in\mathbb{Z}$ 有
$$F^{(s)}_p(D(n)/L)\subset F^{(1)}_p(D(n)/L),\qquad F^{(1)}_p(D(n)/L)\subset F^{(s)}_{sp}(D(n)/L).$$
这进而蕴含与 7.8 类似的如下引理。

\begin{lemma}
对任意 $p\in\mathbb{Z}$，有
$$\dim_k F^{(s)}_p(D(n)/L)\le\dim_k F^{(1)}_p(D(n)/L)$$
以及
$$\dim_k F^{(1)}_p(D(n)/L)\le\dim_k F^{(s)}_{sp}(D(n)/L).$$
\end{lemma}

\begin{lemma}
设 $L$ 是 $D(n)$ 中的左理想。则对任意 $s\in\mathbb{N}$，
$$d(D(n)/L)=\dim V(\Gr^{(s)}L).$$
\end{lemma}

\begin{proof}
正合列
$$0\longrightarrow L\longrightarrow D(n)\longrightarrow D(n)/L\longrightarrow 0$$
（其中 $D(n)$ 配置滤过 $F^{(s)}D(n)$，$L$ 和 $D(n)/L$ 分别配置诱导滤过 $F^{(s)}L$ 和 $F^{(s)}(D(n)/L)$）导出分次 $A$-模（$A=\Gr^{(s)}D(n)$）的正合列
$$0\longrightarrow\Gr^{(s)}L\longrightarrow\Gr^{(s)}D(n)\longrightarrow\Gr^{(s)}(D(n)/L)\longrightarrow 0.$$
显然作为分次环 $\Gr^{(s)}D(n)=\Gr^{(s)}A$。这意味着对任意 $p\in\mathbb{Z}$，
\begin{align*}
\dim_k F^{(s)}_p(D(n)/L)&=\sum_{q=0}^{p}\bigl(\dim_k F^{(s)}_q(D(n)/L)-\dim_k F^{(s)}_{q-1}(D(n)/L)\bigr)\\
&=\sum_{q=0}^{p}\dim_k\Gr^{(s)}_q(D(n)/L)=\sum_{q=0}^{p}\bigl(\dim_k\Gr^{(s)}_q D(n)-\dim_k\Gr^{(s)}_q L\bigr)\\
&=\sum_{q=0}^{p}\bigl(\dim_k\Gr^{(s)}_q A-\dim_k\Gr^{(s)}_q L\bigr)=\sum_{q=0}^{p}\dim_k\Gr^{(s)}_q(A/\Gr^{(s)}L)\\
&=\dim_k F^{(s)}_p(A/\Gr^{(s)}L).
\end{align*}
这进而利用 7.8 和 7.9 蕴含
$$\dim_k F^{(1)}_p(D(n)/L)\le\dim_k F^{(s)}_{sp}(D(n)/L)=\dim_k F^{(s)}_{sp}(A/\Gr^{(s)}L)\le\dim_k F_{sp}(A/\Gr^{(s)}L)$$
以及
$$\dim_k F_p(A/\Gr^{(s)}L)\le\dim_k F^{(s)}_{sp}(A/\Gr^{(s)}L)=\dim_k F^{(s)}_{sp}(D(n)/L)\le\dim_k F^{(1)}_{sp}(D(n)/L).$$
由于函数 $p\longmapsto\dim_k F^{(1)}_p(D(n)/L)$ 和 $p\longmapsto\dim_k F_p(A/\Gr^{(s)}L)$ 对充分大的 $p\in\mathbb{Z}$ 由多项式表示，这些多项式必有相等的次数。这进而蕴含 $d(D(n)/L)=d(A/\Gr^{(s)}L)$。由于 $\Gr^{(s)}L$ 是 $A/\Gr^{(s)}L$ 的零化子，结论由 4.5 得出。
\end{proof}

对任意 $s\in\mathbb{N}$，记 $\sigma^{(s)}_p(T)$ 为 $T\in F^{(s)}_p D(n)$ 在 $\Gr^{(s)}_p D(n)=A$ 中的投影。此外，对 $A$ 上由多项式次数给出的自然滤过，记 $\sigma_p$ 为把 $p$ 次多项式映为其 $p$ 次齐次分量的映射。

\begin{example}
设 $D=D(1)$，$T\in D$ 由 $T=x^3\partial+\partial^2$ 给出。则 $T$ 的次数等于 $2$，且 $\Symb_2(T)=\xi^2$。故 $\sigma_2(\Symb_2(T))=\xi^2$。

另一方面，有 $\sigma^{(1)}_4(T)=x^3\xi$，$\sigma^{(2)}_5(T)=x^3\xi$，$\sigma^{(3)}_6(T)=x^3\xi+\xi^2$，且当 $s>3$ 时 $\sigma^{(s)}_{2s}(T)=\xi^2$。
\end{example}

因此，对充分大的 $s$，$\sigma^{(s)}(T)$ 变得等于 $\sigma(\Symb(T))$。一般情形也成立，更精确地说有如下结果。

\begin{lemma}
设 $T$ 是 $D(n)$ 中阶数 $\le m$ 的微分算子，其符号 $\Symb_m(T)$ 是 $p$ 次多项式。则存在 $s_0$ 使得对 $s\ge s_0$，
$$\sigma^{(s)}_{p+(s-1)m}(T)=\sigma_p(\Symb_m(T)).$$
\end{lemma}

\begin{proof}
由假设
$$T=\sum_{|J|\le m}c_{I,J}X^I\partial^J.$$
此外可固定 $q_0$ 使得 $c_{I,J}\neq 0$ 蕴含 $|I|\le q_0$。则
$$\Symb_m(T)=\sum_{|J|=m}c_{I,J}X^I\xi^J$$
是 $p$ 次多项式，其首项为
$$\sigma_p(\Symb_m(T))=\sum_{|I|=p-m,\,|J|=m}c_{I,J}X^I\xi^J.$$
另一方面，项 $X^I\partial^J$ 属于 $F^{(s)}_{|I|+s|J|}D(n)$。设 $c_{I,J}\neq 0$。则有以下可能：

(i) $|J|=m$，$|I|=p-m$：$X^I\partial^J$ 属于 $F^{(s)}_{p+(s-1)m}D(n)$。

(ii) $|J|=m$，$|I|<p-m$：$X^I\partial^J$ 属于 $F^{(s)}_{p+(s-1)m-1}D(n)$。

(iii) $m\ge 1$，$|J|<m$，$|I|\le q_0$：$X^I\partial^J$ 属于 $F^{(s)}_{q_0+s(m-1)}D(n)$。此外
$$q_0+s(m-1)=q_0+sm-s=q_0+m-s+(s-1)m.$$
故若 $s\ge s_0=q_0+m-p+1$，则 $q_0+s(m-1)\le p+(s-1)m-1$。由此可知此时微分算子 $X^I\partial^J$ 也属于 $F^{(s)}_{p+(s-1)m-1}D(n)$。

这意味着对 $s\ge s_0$ 有
\begin{align*}
\sigma^{(s)}_{p+(s-1)m}(T)&=\sigma^{(s)}_{p+(s-1)m}\Bigl(\sum_{|I|=p-m,\,|J|=m}c_{I,J}X^I\partial^J\Bigr)\\
&=\sum_{|I|=p-m,\,|J|=m}c_{I,J}X^I\xi^J=\sigma_p(\Symb_m(T)).
\end{align*}
\end{proof}

特别地，可取有限多个 $D_i\in L$（$1\le i\le k$），使得 $\sigma(\Symb(D_i))$ 生成 $\Gr(\Gr L)$。由 7.12，存在 $s$ 使得对任意 $1\le i\le k$ 有 $\sigma(\Symb(D_i))=\sigma^{(s)}(D_i)$。因此有 $\Gr(\Gr L)\subset\Gr^{(s)}L$ 和 $V(\Gr(\Gr L))\supset V(\Gr^{(s)}L)$。这进而蕴含 $\dim V(\Gr(\Gr L))\ge\dim V(\Gr^{(s)}L)$。由 4.8 和 7.10，最终得 $\dim V(\Gr L)\ge d(D(n)/L)$。

反向不等式简单得多。首先设 $D\in F^{(1)}_p(D(n))$。则其次数 $\le p$，且 $\Symb(D)$ 是次数 $\le p$ 的多项式。因此 $F^{(1)}_p(L)$ 只含次数 $\le p$ 的微分算子，且 $\Gr F^{(1)}_p(L)\subset F_p(\Gr L)$。由于 $F^{(1)}_p(L)$ 有限维，且其按微分算子次数的滤过是 Hausdorff 且穷尽的，有
$$\dim_k F^{(1)}_p(L)=\dim_k\Gr F^{(1)}_p(L)\le\dim_k F_p(\Gr L).$$
因此，对任意 $p\in\mathbb{N}$，
\begin{align*}
\dim_k F^{(1)}_p(D(n)/L)&=\dim_k D_p(n)-\dim_k F^{(1)}_p(L)\\
&\ge\dim_k F_p A-\dim_k F_p(\Gr L)=\dim_k F_p(A/\Gr L).
\end{align*}
这进而蕴含 $d(D(n)/L)\ge d(A/\Gr L)$。由于 $\Gr L$ 是 $A/\Gr L$ 的零化子，由 4.5 得 $d(D(n)/L)\ge\dim V(\Gr L)$。

故证明了
$$d(D(n)/L)=\dim V(\Gr L).$$
为证明一般情形的定理，考虑正合列
$$0\longrightarrow M'\longrightarrow M\longrightarrow M''\longrightarrow 0,$$
其中 $M$ 有 $q$ 个生成元，$M'$ 有 $q-1$ 个生成元，$M''$ 是循环模。因此 $M'$ 同构于某左理想 $L$ 的 $D(n)/L$。

由证明的第一部分有 $d(M'')=\dim\operatorname{Ch}(M'')$。此外由归纳假设有 $d(M')=\dim\operatorname{Ch}(M')$。由 3.9 和 7.4 得
\begin{align*}
d(M)&=\max(d(M'),d(M''))=\max(\dim\operatorname{Ch}(M'),\dim\operatorname{Ch}(M''))\\
&=\dim(\operatorname{Ch}(M')\cup\operatorname{Ch}(M''))=\dim\operatorname{Ch}(M).
\end{align*}
这完成 7.7 的证明。

特别地，结合 6.3 和 7.7 得到如下结果。

\begin{theorem}
设 $M$ 是有限生成 $D(n)$-模，$M\neq 0$，$\operatorname{Ch}(M)$ 是其特征簇。则 $\dim\operatorname{Ch}(M)\ge n$。
\end{theorem}
\begin{annote}
$\mathrm{Ch}(M)$ 定义为 $k[x,\xi]$ 中 $\mathrm{gr}M$ 的零化子所定义的闭集，因此它是 $T^*A^n$ 中的锥：$\xi$ 方向齐次。这一点不依赖好滤过的选取，因为两条好滤过的 $\sqrt{\mathrm{Ann}}$ 相同。nLab 的 characteristic variety 条目只列文献、无陈述，故本条不作外部断言。[来源：nLab 条目 characteristic+variety（Last revised 2017-08-29，仅文献列表）；原书 §7, p32–38]
\end{annote}



\section{完整模}

设 $M$ 是非平凡有限生成 $D(n)$-模。则由 7.13，其特征簇 $\operatorname{Ch}(M)$ 的维数 $\ge n$。

若有限生成 $D(n)$-模的特征簇 $\operatorname{Ch}(M)$ 的维数 $\le n$，则称它是\emph{完整（holonomic）的}。因此 $M$ 是完整的，当且仅当 $M=0$ 或 $\dim\operatorname{Ch}(M)=n$。

粗略地说，完整模是具有最小可能特征簇的模。

以下是关于完整模的基本观察。

\begin{theorem}\ 
\begin{enumerate}[label=(\roman*)]
  \item 完整模具有有限长度；
  \item 完整模的子模、商模和扩张都是完整的。
\end{enumerate}
\end{theorem}

\begin{proof}
(ii) 由 3.9 立即得出。

(i) 设 $M$ 是非零完整 $D(n)$-模。则由定义其特征簇 $\operatorname{Ch}(M)$ 的维数等于 $n$。由 7.7，其 Bernstein 维数 $d(M)$ 也等于 $n$。由于 $M$ 有限生成且 $D(n)$ 是诺特环，存在 $M$ 中不同于 $M$ 的极大 $D(n)$-子模 $M'$。因此有正合列
$$0\longrightarrow M'\longrightarrow M\longrightarrow M/M'\longrightarrow 0.$$
由 (ii)，$M'$ 和 $M/M'$ 是完整的，且 $M/M'$ 是不可约 $D(n)$-模。若 $M'\neq 0$，由 3.9 得 $e(M')<e(M)$。因此由 $e(M)$ 上的归纳，$M$ 具有有限长度。
\end{proof}

因此范畴 $\mathcal{M}_{fg}(D(n))$ 的全子范畴 $\operatorname{Hol}(D(n))$ 在取子模、商模和扩张下封闭。此外，若记 $\mathcal{M}_{fl}(D(n))$ 为 $\mathcal{M}_{fg}(D(n))$ 中由有限长度 $D(n)$-模构成的全子范畴，则 $\operatorname{Hol}(D(n))$ 是 $\mathcal{M}_{fl}(D(n))$ 的子范畴。可以证明当 $n>1$ 时 $\operatorname{Hol}(D(n))$ 严格小于 $\mathcal{M}_{fl}(D(n))$。

此外，转置函子和 Fourier 函子把完整模映为完整模。

现在讨论完整模的一些例子。

\begin{example}
设 $\mathcal{O}_n=k[X_1,X_2,\dots,X_n]$。则
$$\mathcal{O}_n=D(n)/\bigl(D(n)(\partial_1,\partial_2,\dots,\partial_n)\bigr)$$
是有限生成 $D(n)$-模。此外，若令 $p<0$ 时 $F_p\mathcal{O}_n=0$，$p\ge 0$ 时 $F_p\mathcal{O}_n=\mathcal{O}_n$，则滤过 $F\mathcal{O}_n$ 是关于 $D(n)$ 次数滤过的好滤过。相应的分次模 $\Gr\mathcal{O}_n$ 满足：$p\neq 0$ 时 $\Gr_p\mathcal{O}_n=0$，$\Gr_0\mathcal{O}_n=k[X_1,X_2,\dots,X_n]$。由此可知 $\Gr\mathcal{O}_n$ 的零化子等于 $k[X_1,X_2,\dots,X_n,\xi_1,\xi_2,\dots,\xi_n]$ 中由 $\xi_1,\xi_2,\dots,\xi_n$ 生成的理想。这意味着 $\operatorname{Ch}(\mathcal{O}_n)=k^n\times\{0\}\subset k^{2n}$。特别地 $\dim\operatorname{Ch}(\mathcal{O}_n)=n$，$\mathcal{O}_n$ 是完整的。此外 $\operatorname{supp}(\mathcal{O}_n)=k^n$，投射 $\pi:k^{2n}\longrightarrow k^n$ 是 $\operatorname{Ch}(\mathcal{O}_n)$ 到 $\mathcal{O}$ 上的双射。

通过求微分可见 $\mathcal{O}_n$ 的任何子模必含常数，因此 $\mathcal{O}_n$ 是不可约的。
\end{example}

\begin{example}
现考虑 $\Delta_n=F(\mathcal{O}_n)$。则有
$$\Delta_n=D(n)/\bigl(D(n)(X_1,X_2,\dots,X_n)\bigr).$$
显然 $\Delta_n$ 是完整的且不可约的。设 $\delta$ 是对应于 $1\in\mathcal{O}_1$ 的向量。则对任意 $1\le i\le n$ 有 $X_i\delta=0$。显然 $\Delta_n$ 由 $\delta^{(I)}=\partial^I\delta$（$I\in\mathbb{Z}_+^n$）张成。设 $F\Delta_n$ 是 $\Delta_n$ 的滤过：$p<0$ 时 $F_p\Delta_n=\{0\}$，$p\ge 0$ 时 $F_p\Delta_n$ 由 $\delta^{(I)}$（$|I|\le p$）张成。记 $\epsilon_i$ 为第 $i$ 个位置为 $1$ 的多重指标 $(0,\dots,0,1,0,\dots,0)$。则由 Fourier 变换的定义，
$$\partial_j\delta^{(I)}=\delta^{(I+\epsilon_j)},\qquad X_j\delta^{(I)}=-i_j\delta^{(I-\epsilon_j)}$$
对所有 $1\le j\le n$ 和 $I\in\mathbb{Z}_+^n$ 成立。这意味着 $F_p\Delta_n$ 是 $k[X_1,X_2,\dots,X_n]$-子模。此外对所有 $1\le i\le n$ 和 $p\in\mathbb{Z}$ 有 $\partial_i F_p\Delta_n\subset F_{p+1}\Delta_n$。故 $F\Delta_n$ 是关于按微分算子次数滤过的 $D(n)$ 的穷尽 $D(n)$-模滤过。设 $\bar\delta^{(I)}$ 为 $\delta^{(I)}$ 在 $\Gr_{|I|}\Delta_n$ 中代表的陪集。则 $\Gr_p\Delta_n$ 由 $\bar\delta^{(I)}$（$I\in\mathbb{Z}_+$，$|I|=p$）张成。显然 $X_i$ 在 $\Gr\Delta_n$ 上作为 $0$ 作用，而 $\xi_i$ 把 $\bar\delta^{(I)}$ 映为 $\bar\delta^{(I+\epsilon_i)}$。因此 $\bar\delta$ 作为 $k[X_1,X_2,\dots,X_n,\xi_1,\xi_2,\dots,\xi_n]$-模生成 $\Gr\Delta_n$，且 $F\Delta_n$ 是好滤过。此外 $\Gr\Delta_n$ 的零化子是由 $X_i$（$1\le i\le n$）生成的理想。故 $\Delta_n$ 的特征簇为 $\operatorname{Ch}(\Delta_n)=\{0\}\times k^n\subset k^{2n}$，支集 $\operatorname{supp}(\Delta_n)=\{0\}\subset k^n$。
\end{example}

现在构造更多完整模。从一个简单的完整性判别法开始。

\begin{lemma}
设 $D(n)$ 配置 Bernstein 滤过。设 $M$ 是 $D(n)$-模，$F M$ 是 $M$ 上穷尽的 $D(n)$-模滤过。若
$$\dim_k F_p M\le\frac{c^n}{n!}\,p+\text{（关于 $p$ 的低阶项）}$$
对所有 $p\in\mathbb{Z}_+$ 成立，则 $M$ 是完整 $D(n)$-模且其长度 $\le c$。特别地，$M$ 是有限生成 $D(n)$-模。
\end{lemma}

\begin{proof}
设 $N$ 是 $M$ 的有限生成 $D(n)$-子模。则 $F M$ 在 $N$ 上诱导穷尽 $D(n)$-模滤过。由 3.6 存在 $N$ 的好滤过 $F'N$ 和 $s\in\mathbb{Z}_+$ 使得对任意 $p\in\mathbb{Z}$ 有 $F'_p N\subset F_{p+s}N$。由此对 $p\in\mathbb{Z}_+$ 有
$$\dim_k F'_p N\le\dim_k F_{p+s}N\le\dim_k F_{p+s}M\le\frac{c^n}{n!}\,p+\text{（关于 $p$ 的低阶项）}.$$
因此 $d(N)\le n$，$N$ 是完整的。若 $N\neq 0$，则 $e(N)\le c$。显然这意味着 $N$ 的长度 $\le e(N)\le c$。由此可知 $M$ 的有限生成 $D(n)$-子模的任意递增序列稳定，且 $M$ 本身有限生成。
\end{proof}

\begin{example}
设 $n=1$，令 $D=D(1)$。对任意 $\alpha\in k$ 考虑 $D$-模 $M_\alpha=D/D(z\partial-\alpha)$。

令 $E=z\partial$。如 5.10 的证明，算子 $(z^pE^q,\ \partial^pE^q;\ p,q\in\mathbb{Z}_+)$ 构成 $D$ 作为 $k$-线性空间的一组基。此外，理想 $D(z\partial-\alpha)$ 由元素 $(z^pE^q(E-\alpha),\ \partial^pE^q(E-\alpha);\ p,q\in\mathbb{Z}_+)$ 张成。故 $M_\alpha$ 由对应于 $(z^p,\partial^p;\ p\in\mathbb{Z}_+)$ 的陪集张成。

显然
$$[E,z]=z\partial z-z^2\partial=z,\qquad [E,\partial]=z\partial^2-\partial z\partial=-\partial.$$
因此
$$Ez=z(E+1),\qquad E\partial=\partial(E-1).$$
这立即蕴含对任意 $n\in\mathbb{Z}_+$，$z^n$ 的陪集是 $E$ 的属于特征值 $\alpha+n$ 的特征向量。另一方面，对任意 $n\in\mathbb{Z}_+$，$\partial^n$ 的陪集是 $E$ 的属于特征值 $\alpha-n$ 的特征向量。因此 $E$ 在 $M_\alpha$ 上的谱等于 $\{\alpha+n;\ n\in\mathbb{Z}\}$，且每个特征值的重数等于 $1$。

$M_\alpha$ 的 Fourier 变换同构于
$$D/D(-\partial z-\alpha)=D/D(z\partial+\alpha+1)=M_{-\alpha-1}.$$

先设 $\alpha\notin\mathbb{Z}$。则 $E$ 是线性同构，$z$ 必为满射。由于 $z$ 把属于特征值 $\alpha+n$ 的特征空间映到属于特征值 $\alpha+n+1$ 的特征空间，$z$ 也是单射。因此可归纳构造一族向量 $z^{\alpha+n}$（$n\in\mathbb{Z}$），使得 $Ez^{\alpha+n}=(\alpha+n)z^{\alpha+n}$ 且 $zz^{\alpha+n}=z^{\alpha+n+1}$。显然这些向量构成 $M_\alpha$ 的一组基。此外对任意 $n\in\mathbb{Z}$，
$$\partial z^{\alpha+n}=\partial zz^{\alpha+n-1}=[\partial,z]z^{\alpha+n-1}+Ez^{\alpha+n-1}=(\alpha+n)z^{\alpha+n-1}.$$
这立即蕴含对任意整数 $p\in\mathbb{Z}$ 有 $M_\alpha\cong M_{\alpha+p}$。

此外，$M$ 的任何非零 $D$-子模在 $E$ 下不变，故含 $E$ 的一个特征向量。这进而蕴含它含某 $z^{\alpha+p}$（$p\in\mathbb{Z}$）。由此可知它含所有 $z^{\alpha+n}$（$n\in\mathbb{Z}$），即等于 $M_\alpha$。故 $M_\alpha$ 是不可约 $D$-模。

定义 $M_\alpha$ 的滤过 $F M_\alpha$：$p<0$ 时 $F_pM_\alpha=\{0\}$；$p\ge 0$ 时 $F_pM_\alpha$ 为 $\{z^{\alpha+n};\ |n|\le p\}$ 的张成。显然 $F M_\alpha$ 是 $M_\alpha$ 关于线性子空间的递增穷尽滤过。此外由上述注记，对任意 $p\in\mathbb{Z}$ 有 $zF_pM_\alpha\subset F_{p+1}M_\alpha$ 和 $\partial F_pM_\alpha\subset F_{p+1}M_\alpha$。因此 $F M_\alpha$ 是关于配置 Bernstein 滤过的 $D$ 的 $D$-模滤过。由于 $p\ge 0$ 时 $\dim_k F_pM_\alpha=2p+1$，由 8.4 知 $M_\alpha$ 是完整的。

为计算其特征簇，考虑另一个滤过 $F M_\alpha$：$n<0$ 时 $F_nM_\alpha=\{0\}$，$n\ge 0$ 时 $F_nM_\alpha$ 由 $\{z^{\alpha+p};\ p\ge -n\}$ 张成。显然这是 $F M_\alpha$ 关于 $z$ 中多项式环的模的穷尽滤过。此外对任意 $p\in\mathbb{Z}_+$ 有 $\partial F_pM_\alpha=F_{p+1}M_\alpha$，这是关于按微分算子次数滤过的 $D$ 的好 $D$-模滤过。分次模 $\Gr M_\alpha$ 是 $\Gr_pM_\alpha$ 的直和，其中 $n<0$ 时 $\Gr_nM_\alpha=0$；$\Gr_0M_\alpha$ 等于 $z^{\alpha+p}$（$p\ge 0$）的张成；$\Gr_pM_\alpha$ 由 $z^{\alpha-p}$ 模 $\{z^{\alpha+q};\ q>-p\}$ 的陪集张成。因此 $z$ 零化 $\Gr_pM_\alpha$（$p\neq 0$），而 $\partial$ 的符号 $\xi$ 零化 $\Gr_0M_\alpha$ 并把 $\Gr_pM_\alpha$ 映到 $\Gr_{p+1}M_\alpha$（$p>0$）上。由此可知 $\Gr M_\alpha$ 的零化子是 $k[z,\xi]$ 中由 $z\xi$ 生成的理想。故特征簇 $\operatorname{Ch}(M_\alpha)$ 是 $k^2$ 中直线 $\{z=0\}$ 和 $\{\xi=0\}$ 的并。

现设 $\alpha\in\mathbb{Z}$。则 $E$ 的特征值是整数。若 $v$ 是 $E$ 的属于特征值 $m\neq 0$ 的非零特征向量，则 $\partial v$ 是 $E$ 的属于特征值 $m-1$ 的特征向量，且 $z\partial v=mv\neq 0$。因此 $z$ 把 $E$ 的属于特征值 $q\neq -1$ 的所有特征空间映到属于特征值 $q+1$ 的特征空间。

先设 $n=-\alpha>0$。则 $z^{n-1}$ 的陪集是 $E$ 的属于特征值 $-1$ 的特征向量。因此 $z$ 把 $E$ 的属于特征值 $-1$ 的特征空间映到属于特征值 $0$ 的特征空间。故此时可选取特征值 $m\in\mathbb{Z}$ 的特征空间的基向量 $v_m$，使得对 $m\in\mathbb{Z}$ 有 $zv_m=v_{m+1}$。有
$$\partial v_m=\partial z v_{m-1}=[\partial,z]v_{m-1}+Ev_{m-1}=mv_{m-1}$$
对所有 $m\in\mathbb{Z}$ 成立。这意味着所有 $M_{-n}$（$n>0$）互相同构。

此外，通过考察 $z$ 和 $\partial$ 的作用，可见向量 $v_m$（$m\in\mathbb{Z}_+$）张成一个同构于 8.2 中 $\mathcal{O}_1$ 的 $D$-子模 $N_{-n}$。特别地 $M_{-n}$ 是可约的。

显然 $zv_{-1}\in N_{-n}$。由此可知 $v_{-1}$ 在 $M_{-n}/N_{-n}$ 中的陪集 $\delta$ 满足 $z\delta=0$。

$E$ 在 $L_{-n}=M_{-n}/N_{-n}$ 上的谱由所有严格负整数构成。因此 $\partial$ 在 $L_{-n}$ 上是单射。故 $\delta^{(m)}=\partial^m\delta$ 是 $E$ 的属于特征值 $-(m+1)$（$m\in\mathbb{Z}_+$）的非零特征向量，即它们与 $v_{-(m+1)}$ 的陪集成比例。显然对所有 $m>0$ 有
$$z\delta^{(m)}=z\partial\delta^{(m-1)}=-m\delta^{(m-1)}.$$
由此可知 $L_{-n}$ 同构于 8.3 中描述的 $D$-模 $\Delta_1$。故有正合列
$$0\longrightarrow\mathcal{O}_1\longrightarrow M_{-n}\longrightarrow\Delta_1\longrightarrow 0,$$
且此正合列不分裂。特别地，所有这些 $D$-模都同构于 $M_{-1}$。

由 Fourier 变换可见 $D$-模 $M_n$（$n\ge 0$）同构于 $M_0$。此外有正合列
$$0\longrightarrow\Delta_1\longrightarrow M_n\longrightarrow\mathcal{O}_1\longrightarrow 0,$$
它也不分裂。

由于 $\mathcal{O}_1$ 和 $\Delta_1$ 由 8.2 和 8.3 是完整的，由 8.1 知 $M_n$（$n\in\mathbb{Z}$）是完整的。此外由 7.4 得
$$\operatorname{Ch}(M_n)=\operatorname{Ch}(\mathcal{O}_1)\cup\operatorname{Ch}(\Delta_1)$$
对所有 $n\in\mathbb{Z}$ 成立。故由 8.2 和 8.3，它们等于 $k^2$ 中直线 $\{z=0\}$ 和 $\{\xi=0\}$ 的并。

由上述例子可见特征簇不能决定相应的 $D$-模。此外，不可约完整 $D(n)$-模的特征簇可以是可约的。
\end{example}

现在把上述例子中 $M_{-1}$ 的构造推广。

设 $M$ 是 $D(n)$-模，$P\in k[X_1,X_2,\dots,X_n]$。则在 $M$ 的局部化 $M_P$ 上可定义 $k$-线性映射 $\partial_i:M_P\to M_P$ 为
$$\partial_i\Bigl(\frac{m}{P^p}\Bigr)=-\frac{p\,\partial_i(P)m}{P^{p+1}}+\frac{\partial_i m}{P^p}$$
对任意 $m\in M$ 和 $p\in\mathbb{Z}_+$。通过直接计算可检验
$$[\partial_i,\partial_j]\Bigl(\frac{m}{P^p}\Bigr)=0,\qquad [\partial_i,x_j]\Bigl(\frac{m}{P^p}\Bigr)=\delta_{ij}\frac{m}{P^p}$$
对任意 $1\le i,j\le n$ 和 $p\in\mathbb{Z}_+$。由 5.11 这在 $M_P$ 上定义了 $D(n)$-模结构。

\begin{proposition}
设 $M$ 是完整 $D(n)$-模，$P\in k[X_1,X_2,\dots,X_n]$。则 $M_P$ 是完整 $D(n)$-模。
\end{proposition}

\begin{proof}
显然可设 $P\neq 0$。设 $F M$ 是 $M$ 上好滤过，使得 $p\le 0$ 时 $F_pM=0$，且 $m=\deg P$。定义 $p<0$ 时 $F_pM_P=0$，
$$F_pM_P=\Bigl\{\frac{v}{P^p}\ \Big|\ v\in F_{(m+1)p}M\Bigr\}\quad(p\in\mathbb{Z}_+).$$
显然 $F_pM_P$（$p\in\mathbb{Z}$）是 $M_P$ 的向量子空间。

设 $w\in F_pM_P$，$p\ge 0$。则对某 $v\in F_{(m+1)p}M$ 有 $w=\frac{Pv}{P^{p+1}}$。由于 $Pv\in F_{(m+1)p+m}M\subset F_{(m+1)(p+1)}M$，可见 $w\in F_{p+1}M_P$。这证明滤过 $F M_P$ 是递增的。

设 $v\in F_qM$。则对任意 $s\in\mathbb{Z}_+$ 有 $\frac{v}{P^p}=\frac{P^sv}{P^{p+s}}$。此外对任意 $s\in\mathbb{Z}_+$ 有 $P^sv\in F_{q+sm}M$。而且当 $s\ge q-(m+1)p$ 时，$(m+1)(p+s)-(q+sm)=s+(m+1)p-q\ge 0$。故
$$P^sv\in F_{q+sm}M\subset F_{(m+1)(p+s)}M,\quad \frac{v}{P^p}\in F_{p+s}M_P.$$
因此滤过 $F M_P$ 是穷尽的。

还需证明它是 $D(n)$-模滤过。首先，对 $v\in F_{(m+1)p}M$，有 $x_iPv\in F_{(m+1)(p+1)}M$，故 $x_i\frac{v}{P^p}=\frac{Pv}{P^{p+1}}x_i\in F_{p+1}M_P$。此外，
$$\partial_i\Bigl(\frac{v}{P^p}\Bigr)=\frac{-p\,\partial_i(P)v+P\,\partial_i v}{P^{p+1}},$$
且 $-p\,\partial_i(P)v+P\,\partial_i v\in F_{(m+1)(p+1)}M$，故 $\partial_i\frac{v}{P^p}\in F_{p+1}M_P$。

因此在 $M_P$ 上构造了穷尽 $D(n)$-模滤过。由于对 $p\in\mathbb{Z}_+$，
$$\dim_k F_pM_P\le\dim_k F_{(m+1)p}M\le e(M)\frac{((m+1)p)^n}{n!}+\text{（关于 $p$ 的低阶项）},$$
由 8.4 知 $M_P$ 是完整的。
\end{proof}

\begin{corollary}
设 $P\in k[X_1,X_2,\dots,X_n]$。则 $k[X_1,X_2,\dots,X_n]_P$ 是完整 $D(n)$-模。
\end{corollary}
\begin{annote}
“完整”（holonomic）指 $\dim\mathrm{Ch}(M)=n$；因 $\mathrm{Ch}(M)$ 总是维数 $\ge n$ 的锥，这是可能的最小值。等价刻画：$M$ 的每个非零元在 $D(n)$ 上的零化子理想的维数…（Gabber 定理）——nLab 直言完整当且仅当特征簇维数为 $n$，此时特征簇既是锥又是 Lagrangian 子簇。这也解释了为什么完整模“几乎”是有限维的：解空间有限维（Riemann–Hilbert 对应的 D-模侧）。[来源：nLab 条目 holonomic+D-module（Last revised 2014-02-16）；原书 §8, p38–40]
\end{annote}



\section{外张量积}

以下设 $X=k^n$，$Y=k^m$，并以 $D_X$、$D_Y$ 记相应的多项式系数微分算子代数。则可考虑代数 $D_X\boxtimesext D_Y$，它作为 $k$-向量空间等于 $D_X\otimes_k D_Y$，乘法定义为 $(T\otimes S)(T'\otimes S')=TT'\otimes SS'$（$T,T'\in D_X$，$S,S'\in D_Y$）。称 $D_X\boxtimesext D_Y$ 为 $D_X$ 和 $D_Y$ 的\emph{外张量积}。

以下结果显然。

\begin{lemma}
$D_X\boxtimesext D_Y=D_{X\times Y}$。
\end{lemma}

若 $M$、$N$ 分别是 $D_X$-、$D_Y$-模，可定义 $D_{X\times Y}$-模 $M\boxtimesext N$，它作为 $k$-向量空间等于 $M\otimes_k N$，$D_X\boxtimesext D_Y=D_{X\times Y}$ 的作用由 $(T\otimes S)(m\otimes n)=Tm\otimes Sn$ 给出（$T\in D_X$，$S\in D_Y$，$m\in M$，$n\in N$）。

\begin{lemma}
设 $M$ 是有限生成 $D_X$-模，$N$ 是有限生成 $D_Y$-模。则 $M\boxtimesext N$ 是有限生成 $D_{X\times Y}$-模。
\end{lemma}

\begin{proof}
设 $e_1,e_2,\dots,e_p$ 和 $f_1,f_2,\dots,f_q$ 分别是 $M$ 和 $N$ 的生成元。则对任意 $m\in M$ 和 $n\in N$，有 $m=\sum_i T_ie_i$（$T_i\in D_X$），$n=\sum_j S_jf_j$（$S_j\in D_Y$）。这意味着 $m\otimes n=\sum_{i,j}T_ie_i\otimes S_jf_j=\sum_{i,j}(T_i\otimes S_j)(e_i\otimes f_j)$，且 $e_i\otimes f_j$（$1\le i\le p$，$1\le j\le q$）生成 $M\boxtimesext N$。
\end{proof}

本节的主要目标是证明如下结果。

\begin{theorem}
设 $M$ 是有限生成 $D_X$-模，$N$ 是有限生成 $D_Y$-模。则
$$d(M\boxtimesext N)=d(M)+d(N).$$
\end{theorem}

这有如下重要推论。

\begin{corollary}
设 $M$ 是完整 $D_X$-模，$N$ 是完整 $D_Y$-模。则 $M\boxtimesext N$ 是完整 $D_{X\times Y}$-模。
\end{corollary}

设 $D_X$ 和 $D_Y$ 配置 Bernstein 滤过。设 $M$、$N$ 分别是带好滤过 $F M$、$F N$ 的有限生成 $D_X$-、$D_Y$-模。定义 $M\boxtimesext N$ 上的\emph{乘积滤过}为
$$F_j(M\boxtimesext N)=\sum_{p+q=j}F_pM\otimes_k F_qN$$
对任意 $j\in\mathbb{Z}$。显然 $D_X\boxtimesext D_Y=D_{X\times Y}$ 上的乘积滤过与 Bernstein 滤过一致。因此 $F(M\boxtimesext N)$ 是穷尽 Hausdorff 的 $D_{X\times Y}$-模滤过。

为证明此滤过是好的，需要线性代数中的一些准备。从如下引理开始。

\begin{lemma}
设 $M,M',N,N'$ 是 $k$ 上线性空间，$\varphi:M\to M'$ 和 $\psi:N\to N'$ 是线性映射。则它们定义线性映射 $\varphi\otimes\psi:M\otimes_k N\to M'\otimes_k N'$。有
\begin{enumerate}[label=(\roman*)]
  \item $\operatorname{im}(\varphi\otimes\psi)=\operatorname{im}\varphi\otimes\operatorname{im}\psi$；
  \item $\ker(\varphi\otimes\psi)=\ker\varphi\otimes_k N+M\otimes_k\ker\psi$。
\end{enumerate}
\end{lemma}

\begin{proof}
(i) 由定义显然。

(ii) 由 (i)，为证 (ii) 可设 $\varphi$ 和 $\psi$ 是满射。这时有短正合列
$$0\longrightarrow M''\longrightarrow M\xrightarrow{\ \varphi\ }M'\longrightarrow 0,$$
其中 $M''=\ker\varphi$，以及
$$0\longrightarrow N''\longrightarrow N\xrightarrow{\ \psi\ }N'\longrightarrow 0,$$
其中 $N''=\ker\psi$。

显然 $\varphi\otimes\psi=(\varphi\otimes\operatorname{id}_{N'})\circ(\operatorname{id}_M\otimes\psi)$。由于与 $N'$ 做张量是正合的，第一正合列蕴含序列
$$0\longrightarrow M''\otimes_k N'\longrightarrow M\otimes_k N'\xrightarrow{\ \varphi\otimes\operatorname{id}\ }M'\otimes_k N'\longrightarrow 0$$
是正合的。故 $\ker(\varphi\otimes\operatorname{id}_{N'})=M''\otimes_k N'=\ker\varphi\otimes_k N'$。因此 $M\otimes_k N$ 中元素 $z$ 属于 $\varphi\otimes\psi$ 的核，当且仅当 $(\operatorname{id}_M\otimes\psi)(z)$ 属于 $\ker\varphi\otimes_k N'$。

由于
$$0\longrightarrow M\otimes_k N''\longrightarrow M\otimes_k N\xrightarrow{\ \operatorname{id}_M\otimes\psi\ }M\otimes_k N'\longrightarrow 0$$
也正合，$\ker\varphi\otimes_k N$ 满射到 $\ker\varphi\otimes_k N'$ 上，且 $\ker(\operatorname{id}_M\otimes\psi)=M\otimes_k N''=M\otimes_k\ker\psi$。因此 $z$ 属于 $\varphi\otimes\psi$ 的核当且仅当
$$z\in\ker\varphi\otimes_k N+M\otimes_k\ker\psi.$$
\end{proof}

\begin{lemma}
设 $X_1,X_2,\dots,X_n$ 是张成线性空间 $X$ 的线性子空间。若对 $1\le i\le n$ 有
$$X_i\cap\sum_{j\neq i}X_j=\{0\},$$
则线性空间 $X$ 是 $X_1,X_2,\dots,X_n$ 的直和。
\end{lemma}

\begin{proof}
设 $x_i\in X_i$（$1\le i\le n$）满足 $x_1+x_2+\cdots+x_n=0$。则 $x_i=-\sum_{j\neq i}x_j\in X_i\cap\sum_{j\neq i}X_j$，由假设对每个 $1\le i\le n$ 它等于 $0$。
\end{proof}

现在描述 $\Gr(M\boxtimesext N)$。设 $j\in\mathbb{Z}$。若 $p+q=j$，有良定义的 $k$-线性映射 $F_pM\otimes_k F_qN\to F_j(M\boxtimesext N)$，故有良定义的 $k$-线性映射 $F_pM\otimes_k F_qN\to\Gr_j(M\boxtimesext N)$。由 9.5，自然映射 $F_pM\otimes_k F_qN\to\Gr_pM\otimes_k\Gr_qN$ 的核为 $F_{p-1}M\otimes_k F_qN+F_pM\otimes_k F_{q-1}N$，即含于 $F_{j-1}(M\boxtimesext N)$。故线性映射 $F_pM\otimes_k F_qN\to\Gr_j(M\boxtimesext N)$ 通过 $\Gr_pM\otimes_k\Gr_qN$ 分解。这导出线性映射
$$\pi:\bigoplus_{p+q=j}\Gr_pM\otimes_k\Gr_qN\longrightarrow\Gr_j(M\boxtimesext N).$$
显然由其构造此映射是满射。此外它在直和中每个加项 $\Gr_pM\otimes_k\Gr_qN$ 上的限制是单射。设 $X_{p,q}$ 为 $\Gr_pM\otimes_k\Gr_qN$ 在 $\Gr_j(M\boxtimesext N)$ 中的像。由于
$$F_{p-1}M\otimes_k F_qN+F_pM\otimes_k F_{q-1}N=(F_pM\otimes_k F_qN)\cap\Bigl(\sum_{\substack{p'+q'=j\\ p'\neq p,\ q'\neq q}}F_{p'}M\otimes_k F_{q'}N\Bigr),$$
可见
$$X_{p,q}\cap\Bigl(\sum_{\substack{p'+q'=j\\ p'\neq p,\ q'\neq q}}X_{p',q'}\Bigr)=\{0\}.$$
故由 9.6 此映射是同构。这意味着对任意 $j\in\mathbb{Z}$，
$$\Gr_j(M\boxtimesext N)=\bigoplus_{p+q=j}\Gr_pM\otimes_k\Gr_qN.$$

若类似地定义代数 $\Gr D_X\boxtimesext\Gr D_Y$，其分次由总次数给出，则 $\Gr D_X\boxtimesext\Gr D_Y=\Gr D_{X\times Y}$。此外，$\Gr M\boxtimesext\Gr N$ 成为分次 $\Gr D_{X\times Y}$-模，由上述讨论同构于 $\Gr(M\boxtimesext N)$。由于滤过 $F M$ 和 $F N$ 是好的，由 3.1，$\Gr M$ 和 $\Gr N$ 分别是有限生成 $\Gr D_X$-、$\Gr D_Y$-模。由 9.2 的类似结论，$\Gr(M\boxtimesext N)$ 是有限生成 $\Gr D_{X\times Y}$-模。这意味着乘积滤过是 $M\boxtimesext N$ 上的好滤过。

设
$$P(M,t)=\sum_{p\in\mathbb{Z}}\dim_k(\Gr_p M)\,t^p,\qquad P(N,t)=\sum_{q\in\mathbb{Z}}\dim_k(\Gr_q N)\,t^q$$
为 $\Gr M$ 和 $\Gr N$ 的 Poincaré 级数。则
\begin{align*}
P(M,t)P(N,t)&=\sum_{p\in\mathbb{Z}}\sum_{q\in\mathbb{Z}}\dim_k(\Gr_p M)\dim_k(\Gr_q N)\,t^{p+q}\\
&=\sum_{j\in\mathbb{Z}}\Bigl(\sum_{p+q=j}\dim_k(\Gr_p M)\dim_k(\Gr_q N)\Bigr)t^j\\
&=\sum_{j\in\mathbb{Z}}\Bigl(\sum_{p+q=j}\dim_k(\Gr_p M\otimes_k\Gr_q N)\Bigr)t^j\\
&=\sum_{j\in\mathbb{Z}}\dim_k\Gr_j(M\boxtimesext N)\,t^j=P(M\boxtimesext N,t)
\end{align*}
是 $M\boxtimesext N$ 的 Poincaré 级数。因此 $P(M\boxtimesext N,t)$ 在 $1$ 处极点的阶数是 $P(M,t)$ 和 $P(N,t)$ 极点阶数之和。由 1.5 可见这立即蕴含 9.3。

也可通过考虑特征簇推导 9.3。把 $D_X$、$D_Y$、$D_{X\times Y}$ 看作按微分算子次数滤过的环。设 $M$、$N$ 是带好滤过 $F M$、$F N$ 的有限生成 $D_X$-、$D_Y$-模。如上定义 $M\boxtimesext N$ 上的 $D_{X\times Y}$-模滤过 $F(M\boxtimesext N)$。则如同上述论证，$F(M\boxtimesext N)$ 是 $M\boxtimesext N$ 的好滤过。设 $I$ 为 $\Gr M$ 在 $\Gr D_X$ 中的零化子，$J$ 为 $\Gr N$ 在 $\Gr D_Y$ 中的零化子。则由 9.5，$\Gr(M\boxtimesext N)$ 的零化子等于 $\Gr D_{X\times Y}=\Gr D_X\boxtimesext\Gr D_Y$ 中的理想 $I\otimes_k\Gr D_Y+\Gr D_X\otimes_k J$。

可把 $\Gr D_X$ 与多项式环 $k[x_1,\dots,x_n,\xi_1,\dots,\xi_n]$ 等同，把 $\Gr D_Y$ 与
$k[y_1,\dots,y_m,\eta_1,\dots,\eta_m]$ 等同。此外可把 $\Gr D_{X\times Y}$ 与多项式环
$$k[x_1,\dots,x_n,y_1,\dots,y_m,\xi_1,\dots,\xi_n,\eta_1,\dots,\eta_m]$$
等同。则 $\Gr(M\boxtimesext N)$ 的零化子对应于该环中由 $I$ 和 $J$ 的像生成的理想。若定义映射
$q:k^{2n}\times k^{2m}\to k^{2(n+m)}$ 为
\begin{multline*}
q(x_1,\dots,x_n,\xi_1,\dots,\xi_n,y_1,\dots,y_m,\eta_1,\dots,\eta_m)\\
=(x_1,\dots,x_n,y_1,\dots,y_m,\xi_1,\dots,\xi_n,\eta_1,\dots,\eta_m),
\end{multline*}
则有如下结果。

\begin{theorem}
设 $M$ 和 $N$ 是有限生成 $D_X$-、$D_Y$-模。则
$$\operatorname{Ch}(M\boxtimesext N)=q\bigl(\operatorname{Ch}(M)\times\operatorname{Ch}(N)\bigr).$$
\end{theorem}

这进而蕴含 $\dim\operatorname{Ch}(M\boxtimesext N)=\dim\operatorname{Ch}(M)+\dim\operatorname{Ch}(N)$，由 7.7 得到 9.3 的另一个证明。

利用 7.1 的证明并如上论证，或利用 7.5，也可得到如下结果。

\begin{proposition}
设 $M$ 和 $N$ 是有限生成 $D_X$-、$D_Y$-模。则
$$\operatorname{supp}(M\boxtimesext N)=\operatorname{supp}(M)\times\operatorname{supp}(N).$$
\end{proposition}

\section{逆像}

设 $X=k^n$，$Y=k^m$，以 $x_1,x_2,\dots,x_n$ 和 $y_1,y_2,\dots,y_m$ 分别记 $X$ 和 $Y$ 上的典范坐标函数。以 $R(X)=k[x_1,x_2,\dots,x_n]$ 和 $R(Y)=k[y_1,y_2,\dots,y_m]$ 分别记 $X$ 和 $Y$ 上正则函数环。

设 $F:X\longrightarrow Y$ 是多项式映射，即
$$F(x_1,x_2,\dots,x_n)=\bigl(F_1(x_1,\dots,x_n),F_2(x_1,\dots,x_n),\dots,F_m(x_1,\dots,x_n)\bigr),$$
其中 $F_i\in R(X)$。则 $F$ 定义环同态 $\varphi_F:R(Y)\to R(X)$：$\varphi_F(P)=P\circ F$（$P\in R(Y)$）。因此可把 $R(X)$ 看作 $R(Y)$-模。于是可定义函子 $F^*$，从 $R(Y)$-模范畴 $\mathcal{M}(R(Y))$ 到 $R(X)$-模范畴 $\mathcal{M}(R(X))$，由
$$F^*(N)=R(X)\otimes_{R(Y)}N$$
给出（对任意 $R(Y)$-模 $N$）。显然 $F^*:\mathcal{M}(R(Y))\to\mathcal{M}(R(X))$ 是右正合函子。我们称它为从 $\mathcal{M}(R(Y))$ 到 $\mathcal{M}(R(X))$ 的\emph{逆像函子}。

现在想把此函子扩充到 $D$-模。以 $D_X$、$D_Y$ 分别记 $X$、$Y$ 上多项式系数微分算子代数。若 $N$ 是左 $D_Y$-模，想在逆像 $F^*(N)$ 上定义 $D_X$-模结构。（如 §6 开头所述，转置函子是左 $D$-模范畴与右 $D$-模范畴的等价，故可类似处理右模。）先考虑双线性映射
$$(P,v)\longmapsto\frac{\partial P}{\partial x_i}\otimes v+\sum_{j=1}^m P\frac{\partial F_j}{\partial x_i}\otimes\frac{\partial}{\partial y_j}v,$$
从 $R(X)\times N$ 到 $R(X)\otimes_{R(Y)}N$。由于对任意 $Q\in R(Y)$，
\begin{align*}
&\frac{\partial P\,(Q\circ F)}{\partial x_i}\otimes v+\sum_j P\frac{\partial F_j}{\partial x_i}(Q\circ F)\otimes\frac{\partial}{\partial y_j}v\\
&=\frac{\partial P}{\partial x_i}\otimes Qv+\sum_j P\Bigl(\frac{\partial Q}{\partial y_j}\circ F\Bigr)\frac{\partial F_j}{\partial x_i}\otimes v+\sum_j P(Q\circ F)\frac{\partial F_j}{\partial x_i}\otimes\frac{\partial}{\partial y_j}v\\
&=\frac{\partial P}{\partial x_i}\otimes Qv+\sum_j P\frac{\partial F_j}{\partial x_i}\otimes\Bigl(\frac{\partial Q}{\partial y_j}v+Q\frac{\partial}{\partial y_j}v\Bigr)\\
&=\frac{\partial P}{\partial x_i}\otimes Qv+\sum_j P\frac{\partial F_j}{\partial x_i}\otimes\frac{\partial}{\partial y_j}(Qv),
\end{align*}
此映射通过 $F^*(N)$ 的线性自同态分解，记为 $\dfrac{\partial}{\partial x_i}$。直接计算得
$$\Bigl[\frac{\partial}{\partial x_i},\frac{\partial}{\partial x_j}\Bigr](P\otimes v)=0,\qquad \Bigl[\frac{\partial}{\partial x_i},x_j\Bigr](P\otimes v)=\delta_{ij}(P\otimes v),$$
故由 5.11，$F^*(N)$ 具有自然的左 $D_X$-模结构。

其结构可用另一方式描述。令
$$D_{X\to Y}=F^*(D_Y)=R(X)\otimes_{R(Y)}D_Y.$$
则如上所述，$D_{X\to Y}$ 有左 $D_X$-模结构。但它也有由 $D_Y$ 中右乘给出的右 $D_Y$-模结构。这两个作用显然可交换，故 $D_{X\to Y}$ 是（左 $D_X$，右 $D_Y$）-双模。此外，对任意 $D_Y$-模 $N$ 有
$$F^*(N)=R(X)\otimes_{R(Y)}N=\bigl(R(X)\otimes_{R(Y)}D_Y\bigr)\otimes_{D_Y}N=D_{X\to Y}\otimes_{D_Y}N,$$
且 $D_X$ 在 $F^*(N)$ 上的作用由最后一式中第一个因子上的作用给出。

我们记此 $D_X$-模为 $F^+(N)$，称为 $D_Y$-模 $N$ 的\emph{逆像}。

显然逆像函子 $F^+$ 是从 $\mathcal{M}(D_Y)$ 到 $\mathcal{M}^L(D_X)$ 的右正合函子。它的左导出函子 $L^iF^+$ 对左 $D_Y$-模 $N$ 由
$$L^iF^+(N)=\operatorname{Tor}^{-i}_{D_Y}(D_{X\to Y},N)$$
给出。

设 $\operatorname{For}$ 为从 $D_X$-模范畴（相应地从 $D_Y$-模范畴）到 $R(X)$-模范畴（相应地 $R(Y)$-模范畴）的遗忘函子。则如下函子图交换：
$$\begin{tikzcd}
\mathcal{M}(D_Y) \arrow[r,"F^+"] \arrow[d,"\operatorname{For}"'] & \mathcal{M}(D_X) \arrow[d,"\operatorname{For}"]\\
\mathcal{M}(R(Y)) \arrow[r,"F^*"'] & \mathcal{M}(R(X))
\end{tikzcd}$$

我们断言对左导出函子也有类似陈述，即：

\begin{proposition}
对任意 $i\in\mathbb{Z}$，如下函子图交换：
$$\begin{tikzcd}
\mathcal{M}(D_Y) \arrow[r,"L^iF^+"] \arrow[d,"\operatorname{For}"'] & \mathcal{M}(D_X) \arrow[d,"\operatorname{For}"]\\
\mathcal{M}(R(Y)) \arrow[r,"L^iF^*"'] & \mathcal{M}(R(X))
\end{tikzcd}$$
\end{proposition}

\begin{proof}
设 $F^\cdot$ 是 $D_Y$-模 $N$ 由自由 $D_Y$-模构成的一个左分解。由于自由 $D_Y$-模由 5.10 也是自由 $R(Y)$-模，由上述注记，对任意 $i\in\mathbb{Z}$ 有
$$\operatorname{For}(L^iF^+(N))=\operatorname{For}(H^i(F^+(F^\cdot)))=H^i(\operatorname{For}(F^+(F^\cdot)))=H^i(F^*(\operatorname{For}F^\cdot))=L^iF^*(\operatorname{For}N).$$
\end{proof}

现在研究导出逆像对态射复合的行为。先需要一个无环性结果。

\begin{lemma}
设 $P$ 是射影左 $D_Y$-模。则 $F^*(P)$ 是射影 $R(X)$-模。
\end{lemma}

\begin{proof}
设 $P$ 是射影 $D_Y$-模。则它是自由 $D_Y$-模 $(D_Y)^{(I)}$ 的直和因子。这意味着 $F^+(P)$ 是 $F^+(D_Y^{(I)})$ 的直和因子。由于 $D_Y$ 是自由 $R(Y)$-模，$\operatorname{For}(F^+(D_Y^{(I)}))=R(X)\otimes_{R(Y)}D_Y^{(I)}$ 是自由 $R(X)$-模。
\end{proof}

\begin{theorem}
设 $X=k^n$，$Y=k^m$，$Z=k^p$，$F:X\to Y$ 和 $G:Y\to Z$ 是多项式映射。则
\begin{enumerate}[label=(\roman*)]
  \item 逆像函子 $(G\circ F)^+$（从 $\mathcal{M}^L(D_Z)$ 到 $\mathcal{M}^L(D_X)$）同构于 $F^+\circ G^+$；
  \item 对任意左 $D_Z$-模 $N$，存在谱序列，其 $E_2$ 项为 $E_2^{pq}=L^pF^+(L^qG^+(N))$，收敛于 $L^{p+q}(G\circ F)^+(N)$。
\end{enumerate}
\end{theorem}

\begin{proof}
(i) 先考虑多项式环结构。对任意 $D_Z$-模 $N$，
$$(G\circ F)^*(N)=R(X)\otimes_{R(Z)}N=R(X)\otimes_{R(Y)}(R(Y)\otimes_{R(Z)}N)=F^*(G^*(N)).$$
另一方面，对任意 $P\in R(X)$ 和 $v\in N$，
\begin{align*}
\frac{\partial}{\partial x_i}(P\otimes v)&=\frac{\partial}{\partial x_i}\bigl(P\otimes(1\otimes v)\bigr)\\
&=\frac{\partial P}{\partial x_i}\otimes(1\otimes v)+\sum_{j=1}^m P\frac{\partial F_j}{\partial x_i}\otimes\frac{\partial}{\partial y_j}(1\otimes v)\\
&=\frac{\partial P}{\partial x_i}\otimes v+\sum_{j=1}^m P\frac{\partial F_j}{\partial x_i}\otimes\Bigl(\sum_{k=1}^p\frac{\partial G_k}{\partial y_j}\otimes\frac{\partial}{\partial z_k}v\Bigr)\\
&=\frac{\partial P}{\partial x_i}\otimes v+\sum_{k=1}^p\sum_{j=1}^m P\Bigl(\frac{\partial F_j}{\partial x_i}\frac{\partial G_k}{\partial y_j}\Bigr)\circ F\otimes\frac{\partial}{\partial z_k}v\\
&=\frac{\partial P}{\partial x_i}\otimes v+\sum_{k=1}^p P\frac{\partial(G_k\circ F)}{\partial x_i}\otimes\frac{\partial}{\partial z_k}v.
\end{align*}
故 $D_X$ 作用一致。

(ii) 由 10.1 和 10.2，对任意射影 $D_Z$-模 $P$，逆像 $G^+(P)$ 是 $F^+$-无环的。因此结论由 Grothendieck 谱序列得出。
\end{proof}

这立即有如下推论。

\begin{corollary}
设 $X=k^n$，$Y=k^m$，$Z=k^p$，$F:X\to Y$ 和 $G:Y\to Z$ 是多项式映射。则
\begin{enumerate}[label=(\roman*)]
  \item $D_{X\to Z}=D_{X\to Y}\otimes_{D_Y}D_{Y\to Z}$；
  \item $\operatorname{Tor}^{D_Y}_j(D_{X\to Y},D_{Y\to Z})=0$ 对 $j\in\mathbb{N}$。
\end{enumerate}
\end{corollary}

\begin{proof}
(i) 由 10.3.(i) 有
$$D_{X\to Z}=(G\circ F)^+(D_Z)=F^+(G^+(D_Z))=F^+(D_{Y\to Z})=D_{X\to Y}\otimes_{D_Y}D_{Y\to Z}.$$

(ii) 如在 10.3.(ii) 的证明中所述，由 10.1 和 10.2 可见 $D_{Y\to Z}=G^+(D_Z)$ 是 $F^+$-无环的。故对 $j>0$，
$$0=L^{-j}F^+(G^+(D_Z))=\operatorname{Tor}^{-j}_{D_Y}(D_{X\to Y},D_{Y\to Z}).$$
\end{proof}

现在考虑两个简单的例子。首先设 $p$ 是 $X\times Y$ 到 $Y$ 的投影，由 $p(x,y)=y$（$x\in X$，$y\in Y$）定义。则众所周知，$R(X\times Y)=R(X)\boxtimesext R(Y)$。因此对 $R(Y)$-模 $N$ 有
$$p^*(N)=R(X\times Y)\otimes_{R(Y)}N=\bigl(R(X)\boxtimesext R(Y)\bigr)\otimes_{R(Y)}N=R(X)\boxtimesext N$$
作为 $R(X\times Y)=R(X)\boxtimesext R(Y)$ 上的模。另一方面，若 $N$ 是 $D_Y$-模，则立即看出 $\frac{\partial}{\partial x_i}$ 和 $\frac{\partial}{\partial y_j}$ 的作用也一致，即
$$p^+(N)=R(X)\boxtimesext N.$$
由 9.3 和 9.4 立即得到如下结果。

\begin{proposition}
设 $p:X\times Y\longrightarrow Y$ 是典范投影。则
\begin{enumerate}[label=(\roman*)]
  \item $p^+$ 是从 $\mathcal{M}^L(D_Y)$ 到 $\mathcal{M}^L(D_{X\times Y})$ 的正合函子；
  \item 对任意左 $D_Y$-模 $N$，$p^+(N)=R(X)\boxtimesext N$；
  \item 若 $N$ 是有限生成左 $D_Y$-模，则 $p^+(N)$ 是有限生成 $D_{X\times Y}$-模；
  \item 对任意有限生成左 $D_Y$-模 $N$，$d(p^+(N))=d(N)+n$。
\end{enumerate}
特别地，有限生成左 $D_Y$-模 $N$ 是完整的当且仅当 $p^+(N)$ 是完整的。
\end{proposition}

现在考虑另一个例子。设 $i$ 是 $X$ 到 $X\times Y$ 的典范注入，由 $i(x)=(x,0)$（$x\in X$）给出。则
$$D_{X\to X\times Y}=i^*(D_{X\times Y})=R(X)\otimes_{R(X)\boxtimesext R(Y)}(D_X\boxtimesext D_Y)=D_X\boxtimesext D_Y/\bigl((y_1,y_2,\dots,y_m)D_Y\bigr),$$
其中 $D_X$ 显然通过第一因子上的左乘作用，而 $D_{X\times Y}=D_X\boxtimesext D_Y$ 通过右乘作用。

此外设 $m=1$。则有正合列
$$0\longrightarrow D_Y\xrightarrow{\ y_1\ }D_Y\longrightarrow D_Y/y_1D_Y\longrightarrow 0,$$
其中第二个箭头由左乘 $y_1$ 给出。与 $D_X$ 作张量，得到左 $D_X$-模（对左乘）和右 $D_{X\times Y}$-模（对右乘）的短正合列
$$0\longrightarrow D_{X\times Y}\xrightarrow{\ y_1\ }D_{X\times Y}\longrightarrow D_{X\to X\times Y}\longrightarrow 0.$$
因此可把此正合列的前两项看作 $D_{X\to X\times Y}$ 的、由作为左 $D_{X\times Y}$-模自由的（左 $D_X$，右 $D_{X\times Y}$）-双模构成的左分解。故对 $D_{X\times Y}$-模 $N$，复形
$$0\longrightarrow N\xrightarrow{\ y_1\ }N\longrightarrow 0$$
的上同调计算导出逆像。特别地，有如下引理。

\begin{lemma}
设 $\dim Y=1$，$i$ 是 $X$ 到 $X\times Y$ 的典范注入。则对任意 $D_{X\times Y}$-模 $N$ 有
\begin{enumerate}[label=(\roman*)]
  \item $i^+(N)=\operatorname{coker}y_1$；
  \item $L^{-1}i^+(N)=\ker y_1$；
  \item 当 $p\neq 0,-1$ 时 $L^pi^+(N)=0$。
\end{enumerate}
特别地，$i^+$ 的左同调维数 $\le 1$。
\end{lemma}

最后陈述对任意 $Y$ 有显然的推广。

\begin{lemma}
设 $i$ 是 $X$ 到 $X\times Y$ 的典范注入。则 $i^+$ 的左同调维数 $\le\dim Y$。
\end{lemma}

\begin{proof}
对 $\dim Y$ 归纳。$\dim Y=1$ 时结论已证。可把 $Y$ 表为 $Y=Y'\times Y''$，其中 $Y'=k^{m-1}$，$Y''=k$。以 $i'$ 记 $X$ 到 $X\times Y'$ 的典范包含，以 $j$ 记 $X\times Y'$ 到 $X\times Y'\times Y''=X\times Y$ 的典范包含。则 $i=j\circ i'$。此外由 10.6，$j^+$ 的左同调维数 $\le 1$，而由归纳假设，$i'^+$ 的左同调维数 $\le\dim Y'$。因此由 10.3.(ii) 的 Grothendieck 谱序列断定，导出逆像 $L^{-p}i^+$ 对 $p\ge\dim Y'+1=\dim Y$ 为零。
\end{proof}

设 $F:X\longrightarrow X$ 是 $X$ 的同构，$G$ 是其逆。则由 $\alpha(f)=f\circ F$ 定义的映射 $\alpha:R(X)\to R(X)$ 是环 $R(X)$ 的自同构。其逆为 $\beta$：$\beta(f)=f\circ G$（$f\in R(X)$）。若 $M$ 是 $R(X)$-模，则作为 $k$ 上的线性空间，$F^*(M)$ 通过映射 $\varphi:m\longmapsto1\otimes m$ 同构于 $M$。另一方面，对 $f\in R(X)$ 有
$$f\varphi(m)=f\otimes m=f\circ G\circ F\otimes m=1\otimes(f\circ G)m=\varphi(\beta(f)m)$$
（$m\in M$），即 $R(X)$-模 $F^*(M)$ 同构于 $M$，其 $R(X)$-模结构由 $(f,m)\longmapsto\beta(f)m$（$f\in R(X)$，$m\in M$）给出。

现在想给出 $F^+(M)$ 的类似描述。先要把自同构 $\beta$ 扩充到 $D_X$。

设 $T$ 是 $X$ 上的微分算子，令 $\tilde\beta(T)(f)=\beta(T\alpha(f))$（$f\in R(X)$）。显然 $\tilde\beta(T)$ 是 $R(X)$ 的 $k$-线性自同态。此外 $T\longmapsto\tilde\beta(T)$ 是线性映射。又对 $D_X$ 中两个微分算子 $T$ 和 $S$，对一切 $f\in R(X)$ 有
\begin{align*}
\tilde\beta(TS)(f)&=\beta(TS\alpha(f))=\beta\bigl(T\alpha(\beta(S\alpha(f)))\bigr)\\
&=\beta\bigl(T\alpha(\tilde\beta(S)(f))\bigr)=\tilde\beta(T)\bigl(\tilde\beta(S)(f)\bigr),
\end{align*}
即 $\tilde\beta$ 是从 $k$-代数 $D_X$ 到 $R(X)$ 的 $k$-线性自同态代数的同态。由于对 $g\in R(X)$ 有
$$\tilde\beta(g)f=\beta(g\alpha(f))=\beta(g)f$$
（$f\in R(X)$），可见 $\tilde\beta$ 扩充了 $R(X)$ 的自同构 $\beta$。这进而蕴含 $\tilde\beta(T)\in D_X$（$T\in D_X$），即 $\tilde\beta$ 是 $D_X$ 的自同构，它扩充了 $R(X)$ 的自同构 $\beta$。因此可简单地仍记为 $\beta$。

设 $1\le i\le n$。则
\begin{align*}
\beta(\partial_i)(f)&=\beta(\partial_i\alpha(f))=\beta\bigl(\partial_i(f\circ F)\bigr)\\
&=\beta\Bigl(\sum_{j=1}^n\bigl((\partial_jf)\circ F\bigr)\partial_iF_j\Bigr)=\sum_{j=1}^n\bigl((\partial_iF_j)\circ G\bigr)\partial_jf=\Bigl(\sum_{j=1}^n\beta(\partial_iF_j)\partial_j\Bigr)(f).
\end{align*}

现在考虑与映射 $F$ 相伴的双模 $D_{X\to X}$。线性映射 $\varphi:f\otimes T\longmapsto\beta(f)T$ 把它与 $D_X$ 等同。由右乘给出的 $D_X$-模结构相同。另一方面，对任意 $T\in D_X$ 和 $1\le i\le n$ 有
$$\varphi\bigl(\partial_i(1\otimes T)\bigr)=\varphi\Bigl(\sum_{j=1}^n\partial_iF_j\otimes\partial_jT\Bigr)=\sum_{j=1}^n\beta(\partial_iF_j)\partial_jT=\beta(\partial_i)\varphi(1\otimes T).$$
因此双模 $D_{X\to X}$ 同构于 $D_X$，其右作用为右乘，左作用为 $\beta$ 与左乘的复合。这进而蕴含 $F^+(M)$ 同构于 $M$，其 $D_X$-模结构由 $(T,m)\longmapsto\beta(T)m$（$T\in D_X$，$m\in M$）给出。

因此由 3.10 得到如下结果。

\begin{lemma}
设 $F:X\longrightarrow X$ 是 $X$ 的同构。
\begin{enumerate}[label=(\roman*)]
  \item 设 $M$ 是 $D_X$-模。则作为线性空间 $F^+(M)$ 等于 $M$，其 $D_X$-作用由 $(T,m)\longmapsto\beta(T)m$（$T\in D_X$，$m\in M$）给出。
  \item 函子 $F^+:\mathcal{M}^L(D_X)\to\mathcal{M}^L(D_X)$ 是正合的。
  \item 函子 $F^+$ 把有限生成 $D_X$-模映为有限生成 $D_X$-模。若 $M$ 是有限生成 $D_X$-模，则 $d(F^+(M))=d(M)$。
\end{enumerate}
特别地，$F^+$ 把完整模映为完整模。
\end{lemma}

可通过描述有限生成 $D_X$-模 $M$ 的特征簇 $\operatorname{Ch}(F^+(M))$ 使上述陈述更精确。首先，由上述计算可见 $D_X$ 的自同构 $\beta$ 诱导 $\Gr D_X=k[X_1,X_2,\dots,X_n,\xi_1,\xi_2,\dots,\xi_n]$ 的自同构 $\Gr\beta$，它由
$$X_i\longmapsto\beta(X_i)=G_i,\qquad \xi_i\longmapsto\sum_{j=1}^n\beta(\partial_iF_j)\xi_j=\sum_{j=1}^n\bigl((\partial_iF_j)\circ G\bigr)\xi_j$$
（$1\le i\le n$）定义。

现在想用更几何的语言描述此构造。若 $x=(x_1,x_2,\dots,x_n)$ 是 $X=k^n$ 中的点，则通过 $df(x)\longmapsto\bigl((\partial_1f)(x),(\partial_2f)(x),\dots,(\partial_nf)(x)\bigr)$ 把 $x$ 处的余切空间 $T_x^*(X)$ 与 $k^n$ 等同。因此通过 $(x,df(x))\longmapsto\bigl(x_1,\dots,x_n,(\partial_1f)(x),\dots,(\partial_nf)(x)\bigr)$（$x\in X$）可把 $X$ 的余切丛 $T^*(X)$ 与 $k^{2n}$ 等同。设 $F:X\longrightarrow X$ 是 $X$ 的同构，$G$ 是其逆。则映射 $G$ 把 $X$ 中的点 $x$ 映为 $G(x)$，而 $F$ 把 $G(x)$ 映为 $x$。它们的微分 $T_x(G)$ 和 $T_{G(x)}(F)$ 是切空间 $T_x(X)$ 与 $T_{G(x)}(X)$ 之间互逆的线性同构。因此它们的伴随 $T_x(G)^*:T_{G(x)}^*(X)\to T_x^*(X)$ 和 $T_{G(x)}(F)^*:T_x^*(X)\to T_{G(x)}^*(X)$ 是互逆的线性同构。这进而蕴含可定义 $X$ 的余切丛 $T^*(X)$ 的同构 $\gamma$，由 $(x,\xi)\longmapsto\bigl(G(x),T_{G(x)}(F)^*\xi\bigr)$（$\xi\in T_x^*(X)$，$x\in X$）给出。若把 $T^*(X)$ 与 $k^{2n}$ 等同，则由考察上述公式可见 $(\Gr\beta)(P)=P\circ\gamma$（$P\in k[X_1,X_2,\dots,X_n,\xi_1,\xi_2,\dots,\xi_n]$）。

设 $M$ 是带好滤过 $F M$ 的有限生成 $D_X$-模。则可把 $F^+(M)$ 实现为按上述方式作用的 $M$。显然 $F M$ 是如上实现的 $F^+(M)$ 的好滤过。因此可把 $\Gr F^+(M)$ 与 $\Gr M$ 等同，其作用由 $(Q,m)\longmapsto(\Gr\beta)(Q)m$（$Q\in k[X_1,X_2,\dots,X_n,\xi_1,\xi_2,\dots,\xi_n]$，$m\in\Gr M$）给出。故 $Q$ 属于 $\Gr F^+(M)$ 的零化子当且仅当 $(\Gr\beta)(Q)$ 属于 $\Gr M$ 的零化子。若 $I$ 是 $\Gr F^+(M)$ 的零化子，则 $(\Gr\beta)I$ 是 $\Gr M$ 的零化子。因此 $(x,\xi)\in\operatorname{Ch}(F^+(M))$ 当且仅当 $\gamma^{-1}(x,\xi)\in\operatorname{Ch}(M)$。

\begin{lemma}
设 $M$ 是有限生成 $D_X$-模。则
$$\operatorname{Ch}(F^+(M))=\gamma(\operatorname{Ch}(M)).$$
\end{lemma}

最后，这允许给出逆像函子左同调维数的一个估计。

\begin{theorem}
设 $X=k^n$，$Y=k^m$，$F:X\longrightarrow Y$ 是多项式映射。则 $F^+$ 的左同调维数 $\le\dim Y$。
\end{theorem}

\begin{proof}
为证明此陈述，我们使用图像构造。设 $i:X\to X\times Y$ 是由 $i(x)=(x,0)$（$x\in X$）给出的映射。设 $\Phi:X\times Y\to X\times Y$ 是由 $\Phi(x,y)=(x,y+F(x))$（$x\in X$，$y\in Y$）给出的映射。最后，设 $p:X\times Y\to Y$ 是由 $p(x,y)=y$（$x\in X$，$y\in Y$）给出的投影。则 $F=p\circ\Phi\circ i$。此外，$\Phi$ 是 $X\times Y$ 的同构，其逆为 $(x,y)\longmapsto(x,y-F(x))$。

由 10.3，$F^+=i^+\circ\Phi^+\circ p^+$。此外由 10.5 和 10.8，函子 $p^+$ 和 $\Phi^+$ 正合。因此对一切 $q\in\mathbb{Z}$ 有 $L^qF^+=L^qi^+\circ\Phi^+\circ p^+$。由 10.7 断定，当 $q<-\dim Y$ 时 $L^qF^+=0$。
\end{proof}
\begin{annote}
本节以下六条（投影空间例子、$i_+$ 的同调维数、自同构情形的 $\beta$、特征簇的像、以及左同调维数 $\le\dim Y$ 的图像构造证明）在初版译本中整段缺失，本次校订据英文原书 p49–p52 补译。若与旧版页边编号不合，以此版为准。[来源：原书 §10.5–10.10, p49–p52；校订日志 chap1.md]
\end{annote}



\section{直像}

设 $X=k^n$，$Y=k^n$，$F:X\to Y$ 是如上节的多项式映射。与 $F$ 的复合定义自然环同态 $\widehat F:R(Y)\to R(X)$。此同态进而定义从 $R(X)$-模范畴到 $R(Y)$-模范畴的函子 $F_*$。对任意 $R(X)$-模 $M$，定义 $F_*(M)$ 为作为 $k$-线性空间等于 $M$ 的模，$R(Y)$ 的作用由 $(f,m)\longmapsto\widehat F(f)\cdot m$ 给出（$f\in R(Y)$，$m\in M$）。函子 $F_*:\mathcal{M}(R(X))\to\mathcal{M}(R(Y))$ 称为\emph{直像函子}。显然 $F_*$ 是正合函子。

遗憾的是，若 $M$ 是 $D_X$-模，直像 $F_*(M)$ 一般不允许 $D_Y$-模结构。例如考虑 $X=\{0\}$ 到 $Y=k$ 的包含 $i$，$D_X=R(X)=k$，$D_Y$ 是一变量多项式系数微分算子代数。$D_X$-模范畴就是 $k$-线性空间范畴。由 6.2，非零有限维 $D_X$-模 $M$ 的逆像不能有 $D_Y$-模结构。因此 $D$-模的直像不会像逆像那样与正则函数环上模的直像相关。

若对 $D_{X\to Y}$ 的两个作用都取转置，得到（左 $D_Y$，右 $D_X$）-双模 $D_{Y\leftarrow X}$。这允许对任意左 $D_X$-模 $M$ 定义左 $D_Y$-模
$$F_+(M)=D_{Y\leftarrow X}\otimes_{D_X}M.$$
显然 $F_+$ 是从 $\mathcal{M}^L(D_X)$ 到 $\mathcal{M}^L(D_Y)$ 的右正合函子。称之为\emph{直像函子}。$F_+$ 的左导出函子 $L^iF_+$ 对左 $D_X$-模 $M$ 由
$$L^{-j}F_+(M)=\operatorname{Tor}^{D_X}_j(D_{Y\leftarrow X},M)$$
给出。

设 $X=k^n$，$Y=k^m$，$Z=k^p$，$F:X\to Y$ 和 $G:Y\to Z$ 是多项式映射。若转置作用，10.4 蕴含如下陈述：
$$D_{Z\leftarrow X}=D_{Z\leftarrow Y}\otimes_{D_Y}D_{Y\leftarrow X},$$
以及
$$\operatorname{Tor}^{D_Y}_j(D_{Z\leftarrow Y},D_{Y\leftarrow X})=0\quad\text{对 }j\in\mathbb{N}.$$

若 $P$ 是射影左 $D_X$-模，$P\oplus Q=D_X^{(I)}$（对某左 $D_X$-模 $Q$ 和某 $I$）。因此 $F_+(P)\oplus F_+(Q)=F_+(D_X^{(I)})=(D_{Y\leftarrow X})^{(I)}$。这蕴含如下结果。

\begin{lemma}
设 $P$ 是射影左 $D_X$-模。则
$$\operatorname{Tor}^{D_Y}_j(D_{Z\leftarrow Y},F_+(P))=0\quad\text{对 }j\in\mathbb{N}.$$
\end{lemma}

\begin{theorem}
设 $X=k^n$，$Y=k^m$，$Z=k^p$，$F:X\to Y$ 和 $G:Y\to Z$ 是多项式映射。则
\begin{enumerate}[label=(\roman*)]
  \item 直像函子 $(G\circ F)_+$（从 $\mathcal{M}^L(D_X)$ 到 $\mathcal{M}^L(D_Z)$）同构于 $G_+\circ F_+$；
  \item 对任意左 $D_X$-模 $M$，存在谱序列，其 $E_2$ 项为 $E_2^{pq}=L^pG_+(L^qF_+(M))$，收敛于 $L^{p+q}(G\circ F)_+(M)$。
\end{enumerate}
\end{theorem}

\begin{proof}
(i) 对任意左 $D_X$-模 $M$，由 10.4.(i) 有
\begin{align*}
(G\circ F)_+(M)&=D_{Z\leftarrow X}\otimes_{D_X}M
=(D_{Z\leftarrow Y}\otimes_{D_Y}D_{Y\leftarrow X})\otimes_{D_X}M\\
&=D_{Z\leftarrow Y}\otimes_{D_Y}(D_{Y\leftarrow X}\otimes_{D_X}M)
=D_{Z\leftarrow Y}\otimes_{D_Y}F_+(M)=G_+(F_+(M)).
\end{align*}

(ii) 由 11.1，对任意射影 $D_X$-模 $P$，直像 $F_+(P)$ 是 $G_+$-无环的。因此结论由 Grothendieck 谱序列得出。
\end{proof}

现在考虑一个简单的例子。设 $i$ 为由 $i(x)=(x,0)$（$x\in X$）给出的 $X$ 到 $X\times Y$ 的典范注入。则
$$D_{X\to X\times Y}=i_+(D_{X\times Y})=i_+(D_X\boxtimesext D_Y)=D_X\boxtimesext D_Y/\bigl((y_1,y_2,\dots,y_m)D_Y\bigr)$$
以及
$$D_{X\times Y\leftarrow X}=D_X\boxtimesext D_Y/\bigl(D_Y(y_1,y_2,\dots,y_m)\bigr).$$
这意味着对任意左 $D_X$-模 $M$，
$$i_+(M)=M\boxtimesext D_Y/\bigl(D_Y(y_1,y_2,\dots,y_m)\bigr).$$
此外，模 $D_Y/(D_Y(y_1,y_2,\dots,y_m))$ 同构于 8.3 中讨论的 $\Delta_m$。

\begin{proposition}
设 $i:X\to X\times Y$ 是由 $i(x)=(x,0)$（$x\in X$）定义的注入。则
\begin{enumerate}[label=(\roman*)]
  \item $i_+$ 是从 $\mathcal{M}^L(D_X)$ 到 $\mathcal{M}^L(D_{X\times Y})$ 的正合函子；
  \item 对任意左 $D_X$-模 $M$，$i_+(M)=M\boxtimesext D_Y/\bigl(D_Y(y_1,y_2,\dots,y_m)\bigr)$；
  \item 若 $M$ 是有限生成 $D_X$-模，则 $i_+(M)$ 是有限生成 $D_{X\times Y}$-模；
  \item 对任意有限生成左 $D_X$-模 $M$，$d(i_+(M))=d(M)+m$。
\end{enumerate}
特别地，有限生成 $D_X$-模 $M$ 是完整的当且仅当 $i_+(M)$ 是完整的。
\end{proposition}

\begin{proof}
已证明 (ii)，它立即蕴含 (i)。如 8.3 中所述，$D_Y/(D_Y(y_1,y_2,\dots,y_m))$ 是不可约完整 $D_Y$-模，故 (iii) 由 9.2 得出。为证 (iv)，先注意由 9.3 有
$$d(i_+(M))=d(M)+d\bigl(D_Y/(D_Y(y_1,y_2,\dots,y_m))\bigr).$$
由于 $D_Y/(D_Y(y_1,y_2,\dots,y_m))$ 是完整的，其维数等于 $m$。
\end{proof}

现在研究投影 $p:X\times Y\to Y$（$p(x,y)=y$，$x\in X$，$y\in Y$）的直像。

先考虑 $\dim X=1$ 的情形。则
$$D_{X\times Y\to Y}=p_+(D_Y)=D_X/D_X(\partial_1)\boxtimesext D_Y.$$
故 $D_{Y\leftarrow X\times Y}=D_X/((\partial_1)D_X)\boxtimesext D_Y$。有正合列
$$0\longrightarrow D_{X\times Y}\xrightarrow{\ \partial_1\ }D_{X\times Y}\longrightarrow D_{Y\leftarrow X\times Y}\longrightarrow 0$$
（（左 $D_Y$，右 $D_{X\times Y}$）-双模），其中第二个箭头表示左乘 $\partial_1$。显然这是 $D_{Y\leftarrow X\times Y}$ 由自由右 $D_{X\times Y}$-模构成的左分解，故对任意 $D_{X\times Y}$-模 $M$，复形
$$\cdots\longrightarrow 0\longrightarrow M\xrightarrow{\ \partial_1\ }M\longrightarrow 0\longrightarrow\cdots$$
的上同调是 $\operatorname{Tor}^{D_{X\times Y}}_\cdot(D_{Y\leftarrow X\times Y},M)=L_\cdot p_+(M)$。由此可知 $L^q p_+(M)=0$ 对 $q\notin\{0,-1\}$。

因此建立了如下结果。

\begin{lemma}
设 $\dim X=1$，$p$ 是 $X\times Y$ 到 $Y$ 的典范投影。则对任意 $D_{X\times Y}$-模 $M$ 有
\begin{enumerate}[label=(\roman*)]
  \item $p_+(M)=\operatorname{coker}\partial_1$；
  \item $L^{-1}p_+(N)=\ker\partial_1$；
  \item $L^q p_+(N)=0$ 对不同于 $0$ 或 $-1$ 的 $q$。
\end{enumerate}
特别地，$p_+$ 的左同调维数 $\le 1$。
\end{lemma}

最后陈述对任意 $X$ 有如下推广。

\begin{lemma}
设 $p$ 是 $X\times Y$ 到 $Y$ 的典范投影。则 $p_+$ 的左同调维数 $\le\dim X$。
\end{lemma}

\begin{proof}
设 $X'=\{x_n=0\}\subset X$，以 $p'$ 记 $X'\times Y$ 到 $Y$ 的典范投影。又以 $p''$ 记 $X\times Y$ 到 $X'\times Y$ 的典范投影。则 $p=p'\circ p''$。故由 11.2.(ii)、11.4 和归纳假设得对 $q<-\dim X$ 有 $L^q p_+(M)=0$。
\end{proof}

设 $F:X\to X$ 是 $X$ 的同构，$G$ 是其逆。如 §10 中，定义 $D_X$ 的自同构 $\alpha$ 和 $\beta$。我们在那里把 $F$ 附带的双模 $D_{X\to X}$ 与 $D_X$ 等同，其作用是右乘和与 $\beta$ 复合的左乘。对它应用 $\alpha$，可见它也同构于 $D_X$，其作用是右乘与 $\alpha$ 复合和左乘。通过应用主反自同构，可见双模 $D_{X\leftarrow X}$ 同构于 $D_X$，其作用是与 $\alpha$ 复合的左乘和右乘。这进而蕴含对任意 $D_X$-模 $M$，直像 $F_+(M)$ 同构于 $M$，其作用由 $(T,m)\longmapsto\alpha(T)m$ 给出。特别地 $F_+(M)=G_+(M)$。

因此，由 10.8 立即推出如下结果。

\begin{lemma}
设 $F:X\to X$ 是 $X$ 的同构，$G:X\to X$ 是其逆。

(i) 设 $M$ 是 $D_X$-模。则 $F_+(M)$ 作为线性空间等于 $M$，$D_X$ 作用由 $(T,m)\longmapsto\alpha(T)m$ 给出（$T\in D_X$，$m\in M$）。

(ii) 函子 $F_+:\mathcal{M}^L(D_X)\to\mathcal{M}^L(D_X)$ 是正合的。

(iii) 函子 $F_+$ 把有限生成 $D_X$-模映为有限生成 $D_X$-模。若 $M$ 是有限生成 $D_X$-模，则 $d(F_+(M))=d(M)$。特别地 $F_+$ 把完整模映为完整模。
\end{lemma}

此外 $F_+=G_+$ 且 $F^+=G^+$，这些函子互为范畴的拟逆等价。

如上节，这允许给出直像函子左同调维数的估计。

\begin{theorem}
设 $X=k^n$，$Y=k^m$，$F:X\to Y$ 是多项式映射。则 $F_+$ 的左同调维数 $\le\dim X$。
\end{theorem}

\begin{proof}
如 10.10 的证明，使用图像构造。设 $i:X\to X\times Y$ 为由 $i(x)=(x,0)$（$x\in X$）给出的态射。设 $\Phi:X\times Y\to X\times Y$ 为由 $\Phi(x,y)=(x,y+F(x))$（$x\in X$，$y\in Y$）给出的态射。最后设 $p:X\times Y\to Y$ 为由 $p(x,y)=y$（$x\in X$，$y\in Y$）给出的投影。则 $F=p\circ\Phi\circ i$。此外 $\Phi$ 是 $X\times Y$ 的同构，其逆为 $(x,y)\longmapsto(x,y-F(x))$。

由 11.2，$F_+=p_+\circ\Phi_+\circ i_+$。此外由 11.3 和 11.6，函子 $i_+$ 和 $\Phi_+$ 是正合的。因此对所有 $q\in\mathbb{Z}$ 有 $L^qF_+=L^qp_+\circ\Phi_+\circ i_+$。由 11.5，对 $q<-\dim X$ 有 $L^qF_+=0$。
\end{proof}

\section{Kashiwara 定理}

设 $X=k^n$，$Y=\{x_n=0\}\subset X$。再令 $Z=\{x_1=x_2=\cdots=x_{n-1}=0\}\cong k$。于是 $X=Y\times Z$，这也意味着 $D_X=D_Y\boxtimesext D_Z$。设 $M$ 是 $D_X$-模，令
$$\Gamma_{[Y]}(M)=\{m\in M\mid x_n^p m=0\ \text{对某 }p\in\mathbb{N}\}.$$

\begin{lemma}
设 $M$ 是 $D_X$-模。则
\begin{enumerate}[label=(\roman*)]
  \item $\Gamma_{[Y]}(M)$ 是 $M$ 的 $D_X$-子模；
  \item $\operatorname{supp}(\Gamma_{[Y]}(M))\subset Y$；
  \item 若 $N$ 是 $M$ 的 $D_X$-子模且 $\operatorname{supp}(N)\subset Y$，则 $N\subset\Gamma_{[Y]}(M)$。
\end{enumerate}
\end{lemma}

\begin{proof}
(i) 设 $m\in\Gamma_{[Y]}(M)$。则对 $1\le i\le n$ 和 $1\le j<n$ 有 $x_i m\in\Gamma_{[Y]}(M)$、$\partial_j m\in\Gamma_{[Y]}(M)$。只需检验 $\partial_n m\in\Gamma_{[Y]}(M)$。对任意 $j\in\mathbb{N}$ 有
$$x_n^{j+1}\partial_n m=[x_n,\partial_n]m+\partial_n x_n m=-(j+1)x_n^j m+\partial_n x_n^{j+1}m.$$
故若 $x_n^j m=0$，则 $x_n^{j+1}\partial_n m=0$。

(ii) 若 $x\notin Y$，则 $x_n\notin\mathfrak{m}_x$，局部化 $\Gamma_{[Y]}(M)_x=0$。

(iii) 设 $N$ 是 $M$ 的 $D_X$-子模且 $\operatorname{supp}(N)\subset Y$。取 $m\in N$，记 $N'$ 为由 $m$ 生成的 $R(X)$-子模。则 $\operatorname{supp}(N')\subset Y$。由于 $N'$ 有限生成，由 4.2 其支集等于其在 $R(X)$ 中零化子 $I$ 决定的簇。由零点定理 $r(I)\supset(x_n)$。这意味着对某 $j\in\mathbb{N}$，$x_n^j$ 零化 $N'$，即 $m\in\Gamma_{[Y]}(M)$。
\end{proof}

因此 $\Gamma_{[Y]}(M)$ 是 $M$ 中支集含于 $Y$ 的最大 $D_X$-子模。

乘以 $x_n$ 定义了 $M$ 作为 $D_Y$-模的自同态。令
$$M_0=\ker x_n\subset\Gamma_{[Y]}(M),\qquad M_1=\operatorname{coker} x_n=M/x_nM.$$
记 $i$ 为 $Y$ 到 $X$ 的自然包含。如 10.6 所述，$i^+(M)=M_1$，$L^{-1}i^+(M)=M_0$，其余逆像全为零。

考虑双加性映射 $D_X\times M_0\to M$。显然它通过 $D_X\otimes_{D_Y}M_0\to M$ 分解。此外由 $M_0$ 的定义，后一态射在 $D_Xx_n\otimes_{D_Y}M_0$ 在 $D_X\otimes_{D_Y}M_0$ 中的像上为零。如 §11 中所述，
$$D_{X\leftarrow Y}=D_Y\boxtimesext D_Z/D_Zx_n=\bigoplus_{j=0}^{\infty}\partial_n^jD_Y.$$
因此上述态射诱导自然的 $D_X$-模态射
$$i_+(M_0)=D_{X\leftarrow Y}\otimes_{D_Y}M_0\longrightarrow M.$$
显然其像含于 $\Gamma_{[Y]}(M)$。容易检验这实际上是函子 $i_+\circ L^{-1}i^+$ 到 $\Gamma_{[Y]}$ 的态射。

本节的关键结果是如下引理。

\begin{lemma}
态射 $i_+(M_0)\to\Gamma_{[Y]}(M)$ 是 $D_X$-模的同构。
\end{lemma}

\begin{proof}
先证此态射是满射。断言对任意 $p\in\mathbb{N}$，
$$\{m\in M\mid x_n^p m=0\}\subset D_X\cdot M_0.$$
$p=1$ 时显然。若 $p>1$ 且 $x_n^p m=0$，则
$$0=\partial_n(x_n^p m)=x_n^{p-1}(pm+x_n\partial_n m),$$
由归纳假设 $pm+x_n\partial_n m\in D_X\cdot M_0$。此外由归纳假设 $x_n m\in D_X\cdot M_0$。这意味着
$$(p-1)m=pm+[x_n,\partial_n]m=pm+x_n\partial_n m-\partial_n x_n m\in D_X\cdot M_0,$$
故 $m\in D_X\cdot M_0$。因此映射是满射。

现证单射。由上述讨论
$$i_+(M_0)=D_{X\leftarrow Y}\otimes_{D_Y}M_0=\bigoplus_{j=0}^{\infty}\partial_n^j M_0.$$
设 $(m_0,\partial_n m_1,\dots,\partial_n^q m_q,0,\dots)$ 是此直和中映为 $0$ 的非零元，且 $q$ 最小，即
$$m_0+\partial_n m_1+\cdots+\partial_n^q m_q=0.$$
则
$$0=x_n\sum_{j=0}^q\partial_n^j m_j=\sum_{j=1}^q[x_n,\partial_n^j]m_j=-\sum_{j=1}^q j\,\partial_n^{j-1}m_j,$$
矛盾。因此映射的核为零。
\end{proof}

\begin{corollary}
$x_n\Gamma_{[Y]}(M)=\Gamma_{[Y]}(M)$。
\end{corollary}

\begin{proof}
由 12.2，$\Gamma_{[Y]}(M)$ 的任意元素形如 $\sum_{j\in\mathbb{Z}_+}\partial_n^j m_j$，其中 $m_j\in M_0$。另一方面，
$$x_n\sum_{j\in\mathbb{Z}_+}\frac{1}{j+1}\partial_n^{j+1}m_j=-\sum_{j\in\mathbb{Z}_+}\partial_n^j m_j.$$
\end{proof}

\begin{corollary}
设 $M$ 是 $D_X$-模。则
\begin{enumerate}[label=(\roman*)]
  \item $\Gamma_{[Y]}(M)$ 是有限生成 $D_X$-模当且仅当 $M_0$ 是有限生成 $D_Y$-模；
  \item $d(\Gamma_{[Y]}(M))=d(M_0)+1$。
\end{enumerate}
特别地，$\Gamma_{[Y]}(M)$ 是完整的当且仅当 $L^{-1}i^+(M)=M_0$ 是完整的。
\end{corollary}

\begin{proof}
(i) 由 12.2 和 11.3.(iii)，若 $M_0$ 有限生成则 $\Gamma_{[Y]}(M)$ 有限生成。反之设 $\Gamma_{[Y]}(M)$ 是有限生成 $D_X$-模。设 $N_j$（$j\in\mathbb{N}$）是 $M_0$ 的 $D_Y$-子模的递增序列。则它们生成 $\Gamma_{[Y]}(M)$ 的 $D_X$-子模 $i_+(N_j)=\bigoplus_{p=0}^{\infty}\partial_n^p N_j$。由于 $\Gamma_{[Y]}(M)$ 有限生成，递增序列 $i_+(N_j)$ 稳定。此外 $N_j$ 是 $i_+(N_j)$ 中 $x_n$ 的核，故序列 $N_j$ 也必稳定。因此 $M_0$ 有限生成。

(ii) 由 12.2 和 11.3.(iv) 得出。
\end{proof}

\begin{corollary}
设 $M$ 是完整 $D_X$-模。则 $M_0$ 是完整 $D_Y$-模。
\end{corollary}

\begin{proof}
若 $M$ 是完整的，则 $\Gamma_{[Y]}(M)$ 也是完整的。故结论由 12.4 得出。
\end{proof}

设 $\mathcal{M}_Y(D_X)$ 为 $\mathcal{M}(D_X)$ 中由支集含于 $Y$ 的 $D_X$-模构成的全子范畴。以 $\mathcal{M}_{fg,Y}(D_X)$ 和 $\operatorname{Hol}_Y(D_X)$ 分别记支集含于 $Y$ 的有限生成、完整 $D_X$-模的相应子范畴。则由 12.1，对 $\mathcal{M}_Y(D_X)$ 中任意 $M$ 有 $M=\Gamma_{[Y]}(M)$。由 10.6 和 12.3 可见对 $\mathcal{M}_Y(D_X)$ 中任意 $M$ 有 $i^+(M)=0$，故 $L^{-1}i^+$ 是从 $\mathcal{M}_Y(D_X)$ 到 $\mathcal{M}(D_Y)$ 的正合函子。另一方面，$i_+$ 定义反方向的正合函子，且由 12.2 复合 $i_+\circ L^{-1}i^+$ 同构于 $\mathcal{M}_Y(D_X)$ 上的恒等函子。也显然 $L^{-1}i^+\circ i_+$ 同构于 $\mathcal{M}(D_Y)$ 上的恒等函子。

这导出如下基本结果。

\begin{theorem}[Kashiwara]
直像函子 $i_+$ 定义范畴 $\mathcal{M}(D_Y)$（相应地 $\mathcal{M}_{fg}(D_Y)$、$\operatorname{Hol}(D_Y)$）与范畴 $\mathcal{M}_Y(D_X)$（相应地 $\mathcal{M}_{fg,Y}(D_X)$、$\operatorname{Hol}_Y(D_X)$）的等价。其逆为函子 $L^{-1}i^+$。
\end{theorem}

\begin{proof}
只需证明括号中的陈述，它们由 12.4 立即得出。
\end{proof}
\begin{annote}
完整模经真态射的直像仍完整（Kashiwara），这是“完整”一词的稳定性来源：它使完整模在六种运算下有封闭性，从而可以谈论 $D^b_{rh}(X)$ 与 $D^b_c(X^{an})$ 的等价（Riemann–Hilbert）。1 维情形（$\dim X=1$）就是 Deligne 关于正则奇点线性微分方程有解的存在性定理。[来源：nLab 条目 Riemann-Hilbert+correspondence（Last revised 2024-08-09）、holonomic+D-module（2014-02-16）；原书 §12, p51–54]
\end{annote}



\section{完整性的保持}

本节证明直像和逆像保持完整模。先给出一个完整性的简单判别法。

设 $X=k^n$，$Y=k^m$，$F:X\to Y$ 是多项式映射。想研究完整模在逆像和直像函子作用下的行为。

首先再次使用图像构造把问题归约到特殊映射。如 10.10 和 11.7 的证明：
$$\begin{tikzcd}
X \arrow[r,"F"] \arrow[d,"i"' ] & Y \arrow[d,"p"]\\
X\times Y \arrow[r,"\Phi"'] & X\times Y
\end{tikzcd}$$
其中 $i(x)=(x,0)$（$x\in X$），$p(x,y)=y$（$x\in X$，$y\in Y$），$\Phi(x,y)=(x,y+F(x))$（$x\in X$，$y\in Y$）。

由 10.5 知 $p^+$ 是正合的且把完整模映为完整模。由 11.3 知 $i_+$ 是正合的且把完整模映为完整模。此外由 10.8 和 11.6 知 $\Phi^+$ 和 $\Phi_+$ 是正合的且把完整模映为完整模。

因此只需研究 $i^+$ 和 $p_+$ 的导出函子。先讨论浸入 $i:X\to X\times Y$。

\begin{lemma}
设 $N$ 是完整 $D_{X\times Y}$-模。则 $D_X$-模 $L^qi^+(N)$（$q\in\mathbb{Z}$）是完整的。
\end{lemma}

由于由 8.1.(ii)，完整模的子模、商模和扩张都是完整的，如 10.7 的证明用谱序列论证可把证明归约到 $\dim Y=1$ 的情形。此时若以 $y$ 记 $Y$ 上自然坐标，考虑 $D_X$-模态射 $N\xrightarrow{\ y\ }N$，则有 $i^+(N)=\operatorname{coker}y$，$L^{-1}i^+(N)=\ker y$，其余导出逆像全为零（如 10.6 所述）。此外若 $N$ 是完整的，由 12.4 知 $L^{-1}i_+(N)$ 是完整的。故只需处理 $i^+(N)$。

\begin{lemma}
设 $N$ 是完整 $D_{X\times Y}$-模。则 $i^+(N)$ 是完整的。
\end{lemma}

\begin{proof}
令 $\bar N=N/\Gamma_{[X]}(N)$。可考虑短正合列
$$0\longrightarrow\Gamma_{[X]}(N)\longrightarrow N\longrightarrow\bar N\longrightarrow 0.$$
由于 $i^+$ 是右正合的，这导出正合列
$$i^+(\Gamma_{[X]}(N))\longrightarrow i^+(N)\longrightarrow i^+(\bar N)\longrightarrow 0.$$
另一方面，由 12.3 知 $i^+(\Gamma_{[X]}(N))=0$。因此自然映射 $i^+(N)\to i^+(\bar N)$ 是同构。

设 $\bar v\in\Gamma_{[X]}(\bar N)\subset\bar N$，以 $v\in N$ 记 $\bar v$ 的代表。则对充分大的 $p\in\mathbb{Z}_+$ 有 $y^p\bar v=0$。因此 $y^p v\in\Gamma_{[X]}(N)$。这进而蕴含对充分大的 $q\in\mathbb{Z}_+$ 有 $y^{p+q}v=y^q(y^pv)=0$。故 $v\in\Gamma_{[X]}(N)$，$\bar v=0$。由此可知 $\Gamma_{[X]}(\bar N)=0$。

此外，若 $N$ 是完整 $D_{X\times Y}$-模，则 $\bar N$ 是完整 $D_{X\times Y}$-模。

因此可从一开始假设 $\Gamma_{[X]}(N)=0$。这意味着乘以 $y$ 在 $N$ 上是单射，$N$ 嵌入其局部化 $N_y$。考虑正合列
$$0\longrightarrow N\longrightarrow N_y\longrightarrow L\longrightarrow 0.$$
由于 $N$ 是完整 $D_{X\times Y}$-模，由 8.6 知 $N_y$ 是完整的。故 $L$ 是完整 $D_{X\times Y}$-模。由上述讨论，这意味着 $L^{-1}i^+(L)$ 是完整 $D_X$-模。

对短正合列应用 $i$ 的逆像长正合列，得到
$$\cdots\to L^{-1}i^+(N_y)\to L^{-1}i^+(L)\to i^+(N)\to i^+(N_y)\to i^+(L)\to 0.$$
由于乘以 $y$ 在 $N_y$ 上可逆，由 10.6 知
$$i^+(N_y)=L^{-1}i^+(N_y)=0.$$
故 $i^+(N)\cong L^{-1}i^+(L)$。由上述讨论断定 $i^+(N)$ 是完整 $D_X$-模。
\end{proof}

因此，由 10.3 得到如下结果。

\begin{theorem}
设 $F:X\to Y$ 是多项式映射，$M$ 是完整 $D_Y$-模。则 $L^qF^+(M)$（$q\in\mathbb{Z}$）是完整 $D_X$-模。
\end{theorem}

现在研究 $p$ 的直像。

\begin{lemma}
设 $M$ 是完整 $D_{X\times Y}$-模。则 $D_Y$-模 $L^qp_+(M)$（$q\in\mathbb{Z}$）是完整的。
\end{lemma}

\begin{proof}
由于由 8.1.(ii)，完整模的子模、商模和扩张都是完整的，如 11.5 的证明用谱序列论证可把证明归约到 $\dim X=1$ 的情形。此时若以 $\partial$ 记关于 $X$ 上坐标 $x$ 的导数，考虑 $D_Y$-模态射 $M\xrightarrow{\ \partial\ }M$，则有 $p_+(M)=\operatorname{coker}\partial$，$L^{-1}p_+(M)=\ker\partial$，其余导出逆像全为零（如 11.4 所述）。通过应用 Fourier 变换得到计算 $F(L_\cdot p_+(M))$ 的复形
$$\cdots\longrightarrow 0\longrightarrow F(M)\xrightarrow{\ x\ }F(M)\longrightarrow 0\longrightarrow\cdots.$$
由 13.2 证明中的论证，可见此复形计算典范包含 $j:Y\to X\times Y$（$j(y)=(0,y)$，$y\in Y$）的逆像。故由上同调由 13.2 是完整的。由 6.4，对所有 $q\in\mathbb{Z}$ 有 $L^q p_+(M)$ 是完整的。
\end{proof}

因此，由 11.2 得到如下结果。

\begin{theorem}
设 $F:X\to Y$ 是多项式映射，$M$ 是完整 $D_X$-模。则 $L^qF_+(M)$（$q\in\mathbb{Z}$）是完整 $D_Y$-模。
\end{theorem}

\begin{remark}
13.4 和 13.5 关于有限生成模的类似陈述不成立。例如令 $X=\{0\}$，$Y=k$，以 $i:X\to Y$ 记自然包含，则逆像 $i^+(D_Y)$ 是 $k$ 上无限维向量空间。类似地，若 $p$ 是 $Y$ 到一点的投影，则 $p_+(D_Y)$ 是 $k$ 上无限维向量空间。
\end{remark}

\chapter{光滑代数簇上的微分算子层}

\section{代数簇上的微分算子}

设 $X$ 是特征零的代数闭域 $k$ 上的仿射代数簇。以 $\mathcal{O}_X$ 记 $X$ 的结构层，以 $R(X)$ 记其整体截面，即 $X$ 上正则函数环。则 $R(X)$ 是交换 $k$-代数，可如第一章 §5 那样定义 $R(X)$ 上 $k$-线性微分算子环 $D(X)$。称此环为 $X$ 上的\emph{微分算子环}。微分算子的阶数定义 $D(X)$ 上一个递增环滤过 $(D_p(X);\ p\in\mathbb{Z})$，它满足第一章 §3 开头的性质 (i)--(v)。

如第一章所讨论，在 $X=k^n$ 的情形已知 $D(X)=D(n)$ 是 $n$ 变量多项式系数微分算子环。

可把 $X$ 实现为某仿射空间 $k^n$（$n\in\mathbb{Z}_+$）中的闭子集。设 $I(X)$ 为 $k[X_1,X_2,\dots,X_n]$ 中所有在 $X$ 上消失的多项式构成的理想。则 $R(X)=k[X_1,X_2,\dots,X_n]/I(X)$，以 $r$ 记 $k[X_1,X_2,\dots,X_n]$ 到 $R(X)$ 上的限制同态。定义
$$A=\{T\in D(n)\mid T(I(X))\subset I(X)\}.$$
显然 $A$ 是 $D(n)$ 的子代数。

设 $T\in A$。则 $T$ 诱导 $k[X_1,X_2,\dots,X_n]/I(X)$ 的线性自同态 $\varphi(T)$。映射 $\varphi$ 是 $A$ 到 $R(X)$ 的所有线性自同态环的同态。显然 $A$ 配置微分算子阶数滤过是一个滤过环。

此外，$k[X_1,X_2,\dots,X_n]$ 是 $A$ 的子环。因此有如下交换图
$$\begin{tikzcd}
k[X_1,\dots,X_n] \arrow[r,"r"] \arrow[d,hook] & R(X) \arrow[d,hook]\\
A \arrow[r,"\varphi"'] & \End_k(R(X))
\end{tikzcd}$$

特别地，对任意多项式 $P\in k[X_1,\dots,X_n]$，$\varphi(P)$ 是 $X$ 上乘以 $r(P)$。设 $T\in A\cap D_p(n)$，$f_0,f_1,\dots,f_p$ 是 $R(X)$ 中元素的 $(p+1)$-元组。则可取 $P_0,P_1,\dots,P_p\in k[X_1,X_2,\dots,X_n]$ 使得 $r(P_i)=f_i$（$0\le i\le p$）。故
$$[[\cdots[[\varphi(T),f_0],f_1]\cdots,f_{p-1}],f_p]=\varphi([[\cdots[[T,P_0],P_1]\cdots,P_{p-1}],P_p])=0.$$
因此 $\varphi(T)$ 是 $X$ 上阶数 $\le p$ 的微分算子。由此可知 $\varphi:A\to D(X)$ 是与微分算子阶数滤过相容的环同态。

此外，
$$J(X)=\{T\in D(n)\mid T(k[X_1,X_2,\dots,X_n])\subset I(X)\}$$
是 $A$ 的双边理想。显然 $J(X)$ 含于 $\varphi$ 的核中。

\begin{lemma}
设 $T\in D(n)$。则下列条件等价：
\begin{enumerate}[label=(\roman*)]
  \item $T\in J(X)$；
  \item $T=\sum_I P_I\partial^I$，其中 $P_I\in I(X)$。
\end{enumerate}
\end{lemma}

\begin{proof}
若 $T$ 的系数 $P_I$ 在 $X$ 上消失，显然微分算子 $T$ 属于 $J(X)$。反之若 $T\in J(X)$，则 $P_0=T(1)$ 在 $X$ 上消失，即 $P_0\in I(X)$。假设 $|I|<m$ 时 $P_I\in I(X)$。则 $T'=\sum_{|I|<m}P_I\partial^I$ 属于 $J(X)$。因此 $T''=T-T'\in J(X)$。另一方面，对任意 $J\in\mathbb{Z}_+^n$，$|J|=m$，
$$T''(X^J)=\Bigl(\sum_{|I|\ge m}P_I\partial^I\Bigr)(X^J)=J!P_J$$
在 $X$ 上消失，即 $P_J\in I(X)$。故由 $m$ 上的归纳断定对所有 $I\in\mathbb{Z}_+^n$ 有 $P_I\in I(X)$。
\end{proof}

记 $D$ 为商环 $A/J(X)$。$A$ 按微分算子阶数的滤过诱导 $D$ 上的商环滤过 $(D_p;\ p\in\mathbb{Z})$，满足 I.3 开头的条件 (i)--(v)。

由于 $J(X)$ 在 $\varphi$ 的核中，它定义 $D$ 到 $D(X)$ 的同态 $\Phi$。显然 $\Phi$ 是与环滤过相容的 $D$ 到 $D(X)$ 的同态。

\begin{proposition}
态射 $\Phi:D\to D(X)$ 是滤过环的同构。
\end{proposition}

先证 $\Phi$ 是单射。设 $T\in A$ 使 $\varphi(T)=0$。这意味着 $\varphi(T)(P+I(X))=T(P)+I(X)=0$，即对任意 $P\in k[X_1,X_2,\dots,X_n]$ 有 $T(P)\in I(X)$。故 $T\in J(X)$，$\Phi$ 是单射。

为证满射，先作如下注记。

\begin{lemma}
设 $p\in\mathbb{Z}_+$，$P_I\in k[X_1,X_2,\dots,X_n]$（$I\in\mathbb{Z}_+^n$，$|I|\le p$）。则存在阶数 $\le p$ 的微分算子 $T\in D(n)$，使得对所有 $|I|\le p$ 的 $I\in\mathbb{Z}_+^n$ 有 $T(X^I)=P_I$。
\end{lemma}

\begin{proof}
显然 $p=0$ 时成立。设 $p>0$ 且结论对 $p-1$ 成立。由归纳假设存在阶数 $\le p-1$ 的微分算子 $T'$，使得对所有 $|I|\le p-1$ 有 $T'(X^I)=P_I$。对所有 $|I|=p$ 的 $I$，记 $T'(X^I)=Q_I\in k[X_1,X_2,\dots,X_n]$。显然 $\partial^J$（$|J|=p$）零化 $X^I$（$|I|\le p-1$），且当 $|I|=|J|=p$ 时 $\partial^J(X^I)=I!\delta_{I,J}$。这意味着若定义
$$T''=\sum_{|J|=p}\frac{P_J-Q_J}{J!}\partial^J,$$
则 $T''$ 零化 $X^I$（$|I|\le p-1$），且
$$T''(X^I)=\Bigl(\sum_{|J|=p}\frac{P_J-Q_J}{J!}\partial^J\Bigr)(X^I)=P_I-Q_I.$$
因此对 $|I|=p$ 有 $(T'+T'')(X^I)=P_I$，对 $|I|<p$ 有 $(T'+T'')(X^I)=T'(X^I)=P_I$。
\end{proof}

现断言：对任意 $T\in D_p(X)$ 和阶数 $\le p$ 的 $S\in D(n)$，若对 $|I|\le p$ 的 $I\in\mathbb{Z}_+^n$ 有 $T(r(X^I))=r(S(X^I))$，则 $T\circ r=r\circ S$。$p=0$ 时无需证明。设 $p>0$。则对 $1\le j\le n$，
\begin{align*}
[T,r(X_j)](r(X^I))&=T r(X_jX^I)-r(X_j)T(r(X^I))\\
&=r(S(X_jX^I))-r(X_jS(X^I))=r([S,X_j](X^I))
\end{align*}
对所有 $|I|\le p-1$ 成立，且 $[T,r(X_j)]$ 和 $[S,X_j]$ 的阶数 $\le p-1$。故由归纳假设 $[T,r(X_j)]\circ r=r\circ[S,X_j]$。特别地对任意 $I\in\mathbb{Z}_+^n$，
\begin{align*}
T(r(X_jX^I))&=[T,r(X_j)](r(X^I))+r(X_j)T(r(X^I))\\
&=r([S,X_j](X^I))+r(X_j)T(r(X^I))\\
&=r(S(X_jX^I))+r(X_j)\bigl[T(r(X^I))-r(S(X^I))\bigr].
\end{align*}
故由 $|I|$ 上的归纳得对所有 $I\in\mathbb{Z}_+^n$ 有 $T(r(X^I))=r(S(X^I))$，这证明断言。

设 $T$ 是 $X$ 上阶数 $\le p$ 的微分算子。则可选 $P_I\in k[X_1,X_2,\dots,X_n]$（$|I|\le p$）使得对所有 $|I|\le p$ 有 $T(r(X^I))=r(P_I)$。由上一条引理存在阶数 $\le p$ 的微分算子 $S\in D(n)$ 使得对所有 $|I|\le p$ 有 $S(X^I)=P_I$。这意味着对 $|I|\le p$ 有 $T(r(X^I))=r(S(X^I))$。

由前面的结果得 $T\circ r=r\circ S$。特别地 $r(S(I(X)))=T(r(I(X)))=0$，即 $S\in A$。显然 $\varphi(S)=T$，$\Phi$ 是满射。这完成 1.2 的证明。

\begin{corollary}
设 $X$ 是仿射代数簇。对任意 $p\in\mathbb{Z}_+$，$D_p(X)$ 作为 $R(X)$-模（对左乘和右乘）是有限生成的。
\end{corollary}

\begin{proof}
可设 $X$ 是 $k^n$ 中的闭子集。由 I.8.9 知 $X=k^n$ 时陈述成立。由于 $k[X_1,X_2,\dots,X_n]$ 是诺特环，$A\cap F_pD(n)$ 作为 $k[X_1,X_2,\dots,X_n]$-模（对左乘和右乘）有限生成，且对任意 $p\in\mathbb{Z}$，$D_p$ 作为 $R(X)$-模（对左乘和右乘）有限生成。结论由 1.2 得出。
\end{proof}

设 $f\in R(X)$，$f\neq 0$，$X_f=\{x\in X\mid f(x)\neq 0\}$ 是 $X$ 中相应的主开集。则 $X_f$ 是仿射簇，且 $R(X_f)=R(X)_f$。以 $r_f$ 记 $R(X)$ 到 $R(X_f)$ 的限制映射。

\begin{proposition}
设 $T\in D_p(X)$。则存在唯一的微分算子 $\bar T\in D_p(X_f)$ 使得如下交换图交换：
$$\begin{tikzcd}
R(X) \arrow[r,"T"] \arrow[d,"r_f"' ] & R(X) \arrow[d,"r_f"]\\
R(X_f) \arrow[r,"\bar T"'] & R(X_f)
\end{tikzcd}$$
\end{proposition}

先证 $\bar T$ 的唯一性。只需证明如下引理。

\begin{lemma}
设 $S\in D(X_f)$ 使得对任意 $g\in r_f(R(X))$ 有 $S(g)=0$。则 $S=0$。
\end{lemma}

\begin{proof}
对 $S$ 的阶数 $p$ 归纳证明。若 $p=0$，$S\in R(X_f)$，条件立即导致 $S=0$。现设 $p>0$。则 $S'=[S,f]\in D_{p-1}(X_f)$ 且它零化 $r_f(R(X))$。故由归纳假设 $S'=0$。这意味着 $S$ 与 $f$ 可交换。取 $h\in R(X_f)$。则存在 $n\in\mathbb{Z}_+$ 使得 $f^nh\in r_f(R(X))$。这意味着 $f^nS(h)=S(f^nh)=0$。由于 $f^{-1}\in R(X_f)$，断定 $S(h)=0$。因此 $S=0$。
\end{proof}

只需证明 $\bar T$ 的存在性。

先讨论 $X=k^n$ 的情形。由于 $D(n)$ 作为 $k$-代数由 $X_i,\partial_i$（$1\le i\le n$）生成，只需对 $T=\partial_i$（$1\le i\le n$）证明 $\bar T$ 的存在性。但导子 $\partial_i$ 扩充到有理函数域 $k(X_1,X_2,\dots,X_n)$ 且满足
$$\partial_i\Bigl(\frac{g}{f^m}\Bigr)=\frac{\partial_i(g)f-mg\,\partial_i(f)}{f^{m+1}}$$
对任意 $g\in k[X_1,X_2,\dots,X_n]$ 和 $m\in\mathbb{Z}_+$。因此它们诱导 $k[X_1,X_2,\dots,X_n]_f$ 的导子。这完成 $D(n)$ 存在性的证明。

只需证明一般情形下 $\bar T$ 的存在性。可设 $X$ 作为闭子集嵌入某 $k^n$。设 $P\in k[X_1,X_2,\dots,X_n]$ 是在 $X$ 上限制为 $f$ 的多项式，以 $U$ 记 $k^n$ 中 $P$ 的零点集的补集（仿射开集）。则 $X\cap U=X_f$。由 1.2 可找到 $S\in A\cap F_pD(n)$ 使得 $\varphi(S)=T$。此微分算子扩充为 $U$ 上阶数 $\le p$ 的微分算子 $\bar S$。

\begin{lemma}
设 $S\in A$。则 $\bar S$ 把 $I(X)_P$ 映入自身。
\end{lemma}

\begin{proof}
对 $S$ 的阶数 $p$ 归纳证明。$p=0$ 时显然。设 $p>0$。则 $S'=[S,P]\in A$ 且阶数 $\le p-1$。故由归纳假设 $\bar S'$ 把 $I(X)_P$ 映入自身。取 $Q\in I(X)$。则
$$\bar S\Bigl(\frac{Q}{P^m}\Bigr)=\bar S'\Bigl(\frac{Q}{P^{m+1}}\Bigr)+P\,\bar S\Bigl(\frac{Q}{P^{m+1}}\Bigr)$$
且由归纳假设
$$\bar S\Bigl(\frac{Q}{P^{m+1}}\Bigr)-P^{-1}\bar S\Bigl(\frac{Q}{P^{m}}\Bigr)\in I(X)_P$$
对任意 $m\in\mathbb{Z}_+$。由 $m$ 上的归纳得对任意 $m\in\mathbb{Z}_+$ 有 $\bar S\frac{Q}{P^m}\in I(X)_P$。
\end{proof}

因此 $\bar S$ 诱导 $k[X_1,X_2,\dots,X_n]_P/I(X)_P=R(X)_f=R(X_f)$ 的线性自同态。如 1.2 之前的讨论，这实际上是 $X_f$ 上的微分算子。此外在 $r_f(R(X))$ 上它与 $T$ 一致。因此构造了 $\bar T$。这完成 1.5 的证明。

设 $X$ 是仿射代数簇，$f\in R(X)$，$f\neq 0$。则由 1.5 有良定义的限制映射 $\rho_f:D(X)\to D(X_f)$。1.5 的唯一性部分蕴含 $\rho_f$ 是环同态，故有如下结果。

\begin{proposition}
映射 $\rho_f:D(X)\to D(X_f)$ 是滤过环的同态。
\end{proposition}

特别地，$\rho_f$ 是 $R(X)$-模（对左乘和右乘）的同态。

\begin{lemma}
设 $D(X)_f$ 是把 $D(X)$ 看作 $R(X)$-模（对左乘）的局部化。则态射 $\rho_f$ 诱导 $D(X)_f$ 到 $D(X_f)$ 上的同构 $\beta_f$。
\end{lemma}

\begin{proof}
先设 $X_f$ 在 $X$ 中稠密。此时自然映射 $r_f:R(X)\to R(X_f)$ 是单射。故 $\beta_f(\frac{1}{f^m}T)=\frac{1}{f^m}\bar T=0$（对某 $T\in D(X)$）蕴含对任意 $g\in R(X)$ 存在 $s\in\mathbb{Z}_+$ 使 $f^sT(g)=0$。因此对所有 $g\in R(X)$ 有 $T(g)=0$，即 $T=0$。由此可知 $\beta_f$ 是单射。

为证 $\beta_f$ 满射，只需证明对任意 $T\in D(X_f)$ 存在 $m\in\mathbb{Z}_+$ 使 $(f^mT)(R(X))\subset R(X)$。对 $T$ 的阶数 $p$ 归纳证明。$p=0$ 时显然。

设 $p>0$。记 $g_1,g_2,\dots,g_n$ 为 $k$-代数 $R(X)$ 的生成元。由归纳假设，存在 $m\in\mathbb{Z}_+$ 使得 $f^m[T,g_i](R(X))\subset R(X)$（$1\le i\le n$），且 $f^mT(1)\in R(X)$。这意味着若 $h\in R(X)$ 满足 $f^mT(h)\in R(X)$，则
$$f^mT(g_ih)=f^m[T,g_i](h)+f^m g_iT(h)\in R(X).$$
利用 $f^mT(1)\in R(X)$ 并对单项式 $g_1^{i_1}g_2^{i_2}\cdots g_n^{i_n}$ 的长度归纳，可见此关系蕴含 $f^mT(R(X))\subset R(X)$。故此时 $\beta_f$ 是同构。

现考虑一般情形。假设 $X_f$ 在 $X$ 中不稠密。断言此时存在 $f'\in R(X)$ 使得 $X_f$ 与 $X_{f'}$ 不交且它们的并在 $X$ 中稠密。首先，若 $X_f$ 在 $X$ 中不稠密，可找到 $a_1\in R(X)$ 使 $a_1\neq 0$ 且在 $X_f$ 上消失。这意味着 $X_f$ 与 $X_{a_1}$ 不交且 $X_{f+a_1}=X_f\cup X_{a_1}$。若 $X_{f+a_1}$ 在 $X$ 中不稠密可重复此构造，由于 $X$ 是诺特拓扑空间，经有限步构造序列 $a_1,a_2,\dots,a_s$，使得 $X_f,X_{a_1},X_{a_2},\dots,X_{a_s}$ 是 $X$ 中两两不交的主开集且并在 $X$ 中稠密。若令 $f'=a_1+a_2+\cdots+a_s$，它显然具有所需性质。

现断言 $D(X_{f+f'})=D(X_f)\oplus D(X_{f'})$。显然 $R(X_{f+f'})=R(X_f)\oplus R(X_{f'})$。设 $\chi,\chi'\in R(X_{f+f'})$ 分别为 $X_f$ 和 $X_{f'}$ 的特征函数。断言对任意 $T\in D_p(X_{f+f'})$ 有 $[T,\chi]=[T,\chi']=0$。$T$ 阶数 $\le 0$ 时成立。对 $T$ 的阶数归纳。设 $p>0$。若 $T$ 的阶数 $\le p$，由归纳假设断言对 $[T,\chi]$ 成立。因此
$$0=[[T,\chi],\chi']=[T,\chi]\chi'-\chi'[T,\chi]=-\chi T\chi'-\chi'T\chi$$
即 $\chi T\chi'=-\chi'T\chi$。右乘 $\chi$ 得 $\chi T\chi'=\chi'T\chi=0$。因此 $T\chi=(\chi+\chi')T\chi=\chi T\chi$，类似地 $\chi'T=\chi'T(\chi+\chi')=\chi'T\chi'$。由对称性也有 $\chi T=\chi T\chi$，最终导致 $[T,\chi]=0$。因此对 $g\in R(X_f)$ 有 $T(g)=T(\chi g)=\chi T(g)$，$T(R(X_f))\subset R(X_f)$。用 $\chi'$ 的类似论证得 $T(R(X_{f'}))\subset R(X_{f'})$。因此 $T$ 诱导 $X_f$ 和 $X_{f'}$ 上的微分算子 $S$ 和 $S'$，且 $T=S\oplus S'$。

由证明的第一部分，$\beta_{f+f'}:D(X)_{f+f'}\to D(X_f)\oplus D(X_{f'})$ 是同构。关于 $\chi$ 局部化得 $\beta_f:D(X)_f\to D(X_f)$ 是同构。
\end{proof}

设 $U$ 是 $X$ 中开集。以 $\mathcal{P}_U$ 记含于 $U$ 的所有主开集按包含排序的族。若 $V,W\in\mathcal{P}_U$ 且 $V\subset W$，有自然环同态 $r_{VW}:D(W)\to D(V)$。显然 $(D(V);r_{VW})$ 是环的逆系统。记其逆极限为 $D(U)$。显然 $\mathcal{D}_X:U\longmapsto D(U)$ 是 $X$ 上环的预层。由 1.9，这是左乘 $\mathcal{O}_X$-模的层。这蕴含如下结果。

\begin{proposition}
$\mathcal{D}_X$ 是 $X$ 上的环层。
\end{proposition}

称 $\mathcal{D}_X$ 为 $X$ 上\emph{局部微分算子层}。

\begin{theorem}
设 $X$ 是仿射簇，$\mathcal{D}_X$ 是 $X$ 上局部微分算子层。则对任意仿射开子集 $U\subset X$ 有 $\Gamma(U,\mathcal{D}_X)=D(U)$。
\end{theorem}

\begin{proof}
$U$ 是 $X$ 的主开集时显然。设 $U$ 是 $X$ 的任意仿射开子集。取 $f\in R(X)$ 使 $X_f\subset U$。则若记 $g=f|_U$，有 $U_g=X_f$。这意味着
$$\Gamma(U_g,\mathcal{D}_U)=D(U_g)=D(X_f)=\Gamma(X_f,\mathcal{D}_X).$$
此外这些同构与限制态射相容。由于主开集 $\{X_f\mid f\in R(X)\}$ 构成 $X$ 拓扑的基，含于 $U$ 的那些构成 $U$ 拓扑的基。此外，由于 $\mathcal{D}_X|_U$ 和 $\mathcal{D}_U$ 是 $U$ 上的层且在拓扑基上一致，它们相等。这意味着 $\Gamma(U,\mathcal{D}_X)=\Gamma(U,\mathcal{D}_U)=D(U)$。
\end{proof}

设 $X$ 是 $k$ 上任意代数簇。对 $X$ 中任意开集 $U$，以 $\mathcal{B}_U$ 记 $U$ 中所有仿射开子集按包含排序的族。若 $V,W\in\mathcal{B}_U$ 且 $V\subset W$，有自然环同态 $r_{VW}:D(W)\to D(V)$。显然 $(D(V);r_{VW})$ 是环的逆系统。记其逆极限为 $D(U)$。同样 $\mathcal{D}_X:U\longmapsto D(U)$ 是 $X$ 上环的预层。

\begin{proposition}
设 $X$ 是 $k$ 上的代数簇。则 $\mathcal{D}_X$ 是 $X$ 上的环层。
\end{proposition}

此结果以及 1.10 是如下引理的特例。设 $\mathcal{C}$ 是这样的范畴：其中任意对象的逆系统都有逆极限。设 $X$ 是拓扑空间，$\mathcal{B}$ 是 $X$ 拓扑的开集基。称 $\mathcal{B}$ 上取值于 $\mathcal{C}$ 的预层 $\mathcal{F}$ 为对象族 $\mathcal{F}(U)$（$U\in\mathcal{B}$）和态射族 $\rho_{VU}:\mathcal{F}(V)\to\mathcal{F}(U)$（对任意满足 $U\subset V$ 的对 $(U,V)$ 定义），满足
\begin{enumerate}[label=(\roman*)]
  \item 对任意 $U\in\mathcal{B}$，$\rho^U_U=\operatorname{id}$；
  \item 对任意 $U,V,W\in\mathcal{B}$ 且 $U\subset V\subset W$，$\rho^W_U=\rho^V_U\circ\rho^W_V$。
\end{enumerate}
则可通过令 $\mathcal{F}'(U)$ 等于所有满足 $V\subset U$ 的 $V\in\mathcal{B}$ 的 $\mathcal{F}(V)$ 的逆极限，定义 $X$ 上的预层 $\mathcal{F}'$。此外对任意 $U\in\mathcal{B}$ 有 $\mathcal{F}'(U)=\mathcal{F}(U)$。

\begin{lemma}
若 $\mathcal{F}$ 满足如下条件，则 $X$ 上的预层 $\mathcal{F}'$ 是 $X$ 上的层：

(F) 对 $U\in\mathcal{B}$ 的任意由 $U_i\in\mathcal{B}$ 构成的覆盖 $(U_i;\ i\in I)$ 以及任意对象 $T\in\mathcal{C}$，把 $f\in\operatorname{Hom}(T,\mathcal{F}(U))$ 映为族 $(\rho^{U_i}_U\circ f)_{i\in I}\in\prod_{i\in I}\operatorname{Hom}(T,\mathcal{F}(U_i))$ 的映射，是 $\operatorname{Hom}(T,\mathcal{F}(U))$ 到所有满足如下条件的族 $(f_i;\ i\in I)$ 的集合的双射：对任意指标对 $(i,j)$ 和任意 $V\in\mathcal{B}$ 且 $V\subset U_i\cap U_j$，有 $\rho^{V}_{U_j}\circ f_i=\rho^V_{U_i}\circ f_j$。
\end{lemma}

\begin{proof}
设 $\mathcal{B}'$ 是含于 $\mathcal{B}$ 的 $X$ 拓扑的另一个基。则可定义另一个预层 $\mathcal{F}''$，它由 $\mathcal{B}'$ 上 $\mathcal{F}$ 定义的预层附着而来。由逆极限的定义，对 $X$ 上任意开集 $U$ 存在 $\mathcal{F}'(U)$ 到 $\mathcal{F}''(U)$ 的典范态射。若 $U\in\mathcal{B}$，此典范态射是从 $\mathcal{F}(U)$ 到 $\mathcal{F}''(U)$ 的态射。断言此态射是同构。事实上由条件 (F)，$\mathcal{F}''(U)$ 到 $\mathcal{F}(V)$（$V\in\mathcal{B}'$，$V\subset U$）的典范态射通过 $\mathcal{F}(U)$ 分解。此外由泛性质，典范态射 $\mathcal{F}''(U)\to\mathcal{F}(U)$ 和前述态射 $\mathcal{F}(U)\to\mathcal{F}''(U)$ 两个顺序的复合都是恒等态射。这证明断言。另一方面，这也蕴含对 $X$ 中任意开集 $U$，态射 $\mathcal{F}''(U)\to\mathcal{F}''(V)=\mathcal{F}(V)$（$V\in\mathcal{B}$，$V\subset U$）满足逆系统 $(F(V);\ V\subset U,\ V\in\mathcal{B})$ 逆极限的条件，故 $\mathcal{F}''(U)=\mathcal{F}'(U)$。

现设 $U$ 是 $X$ 的开子集，$(U_i)$ 是 $U$ 的开覆盖，$\mathcal{B}'$ 是 $\mathcal{B}$ 中含于至少一个 $U_i$ 的元素构成的子族。显然 $\mathcal{B}'$ 是 $U$ 拓扑的基，故 $\mathcal{F}'(U)$（相应地 $\mathcal{F}'(U_i)$）是 $V\in\mathcal{B}'$ 且 $V\subset U$（相应地 $V\subset U_i$）的 $\mathcal{F}(V)$ 的逆极限。由逆极限的定义现在得出 $\mathcal{F}'$ 是层。
\end{proof}

设 $U$ 是 $X$ 中开集，$T\in\mathcal{D}_X(U)$。若对任意仿射开集 $V\subset U$，微分算子 $r^V_U(T)$ 阶数 $\le p$，则称 $T$ \emph{阶数 $\le p$}。这在 $\mathcal{D}_X(U)$ 上定义递增滤过 $F\mathcal{D}_X(U)$。

\begin{lemma}
$\mathcal{D}_X(U)$ 上的滤过 $F\mathcal{D}_X(U)$ 是穷尽的。
\end{lemma}

\begin{proof}
设 $T\in\mathcal{D}_X(U)$。由于 $U$ 拟紧，可找到由仿射开集构成的 $U$ 的有限开覆盖 $(U_i;\ 1\le i\le s)$。取 $p\in\mathbb{Z}$ 使 $T$ 在覆盖各元素上的限制阶数 $\le p$。设 $V$ 是 $U$ 的任意仿射开子集，$S=r^V_U(T)$。断言 $S$ 阶数 $\le p$。取 $f_0,f_1,\dots,f_p\in R(V)$。则 $R=[\cdots[[S,f_0],f_1],\dots,f_p]$ 是 $V$ 上的微分算子，其在 $V\cap U_i$（$1\le i\le s$）上的限制均为零。这意味着 $R=0$，$S$ 阶数 $\le p$。
\end{proof}

因此此滤过满足 I.3 开头的性质 (i)--(v)。显然这样得到 $X$ 上局部微分算子层 $\mathcal{D}_X$ 由 $k$-向量子层构成的滤过 $F\mathcal{D}_X$。称之为\emph{按微分算子阶数的滤过}。在 $X$ 上任意仿射开集 $U$ 上有 $F_p\mathcal{D}_X(U)=D_p(U)$（$p\in\mathbb{Z}$）。因此可考虑分次环层 $\Gr\mathcal{D}_X$。它是交换环层且 $\Gr_0\mathcal{D}_X=\mathcal{O}_X$。

\begin{theorem}
设 $X$ 是 $k$ 上的代数簇。则
\begin{enumerate}[label=(\roman*)]
  \item 层 $\mathcal{D}_X$ 是拟凝聚 $\mathcal{O}_X$-模（对左乘和右乘）；
  \item 层 $F_p\mathcal{D}_X$（$p\in\mathbb{Z}$）是凝聚 $\mathcal{O}_X$-模（对左乘和右乘）；
  \item 层 $\Gr_p\mathcal{D}_X$（$p\in\mathbb{Z}$）是凝聚 $\mathcal{O}_X$-模。
\end{enumerate}
\end{theorem}

\begin{proof}
由于断言是局部的，可设 $X$ 是仿射的。对左乘，(i) 由前述结果得出，(ii) 由 1.3 得出，(iii) 由 (ii) 得出。由于 $\Gr_p\mathcal{D}_X$ 上左乘和右乘定义相同的 $\mathcal{O}_X$-模结构，(iii) 成立。另一方面，
$$0\longrightarrow F_{p-1}\mathcal{D}_X\longrightarrow F_p\mathcal{D}_X\longrightarrow\Gr_p\mathcal{D}_X\longrightarrow 0$$
对右乘和左乘都是正合列，故由 $p$ 上的归纳得右乘的 (ii)。由于 $\mathcal{D}_X$ 是 $F_p\mathcal{D}_X$（$p\in\mathbb{Z}$）的直极限，(i) 成立。
\end{proof}

设 $X$ 是代数簇。对任意仿射开集 $U\subset X$ 记 $T_X(U)=\Der_k(R(U))$。由 I.8.2.(iii) 有 $D_1(U)=R(U)\oplus T_X(U)$。设 $V$ 是 $U$ 的仿射开子集。则对任意 $T\in T_X(U)$ 有 $r^V_U(T)(1)=T(1)=0$，故 $r^V_U(T)\in T_X(V)$，且限制映射与此直和分解相容。这意味着赋值 $U\longmapsto T_X(U)$ 定义 $X$ 中所有仿射开集构成的基 $\mathcal{B}$ 上的预层。记 $X$ 上相应的预层为 $\mathcal{T}_X$。由于 $F_1\mathcal{D}_X=\mathcal{O}_X\oplus\mathcal{T}_X$ 且 $F_1\mathcal{D}_X$ 和 $\mathcal{O}_X$ 都是层，$\mathcal{T}_X$ 是 $X$ 上的层。称之为 $X$ 的\emph{切层}。其在开集 $U\subset X$ 上的局部截面称为 $U$ 上的\emph{局部向量场}。显然它有自然的 $\mathcal{O}_X$-模结构，作为这样的模它同构于 $\Gr_1\mathcal{D}_X$。由前面的 (iii) 得出如下结果。

\begin{proposition}
设 $X$ 是 $k$ 上的代数簇。则
\begin{enumerate}[label=(\roman*)]
  \item $X$ 的切层 $\mathcal{T}_X$ 是凝聚 $\mathcal{O}_X$-模；
  \item $F_1\mathcal{D}_X=\mathcal{O}_X\oplus\mathcal{T}_X$。
\end{enumerate}
\end{proposition}

显然若 $T,T'$ 是 $U$ 上两个向量场，其换位子 $[T,T']$ 是 $U$ 上的向量场。因此 $\mathcal{T}_X$ 是 $k$ 上 Lie 代数的层。
\begin{annote}
相对微分算子环 $\mathcal D_{X/S}$ 的一阶部分正是 $\Omega_{X/S}=\Omega_{S/\mathbb Z}\otimes\mathcal O_X$ 的对偶：$\mathcal D_{X/S,\le1}=\mathcal O_X\oplus\Omega_{X/S}^{\vee}$ 只在光滑时成立（一般时左端是 $\mathcal O_X\oplus\mathcal T_{X/S}$ 的子层）。记 $\Omega_{S/R}$ 的定义与“模 $R$ 的微分”全称。[来源：Stacks Project Tag 07BK（Definition 10.131.2, module of Kähler differentials）；原书 §1, p55–62]
\end{annote}



\section{代数簇的光滑点}

设 $X$ 是特征零代数闭域 $k$ 上的代数簇。以 $\mathcal{O}_X$ 记其结构层。设 $x\in X$，以 $\mathcal{O}_x=\mathcal{O}_{X,x}$ 记 $\mathcal{O}_X$ 在 $x$ 处的茎。则 $\mathcal{O}_x$ 是诺特局部环，其极大理想 $\mathfrak{m}_x$ 由在 $x$ 处消失的函数芽构成。

\begin{lemma}
设 $x\in X$。则 $d(\mathcal{O}_{X,x})=\dim_x X$。
\end{lemma}

\begin{proof}
由于断言是局部的，可设簇 $X$ 是某 $k^n$ 的闭子集。则限制映射定义 $k[X_1,X_2,\dots,X_n]$ 到 $R(X)$ 上的满同态，其核 $I$ 由所有在 $X$ 上消失的多项式构成。可考虑 $k[X_1,X_2,\dots,X_n]$-模的正合列
$$0\longrightarrow I\longrightarrow k[X_1,X_2,\dots,X_n]\longrightarrow R(X)\longrightarrow 0$$
及其在 $x$ 处的局部化
$$0\longrightarrow I_x\longrightarrow k[X_1,X_2,\dots,X_n]_x\longrightarrow R(X)_x\longrightarrow 0.$$
这把 $\mathcal{O}_{k^n,x}=k[X_1,X_2,\dots,X_n]_x$ 模 $I_x$ 的商等同于 $\mathcal{O}_{X,x}=R(X)_x$。此外若以 $\mathcal{M}_x$ 记 $k[X_1,X_2,\dots,X_n]$ 中由 $X_i-x_i$（$1\le i\le n$）生成的极大理想，可见商态射把其局部化 $(\mathcal{M}_x)_x$ 映到 $\mathfrak{m}_x$ 上。这意味着 $\mathcal{O}_{X,x}$ 上的滤过 $(\mathfrak{m}_x^p;\ p\in\mathbb{Z}_+)$ 与 $\mathcal{O}_{X,x}$ 作为 $\mathcal{O}_{k^n,x}$-模的滤过 $((\mathcal{M}_x)^p_x\mathcal{O}_{X,x};\ p\in\mathbb{Z}_+)$ 一致。因此局部环 $\mathcal{O}_{X,x}$ 的维数 $d(\mathcal{O}_{X,x})$ 等于模 $\mathcal{O}_{X,x}$ 在局部环 $\mathcal{O}_{k^n,x}$ 上的维数。由 I.4.2 和 I.4.6 得
$$d(\mathcal{O}_{X,x})=\dim_x(\operatorname{supp}(R(X)))=\dim_x V(I)=\dim_x X.$$
\end{proof}

称向量空间 $T_x^*(X)=\mathfrak{m}_x/\mathfrak{m}_x^2$ 为 $X$ 在 $x$ 处的\emph{余切空间}，其线性对偶 $T_x(X)$ 为 $X$ 在 $x$ 处的\emph{切空间}。

\begin{proposition}
设 $x\in X$。则切空间 $T_x(X)$ 有限维且
$$\dim_k T_x(X)\ge\dim_x X.$$
\end{proposition}

\begin{proof}
这由前面的结果及 I.2.8 立即得出。
\end{proof}

设 $f\in\mathcal{O}_x$。则 $f-f(x)\in\mathfrak{m}_x$，以 $df(x)$ 记其在 $T_x^*(X)$ 中的像。

\begin{lemma}
线性映射 $d:\mathcal{O}_x\to T_x^*(X)$ 满足
$$d(fg)(x)=f(x)\,dg(x)+g(x)\,df(x)$$
对任意 $f,g\in\mathcal{O}_x$。
\end{lemma}

\begin{proof}
有
\begin{align*}
d(fg)(x)&=fg-f(x)g(x)+\mathfrak{m}_x^2\\
&=fg-f(x)g(x)-(f-f(x))(g-g(x))+\mathfrak{m}_x^2\\
&=g(x)(f-f(x))+f(x)(g-g(x))+\mathfrak{m}_x^2\\
&=f(x)\,dg(x)+g(x)\,df(x).
\end{align*}
\end{proof}

例如，若 $X=k^n$，则对任意芽 $f\in k[X_1,X_2,\dots,X_n]_x$ 有
$$df(x)=\sum_{i=1}^n(\partial_i f)(x)\,dX_i(x),$$
且 $(dX_1(x),dX_2(x),\dots,dX_n(x))$ 构成 $T_x^*(k^n)$ 的一组基。因此把向量 $(\xi_1,\xi_2,\dots,\xi_n)\in k^n$ 映为切向量 $f\longmapsto\sum_i\xi_i(\partial_i f)(x)$ 的映射是 $k^n$ 到 $T_x(k^n)$ 的同构。

设 $X$ 和 $Y$ 是 $k$ 上两个代数簇，$\varphi$ 是 $X$ 到 $Y$ 的态射。则对任意 $x\in X$，它诱导态射 $\varphi_x:\mathcal{O}_{Y,\varphi(x)}\to\mathcal{O}_{X,x}$，定义为 $\varphi_x(f)=f\circ\varphi$（$f\in\mathcal{O}_{Y,\varphi(x)}$）。显然 $\varphi_x(\mathfrak{m}_{\varphi(x)})\subset\mathfrak{m}_x$，这蕴含 $\varphi_x(\mathfrak{m}^2_{\varphi(x)})\subset\mathfrak{m}_x^2$，得到线性映射 $T^*_x(\varphi):T^*_{\varphi(x)}(Y)\to T^*_x(X)$。若 $f\in\mathcal{O}_{Y,\varphi(x)}$，有
$$T^*_x(\varphi)(df(\varphi(x)))=d(\varphi_x(f))(x)=d(f\circ\varphi)(x).$$
$T^*_x(\varphi)$ 的转置 $T_x(\varphi):T_x(X)\to T_{\varphi(x)}(Y)$ 称为 $\varphi$ 在 $x$ 处的\emph{切线性映射}。设 $\xi\in T_x(X)$，$f\in\mathcal{O}_{Y,\varphi(x)}$，则
$$\bigl(T_x(\varphi)(\xi)\bigr)(df(\varphi(x)))=\xi\bigl(T^*_x(\varphi)(df(\varphi(x)))\bigr)=\xi(d(f\circ\varphi)(x)).$$

\begin{lemma}\ 
\begin{enumerate}[label=(\roman*)]
  \item 设 $X,Y,Z$ 是代数簇，$\alpha:X\to Y$，$\beta:Y\to Z$ 是代数簇的态射。设 $x\in X$。则
  $$T_x(\beta\circ\alpha)=T_{\alpha(x)}(\beta)\circ T_x(\alpha).$$
  \item 设 $Y$ 是 $X$ 的子簇，$j:Y\to X$ 是典范包含。则对任意 $y\in Y$，$T_y(j):T_y(Y)\to T_y(X)$ 是单射。
\end{enumerate}
\end{lemma}

\begin{proof}
(i) 由定义得出。

(ii) 断言是局部的，故可设 $Y$ 在 $X$ 中闭且 $X$ 仿射。设 $I$ 是 $R(X)$ 中所有在 $Y$ 上消失的函数构成的理想。则 $\mathcal{O}_{Y,x}$ 是环 $R(X)/I$ 在 $x$ 处的局部化，由局部化的正合性，它是 $\mathcal{O}_{X,x}$ 的商。这意味着线性映射 $T^*_x(j):T^*_x(X)\to T^*_x(Y)$ 是满射，其转置 $T_x(j)$ 是单射。
\end{proof}

现设 $X$ 是某 $k^n$ 的闭子空间。由上述讨论和前一个引理 (ii)，自然包含 $j:X\to k^n$ 的切线性映射 $T_x(j)$ 把 $T_x(X)$ 等同于 $k^n$ 的线性子空间。以下结果精确刻画此子空间。如前以 $I$ 记 $k[X_1,X_2,\dots,X_n]$ 中所有在 $X$ 上消失的多项式的理想。

\begin{lemma}
对任意 $x\in X$ 有
$$T_x(X)=\Bigl\{(\xi_1,\xi_2,\dots,\xi_n)\in k^n\ \Big|\ \sum_{i=1}^n\xi_i(\partial_iP)(x)=0,\ P\in I\Bigr\}.$$
\end{lemma}

\begin{proof}
由定义和前一引理 (ii) 证明中的讨论，$T_x(X)$ 是 $T^*_x(j):k^n\to T^*_x(X)$ 的核的正交补。另一方面，$\ker T^*_x(j)=\{df(x)\mid f\in I_x\}$。$I_x$ 中任意芽 $f$ 是有理函数 $\frac{P}{Q}$ 的芽，其中 $Q(x)\neq 0$，$P\in I$。因此 $df(x)=\frac{1}{Q(x)}dP(x)$，且 $\{df(x)\mid f\in I_x\}=\{dP(x)\mid P\in I\}$。
\end{proof}

现考虑代数簇 $X$ 上的函数 $x\longmapsto\dim_kT_x(X)$。

\begin{proposition}
代数簇 $X$ 上的函数 $x\longmapsto\dim_k T_x(X)$ 是上半连续的。
\end{proposition}

\begin{proof}
断言是局部的，故可设 $X$ 是某 $k^n$ 的闭子空间。由前一结果可把切空间 $T_x(X)$ 等同于
$$\Bigl\{(\xi_1,\dots,\xi_n)\in k^n\ \Big|\ \sum_{i=1}^n\xi_i(\partial_iP)(x)=0,\ P\in I\Bigr\}.$$
若 $\dim T_x(X)=p$，则存在多项式 $P_1,P_2,\dots,P_{n-p}\in I$ 使得矩阵 $[(\partial_iP_j)(x)]$ 的秩为 $n-p$。这意味着在 $x$ 的某邻域 $U$ 中其秩等于 $n-p$。特别地对 $y\in U\cap X$ 有 $\dim_kT_y(X)\le p$。
\end{proof}

若 $\dim_kT_x(X)=\dim_xX$，则称点 $x\in X$ 是\emph{光滑的}。换句话说，$x\in X$ 光滑当且仅当局部环 $\mathcal{O}_x$ 是正则的。

\begin{theorem}
设 $X$ 是 $k$ 上的代数簇。则
\begin{enumerate}[label=(\roman*)]
  \item $X$ 的所有光滑点之集在 $X$ 中开且稠密；
  \item 光滑点 $x\in X$ 含于 $X$ 的唯一不可约分支中。
\end{enumerate}
\end{theorem}

\begin{proof}
第二陈述由 I.2.10 立即得出。记 $V_1,V_2,\dots,V_r$ 为 $X$ 的不可约分支，令 $Y=\bigcup_{i\neq j}V_i\cap V_j$。则 $Y$ 是 $X$ 的闭子集，其补在 $X$ 中稠密。此外由 (ii)，$Y$ 不含 $X$ 的光滑点。故可设 $X$ 是其不可约分支的不交并。这把证明归约到不可约簇的情形。

由前面的结果，光滑点之集在 $X$ 中开，证明归约为证明不可约仿射簇 $X$ 中光滑点的存在性。

设 $n=\dim X$。由 Noether 规范化引理，可找到 $f_1,f_2,\dots,f_n\in R(X)$ 使得由 $P\longmapsto P(f_1,\dots,f_n)$ 定义的 $k[X_1,X_2,\dots,X_n]$ 到 $R(X)$ 的同态是单射，且 $R(X)$ 在其像 $B$ 上是整的。几何上，这定义了 $X$ 到 $k^n$ 上的满有限态射 $p$。有理函数域 $R(X)$ 是 $B$ 的分式域 $L$（同构于 $k(X_1,X_2,\dots,X_n)$）的代数扩张。由于 $R(X)$ 是有限生成 $k$-代数，$R(X)$ 由 $R(X)$ 的有限多个元素在 $L$ 上生成。这意味着域 $R(X)$ 是 $L$ 的有限扩张。此外由于 $k$ 特征为零，由本原元定理，存在元素 $g\in R(X)$ 在 $L$ 上生成 $R(X)$。

先断言可设 $g\in R(X)$（原文本处疑有笔误）。设 $S=B-\{0\}$ 为乘法系，$S^{-1}R(X)$ 为相应的分式环。断言 $S^{-1}R(X)=R(X)$（此处指分式域）。取 $f\in S^{-1}R(X)$。由于 $R(X)$ 是 $L$ 的代数扩张，$f$ 在 $L$ 上代数，即存在 $b_1,b_2,\dots,b_n\in L$ 使
$$f^n+b_1f^{n-1}+\cdots+b_{n-1}f+b_n=0,\qquad b_n\neq 0.$$
这意味着
$$\frac{1}{f}=-\frac{1}{b_n}(f^{n-1}+b_1f^{n-2}+\cdots+b_{n-1})\in S^{-1}R(X).$$
因此 $S^{-1}R(X)$ 是含 $L$ 和 $R(X)$ 的域，故等于 $R(X)$（分式域）。由此可知任一本原元形如 $g=\frac{h}{b}$，其中 $h\in R(X)$，$b\in B-\{0\}$。这意味着 $h\in R(X)$ 也是本原元。

因此在以下固定本原元 $g\in R(X)$。设 $\Omega$ 为 $R(X)$ 的代数闭包，则 $\Omega$ 也是 $L$ 的代数闭包。设 $s$ 为 $R(X)$ 在 $L$ 上的次数。则 $g$ 在 $\Omega$ 的 $L$-自同构群 $\operatorname{Aut}_L(\Omega)$ 作用下的轨道由 $s$ 个元素 $g_0=g,g_1,\dots,g_{s-1}$ 组成。由于 $g$ 在 $B$ 上整且 $\operatorname{Aut}_L(\Omega)$ 固定 $B$，可见 $g_i$（$0\le i\le s-1$）在 $B$ 上整。令
$$V(\lambda_1,\lambda_2,\dots,\lambda_s)=\prod_{i<j}(\lambda_i-\lambda_j).$$
则对任意 $T\in\operatorname{Aut}_L(\Omega)$ 有
$$T\bigl(V(g_0,g_1,\dots,g_{s-1})^2\bigr)=V(g_0,g_1,\dots,g_{s-1})^2.$$
因此 $D=V(g_0,g_1,\dots,g_{s-1})^2$ 属于 $L$。此外它在 $B$ 上整。由于 $B$ 整闭，断定 $D\in B$。

显然 $1,g,g^2,\dots,g^{s-1}$ 构成向量空间 $R(X)$ 在 $L$ 上的一组基。因此对任意 $h\in R(X)$ 存在 $a_0,a_1,\dots,a_{s-1}\in L$ 使
$$h=\sum_{i=0}^{s-1}a_i g^i.$$
若令
$$h_j=\sum_{i=0}^{s-1}a_i g_j^i\quad(0\le j\le s-1),$$
可见 $h_0=h,h_1,\dots,h_{s-1}$ 在 $B$ 上整。由 Cramer 法则
$$a_j=\frac{\Delta_j}{V(g_0,g_1,\dots,g_{s-1})},$$
其中 $\Delta_j$ 是把 Vandermonde 行列式第 $j$ 列换为 $(h_0,h_1,\dots,h_{s-1})^T$ 所得行列式。这导出
$$a_j=\frac{\Delta_jV(g_0,g_1,\dots,g_{s-1})}{D}.$$
由于 $B$ 整闭，$a_jD\in L$，且 $\Delta_jV(g_0,g_1,\dots,g_{s-1})$ 在 $B$ 上整，断定 $a_jD\in B$（$0\le j\le s-1$）。这意味着
$$h=\sum_{i=0}^{s-1}a_iD\frac{g^i}{D}.$$
因此 $\frac{g^i}{D}$（$0\le i\le s-1$）生成 $R(X)$ 的一个含 $R(X)$ 的 $B$-子模（此处按原文）。

设 $Y=\{y\in k^n\mid D(y)\neq 0\}$，$X'=\{x\in X\mid (D\circ p)(x)\neq 0\}$。则 $p$ 把 $X'$ 映到 $Y$ 上。取 $y\in Y$ 和 $x\in X'$ 使 $p(x)=y$。由构造，$D$ 是 $g$ 在 $L$ 上的极小多项式 $\mu$ 的判别式。此外
$$\mu(X)=X^s+c_1X^{s-1}+\cdots+c_{s-1}X+c_s=\prod_{i=0}^{s-1}(X-g_i).$$
因此其系数 $c_j$（$1\le j\le s$）在 $B$ 上整，且由于 $B$ 整闭，它们属于 $B$。由于 $D$ 是 $g_i$ 的对称多项式，它是初等对称多项式 $c_1,c_2,\dots,c_s$ 的多项式。特别地 $D(y)$ 是多项式 $P(X)=X^s+c_1(y)X^{s-1}+\cdots+c_{s-1}(y)X+c_s(y)$ 的判别式。由于 $D(y)\neq 0$，此多项式在 $k$ 中有 $s$ 个不同的根。显然其在 $k$ 中的一个根是 $g(x)$。这意味着其导数 $P'$ 满足
$$P'(g(x))=sg(x)^{s-1}+(s-1)c_1(y)g(x)^{s-2}+\cdots+c_{s-1}(y)\neq 0.$$
因此若 $T\in T_x(X)$ 属于 $T_x(p)$ 的核，则
\begin{align*}
0&=T\bigl(d(g^s+(c_1\circ p)g^{s-1}+\cdots+(c_{s-1}\circ p)g+c_s\circ p)(x)\bigr)\\
&=P'(g(x))\,T(dg(x))+T\bigl(g(x)^{s-1}d(c_1\circ p)(x)+\cdots+g(x)d(c_{s-1}\circ p)(x)+d(c_s\circ p)(x)\bigr)\\
&=P'(g(x))\,T(dg(x))\\
&\quad +g(x)^{s-1}T_x(p)(T)(dc_1(y))+\cdots+g(x)T_x(p)(T)(dc_{s-1}(y))\\
&\quad +T_x(p)(T)(dc_s(y))\\
&=P'(g(x))\,T(dg(x)),
\end{align*}
这导致 $T(dg(x))=0$。现取任意 $h\in R(X)$。则由上述讨论，$h=\sum_{i=0}^{s-1}(b_i\circ p)g^i$，其中 $b_i\in R(Y)$。因此
\begin{align*}
T(dh(x))&=\sum_{i=0}^{s-1}g(x)^iT(d(b_i\circ p)(x))
=\sum_{i=0}^{s-1}g(x)^iT_x(p)(T)(db_i(y))=0.
\end{align*}
最后若 $f\in\mathcal{O}_{X,x}$，$f=\frac{h}{a}$，其中 $a,h\in R(X)$，$a(x)\neq 0$，则
$$df(x)=\frac{a(x)dh(x)-h(x)da(x)}{a(x)^2}.$$
因此对任意 $f\in\mathcal{O}_{X,x}$ 有 $T(df(x))=0$，从而断定 $T=0$，$T_x(p)$ 是单射。由于 $T_y(k^n)$ 是 $n$ 维的，可见 $\dim_kT_x(X)\le n$。由前面的结果得 $\dim_kT_x(X)=n$，$x$ 是光滑点。
\end{proof}

代数簇 $X$ 称为\emph{光滑的}，若其所有点都光滑。由前面的结果有如下命题。

\begin{proposition}
设 $X$ 是光滑代数簇。则其不可约分支等于其连通分支。
\end{proposition}

这特别蕴含函数 $x\longmapsto\dim_kX$ 在光滑簇 $X$ 上局部常值。

现在分析 $X$ 的光滑点邻域。断言如下结果。

\begin{theorem}
设 $X$ 是代数簇，$x\in X$ 是光滑点且 $\dim_xX=n$。则存在：
\begin{enumerate}[label=(\roman*)]
  \item $x$ 的仿射开邻域 $U$；
  \item $U$ 上的正则函数 $f_1,f_2,\dots,f_n$ 和向量场 $D_1,D_2,\dots,D_n$，使得 $D_i(f_j)=\delta_{ij}$（$1\le i,j\le n$）。
\end{enumerate}
\end{theorem}

\begin{proof}
由于断言是局部的，可设 $X$ 是嵌入某 $k^m$ 的光滑不可约仿射闭簇。设 $I$ 是 $A=k[X_1,X_2,\dots,X_m]$ 中所有在 $X$ 上消失的多项式的理想。由于 $\dim_kT_x(X)=\dim X=n$，可找到多项式 $P_{n+1},P_{n+2},\dots,P_m\in I$ 使得矩阵 $[(\partial_iP_j)(x)]$ 的秩为 $m-n$。这意味着此矩阵的秩在 $x\in k^m$ 的某邻域 $V$ 中等于 $m-n$，且
$$T_x(X)=\{(\xi_1,\dots,\xi_m)\in k^m\mid \textstyle\sum_{i=1}^m\xi_i(\partial_iP_j)(x)=0,\ n+1\le j\le m\}.$$

以 $J$ 记 $A$ 中由 $P_{n+1},P_{n+2},\dots,P_m$ 生成的理想。先断言 $J_x=I_x$。显然由定义 $J\subset I$。设 $Y$ 为 $J$ 在 $k^m$ 中所有零点的集合，则 $X\subset Y$。有
$$\dim_xY\ge\dim X=\dim_kT_x(X)=\dim_kT_x(Y)\ge\dim_xY.$$
这意味着 $\dim X=\dim_xY$。因此 $X$ 是 $Y$ 的一个不可约分支。另一方面，由于 $\dim_xY=\dim_kT_x(Y)$，$x$ 是 $Y$ 的光滑点且由前面的结果位于 $Y$ 的唯一不可约分支中。这意味着存在 $x$ 在 $k^m$ 中的邻域 $V'\subset V$ 不与 $Y$ 的任何其他不可约分支相交。因此断定 $J_x=I_x$。考虑局部环 $(A/J)_x$，其极大理想为 $\mathfrak{n}_x=\mathfrak{m}_x/J_x$，且对任意 $p\in\mathbb{Z}_+$ 有 $\mathfrak{n}^p_x=\mathfrak{m}^p_x/J_x$。因此局部环 $(A/J)_x$ 的维数等于其作为 $A_x$-模的维数。由 I.4.5 得 $d((A/J)_x)=\dim_x(V(J))=\dim_xY=\dim X=n$。另一方面有有限维向量空间的正合列
$$0\longrightarrow(J_x+\mathfrak{m}_x^2)/\mathfrak{m}_x^2\longrightarrow\mathfrak{m}_x/\mathfrak{m}_x^2\longrightarrow\mathfrak{n}_x/\mathfrak{n}_x^2\longrightarrow 0,$$
其中 $(J_x+\mathfrak{m}_x^2)/\mathfrak{m}_x^2=\{df(x)\mid f\in J_x\}$ 由 $dP_i(x)$（$n+1\le i\le m$）张成，且由上述等同 $\mathfrak{m}_x/\mathfrak{m}_x^2\cong k^m$。这意味着 $\dim_k(\mathfrak{n}_x/\mathfrak{n}_x^2)=n$，$(A/J)_x$ 是正则局部环。由 I.2.10 它是整环，故 $J_x$ 是素的。这最终导致 $J_x=I_x$。

这意味着 $A$-模 $I/J$ 的支集不含 $x$。特别地存在 $g\in A$ 使得由 $g$ 决定的 $k^m$ 中主开集 $V''\subset V'$ 与 $\operatorname{supp}(I/J)$ 不交。因此 $(I/J)_g=0$，$J_g=I_g$。

可找到多项式 $P_1,P_2,\dots,P_n\in A$ 使得矩阵 $[(\partial_iP_j)(x);\ 1\le i,j\le m]$ 可逆。因此必要时改变 $g$，还可设它在主开集 $V''$ 上可逆。以 $Q$ 记此矩阵的逆。则 $Q$ 的矩阵系数属于 $A_g$。因此在 $V''$ 上可定义微分算子 $\delta_i=\sum_{j=1}^mQ_{ij}\partial_j$（$1\le i\le n$）。显然它们满足
$$\delta_iP_j=\sum_{k=1}^mQ_{ik}\partial_kP_j=\delta_{ij}$$
对任意 $1\le j\le m$。由于任意 $f\in J_g$ 可表示为 $f=\sum_{j=n+1}^m h_jP_j$（$h_j\in A_g$），有
$$\delta_i(f)=\delta_i\Bigl(\sum_{j=n+1}^m h_jP_j\Bigr)=\sum_{j=n+1}^m(\delta_i(h_j)P_j+h_j\delta_i(P_j))=\sum_{j=n+1}^m\delta_i(h_j)P_j\in J_g,$$
即 $J_g=I_g$ 在 $\delta_i$（$1\le i\le n$）作用下不变。这意味着 $\delta_i$ 诱导 $U=X\cap V''$ 上的局部向量场，记为 $D_i$（$1\le i\le n$）。此外若以 $f_i$（$1\le i\le n$）记 $P_i$ 在 $U$ 上的限制，可见
$$D_i(f_j)=\delta_i(P_j)=\delta_{ij}.$$
\end{proof}

称 $(f_1,f_2,\dots,f_n;D_1,D_2,\dots,D_n)$ 为 $U\subset X$ 上的一个\emph{坐标系}。

\begin{lemma}
设 $X$ 是代数簇，$x\in X$ 是光滑点且 $\dim_xX=n$。则存在 $x$ 的仿射开邻域 $U$ 和 $U\subset X$ 上的坐标系 $(f_1,f_2,\dots,f_n;D_1,D_2,\dots,D_n)$，使得 $D_1,D_2,\dots,D_n$ 构成 $T_X(U)$ 作为自由 $\mathcal{O}_X(U)$-模的一组基。此外 $[D_i,D_j]=0$（$1\le i,j\le n$）。
\end{lemma}

\begin{proof}
由于任意光滑点位于 $X$ 的唯一不可约分支中，可设前一个定理中的邻域 $U$ 是不可约的。则 $\dim U=n$。以 $R(U)$ 记 $U$ 上有理函数域。由于 $U$ 是 $n$ 维的，$R(U)$ 在 $k$ 上的超越次数等于 $n$。设 $(f_1,f_2,\dots,f_n;D_1,D_2,\dots,D_n)$ 是 $U$ 上的坐标系。

断言 $f_1,f_2,\dots,f_n$ 在 $k$ 上代数无关。否则可找到非零且次数最小的多项式 $P\in k[X_1,X_2,\dots,X_n]$ 满足 $P(f_1,f_2,\dots,f_n)=0$。这意味着
$$0=D_i(P(f_1,f_2,\dots,f_n))=\sum_{j=1}^n(\partial_jP)(f_1,f_2,\dots,f_n)D_i(f_j)=(\partial_iP)(f_1,f_2,\dots,f_n),$$
由 $P$ 次数的最小性，$\partial_iP=0$（$1\le i\le n$）。由于 $k$ 特征为零，断定 $P$ 是常数多项式，显然不可能。

设 $K$ 是 $R(U)$ 中由 $f_1,f_2,\dots,f_n$ 生成的子域。则 $K$ 在 $k$ 上的超越次数也等于 $n$，$R(U)$ 是 $K$ 的代数扩张。

由于 $U$ 上向量场是 $R(U)$ 的导子，它扩充为 $R(U)$ 的导子。另一方面 $R(U)$ 是 $K$ 的代数扩张。故由隐函数微分法，此导子由其限制到 $K$ 唯一确定。可知 $U$ 上向量场完全由它限制到 $R(U)$ 中由 $f_1,f_2,\dots,f_n$ 生成的子环决定。

设 $T\in T_X(U)$，令 $g_i=T(f_i)$（$1\le i\le n$）。则
$$\Bigl(T-\sum_{i=1}^n g_iD_i\Bigr)(f_j)=T(f_j)-\sum_{i=1}^n g_iD_i(f_j)=0,$$
由上述讨论得 $T=\sum_{i=1}^ng_iD_i$，即 $D_1,D_2,\dots,D_n$ 生成 $T_X(U)$。另一方面若 $\sum_{i=1}^nh_iD_i=0$（$h_i\in R(U)$），则
$$0=\Bigl(\sum_{i=1}^nh_iD_i\Bigr)(f_j)=h_j\quad(1\le j\le n).$$
故 $\mathcal{O}_X(U)$-模 $T_X(U)$ 是自由的，$(D_1,D_2,\dots,D_n)$ 是其一组基。

最后对任意 $1\le i,j,k\le n$ 有
$$[D_i,D_j](f_k)=D_i(D_j(f_k))-D_j(D_i(f_k))=0,$$
这意味着 $[D_i,D_j]=0$。
\end{proof}

设 $x\in X$，$T\in T_{X,x}$。则 $T$ 决定局部环 $\mathcal{O}_{X,x}$ 的导子。显然对任意 $f\in\mathfrak{m}_x^2$ 有 $T(f)\in\mathfrak{m}_x$。此外对 $f\in\mathcal{O}_x$，有 $T(f)(x)=T(f-f(x))(x)$，结果只依赖于 $df(x)$。因此映射 $f\longmapsto T(f)(x)$ 通过 $T^*_x(X)$ 分解，定义切向量 $T(x)\in T_x(X)$，满足 $T(x)(df(x))=T(f)(x)$（$f\in\mathcal{O}_x$）。于是构造了 $T_{X,x}$ 到 $T_x(X)$ 的线性映射。显然它决定几何纤维 $\mathcal{O}_x/\mathfrak{m}_x\otimes_{\mathcal{O}_x}T_{X,x}$ 到 $T_x(X)$ 的线性映射。

\begin{proposition}
设 $x$ 是代数簇 $X$ 的光滑点。则 $\mathcal{O}_x/\mathfrak{m}_x\otimes_{\mathcal{O}_x}T_{X,x}$ 到 $T_x(X)$ 的典范映射是向量空间的同构。
\end{proposition}

\begin{proof}
设 $n=\dim_xX$。由前面的引理知 $T_{X,x}$ 是自由 $\mathcal{O}_x$-模。更确切地说，存在 $f_1,f_2,\dots,f_n\in\mathcal{O}_x$ 和 $D_1,D_2,\dots,D_n\in T_{X,x}$，满足 $D_i(f_j)=\delta_{ij}$（$1\le i,j\le n$），使得 $(D_1,D_2,\dots,D_n)$ 是自由 $\mathcal{O}_x$-模 $T_{X,x}$ 的一组基。这意味着 $D_1,D_2,\dots,D_n$ 在 $\mathcal{O}_x/\mathfrak{m}_x\otimes_{\mathcal{O}_x}T_{X,x}$ 中的像构成其作为 $k$-向量空间的基。

另一方面 $D_i(x)$ 满足 $D_i(x)(df_j(x))=\delta_{ij}$（$1\le i,j\le n$），故它们线性无关。由于切空间 $T_x(X)$ 是 $n$ 维的，断定 $(D_i(x),1\le i\le n)$ 是 $T_x(X)$ 的一组基，映射是双射。
\end{proof}

前面的两个结果蕴含如下定理。

\begin{theorem}
设 $X$ 是光滑代数簇。则切层 $\mathcal{T}_X$ 是有限秩局部自由 $\mathcal{O}_X$-模。对任意 $x\in X$，$\mathcal{T}_X$ 的几何纤维自然同构于 $T_x(X)$。
\end{theorem}

设 $X$ 是光滑代数簇，$T(X)=\{(x,\xi)\mid \xi\in T_x(X),\ x\in X\}$。想在 $T(X)$ 上定义自然的代数簇结构。

先设 $X$ 上切层 $\mathcal{T}_X$ 是自由 $\mathcal{O}_X$-模。设 $(T_1,T_2,\dots,T_n)$ 是 $\mathcal{T}_X$ 的一组基。则有 $X\times k^n$ 到 $T(X)$ 上的双射 $\varphi$，定义为
$$\varphi(x,\xi_1,\xi_2,\dots,\xi_n)=\Bigl(x,\sum_{i=1}^n\xi_iT_i(x)\Bigr).$$
通过要求 $\varphi:X\times k^n\to T(X)$ 是同构，可在 $T(X)$ 上定义代数簇结构。设 $(T'_1,T'_2,\dots,T'_n)$ 是自由 $\mathcal{O}_X$-模 $\mathcal{T}_X$ 的另一组基，$\varphi':X\times k^n\to T(X)$ 是相应的映射。则存在元素在 $R(X)$ 中的可逆矩阵 $Q$ 使得 $T'_i=\sum_{j=1}^nQ_{ji}T_j$，这意味着
$$\varphi'(x,\xi_1,\dots,\xi_n)=\Bigl(x,\sum_i\xi_iT'_i(x)\Bigr)=\Bigl(x,\sum_{i,j}\xi_iQ_{ji}(x)T_j(x)\Bigr)=\varphi\bigl(x,Q(x)(\xi_1,\dots,\xi_n)\bigr),$$
且 $\varphi'$ 是同构当且仅当 $\varphi$ 是同构。因此 $X$ 上的代数结构与 $\mathcal{T}_X$ 基的选择无关。

现考虑任意光滑代数簇 $X$。由前面的定理可找到 $X$ 的开覆盖 $(U_1,U_2,\dots,U_s)$ 使得 $\mathcal{T}_X|_{U_i}=\mathcal{T}_{U_i}$ 是自由 $\mathcal{O}_{U_i}$-模。显然 $T(X)$ 是 $T(U_i)$（$1\le i\le s$）的并，由上述讨论每个 $T(U_i)$ 有自然的代数簇结构。此外由于此结构与基的选择无关，可见在交 $T(U_i)\cap T(U_j)$ 上由 $T(U_i)$ 和 $T(U_j)$ 结构诱导的结构相同。这在 $T(X)$ 上定义了一个代数簇结构。称 $T(X)$ 为 $X$ 的\emph{切丛}。有自然映射 $i:X\to T(X)$ 和 $p:T(X)\to X$，定义为 $i(x)=(x,0)$ 和 $p(x,\xi)=x$（$\xi\in T_x(X)$，$x\in X$）。显然这些映射是代数簇的态射。此外有如下显然的结果。

\begin{proposition}
设 $X$ 是光滑代数簇。则
\begin{enumerate}[label=(\roman*)]
  \item 切丛 $T(X)$ 是光滑代数簇，且 $\dim_{(x,\xi)}T(X)=2\dim_xX$；
  \item 纤维化 $p:T(X)\to X$ 是局部平凡的。
\end{enumerate}
\end{proposition}

类似地，若 $X$ 是光滑簇，可在 $X$ 上定义 $\mathcal{O}_X$-模
$$\mathcal{T}^*_X=\mathcal{H}om_{\mathcal{O}_X}(\mathcal{T}_X,\mathcal{O}_X).$$
由于 $\mathcal{T}_X$ 是有限秩局部自由的，$\mathcal{T}^*_X$ 也是。这意味着其在 $x\in X$ 的几何纤维自然同构于余切空间 $T^*_x(X)$。设 $U\subset X$ 是开集，$f\in\mathcal{O}_X(U)$。则它定义 $\mathcal{H}om_{\mathcal{O}_X}(\mathcal{T}_X,\mathcal{O}_X)(U)=\operatorname{Hom}_{\mathcal{O}_X(U)}(\mathcal{T}_X(U),\mathcal{O}_X(U))$ 中由 $T\longmapsto T(f)$ 给出的元素，记为 $df$，称为 $f$ 的\emph{微分}。显然对任意 $x\in U$ 有
$$df(T)(x)=T(f)(x)=T(x)(df(x)),$$
故可把 $df(x)$ 看作由局部截面 $df$ 决定的 $\mathcal{T}^*_X$ 的几何纤维 $T^*_x(X)$ 中的元素。

若 $(f_1,f_2,\dots,f_n;D_1,D_2,\dots,D_n)$ 是充分小的仿射开集 $U$ 上的坐标系，由前面的结果 $(D_1,D_2,\dots,D_n)$ 构成 $\mathcal{T}_X(U)$ 作为自由 $\mathcal{O}_X(U)$-模的一组基。由此可知 $(df_1,df_2,\dots,df_n)$ 是自由 $\mathcal{O}_X(U)$-模 $\mathcal{T}^*_X(U)$ 的对偶基。

如切丛的情形，可定义 $T^*(X)=\{(x,\omega)\mid \omega\in T^*_x(X),\ x\in X\}$ 以及 $T^*(X)$ 上的代数簇结构，使得在带坐标系 $(f_1,f_2,\dots,f_n;D_1,D_2,\dots,D_n)$ 的充分小仿射开集 $U\subset X$ 上，$(df_1,df_2,\dots,df_n)$ 是自由 $\mathcal{O}_X(U)$-模 $\mathcal{T}^*_X(U)$ 的对偶基，且有同构 $U\times k^n\to T^*(U)\subset T^*(X)$，由
$$\varphi^*(x,\xi_1,\xi_2,\dots,\xi_n)=\Bigl(x,\sum_{i=1}^n\xi_i\,df_i(x)\Bigr)$$
给出。此簇称为 $X$ 的\emph{余切丛}。有自然映射 $\iota:X\to T^*(X)$ 和 $\pi:T^*(X)\to X$，定义为 $\iota(x)=(x,0)$ 和 $\pi(x,\omega)=x$（$\omega\in T^*_x(X)$，$x\in X$）。显然这些映射是代数簇的态射。

\begin{proposition}
设 $X$ 是光滑代数簇。则
\begin{enumerate}[label=(\roman*)]
  \item 余切丛 $T^*(X)$ 是光滑代数簇，且 $\dim_{(x,\omega)}T^*(X)=2\dim_xX$；
  \item 纤维化 $\pi:T^*(X)\to X$ 是局部平凡的。
\end{enumerate}
\end{proposition}

最后补充关于光滑点处正则函数芽的“Taylor 级数”的注记。

设 $X$ 是代数簇，$x\in X$ 是光滑点，其局部环 $\mathcal{O}_x$ 是正则局部环。若 $n=\dim_xX$，由前面的结果存在 $x$ 的仿射开邻域 $U$ 和 $U$ 上坐标系 $(f_1,f_2,\dots,f_n;D_1,D_2,\dots,D_n)$，使得 $f_1(x)=f_2(x)=\cdots=f_n(x)=0$ 且 $D_1(x),D_2(x),\dots,D_n(x)$ 构成 $T_x(X)$ 的一组基。因此 $df_1(x),df_2(x),\dots,df_n(x)$ 构成 $T^*_x(X)$ 的一组基，$f_1,f_2,\dots,f_n$ 在正则局部环 $\mathcal{O}_x$ 中定义一个坐标系。故有自然态射 $k[X_1,X_2,\dots,X_n]\to\mathcal{O}_x$，由 $P\longmapsto P(f_1,f_2,\dots,f_n)$ 给出。显然任何非常数项的多项式的像在 $\mathcal{O}_x$ 中可逆。因此这个态射扩充为局部环的态射 $\varphi:A=k[X_1,X_2,\dots,X_n]_0\to\mathcal{O}_x$。由于 $\mathcal{O}_x$ 是正则局部环，由 I.2.9 知 $\Gr\varphi$ 是 $\Gr A$ 到 $\Gr\mathcal{O}_x$ 的同构。由于 $A$ 的滤过是 Hausdorff 的，这意味着 $\varphi:A\to\mathcal{O}_x$ 是单射。故可把 $A$ 看作 $\mathcal{O}_x$ 的子环。自然同态 $A\to\mathcal{O}_x\to\mathcal{O}_x/\mathfrak{m}_x\cong k$ 是满射，因此其核 $A\cap\mathfrak{m}_x$ 是 $A$ 的极大理想，即等于 $\mathfrak{m}$。这意味着对任意 $p\in\mathbb{Z}_+$ 有 $\mathfrak{m}^p\subset A\cap\mathfrak{m}_x^p$。因此对所有 $p\in\mathbb{Z}_+$ 有自然映射 $A/\mathfrak{m}^p\to\mathcal{O}_x/\mathfrak{m}_x^p$，且图
$$\begin{tikzcd}
0 \arrow[r] & \Gr_pA \arrow[r] \arrow[d] & A/\mathfrak{m}^{p+1} \arrow[r] \arrow[d] & A/\mathfrak{m}^p \arrow[r] \arrow[d] & 0\\
0 \arrow[r] & \Gr_p\mathcal{O}_x \arrow[r] & \mathcal{O}_x/\mathfrak{m}_x^{p+1} \arrow[r] & \mathcal{O}_x/\mathfrak{m}_x^p \arrow[r] & 0
\end{tikzcd}$$
交换。由于两行正合且第一个竖直箭头是同构，若最后一个也是同构，由五引理中间箭头是同构。因此由 $p$ 上的归纳断定 $A/\mathfrak{m}^p\to\mathcal{O}_x/\mathfrak{m}_x^p$ 对所有 $p\in\mathbb{Z}_+$ 是同构。由此可知对所有 $p\in\mathbb{Z}_+$ 有 $\mathfrak{m}^p=A\cap\mathfrak{m}_x^p$。此外对任意 $p\in\mathbb{Z}_+$ 有 $A+\mathfrak{m}_x^p=\mathcal{O}_x$。故 $A$ 在 $\mathcal{O}_x$ 中稠密，$\mathfrak{m}_x$-进拓扑在 $A$ 上诱导 $\mathfrak{m}$-进拓扑。这意味着 $A$ 和 $\mathcal{O}_x$ 的完备化环同构。由于 $A$ 的完备化 $\widehat A$ 是形式幂级数环 $k[[X_1,X_2,\dots,X_n]]$，可见 $\mathcal{O}_x$ 可等同于 $k[[X_1,X_2,\dots,X_n]]$ 的子环。

设 $\widehat{\mathfrak{m}}$ 是局部环 $\widehat A=k[[X_1,X_2,\dots,X_n]]$ 的极大理想，即
$$\widehat{\mathfrak{m}}=\Bigl\{F\ \Big|\ F=\sum_{|I|\ge1}a_IX^I\Bigr\}.$$
这意味着
$$\widehat{\mathfrak{m}}^p=\Bigl\{F\ \Big|\ F=\sum_{|I|\ge p}a_IX^I\Bigr\}.$$
设 $\bar X_1,\bar X_2,\dots,\bar X_n$ 是 $X_1,X_2,\dots,X_n$ 在 $\widehat{\mathfrak{m}}/\widehat{\mathfrak{m}}^2$ 中的像。则立即看到 $(\bar X_i;1\le i\le n)$ 是 $\widehat A$ 中的坐标系。此外自然同态 $k[\bar X_1,\bar X_2,\dots,\bar X_n]$ 到 $\Gr\widehat A$ 是同构。因此 $\widehat A=k[[X_1,X_2,\dots,X_n]]$ 是正则局部环。

显然有 $\mathfrak{m}=A\cap\widehat{\mathfrak{m}}$ 和 $\mathfrak{m}^p\subset\widehat{\mathfrak{m}}^p$（$p\in\mathbb{Z}_+$）。$A$ 到 $\widehat A$ 的自然包含诱导 $\Gr A$ 到 $\Gr\widehat A$ 的同构。因此如前有交换图
$$\begin{tikzcd}
0 \arrow[r] & \Gr_pA \arrow[r] \arrow[d] & A/\mathfrak{m}^{p+1} \arrow[r] \arrow[d] & A/\mathfrak{m}^p \arrow[r] \arrow[d] & 0\\
0 \arrow[r] & \Gr_p\widehat A \arrow[r] & \widehat A/\widehat{\mathfrak{m}}^{p+1} \arrow[r] & \widehat A/\widehat{\mathfrak{m}}^p \arrow[r] & 0
\end{tikzcd}$$
且由归纳再次断定竖直箭头是同构。这意味着 $\mathfrak{m}^p=A\cap\widehat{\mathfrak{m}}^p$，即
$$\mathfrak{m}^p=\Bigl\{f\in A\ \Big|\ f=\sum_{|I|\ge p}a_IX^I\Bigr\}.$$

设 $T$ 是 $U$ 上向量场。则它诱导 $\mathcal{O}_x$ 的导子。由 $p$ 上的归纳可见 $T(\mathfrak{m}_x^p)\subset\mathfrak{m}^{p-1}_x$（$p\in\mathbb{N}$）。因此 $T$ 在 $\mathcal{O}_x$ 的 $\mathfrak{m}_x$-进拓扑中连续，并扩充为 $\mathcal{O}_x$ 完备化的连续导子。另一方面，多项式环 $k[X_1,X_2,\dots,X_n]$ 在形式幂级数环 $k[[X_1,X_2,\dots,X_n]]$ 中稠密，故 $k[[X_1,X_2,\dots,X_n]]$ 的任意连续导子完全由其在 $X_i$（$1\le i\le n$）上的作用决定。由于 $D_i(f_j)=\delta_{ij}$（$1\le i,j\le n$），这意味着在上述同构下 $D_i$ 对应于 $\partial_i$（$1\le i\le n$）。

任意形式幂级数 $F\in k[[X_1,X_2,\dots,X_n]]$ 可唯一写成其 Taylor 级数
$$F=\sum_{I\in\mathbb{Z}_+^n}\frac{(\partial^IF)(0)}{I!}\,X^I.$$
这与上述讨论一起立即得出如下结果。

\begin{lemma}
设 $x$ 是代数簇 $X$ 的光滑点，$(f_1,f_2,\dots,f_n;D_1,D_2,\dots,D_n)$ 是 $x$ 邻域中的坐标系。则对任意 $f\in\mathcal{O}_x$ 和 $p\in\mathbb{Z}_+$，下列条件等价：
\begin{enumerate}[label=(\roman*)]
  \item $f\in\mathfrak{m}_x^p$；
  \item $(D^If)(x)=0$ 对所有 $|I|<p$ 的 $I\in\mathbb{Z}_+^n$。
\end{enumerate}
特别地，$(D^If)(x)=0$ 对所有 $I\in\mathbb{Z}_+^n$ 蕴含 $f=0$。
\end{lemma}
\begin{annote}
光滑点 $x\in X$ 的定义是 $\Omega_{X/k}$ 在该处局部自由、且 $\dim T_xX=\dim_xX$；只有在此处才有 $\mathrm{gr}\mathcal D_X\cong\mathrm{Sym}\,\mathcal T_X$（讲义 2.x），从而特征簇落在 $T^*X$ 且维数比较可用。一般点上没有这个同构，只能退回 $\mathcal O_X$-凝聚语言，这正是第三章引入凝聚化函子的动机。[来源：原书 §2–§3, p62–76；译者按]
\end{annote}



\section{光滑簇上的微分算子层}

设 $X$ 是特征零代数闭域 $k$ 上的光滑代数簇。以 $\mathcal{D}_X$ 记 $X$ 上局部微分算子层，以 $F\mathcal{D}_X$ 记按微分算子阶数的滤过，以 $\Gr\mathcal{D}_X$ 记相应的分次环层。

先描述 $\Gr\mathcal{D}_X$ 的结构。设 $U$ 是 $X$ 中仿射开集。则由定义 $\Gamma(U,\mathcal{D}_X)=D(U)$。如 I.8，对任意 $p\in\mathbb{Z}_+$ 和 $T\in D_p(U)$ 可定义映射 $\sigma_p(T):R(U)^p\to R(U)$：
$$\sigma_p(T)(f_1,f_2,\dots,f_p)=[[\cdots[[T,f_1],f_2],\dots,f_{p-1}],f_p].$$
如 I.8.3 中所证，此映射是对称 $k$-多重线性的，且 $\sigma_p(T)=0$ 当且仅当 $T\in D_{p-1}(U)$。此外对任意 $1\le i\le p$，映射
$$f\longmapsto\sigma_p(T)(f_1,f_2,\dots,f_{i-1},f,f_{i+1},\dots,f_p)$$
是 $U$ 上的向量场。由于 $\sigma_p(T)$ 对称，为证此点可设 $i=p$。显然这是 $U$ 上阶数 $\le 1$ 的微分算子且在常数上为零，故由 1.16.(ii) 它是向量场。因此 $\sigma_p(T)(f_1,f_2,\dots,f_p)(x)$ 只依赖于 $f_i$ 在 $x$ 处的微分 $df_i(x)$（$1\le i\le p$）。由此可知可在 $U$ 的余切丛 $T^*(U)$ 上定义函数 $\Symb_p(T)$：
$$\Symb_p(T)(x,\omega)=\frac{1}{p!}\sigma_p(T)(f,f,\dots,f)(x),$$
其中 $f\in R(U)$ 满足 $df(x)=\omega$。

\begin{lemma}\ 
\begin{enumerate}[label=(\roman*)]
  \item 函数 $\Symb_p(T)$ 在 $T^*(U)$ 上正则；
  \item 对固定的 $x\in U$，函数 $\Symb_p(T)$ 是 $T^*_x(X)$ 上 $p$ 次齐次多项式。
\end{enumerate}
\end{lemma}

\begin{proof}
由于断言是局部的，由 2.10 可设 $U$ 充分小，使得 $U$ 上存在坐标系 $(f_1,f_2,\dots,f_n;D_1,D_2,\dots,D_n)$，且映射 $(x,\xi_1,\xi_2,\dots,\xi_n)\longmapsto(x,\sum_{i=1}^n\xi_i\,df_i(x))$ 是 $U\times k^n$ 到 $T^*(U)$ 上的同构。另一方面，
$$(x,\xi_1,\xi_2,\dots,\xi_n)\longmapsto\frac{1}{p!}\sigma_p(T)\Bigl(\sum_i\xi_if_i,\sum_i\xi_if_i,\dots,\sum_i\xi_if_i\Bigr)(x)$$
是 $U\times k^n$ 上的正则函数，这意味着 $\Symb_p(T)$ 在 $T^*(U)$ 上正则。第二陈述显然。
\end{proof}

称函数 $\Symb_p(T)$ 为微分算子 $T$ 的 \emph{$p$-主符号}。

设 $\pi:T^*(X)\to X$ 是自然投射，定义为 $\pi(x,\omega)=x$（$\omega\in T^*_x(X)$，$x\in X$）。由于 $\pi$ 是局部平凡纤维化且 $x\in X$ 处的纤维是 $T^*_x(X)$，可见 $T^*_x(X)$ 上多项式的齐次次数的自然分次在直像层 $\pi_*(\mathcal{O}_{T^*(X)})$ 上诱导分次环层结构。显然符号映射 $\Symb_p$ 定义层 $F_p\mathcal{D}_X$ 到 $\Gr_p\pi_*(\mathcal{O}_{T^*(X)})$ 的同态，同名记之。它在子层 $F_{p-1}\mathcal{D}_X$ 上为零，故它定义层 $\Gr_p\mathcal{D}_X$ 到 $\pi_*(\mathcal{O}_{T^*(X)})$ 的第 $p$ 次齐次分量的同态。设 $\Symb:\Gr\mathcal{D}_X\to\pi_*(\mathcal{O}_{T^*(X)})$ 为相应的分次层同态。

\begin{theorem}
符号映射 $\Symb:\Gr\mathcal{D}_X\to\pi_*(\mathcal{O}_{T^*(X)})$ 是分次 $\mathcal{O}_X$-代数层的同构。
\end{theorem}

此结果的证明分几步。先证明符号映射是 $k$-代数层的同态。

\begin{lemma}
设 $U$ 是 $X$ 的开子集，$T,S\in\mathcal{D}_X(U)$ 分别是阶数 $\le p$ 和 $\le q$ 的算子。则
$$\Symb_{p+q}(TS)=\Symb_p(T)\,\Symb_q(S).$$
\end{lemma}

\begin{proof}
设 $f\in\mathcal{O}_X(U)$，定义映射 $\tau:\mathcal{D}_X(U)\to\mathcal{D}_X(U)$ 为 $\tau(T)=[T,f]$。则
$$\tau(TS)=[TS,f]=TSf-fTS=[T,f]S+T[S,f]=\tau(T)S+T\tau(S).$$
故对任意 $k\in\mathbb{Z}_+$，
$$\tau^k(TS)=\sum_{i=0}^k\binom{k}{i}\tau^{k-i}(T)\,\tau^i(S).$$
这意味着若固定 $x\in X$ 和 $\omega\in T^*_x(X)$ 使 $df(x)=\omega$，有
\begin{align*}
\Symb_{p+q}(TS)(x,\omega)&=\frac{1}{(p+q)!}\sigma_{p+q}(T)(f,f,\dots,f)(x)=\frac{1}{(p+q)!}\tau^{p+q}(TS)(x)\\
&=\frac{1}{p!\,q!}\tau^p(T)(x)\,\tau^q(S)(x)=\Symb_p(T)(x,\omega)\,\Symb_q(S)(x,\omega).
\end{align*}
\end{proof}

显然，由于纤维化 $\pi:T^*(X)\to X$ 是局部平凡的，分次环层 $\pi_*(\mathcal{O}_{T^*(X)})$ 的第零次齐次分量等于 $\mathcal{O}_X$，且符号映射 $\Symb_0$ 是恒等映射。另一方面，第一个齐次分量自然同构于 $\mathcal{T}_X$。此外对 $x$ 邻域上的任意向量场 $T$ 和正则函数 $f$ 有
$$\Symb_1(T)(x,df(x))=[T,f](x)=T(f)(x)=T(x)(df(x)).$$
由于由 1.16.(ii) 有 $F_1\mathcal{D}_X=\mathcal{O}_X\oplus\mathcal{T}_X$，断定 $\Symb_1$ 是 $\Gr_1\mathcal{D}_X$ 到 $\pi_*(\mathcal{O}_{T^*(X)})$ 第一个齐次分量的同构。由 $\pi$ 的局部平凡性，分次环层 $\pi_*(\mathcal{O}_{T^*(X)})$ 由其第零和第一次齐次分量生成。因此符号映射的像等于 $\pi_*(\mathcal{O}_{T^*(X)})$。只需证明其核为零。

\begin{lemma}
设 $T\in F_p\mathcal{D}_X(U)$。则 $\Symb_p(T)=0$ 当且仅当 $T$ 阶数 $\le p-1$。
\end{lemma}

\begin{proof}
断言是局部的，故可设 $U$ 仿射。对 $p$ 归纳证明。$p=0$ 时显然，故设 $p>0$。固定 $f\in\mathcal{O}_X(U)$，则 $[T,f]$ 是阶数 $\le p-1$ 的微分算子。取 $x\in U$ 和 $\omega\in T^*_x(X)$，令 $\eta=df(x)$。必要时缩小 $U$，可取 $g\in\mathcal{O}_X(U)$ 使 $dg(x)=\omega$。对任意 $h\in\mathcal{O}_X(U)$ 定义映射 $\tau_h:\mathcal{D}_X(U)\to\mathcal{D}_X(U)$ 为 $\tau_h(T)=[T,h]$。则对任意 $\lambda\in k$，
$$\tau_{f+\lambda g}(T)=[T,f+\lambda g]=[T,f]+\lambda[T,g]=\tau_f(T)+\lambda\tau_g(T).$$
由于 $\tau_f$ 和 $\tau_g$ 可交换，对任意 $k\in\mathbb{Z}_+$ 有
$$\tau_{f+\lambda g}^k(T)=\sum_{i=0}^k\binom{k}{i}\lambda^i\tau_g^i(\tau_f^{k-i}(T)).$$
由假设，映射
$$\lambda\longmapsto\Symb_p(T)(x,\eta+\lambda\omega)=\frac{1}{p!}\tau_{f+\lambda g}^p(T)(x)$$
在 $k$ 上恒为零。由于 $k$ 无限，对 $1\le i\le p$ 有 $\tau_g^{p-i}(\tau_f^i(T))(x)=0$。特别地，
$$\Symb_{p-1}([T,f])(x,\omega)=\frac{1}{(p-1)!}\tau_g^{p-1}(\tau_f(T))(x)=\frac{1}{(p-1)!}\tau_g^{p-1}([T,f])(x)=0$$
对任意 $\omega\in T^*_x(X)$ 成立。由于 $x\in U$ 任意，由归纳假设 $[T,f]$ 阶数 $\le p-2$。这意味着 $T$ 阶数 $\le p-1$。
\end{proof}

这完成定理 3.2 的证明。

\begin{proposition}
光滑簇 $X$ 上局部微分算子层 $\mathcal{D}_X$ 是局部自由 $\mathcal{O}_X$-模（对左乘和右乘）。

更确切地说，每点 $x\in X$ 有仿射开邻域 $U$ 和 $U$ 上的坐标系 $(f_1,f_2,\dots,f_n;D_1,D_2,\dots,D_n)$，使得
\begin{enumerate}[label=(\roman*)]
  \item 对任意 $I,J\in\mathbb{Z}_+^n$ 有 $D^I\circ D^J=D^{I+J}$；
  \item $(D^I;\ I\in\mathbb{Z}_+^n,\ |I|\le p)$ 是自由 $\mathcal{O}_X(U)$-模 $F_p\mathcal{D}_X(U)$（对左乘和右乘）的一组基；
  \item $(D^I;\ I\in\mathbb{Z}_+^n)$ 是自由 $\mathcal{O}_X(U)$-模 $\mathcal{D}_X(U)$（对左乘和右乘）的一组基。
\end{enumerate}
\end{proposition}

\begin{proof}
设 $U$ 是 $x$ 的邻域，$(f_1,f_2,\dots,f_n;D_1,D_2,\dots,D_n)$ 是如 2.10 中 $U$ 上的坐标系。则对 $1\le i,j\le n$ 有 $[D_i,D_j]=0$，(i) 成立。

记 $\xi_i=\Symb_1(D_i)$（$1\le i\le n$）。则 $\pi_*(\mathcal{O}_{T^*(X)})(U)$ 是以 $(\xi^I;\ I\in\mathbb{Z}_+^n)$ 为基的自由 $\mathcal{O}_X(U)$-模，其齐次分量是自由 $\mathcal{O}_X(U)$-模。由正合列
$$0\longrightarrow F_{p-1}\mathcal{D}_X\longrightarrow F_p\mathcal{D}_X\longrightarrow\Gr_p\mathcal{D}_X\longrightarrow 0$$
及定理 3.2，由 $p$ 上的归纳断定 $F_p\mathcal{D}_X(U)$ 是自由 $\mathcal{O}_X(U)$-模，且由 $(D^I;\ I\in\mathbb{Z}_+^n,\ |I|\le p)$ 生成。若 $\sum_{|I|\le p}f_ID^I=0$，取 $p$ 次主符号得 $|I|=p$ 时 $f_I=0$ 且 $\sum_{|I|\le p-1}f_ID^I=0$。故由递降归纳得 $|I|\le p$ 时 $f_I=0$。这证明 (ii) 和 (iii)。
\end{proof}

\begin{proposition}
设 $X$ 是特征零代数闭域上的光滑仿射簇。则
\begin{enumerate}[label=(\roman*)]
  \item $\Gr D(X)$ 是诺特环；
  \item $\Gr D(X)$ 是由 $\Gr_1D(X)$ 生成的 $R(X)$-代数。
\end{enumerate}
\end{proposition}

\begin{proof}
由于 $\pi:T^*(X)\to X$ 是局部平凡纤维化且纤维是向量空间，断定 $\pi$ 是仿射态射。这意味着 $T^*(X)$ 是仿射簇。故
$$\Gr D(X)=\Gamma(X,\Gr\mathcal{D}_X)\cong\Gamma(X,\pi_*(\mathcal{O}_{T^*(X)}))=\Gamma(T^*(X),\mathcal{O}_{T^*(X)})=R(T^*(X))$$
是有限生成 $k$-代数和诺特环。此外由于 $\pi_*(\mathcal{O}_{T^*(X)})$ 作为 $\mathcal{O}_X$-代数由第一次齐次分量生成，对任意 $p\in\mathbb{Z}_+$，$\pi_*(\mathcal{O}_{T^*(X)})_1\otimes_{\mathcal{O}_X}\pi_*(\mathcal{O}_{T^*(X)})_p\to\pi_*(\mathcal{O}_{T^*(X)})_{p+1}$ 是满射。由于 $X$ 仿射，这意味着整体截面的相应态射是满射。因此 $\Gr D(X)$ 作为 $R(X)$-代数由 $\Gr_1D(X)$ 生成。
\end{proof}

\begin{theorem}
设 $X$ 是特征零代数闭域上的光滑仿射簇。则
\begin{enumerate}[label=(\roman*)]
  \item $D(X)$ 是左诺特且右诺特的环；
  \item 环 $D(X)$ 由 $R(X)$ 和 $X$ 上整体向量场生成。
\end{enumerate}
\end{theorem}

\begin{proof}
由前面的命题，$D(X)$ 是满足 I.3 开头条件 (i)--(vii) 的滤过环。故 (i) 由 I.3.4 得出。

设 $A$ 是 $D(X)$ 中由 $R(X)$ 和 $X$ 上整体向量场生成的子环。设 $F A$ 是 $A$ 上的诱导滤过。则有 $\Gr A$ 到 $\Gr D(X)$ 的单射同态，由前面的命题 (ii) 它也是满射。这意味着 $A=D(X)$。
\end{proof}

\chapter{光滑代数簇上微分算子层上的模}

\section{拟凝聚 $\mathcal{D}_X$-模}

设 $X$ 是拓扑空间，$A$ 是 $X$ 上带单位的环层。以 $\mathcal{M}(A)$ 记 $X$ 上 $A$-模层的范畴，它是 Abel 范畴。以 $A=\Gamma(X,A)$ 记 $A$ 的整体截面环，以 $\mathcal{M}(A)$ 记 $A$-模的范畴。则有自然的加性整体截面函子 $\Gamma=\Gamma(X,-):\mathcal{M}(A)\to\mathcal{M}(A)$。此外存在函子 $\operatorname{Hom}_A(A,-)$ 到函子 $\Gamma(X,-)$ 的自然同构，它把态射 $T:A\to V$ 映为 $T(1_X)\in\Gamma(X,V)$。

还可定义局部化函子 $\Delta:\mathcal{M}(A)\to\mathcal{M}(A)$，由
$$\Delta(V)=A\otimes_A V$$
给出。显然 $\Delta$ 是加性函子，且右正合。

由标准论证，对任意 $V\in\mathcal{M}(A)$ 和 $W\in\mathcal{M}(A)$ 有
$$\operatorname{Hom}_A(A\otimes_A V,W)=\operatorname{Hom}_A(V,\operatorname{Hom}_A(A,W)).$$
故
$$\operatorname{Hom}_A(\Delta(V),W)=\operatorname{Hom}_A(V,\Gamma(X,W)),$$
即 $\Delta$ 是整体截面函子 $\Gamma$ 的左伴随函子。

特别地，存在伴随态射 $\varphi$（从恒等函子到 $\Gamma\circ\Delta$）和 $\psi$（从 $\Delta\circ\Gamma$ 到恒等函子）。

现考虑 $X$ 是代数簇、$A=\mathcal{O}_X$ 是 $X$ 上结构层的特殊情形。此时如前以 $R(X)$ 记 $X$ 上正则函数环。

若 $X$ 是仿射的，称 $\mathcal{M}(\mathcal{O}_X)$ 中的 $V$ 是\emph{拟凝聚 $\mathcal{O}_X$-模}，若存在 $R(X)$-模 $V$ 使得 $V\cong\Delta(V)$。

若 $X$ 是任意代数簇，称 $V$ 是拟凝聚 $\mathcal{O}_X$-模，若每点 $x\in X$ 有仿射开邻域 $U$ 使得 $V|_U$ 是拟凝聚 $\mathcal{O}_U$-模。拟凝聚 $\mathcal{O}_X$-模构成 $\mathcal{M}(\mathcal{O}_X)$ 的全子范畴，记为 $\mathcal{M}_{qc}(\mathcal{O}_X)$。可以检验 $\mathcal{M}_{qc}(\mathcal{O}_X)$ 是 Abel 范畴。

显然 $\Delta$ 是从 $\mathcal{M}(R(X))$ 到 $\mathcal{M}_{qc}(\mathcal{O}_X)$ 的函子。若 $X$ 仿射，由 Serre 定理 $\Delta:\mathcal{M}(R(X))\to\mathcal{M}_{qc}(\mathcal{O}_X)$ 是范畴的等价，$\Gamma:\mathcal{M}_{qc}(\mathcal{O}_X)\to\mathcal{M}(R(X))$ 是其拟逆。

设 $\mathcal{D}_X$ 是 $X$ 上微分算子层。则有自然同态 $\iota:\mathcal{O}_X\to\mathcal{D}_X$，它定义从 $\mathcal{M}(\mathcal{D}_X)$ 到 $\mathcal{M}(\mathcal{O}_X)$ 的遗忘函子。称 $\mathcal{D}_X$-模 $V$ 是拟凝聚的，若它是拟凝聚 $\mathcal{O}_X$-模。

以 $\mathcal{M}_{qc}(\mathcal{D}_X)$ 记由拟凝聚 $\mathcal{D}_X$-模构成的 $\mathcal{M}(\mathcal{D}_X)$ 的全子范畴。则 $\mathcal{M}_{qc}(\mathcal{D}_X)$ 是 Abel 范畴。

再设 $X$ 是仿射簇。以 $D_X=\Gamma(X,\mathcal{D}_X)$ 记 $X$ 上整体微分算子环。则由 Serre 定理 $\Gamma:\mathcal{M}_{qc}(\mathcal{D}_X)\to\mathcal{M}(D_X)$ 是正合函子。此外 $\Delta$ 是从 $\mathcal{M}(D_X)$ 到 $\mathcal{M}_{qc}(\mathcal{D}_X)$ 的函子。有如下 Serre 定理的类似物。

\begin{theorem}
设 $X$ 是仿射簇。则 $\Gamma:\mathcal{M}_{qc}(\mathcal{D}_X)\to\mathcal{M}(D_X)$ 是范畴的等价。局部化函子 $\Delta:\mathcal{M}(D_X)\to\mathcal{M}_{qc}(\mathcal{D}_X)$ 是其拟逆。
\end{theorem}

\begin{proof}
设 $V\in\mathcal{M}(D_X)$。则存在 $D_X$-模的正合列 $D_X^{(I)}\to D_X^{(J)}\to V\to 0$，应用 $\Delta$ 后得到 $\mathcal{D}_X$-模的正合列 $\mathcal{D}_X^{(I)}\to\mathcal{D}_X^{(J)}\to\Delta(V)\to 0$。函子 $\Gamma\circ\Delta$ 是 $\mathcal{M}(D_X)$ 到自身的右正合函子。此外对任意 $V\in\mathcal{M}(D_X)$ 有伴随态射 $\varphi_V:V\to\Gamma(X,\Delta(V))$。断言它是同构。显然对任意自由 $D_X$-模 $F$，$\varphi_F:F\to\Gamma(X,\Delta(F))$ 是同构。因此取 $D_X$-模的正合列 $D_X^{(I)}\to D_X^{(J)}\to V\to 0$，得到如下 $D_X$-模的交换图
$$\begin{tikzcd}
D_X^{(I)} \arrow[r] \arrow[d,"\varphi_{D_X^{(I)}}"'] & D_X^{(J)} \arrow[r] \arrow[d,"\varphi_{D_X^{(J)}}"'] & V \arrow[r] \arrow[d,"\varphi_V"] & 0\\
D_X^{(I)} \arrow[r] & D_X^{(J)} \arrow[r] & \Gamma(X,\Delta(V)) \arrow[r] & 0
\end{tikzcd}$$
其行正合且前两个竖直箭头是同构，因此 $\varphi_V$ 是同构。

现考虑另一伴随态射。对任意拟凝聚 $\mathcal{D}_X$-模 $V$，存在自然态射 $\psi_V:\Delta(\Gamma(X,V))\to V$。断言它是同构。考虑拟凝聚 $\mathcal{D}_X$-模的正合列
$$0\longrightarrow K\longrightarrow\Delta(\Gamma(X,V))\xrightarrow{\ \psi_V\ }V\longrightarrow C\longrightarrow 0.$$
由于 $\Gamma(X,-)$ 正合，得到 $D_X$-模的正合列
$$0\to\Gamma(X,K)\to\Gamma(X,\Delta(\Gamma(X,V)))\to\Gamma(X,V)\to\Gamma(X,C)\to 0.$$
由证明的第一部分，中间箭头是同构，故 $\Gamma(X,K)=\Gamma(X,C)=0$。由 Serre 定理最终断定 $K=C=0$。
\end{proof}

回顾 $X$ 上层 $F$ 的\emph{支集} $\operatorname{supp}(F)$ 是使 $F|_U=\{0\}$ 的最大开集 $U$ 的补集，因此层的支集是闭的。设 $U$ 是 $X$ 中开集，$s$ 是 $F$ 在 $U$ 上的局部截面，则 $\operatorname{supp}(s)$ 是使 $s|_V=0$ 的最大开集 $V\subset U$ 的补集。

\begin{lemma}
对任意 $s\in F(U)$，有
$$\operatorname{supp}(s)=\{x\in U\mid s_x\neq 0\}.$$
\end{lemma}

\begin{proof}
显然若 $x\notin\operatorname{supp}(s)$，存在 $x$ 的开邻域 $V$ 使 $s|_V=0$ 且 $s_x=0$。反之若 $s_x=0$，存在 $x$ 的开邻域 $V\subset U$ 使 $s|_V=0$ 且 $x\notin\operatorname{supp}(s)$。
\end{proof}

此外有如下结果。

\begin{proposition}
设 $F$ 是 $X$ 上的层。则 $\operatorname{supp}(F)$ 是集合 $\{x\in X\mid F_x\neq 0\}$ 的闭包。
\end{proposition}

\begin{proof}
显然 $\operatorname{supp}(F)$ 必含其所有截面的支集，故 $\{x\in X\mid F_x\neq 0\}$ 必含于支集中。由于 $F$ 的支集是闭的，它也含 $\{x\in X\mid F_x\neq 0\}$ 的闭包。设 $x$ 在此闭包之外，则存在 $x$ 的开邻域 $U$ 使对所有 $y\in U$ 有 $F_y=0$。但这意味着 $F|_U=0$ 且 $U\cap\operatorname{supp}(F)=\varnothing$，故 $x\notin\operatorname{supp}(F)$。
\end{proof}

设 $X=k^n$，$V$ 是 $D(n)$-模。则 $\mathcal{V}=\Delta(V)$ 是拟凝聚 $\mathcal{D}_X$-模。此外对任意 $x\in X$ 有 $\mathcal{V}_x=V_x$。因此由上述命题，$\operatorname{supp}(\mathcal{V})$ 等于 I.4 意义下 $\operatorname{supp}(V)$ 在 $X$ 中的闭包。

\begin{proposition}
设
$$0\longrightarrow F_1\longrightarrow F_2\longrightarrow F_3\longrightarrow 0$$
是 $X$ 上层短正合列。则
$$\operatorname{supp}(F_2)=\operatorname{supp}(F_1)\cup\operatorname{supp}(F_3).$$
\end{proposition}

\begin{proof}
由假设，对所有 $x\in X$ 序列
$$0\longrightarrow F_{1,x}\longrightarrow F_{2,x}\longrightarrow F_{3,x}\longrightarrow 0$$
正合。因此
$$\{x\in X\mid F_{2,x}\neq 0\}=\{x\in X\mid F_{1,x}\neq 0\}\cup\{x\in X\mid F_{3,x}\neq 0\}.$$
故取闭包并用上述命题即得结论。
\end{proof}
\begin{annote}
$\mathcal D_X=\varinjlim\mathcal D_{X,\le d}$ 是 $\mathcal O_X$-拟凝聚层；它在光滑点处局部同构于 $\mathcal O_X\otimes\mathrm{Sym}\,\mathcal T_X$ 的完成，但整体上不是 $\mathcal O_X$-凝聚层。讲义此后凡说“$\mathcal D_X$-模”都指左模，右模与双边模的换边靠 $\omega_X$ 与 $\mathcal D_X$ 的自反性。[来源：Stacks Project Tag 01BD（Section 17.10, Quasi-coherent modules）、Tag 01BU（Section 17.12, Coherent modules）；原书 §1, p77–86]
\end{annote}



\section{凝聚 $\mathcal{D}_X$-模}

现设 $X$ 是光滑仿射簇，则 $D_X$ 是诺特环。因此 $\mathcal{M}(\mathcal{D}_X)$ 中由有限生成 $\mathcal{D}_X$-模构成的全子范畴 $\mathcal{M}_{fg}(\mathcal{D}_X)$ 是 Abel 范畴。称 $V$ 是\emph{凝聚 $\mathcal{D}_X$-模}，若 $V\cong\Delta(V_0)$，其中 $V_0$ 是某有限生成 $D_X$-模。以 $\mathcal{M}_{coh}(\mathcal{D}_X)$ 记由凝聚 $\mathcal{D}_X$-模构成的 $\mathcal{M}(\mathcal{D}_X)$ 的全子范畴。显然 $\Gamma$ 把 $\mathcal{M}_{coh}(\mathcal{D}_X)$ 映入 $\mathcal{M}_{fg}(\mathcal{D}_X)$，$\Delta$ 把 $\mathcal{M}_{fg}(\mathcal{D}_X)$ 映入 $\mathcal{M}_{coh}(\mathcal{D}_X)$。因此 $\Gamma:\mathcal{M}_{coh}(\mathcal{D}_X)\to\mathcal{M}_{fg}(\mathcal{D}_X)$ 是范畴的等价，$\Delta:\mathcal{M}_{fg}(\mathcal{D}_X)\to\mathcal{M}_{coh}(\mathcal{D}_X)$ 是其拟逆。故此时可把凝聚性看作有限生成性的层化版本。

\begin{lemma}
设 $X$ 是光滑仿射簇，$V$ 是拟凝聚 $\mathcal{D}_X$-模。则下列条件等价：
\begin{enumerate}[label=(\roman*)]
  \item $V$ 是凝聚 $\mathcal{D}_X$-模；
  \item 对任意 $x\in X$ 存在 $x$ 的开邻域 $U$ 和正合列 $\mathcal{D}_U^p\to\mathcal{D}_U^q\to V|_U\to 0$。
\end{enumerate}
\end{lemma}

\begin{proof}
(i)$\Rightarrow$(ii) 设 $V$ 凝聚，则 $V\cong\Delta(V_0)$，其中 $V_0$ 是有限生成 $D_X$-模。由于 $D_X$ 诺特，存在正合列 $D_X^p\to D_X^q\to V_0\to 0$（$p,q\in\mathbb{Z}_+$）。局部化得正合列 $\mathcal{D}_X^p\to\mathcal{D}_X^q\to\Delta(V_0)\to 0$。由于 $V\cong\Delta(V_0)$，存在正合列 $\mathcal{D}_X^p\to\mathcal{D}_X^q\to V\to 0$。故对任意 $x\in X$ 可取 $U=X$。

(ii)$\Rightarrow$(i) 存在 $f\in R(X)$ 使 $f(x)\neq 0$ 且 $X_f\subset U$。故缩小 $U$ 可设为主开集。则 (ii) 蕴含序列 $\mathcal{D}_U^p\to\mathcal{D}_U^q\to\Gamma(U,V)\to 0$ 正合，即 $\Gamma(U,V)$ 是有限生成 $D_U$-模。现 $\Gamma(U,V)=\Gamma(X,V)_f$，且存在 $v_1,\dots,v_n\in\Gamma(X,V)$ 使其在 $U$ 上的限制作为 $D_U$-模生成 $V|_U$。所有这些主开集构成 $X$ 的开覆盖。由于 $X$ 拟紧，可取有限子覆盖，从而存在 $w_1,\dots,w_m\in\Gamma(X,V)$ 使得每个茎 $V_x$ 作为 $\mathcal{D}_{X,x}$-模由它们的像生成。因此有满射 $\mathcal{D}_X^m\to V$，进而有满射 $D_X^m\to\Gamma(X,V)$，$\Gamma(X,V)$ 是有限生成 $D_X$-模。故 $V$ 凝聚。
\end{proof}

设 $X$ 是任意光滑代数簇。称 $X$ 上拟凝聚 $\mathcal{D}_X$-模 $V$ 是凝聚的，若对任意 $x\in X$ 存在 $x$ 的开邻域 $U$ 和正合列 $\mathcal{D}_U^p\to\mathcal{D}_U^q\to V|_U\to 0$。由引理 2.1，此定义与仿射簇情形的先前定义一致。此外前面的结果蕴含如下命题。

\begin{proposition}
设 $V$ 是光滑代数簇 $X$ 上的拟凝聚 $\mathcal{D}_X$-模。则下列条件等价：
\begin{enumerate}[label=(\roman*)]
  \item $V$ 是凝聚 $\mathcal{D}_X$-模；
  \item 对任意仿射开子集 $U\subset X$，限制 $V|_U$ 是凝聚 $\mathcal{D}_U$-模；
  \item 对 $X$ 的仿射开覆盖 $(U_1,\dots,U_n)$，限制 $V|_{U_i}$（$1\le i\le n$）是凝聚 $\mathcal{D}_{U_i}$-模。
\end{enumerate}
\end{proposition}

以 $\mathcal{M}_{coh}(\mathcal{D}_X)$ 记 $\mathcal{M}_{qc}(\mathcal{D}_X)$ 中由凝聚 $\mathcal{D}_X$-模构成的全子范畴。则 $\mathcal{M}_{coh}(\mathcal{D}_X)$ 是 Abel 范畴。

对凝聚 $\mathcal{D}_X$-模可改进前面的支集结果。

\begin{proposition}
设 $V$ 是凝聚 $\mathcal{D}_X$-模。则
$$\operatorname{supp}(V)=\{x\in X\mid V_x\neq 0\}.$$
\end{proposition}

\begin{proof}
由前面的命题，只需证 $\{x\in X\mid V_x\neq 0\}$ 是闭集。设 $y$ 是其闭包中一点，$U$ 是 $y$ 的仿射邻域。则 $V(U)$ 是有限生成 $D_U$-模。设 $s_1,\dots,s_n$ 是 $V(U)$ 中作为 $D_U$-模生成它的截面，则这些截面也作为 $\mathcal{D}_U$-模生成 $V|_U$。令 $Z=\bigcup_{i=1}^n\operatorname{supp}(s_i)$，则 $Z$ 是 $U$ 中含于 $\{x\in X\mid V_x\neq 0\}$ 的闭子集。假设 $y\notin Z$，则存在 $y$ 的开邻域 $W\subset U$ 使 $s_1,\dots,s_n$ 在 $W$ 上消失。由此可知 $V|_W=0$，$y\notin\operatorname{supp}(V)$，与前面的命题矛盾。故 $y\in Z$。
\end{proof}

设 $U$ 是 $X$ 中开集，$F$ 是 $\mathcal{O}_X$-模，$G$ 是 $F|_U$ 的 $\mathcal{O}_U$-子模。以 $\bar G$ 记由
$$\bar G(V)=\{s\in F(V)\mid s|_{V\cap U}\in G(V\cap U)\}$$
定义的 $F$ 的子层。显然 $\bar G$ 是 $\mathcal{O}_X$-子模，称为 $G$ 的\emph{典范扩张}。

\begin{lemma}
设 $F$ 是拟凝聚 $\mathcal{O}_X$-模，$G$ 是 $F|_U$ 的拟凝聚 $\mathcal{O}_U$-子模。则 $G$ 的典范扩张 $\bar G$ 是拟凝聚 $\mathcal{O}_X$-模。
\end{lemma}

\begin{proof}
设 $i$ 是 $U$ 到 $X$ 的自然包含。以 $H$ 记 $F|_U$ 模 $G$ 的商，则 $H$ 是拟凝聚 $\mathcal{O}_U$-模。考虑拟凝聚 $\mathcal{O}_X$-模的自然态射 $\alpha:i_*(F|_U)\to i_*(H)$。它与典范态射 $F\to i_*(F|_U)$ 的复合定义拟凝聚 $\mathcal{O}_X$-模态射 $\varphi:F\to i_*(H)$。故其核拟凝聚，且
\begin{align*}
\ker\varphi(V)&=\{s\in F(V)\mid \varphi_V(s)=0\}\\
&=\{s\in F(V)\mid \alpha_{V\cap U}(s|_{V\cap U})=0\}\\
&=\{s\in F(V)\mid s|_{V\cap U}\in G(V\cap U)\}=\bar G(V).
\end{align*}
\end{proof}

设 $V$ 是拟凝聚 $\mathcal{D}_X$-模，$W$ 是 $V|_U$ 的拟凝聚 $\mathcal{D}_U$-子模。则 $W$ 的典范扩张 $\bar W$ 是 $V$ 的拟凝聚 $\mathcal{D}_X$-子模。

非零拟凝聚 $\mathcal{D}_X$-模 $V$ 称为\emph{不可约的}，若 $V$ 的任意拟凝聚 $\mathcal{D}_X$-子模 $W$ 或为 $\{0\}$ 或等于 $V$。

\begin{lemma}
设 $U$ 是 $X$ 中开集，$V$ 是不可约拟凝聚 $\mathcal{D}_X$-模。则 $V|_U$ 或为不可约拟凝聚 $\mathcal{D}_U$-模，或为零。
\end{lemma}

\begin{proof}
设 $V|_U\neq 0$。设 $W$ 是 $V|_U$ 的拟凝聚 $\mathcal{D}_U$-子模。以 $\bar W$ 记它到 $V$ 的 $\mathcal{D}_X$-子模的典范扩张。由于 $V$ 不可约，$\bar W$ 或为 $V$ 或为 $0$。这意味着 $W$ 或为 $V|_U$ 或为 $0$。
\end{proof}

特别地有如下推论。

\begin{proposition}
设 $V$ 是不可约拟凝聚 $\mathcal{D}_X$-模。则 $V$ 是凝聚的。
\end{proposition}

\begin{proof}
设 $U$ 是 $X$ 中仿射开集。则由前面的引理 $V|_U$ 或不可约或为 $0$。若 $V|_U$ 不可约，由 1.1，$\Gamma(U,V)$ 必为不可约 $D_U$-模。故 $V|_U$ 是凝聚 $\mathcal{D}_U$-模。结论由前面的命题得出。
\end{proof}

\begin{proposition}
设 $V$ 是不可约拟凝聚 $\mathcal{D}_X$-模。则支集 $\operatorname{supp}(V)$ 是 $X$ 的不可约闭子簇。
\end{proposition}

\begin{proof}
由定义 $\operatorname{supp}(V)$ 是 $X$ 的闭子簇。先断言 $\operatorname{supp}(V)$ 连通。假设 $\operatorname{supp}(V)$ 是 $X$ 的两个闭子簇 $Z_1$ 和 $Z_2$ 的不交并，且 $Z_1\neq\varnothing$。令 $U=X-Z_1$，以 $W$ 记 $V|_U$ 的零 $\mathcal{D}_U$-子模的典范扩张。由于 $V$ 不可约，$W$ 或为 $V$ 或为 $0$。取 $x\in Z_1$ 和不与 $Z_2$ 相交的 $x$ 的仿射开邻域 $W'$。则 $V|_{W'}$ 的支集等于 $Z_1\cap W'=(X-U)\cap W'=W'-(W'\cap U)$。另一方面，
$$\Gamma(W',W)=\{s\in\Gamma(W',V)\mid s|_{W'\cap U}=0\}=\Gamma(W',V),$$
故 $W\neq 0$，从而 $W=V$，$V|_U=0$，故 $Z_2=\varnothing$。因此 $\operatorname{supp}(V)$ 连通。

现证 $\operatorname{supp}(V)$ 不可约。反设不然。设 $Z_1$ 是 $\operatorname{supp}(V)$ 的一个不可约分支，$Z_2$ 是其余所有不可约分支的并，则 $\operatorname{supp}(V)=Z_1\cup Z_2$。令 $Z=Z_1\cap Z_2$，$U=X-Z$，则 $V|_U\neq 0$。由前面的引理它是不可约 $\mathcal{D}_U$-模。显然其支集等于 $(Z_1\cup Z_2)-(Z_1\cap Z_2)=(Z_1-(Z_1\cap Z_2))\cup(Z_2-(Z_1\cap Z_2))$。由前述结果此空间必连通，故 $Z_2-(Z_1\cap Z_2)=\varnothing$。由此可知 $Z_2\subset Z_1$，矛盾。
\end{proof}

拟凝聚 $\mathcal{D}_X$-模 $V$ 称为\emph{有限长度的}，若它有由拟凝聚 $\mathcal{D}_X$-模构成的有限递增滤过
$$\{0\}=V_0\subsetneq V_1\subsetneq\cdots\subsetneq V_n=V,$$
使得 $V_p/V_{p-1}$ 是不可约 $\mathcal{D}_X$-模。数 $\ell(V)=n$ 称为 $V$ 的\emph{长度}。显然由长度上的归纳和前面的命题，任意有限长度的拟凝聚 $\mathcal{D}_X$-模是凝聚的。

\begin{lemma}
设 $V$ 是拟凝聚 $\mathcal{D}_X$-模。则下列条件等价：
\begin{enumerate}[label=(\roman*)]
  \item $V$ 具有有限长度；
  \item 对任意开子集 $U\subset X$，限制 $V|_U$ 具有有限长度；
  \item 存在 $X$ 的开覆盖 $(U_i;\ 1\le i\le n)$ 使得 $V|_{U_i}$（$1\le i\le n$）具有有限长度。
\end{enumerate}
\end{lemma}

\begin{proof}
显然由前面的引理 (i) 蕴含 (ii)，(ii) 蕴含 (iii)。

为证 (iii) 蕴含 (i)，对 $\sum_{j=1}^n\ell(V|_{U_j})$ 归纳。若此和为 $0$，则对所有 $j$ 有 $V|_{U_j}=0$，$V=0$。设此和严格为正。若 $V$ 不可约则已完成。否则存在非平凡拟凝聚 $\mathcal{D}_X$-子模 $U$，即有正合列
$$0\longrightarrow U\longrightarrow V\longrightarrow W\longrightarrow 0,$$
其中 $U$ 和 $W$ 均非零。由于对 $1\le j\le n$ 有 $\ell(V|_{U_j})=\ell(U|_{U_j})+\ell(W|_{U_j})$，可见
$$\sum_{j=1}^n\ell(V|_{U_j})=\sum_{j=1}^n\ell(U|_{U_j})+\sum_{j=1}^n\ell(W|_{U_j}),$$
且右边两项均不为零。故归纳假设可用于两者，由此可知 $U$ 和 $W$ 具有有限长度，从而 $V$ 具有有限长度。
\end{proof}

\section{特征簇}

本节把特征簇的构造推广到任意凝聚 $\mathcal{D}_X$-模。

先设 $X$ 是光滑仿射簇，$D_X$ 是 $X$ 上相应的微分算子环。由 II.3.7，$D_X$ 是左诺特且右诺特的环。此外它有按微分算子阶数的自然滤过 $F D_X$，由 II.3.6，作为配置这样滤过的滤过环满足 I.3 的公理。

设 $\pi:T^*(X)\to X$ 是 $X$ 的余切丛。由于 $\pi$ 是局部平凡纤维化，态射 $\pi$ 是仿射的，因此 $T^*(X)$ 是仿射簇。如 II.3.6 的证明中所述，有
$$\operatorname{Gr}D_X=\Gamma(X,\operatorname{Gr}\mathcal{D}_X)=\Gamma(X,\pi_*(\mathcal{O}_{T^*(X)}))=\Gamma(T^*(X),\mathcal{O}_{T^*(X)})=R(T^*(X)).$$

任意有限生成 $D_X$-模 $V_0$ 有好 $D_X$-模滤过 $F V_0$，且 $\Gr V_0$ 是 $\Gr D_X=R(T^*(X))$ 上的有限生成模。设 $I$ 是 $\Gr V_0$ 的零化子。则由 I.10.2，根 $r(I)$ 不依赖于 $V_0$ 上好滤过的选择。称之为 $V_0$ 的\emph{特征理想}，记为 $J(V_0)$。$J(V_0)$ 在 $T^*(X)$ 中的零点集称为 $V_0$ 的\emph{特征簇}，记为 $\operatorname{Ch}(V_0)$。这些定义与 I.10 中 $k^n$ 上微分算子模的定义一致。

现在把这些概念层化。设 $V$ 是 $X$ 上凝聚 $\mathcal{D}_X$-模，则称 $V$ 的特征簇 $\operatorname{Ch}(V)$ 为 $D_X$-模 $\Gamma(X,V)$ 的特征簇。

称 $V$ 由凝聚 $\mathcal{O}_X$-子模构成的递增 $\mathcal{D}_X$-模滤过 $F V$ 是\emph{好的}，若
\begin{enumerate}[label=(\roman*)]
  \item 对充分负的 $n\in\mathbb{Z}$，$F_nV=\{0\}$；
  \item 滤过 $F V$ 是穷尽的（即 $\bigcup_{n\in\mathbb{Z}}F_nV=V$）；
  \item 滤过 $F V$ 是稳定的，即存在 $m_0\in\mathbb{Z}$ 使得对所有 $n\in\mathbb{Z}_+$ 和 $m\ge m_0$ 有 $F_nD_X\,F_mV=F_{m+n}V$。
\end{enumerate}
设 $F V$ 是 $V$ 的好滤过，则 $\Gamma(X,F_pV)$ 是 $\Gamma(X,V)$ 的有限生成 $R(X)$-子模，且 $(\Gamma(X,F_pV);\ p\in\mathbb{Z})$ 是 $D_X$-模 $\Gamma(X,V)$ 的好滤过。

\begin{lemma}
设 $V$ 是凝聚 $\mathcal{D}_X$-模。则
\begin{enumerate}[label=(\roman*)]
  \item $V$ 有好滤过；
  \item 映射 $(F_pV;\ p\in\mathbb{Z})\longmapsto(\Gamma(X,F_pV);\ p\in\mathbb{Z})$ 是 $V$ 的好滤过之集到 $\Gamma(X,V)$ 的好滤过之集的双射。
\end{enumerate}
\end{lemma}

\begin{proof}
(i) $D_X$-模 $\Gamma(X,V)$ 有限生成，故它有好滤过 $F\Gamma(X,V)$。由 Serre 定理 $V=\Delta(\Gamma(X,V))$，$F_pV=\Delta(F_p\Gamma(X,V))$ 自然等同于 $V$ 的凝聚 $\mathcal{O}_X$-子模。直接检验可知 $F V$ 是 $V$ 上的好滤过。

(ii) 由 Serre 定理立即得出。
\end{proof}

设 $F V$ 是 $V$ 的好滤过，则 $\Gr V$ 是 $\Gr\mathcal{D}_X=\pi_*(\mathcal{O}_{T^*(X)})$ 上的模。此外由于 $X$ 仿射，有
$$\Gamma(X,\Gr V)=\Gr\Gamma(X,V),$$
其中 $\Gamma(X,V)$ 配置好滤过 $\{\Gamma(X,F_pV),\ p\in\mathbb{Z}\}$。由于 $\Gr\Gamma(X,V)$ 是有限生成 $R(T^*(X))$-模，由 I.4.2 对 $A=R(T^*(X))$ 的显然推广得
$$\operatorname{Ch}(V)=\operatorname{Ch}(\Gamma(X,V))=\operatorname{supp}(\Gr\Gamma(X,V))=\operatorname{supp}(\Gamma(X,\Gr V)).$$
这意味着特征簇的构造本质上是局部的，即有如下结果。设 $U$ 是 $X$ 中开集，则可把 $T^*(U)$ 看作 $T^*(X)$ 的开子集 $\pi^{-1}(U)=\{(x,\omega)\in T^*(X)\mid \omega\in T_x^*(X),\ x\in U\}$。

\begin{lemma}
设 $U$ 是 $X$ 中仿射开集。则
$$\operatorname{Ch}(V|_U)=\operatorname{Ch}(V)\cap\pi^{-1}(U).$$
\end{lemma}

现考虑一般情形。设 $X$ 是光滑代数簇，$\mathcal{D}_X$ 是 $X$ 上微分算子层，$\pi:T^*(X)\to X$ 是余切丛。设 $V$ 是凝聚 $\mathcal{D}_X$-模。

\begin{lemma}
存在唯一的 $T^*(X)$ 闭子簇 $M$，使得对任意仿射开集 $U\subset X$ 有 $\operatorname{Ch}(V|_U)=M\cap\pi^{-1}(U)$。
\end{lemma}

\begin{proof}
设 $U$ 和 $W$ 是 $X$ 的两个仿射开子集。则由前面的引理，
$$\operatorname{Ch}(V|_U)\cap\pi^{-1}(U\cap W)=\operatorname{Ch}(V|_{U\cap W})=\operatorname{Ch}(V|_W)\cap\pi^{-1}(U\cap W).$$
因此由所有满足如下条件的对 $(x,\omega)$（$\omega\in T_x^*(X)$）构成的集合 $M$ 是良定义的且具有所需性质：对 $x$ 的某仿射开邻域 $U$ 有 $(x,\omega)\in\operatorname{Ch}(V|_U)$。
\end{proof}

上一引理描述的簇 $M$ 即为 $V$ 的特征簇 $\operatorname{Ch}(V)$。

显然凝聚 $\mathcal{D}_X$-模好滤过的定义即使 $X$ 不仿射也有意义。现证明这种滤过的存在性。

先需要一个辅助结果。设 $X$ 是代数簇。证明凝聚 $\mathcal{O}_X$-子模可从开子簇扩张。

\begin{proposition}
设 $F$ 是 $X$ 上拟凝聚 $\mathcal{O}_X$-模，$G$ 是 $F|_U$ 的凝聚 $\mathcal{O}_U$-子模。则存在 $F$ 的凝聚 $\mathcal{O}_X$-子模 $G'$ 使得 $G'|_U=G$。
\end{proposition}

\begin{proof}
先设 $X$ 仿射。设 $\bar G$ 是 $G$ 的典范扩张。则 $\Gamma(X,\bar G)$ 是有限生成 $\Gamma(X,\mathcal{O}_X)$-模的递增族的直接极限。设 $\{H_i;\ i\in I\}$ 是这些子模的局部化。则它们构成 $F$ 的凝聚 $\mathcal{O}_X$-子模的递增系，其直接极限为 $\bar G$。由于 $\bar G|_U=G$ 凝聚，系 $\{H_i|_U;\ i\in I\}$ 稳定，即存在 $i_0\in I$ 使得对 $i\ge i_0$ 有 $H_i|_U=G$。

现考虑一般情形。对 $X$ 的有限仿射开覆盖的基数归纳证明。设 $(V_i;\ 1\le i\le n)$ 是 $X$ 的仿射开覆盖，$Y=\bigcup_{i=1}^{n-1}V_i$。则由归纳假设存在 $F|_Y$ 的凝聚 $\mathcal{O}_Y$-子模 $H$ 使得 $H|_{Y\cap U}=G|_{Y\cap U}$。$H$ 到限制于 $U$ 的 $F$ 的子模的典范扩张 $\bar H$ 含 $G$。应用证明的第一部分，存在 $\bar H|_{V_n}$ 的凝聚子模 $K$ 使得 $K|_{V_n\cap U}=G|_{V_n\cap U}$。设 $G'$ 为 $K$ 到 $\bar H$ 的子模的典范扩张。则 $G'|_U$ 含 $G$。此外 $G'|_{Y\cap U}\subset\bar H|_{Y\cap U}=G|_{Y\cap U}$，即 $G'|_{Y\cap U}=G|_{Y\cap U}$。又 $G'|_{V_n\cap U}=K|_{V_n\cap U}=G|_{V_n\cap U}$。因此 $G'|_U=G|_U$。另一方面，由于 $G'|_Y\subset\bar H|_Y=H$ 凝聚且 $G'|_{V_n}=K$ 也凝聚，$G'$ 是 $F$ 的凝聚 $\mathcal{O}_X$-子模。
\end{proof}

设 $X$ 是 $k$ 上的光滑代数簇，$\mathcal{D}_X$ 是 $X$ 上微分算子层。

\begin{theorem}
设 $V$ 是凝聚 $\mathcal{D}_X$-模。则 $V$ 有好滤过。
\end{theorem}

\begin{proof}
先断言存在凝聚 $\mathcal{O}_X$-子模 $U'\subset V$ 使得态射 $\mathcal{D}_X\otimes_{\mathcal{O}_X}U'\to V$ 是满射。设 $(U_i;\ 1\le i\le n)$ 是 $X$ 的仿射开覆盖。则对每个 $1\le i\le n$，$\Gamma(U_i,V)$ 是有限生成 $D_{U_i}$-模。由前面的命题，存在 $V$ 的凝聚 $\mathcal{O}_X$-子模 $G_i$ 使得 $\Gamma(U_i,G_i)$ 作为 $D_{U_i}$-模生成 $\Gamma(U_i,V)$。因此它们的和具有所需性质。

现可定义 $F_nV$ 为 $F_n\mathcal{D}_X\otimes_{\mathcal{O}_X}U'$ 在态射 $\mathcal{D}_X\otimes_{\mathcal{O}_X}U'\to V$ 下的像。显然 $F V$ 是 $V$ 的好滤过。
\end{proof}

设 $V$ 是凝聚 $\mathcal{D}_X$-模，$F V$ 是 $V$ 的好滤过。设 $\Gr V$ 是相应的分次 $\pi_*(\mathcal{O}_{T^*(X)})$-模。则对任意开集 $U\subset X$，滤过 $F_p(V|_U)=(F_pV)|_U$（$p\in\mathbb{Z}$）是 $V|_U$ 的好滤过。此外有 $\Gr V|_U=\Gr(V|_U)$。因此在 $X$ 的仿射开集 $U$ 上有 $\Gamma(U,\Gr V)=\Gr\Gamma(U,V)$。如前所述，簇 $T^*(U)$ 是仿射的且 $\Gamma(U,\pi_*(\mathcal{O}_{T^*(X)}))=R(T^*(U))$。故若把 $\Gamma(U,\Gr V)$ 作为 $R(T^*(U))$-模局部化，得到 $U$ 上唯一的 $\mathcal{O}_{T^*(U)}$-模 $\widetilde{V}_U$，具有性质 $\pi_*(\widetilde{V}_U)=\Gr V|_U$。由于 $\Gamma(U,\Gr V)$ 是有限生成 $R(T^*(U))$-模，$\widetilde{V}_U$ 是凝聚 $\mathcal{O}_{T^*(U)}$-模。把 $\widetilde{V}_U$ 粘在一起，得到 $T^*(X)$ 上唯一的凝聚 $\mathcal{O}_{T^*(X)}$-模 $\widetilde V$，具有性质 $\pi_*(\widetilde V)=\Gr V$。

这立即蕴含光滑仿射簇上凝聚模特征簇公式的如下推广。

\begin{proposition}
设 $V$ 是凝聚 $\mathcal{D}_X$-模，$F V$ 是 $V$ 的好滤过，$\Gr V$ 是相应的分次 $\pi_*(\mathcal{O}_{T^*(X)})$-模。则
\begin{enumerate}[label=(\roman*)]
  \item 存在唯一的 $T^*(X)$ 上凝聚 $\mathcal{O}_{T^*(X)}$-模 $\widetilde V$，使得 $\pi_*(\widetilde V)=\Gr V$；
  \item $\operatorname{Ch}(V)=\operatorname{supp}(\widetilde V)$。
\end{enumerate}
\end{proposition}

以下结果是 I.10.4 的推广。由于断言是局部的，可在仿射开集上检验，那里的论证与 I.10.4 相同。

\begin{proposition}
设
$$0\longrightarrow V_1\longrightarrow V_2\longrightarrow V_3\longrightarrow 0$$
是凝聚 $\mathcal{D}_X$-模的短正合列。则
$$\operatorname{Ch}(V_2)=\operatorname{Ch}(V_1)\cup\operatorname{Ch}(V_3).$$
\end{proposition}

设 $M$ 是 $X$ 的余切簇 $T^*(X)$ 的子簇。若 $(x,\omega)\in M$ 蕴含对所有 $\lambda\in k$ 有 $(x,\lambda\omega)\in M$，则称 $M$ 是\emph{锥形子簇}。

I.10.3 和 I.10.7 的推广由下列陈述给出。由于断言是局部的，只需在 II.2.10 意义下“小”的仿射开集 $U$ 上检验。此时 $\pi^{-1}(U)\cong U\times k^n$（$n=\dim_xX$），$R(T^*(U))\cong R(U)\otimes_kk[\xi_1,\dots,\xi_n]=R(U)[\xi_1,\dots,\xi_n]$，且分次是多项式环的自然分次。$\Gr\Gamma(U,V)$ 的零化子是 $R(U)[\xi_1,\dots,\xi_n]$ 中的齐次理想。在此情形下 I.10.3 和 I.10.7 证明的必要修改是直接的。

\begin{proposition}
设 $V$ 是 $X$ 上凝聚 $\mathcal{D}_X$-模。则特征簇 $\operatorname{Ch}(V)$ 是 $T^*(X)$ 的锥形子簇。
\end{proposition}

\begin{theorem}
设 $V$ 是 $X$ 上凝聚 $\mathcal{D}_X$-模。则
$$\pi(\operatorname{Ch}(V))=\operatorname{supp}(V).$$
\end{theorem}

\section{Coherentor（凝聚化函子）}

本节介绍代数簇上 $\mathcal{O}$-模的一些基本定义和结果。

先设 $X$ 是代数闭域 $k$ 上的任意代数簇。以 $\mathcal{M}(\mathcal{O}_X)$ 记 $X$ 上 $\mathcal{O}_X$-模的范畴，以 $\mathcal{M}_{qc}(\mathcal{O}_X)$ 记其拟凝聚 $\mathcal{O}_X$-模的全子范畴。整体截面函子 $\Gamma(X,-)$ 是从 $\mathcal{M}(\mathcal{O}_X)$ 到 $X$ 上正则函数环 $R(X)$ 上模范畴 $\mathcal{M}(R(X))$ 的左正合函子。由文献 [1] III.2.7，$\Gamma(X,-)$ 的右同调维数 $\le\dim X$。

\begin{lemma}
从 $\mathcal{M}_{qc}(\mathcal{O}_X)$ 到 $\mathcal{M}(\mathcal{O}_X)$ 的遗忘函子 $\operatorname{For}$ 有右伴随函子 $Q_X:\mathcal{M}(\mathcal{O}_X)\to\mathcal{M}_{qc}(\mathcal{O}_X)$。
\end{lemma}

\begin{proof}
只需证明对任意 $W\in\mathcal{M}(\mathcal{O}_X)$ 存在 $Q(W)\in\mathcal{M}_{qc}(\mathcal{O}_X)$ 使得
$$\operatorname{Hom}_{\mathcal{O}_X}(V,W)=\operatorname{Hom}_{\mathcal{O}_X}(V,Q(W))$$
对任意 $V\in\mathcal{M}_{qc}(\mathcal{O}_X)$ 成立。

先设 $X$ 是仿射簇。则 $\Gamma(X,-)$ 是 $\mathcal{M}_{qc}(\mathcal{O}_X)$ 与 $\mathcal{M}(R(X))$ 的范畴等价。此外对任意 $R(X)$-模 $M$，以 $\widetilde M=\mathcal{O}_X\otimes_{R(X)}M$ 记其局部化。则 $M\to\widetilde M$ 是从 $\mathcal{M}(R(X))$ 到 $\mathcal{M}_{qc}(\mathcal{O}_X)$ 的正合函子，且是 $\Gamma(X,-)$ 的拟逆。因此对任意 $V\in\mathcal{M}_{qc}(\mathcal{O}_X)$，
\begin{align*}
\operatorname{Hom}_{\mathcal{O}_X}(V,W)&=\operatorname{Hom}_{\mathcal{O}_X}(\widetilde{\Gamma(X,V)},W)\\
&=\operatorname{Hom}_{\mathcal{O}_X}(\mathcal{O}_X\otimes_{R(X)}\Gamma(X,V),W)=\operatorname{Hom}_{R(X)}(\Gamma(X,V),\operatorname{Hom}_{\mathcal{O}_X}(\mathcal{O}_X,W))\\
&=\operatorname{Hom}_{R(X)}(\Gamma(X,V),\Gamma(X,W))=\operatorname{Hom}_{\mathcal{O}_X}(V,\widetilde{\Gamma(X,W)}),
\end{align*}
此时 $Q_X(W)=\widetilde{\Gamma(X,W)}$。

现设 $U$ 是仿射簇 $X$ 的开仿射子集，$i:U\to X$ 是自然浸入。则对任意 $V\in\mathcal{M}_{qc}(\mathcal{O}_X)$ 和 $W\in\mathcal{M}(\mathcal{O}_U)$，
$$\operatorname{Hom}_{\mathcal{O}_X}(V,i_*(W))=\operatorname{Hom}_{\mathcal{O}_U}(V|_U,W)=\operatorname{Hom}_{\mathcal{O}_U}(V|_U,Q_U(W))=\operatorname{Hom}_{\mathcal{O}_X}(V,i_*(Q_U(W))).$$
由于直像保持拟凝聚性（文献 [1] II.5.8），$i_*(Q_U(W))\in\mathcal{M}_{qc}(\mathcal{O}_X)$，且
$$Q_X(i_*(W))=i_*(Q_U(W)),$$
即函子 $Q_X\circ i_*$ 与 $i_*\circ Q_U$ 同构。

现考虑一般情形。设 $X$ 是任意簇，$U=(U_1,U_2,\dots,U_n)$ 是其仿射开覆盖。以 $f_i:U_i\to X$，$f_{ij}:U_i\cap U_j\to X$ 记自然浸入。由相应的 Čech 分解，对任意 $W\in\mathcal{M}(\mathcal{O}_X)$ 有正合列
$$0\longrightarrow W\longrightarrow\bigoplus_{i=1}^n f_{i*}(W|_{U_i})\longrightarrow\bigoplus_{i<j}f_{ij*}(W|_{U_i\cap U_j}).$$
固定 $1\le i\le n$，以 $g_{ij}:U_i\cap U_j\to U_i$ 记自然包含（$j\neq i$）。则 $f_{ij}=f_i\circ g_{ij}$，且 $f_{i*}(W|_{U_i})\to\bigoplus_{i\neq j}f_{ij*}(W|_{U_i\cap U_j})$ 是把 $f_i$ 的直像应用于 $W|_{U_i}\to\bigoplus_{j\neq i}g_{ij*}(W|_{U_i\cap U_j})$ 得到的。由于 $U_i$ 和 $U_i\cap U_j$ 是开仿射子簇，对此态射应用 $Q_{U_i}$ 并用前段结果，得到态射 $Q_{U_i}(W|_{U_i})\to\bigoplus_{j\neq i}g_{ij*}(Q_{U_i\cap U_j}(W|_{U_i\cap U_j}))$，使得对任意 $V\in\mathcal{M}_{qc}(\mathcal{O}_X)$ 图
$$\begin{tikzcd}
\operatorname{Hom}_{\mathcal{O}_X}(V,f_{i*}(W|_{U_i})) \arrow[r] \arrow[d] & \operatorname{Hom}_{\mathcal{O}_X}(V,f_{ij*}(W|_{U_i\cap U_j})) \arrow[d]\\
\operatorname{Hom}_{\mathcal{O}_{U_i}}(V|_{U_i},W|_{U_i}) \arrow[r] \arrow[d] & \operatorname{Hom}_{\mathcal{O}_{U_i}}(V|_{U_i},g_{ij*}(W|_{U_i\cap U_j})) \arrow[d]\\
\operatorname{Hom}_{\mathcal{O}_{U_i}}(V|_{U_i},Q_{U_i}(W|_{U_i})) \arrow[r] & \operatorname{Hom}_{\mathcal{O}_{U_i}}(V|_{U_i},g_{ij*}(Q_{U_i\cap U_j}(W|_{U_i\cap U_j})))\\
\operatorname{Hom}_{\mathcal{O}_X}(V,f_{i*}(Q_{U_i}(W|_{U_i}))) \arrow[u] \arrow[r] & \operatorname{Hom}_{\mathcal{O}_X}(V,f_{ij*}(Q_{U_i\cap U_j}(W|_{U_i\cap U_j}))) \arrow[u]
\end{tikzcd}$$
交换。Čech 分解 $C^\cdot(U,W)$ 的第一个微分
$$d^1:\bigoplus_{i=1}^nf_{i*}(W|_{U_i})\longrightarrow\bigoplus_{i<j}f_{ij*}(W|_{U_i\cap U_j})$$
由此对应决定态射
$$\delta:\bigoplus_{i=1}^nf_{i*}(Q_{U_i}(W|_{U_i}))\longrightarrow\bigoplus_{i<j}f_{ij*}(Q_{U_i\cap U_j}(W|_{U_i\cap U_j})).$$
以 $Q(W)$ 记此态射的核。显然它是拟凝聚 $\mathcal{O}_X$-模。由函子 $\operatorname{Hom}_{\mathcal{O}_X}(V,-)$ 的左正合性得 $\operatorname{Hom}_{\mathcal{O}_X}(V,W)=\operatorname{Hom}_{\mathcal{O}_X}(V,Q(W))$。
\end{proof}

函子 $Q_X:\mathcal{M}(\mathcal{O}_X)\to\mathcal{M}_{qc}(\mathcal{O}_X)$ 称为\emph{凝聚化函子}（coherentor）。

\begin{proposition}
\begin{enumerate}[label=(\roman*)]
  \item 函子 $Q_X:\mathcal{M}(\mathcal{O}_X)\to\mathcal{M}_{qc}(\mathcal{O}_X)$ 是左正合的；
  \item $Q_X$ 把 $\mathcal{M}(\mathcal{O}_X)$ 中的入射对象映为 $\mathcal{M}_{qc}(\mathcal{O}_X)$ 中的入射对象；
  \item 复合 $Q_X\circ\operatorname{For}$ 同构于 $\mathcal{M}_{qc}(\mathcal{O}_X)$ 上的恒等函子；
  \item 设 $U$ 是 $X$ 的开子簇，$i:U\to X$ 是自然包含。则函子 $Q_X\circ i_*$ 与 $i_*\circ Q_U$ 同构。
\end{enumerate}
\end{proposition}

\begin{proof}
(i) 这是任意右伴随函子的性质。

(ii) 这是正合函子的右伴随的性质。

(iii) 由 $Q_X$ 的定义显然。

(iv) 对任意 $V\in\mathcal{M}_{qc}(\mathcal{O}_X)$ 和 $W\in\mathcal{M}(\mathcal{O}_U)$ 有
$$\operatorname{Hom}_{\mathcal{O}_X}(V,i_*(W))=\operatorname{Hom}_{\mathcal{O}_U}(V|_U,W)=\operatorname{Hom}_{\mathcal{O}_U}(V|_U,Q_U(W))=\operatorname{Hom}_{\mathcal{O}_X}(V,i_*(Q_U(W))).$$
由于直像保持拟凝聚性，$i_*(Q_U(W))\in\mathcal{M}_{qc}(\mathcal{O}_X)$，且 $Q_X(i_*(W))=i_*(Q_U(W))$，即函子 $Q_X\circ i_*$ 与 $i_*\circ Q_U$ 同构。
\end{proof}

\begin{theorem}
范畴 $\mathcal{M}_{qc}(\mathcal{O}_X)$ 有充足的入射对象。
\end{theorem}

\begin{proof}
设 $V$ 是拟凝聚 $\mathcal{O}_X$-模。则存在入射 $\mathcal{O}_X$-模 $I$ 和单射 $\epsilon:V\to I$。由于 $Q_X$ 由前面的命题 (i) 左正合，$Q_X(\epsilon):Q_X(V)\to Q_X(I)$ 是单射。此外由 (ii)，$Q_X(I)$ 在 $\mathcal{M}_{qc}(\mathcal{O}_X)$ 中入射。最后由 (iii)，$Q_X(V)\cong V$。
\end{proof}

\begin{lemma}
设 $X$ 是仿射簇，$W\in\mathcal{M}(\mathcal{O}_X)$。则 $R^pQ_X(W)=\widetilde{H^p(X,W)}$。
\end{lemma}

\begin{proof}
在引理 4.1 的证明中已见 $Q_X$ 同构于局部化函子与整体截面函子 $\Gamma(X,-)$ 的复合。设 $I^\cdot$ 是 $W$ 的入射分解。则 $H^p(X,W)=H^p(\Gamma(X,I^\cdot))$。由于局部化函子正合，
$$\widetilde{H^p(X,W)}=H^p(\widetilde{\Gamma(X,I^\cdot)})=H^p(Q_X(I^\cdot))=R^pQ_X(W).$$
\end{proof}

以下引理在凝聚化函子各种性质的证明中至关重要。

\begin{lemma}
设 $V$ 是满足如下条件的 $\mathcal{O}_X$-模：

(V) 对任意仿射开集 $U\subset X$ 和 $p\ge 1$ 有上同调 $H^p(U,V)=0$。

则对任意仿射开集 $U\subset X$，直像层 $i_*(V|_U)$ 是 $Q_X$-无环的。
\end{lemma}

\begin{proof}
先注意高阶直像 $R^pi_*(V|_U)$ 为零。固定 $p>0$。对 $U$ 上任意层 $F$，$R^pi_*(F)$ 是预层 $V\longmapsto H^p(i^{-1}(V),F)=H^p(U\cap V,F)$ 附着的层。若 $V$ 是仿射开集，$U\cap V$ 也是仿射的，故 $H^p(U\cap V,V)=0$。这意味着此预层在所有仿射开集上为零。由于仿射开集构成 $X$ 拓扑的基，相应的层为零。

设 $I^\cdot$ 是 $V|_U$ 在 $\mathcal{M}(\mathcal{O}_X)$ 中的入射分解。由前注复形 $i_*(I^\cdot)$ 是 $i_*(V|_U)$ 的分解。此外由于 $i_*$ 是限制函子 $\mathcal{M}(\mathcal{O}_X)\to\mathcal{M}(\mathcal{O}_U)$ 的右伴随，它把入射映为入射，即 $i_*(I^\cdot)$ 是 $i_*(V|_U)$ 在 $\mathcal{M}(\mathcal{O}_X)$ 中的入射分解。因此
$$R^pQ_X(i_*(V|_U))=H^p(Q_X(i_*(I^\cdot)))=H^p(i_*(Q_U(I^\cdot))),$$
由命题 (iv)。由 4.4 对 $p>0$，
$$H^p(Q_U(I^\cdot))=R^pQ_U(V|_U)=\widetilde{H^p(U,V)}=0.$$
因此 $Q_U(I^\cdot)$ 是由拟凝聚 $\mathcal{O}_U$-模构成的无环复形。由于 $i:U\to X$ 是仿射态射，$i_*$ 是从 $\mathcal{M}_{qc}(\mathcal{O}_U)$ 到 $\mathcal{M}_{qc}(\mathcal{O}_X)$ 的正合函子。故复形 $i_*(Q_U(I^\cdot))$ 也无环，即 $R^pQ_X(i_*(V|_U))=0$（$p>0$）。
\end{proof}

\begin{proposition}
拟凝聚 $\mathcal{O}_X$-模对 $Q_X$ 无环。
\end{proposition}

\begin{proof}
设 $U=(U_0,U_1,\dots,U_n)$ 是 $X$ 的仿射开覆盖。设 $V$ 是拟凝聚 $\mathcal{O}_X$-模，则 $V$ 满足前一引理的条件 (V)。因此 $V$ 的 Čech 复形 $C^\cdot(U,V)$ 中的模都是 $Q_X$-无环的。可知 $C^\cdot(U,V)$ 是 $V$ 的 $Q_X$-无环分解。故由命题 (iii)，
$$R^pQ_X(V)=H^p(Q_X(C^\cdot(U,V)))=H^p(C^\cdot(U,V))=0\quad(p>0).$$
\end{proof}

此结果蕴含如下基本事实。

\begin{theorem}
设 $V,W$ 是两个拟凝聚 $\mathcal{O}_X$-模。则
$$\operatorname{Ext}^\cdot_{\mathcal{M}_{qc}(\mathcal{O}_X)}(V,W)=\operatorname{Ext}^\cdot_{\mathcal{M}(\mathcal{O}_X)}(V,W).$$
\end{theorem}

\begin{proof}
设
$$0\longrightarrow W\longrightarrow I^0\longrightarrow I^1\longrightarrow\cdots\longrightarrow I^n\longrightarrow\cdots$$
是 $W$ 在 $\mathcal{M}(\mathcal{O}_X)$ 中的入射分解。则由命题 (ii)、(iii) 和前一命题，
$$0\longrightarrow W\longrightarrow Q_X(I^0)\longrightarrow Q_X(I^1)\longrightarrow\cdots\longrightarrow Q_X(I^n)\longrightarrow\cdots$$
是 $W$ 在 $\mathcal{M}_{qc}(\mathcal{O}_X)$ 中的入射分解。故对任意 $p\in\mathbb{Z}_+$，
$$\operatorname{Ext}^p_{\mathcal{M}(\mathcal{O}_X)}(V,W)=H^p(\operatorname{Hom}_{\mathcal{O}_X}(V,I^\cdot))=H^p(\operatorname{Hom}_{\mathcal{O}_X}(V,Q_X(I^\cdot)))=\operatorname{Ext}^p_{\mathcal{M}_{qc}(\mathcal{O}_X)}(V,W).$$
\end{proof}

因此可无混淆地记
$$\operatorname{Ext}^\cdot_{\mathcal{O}_X}(V,W)=\operatorname{Ext}^\cdot_{\mathcal{M}_{qc}(\mathcal{O}_X)}(V,W)=\operatorname{Ext}^\cdot_{\mathcal{M}(\mathcal{O}_X)}(V,W)$$
对任意两个拟凝聚 $\mathcal{O}_X$-模 $V,W$。

\begin{theorem}
$Q_X$ 的右同调维数有限。
\end{theorem}

\begin{proof}
固定 $X$ 的仿射开覆盖 $U=(U_0,U_1,\dots,U_n)$。

先观察如下事实：设 $V$ 是满足引理 4.5 条件 (V) 的 $\mathcal{O}_X$-模，则对 $p>n$ 有 $R^pQ_X(V)=0$。

如前证明，由引理 4.5，$V$ 的 Čech 复形 $C^\cdot(U,V)$ 中的模都是 $Q_X$-无环的，故 $C^\cdot(U,V)$ 是 $Q_X$-无环分解。因此 $R^pQ_X(V)=H^p(Q_X(C^\cdot(U,V)))$，由此对 $p>n$ 有 $R^pQ_X(V)=0$。

现建立此结果的如下推广：设 $V$ 是满足如下条件的 $\mathcal{O}_X$-模：

(V$_q$) 对任意仿射开集 $U\subset X$ 和 $p>q$ 有 $H^p(U,V)=0$。

则对 $p>q+n$ 有 $R^pQ_X(V)=0$。

已对 $q=0$ 建立此结果。为证归纳步，对 $p$ 归纳。设结论对某 $q\ge 0$ 成立。设 $V$ 是使 $H^p(U,V)=0$（对 $p>q+1$ 和任意仿射开集 $U\subset X$）的 $\mathcal{O}_X$-模。取短正合列
$$0\longrightarrow V\longrightarrow I\longrightarrow W\longrightarrow 0,$$
其中 $I$ 是入射 $\mathcal{O}_X$-模。则对任意仿射开集 $U\subset X$，限制 $I|_U$ 是入射 $\mathcal{O}_U$-模。因此由上同调长正合列得连接同态 $H^p(U,W)\to H^{p+1}(U,V)$ 对 $p\ge 1$ 是同构。这意味着 $H^p(U,W)=0$（$p>q$），即归纳假设可用于 $W$。由 $Q_X$ 的导出函子的长正合列得连接态射 $R^pQ_X(W)\to R^{p+1}Q_X(V)$ 对 $p\ge 1$ 是同构。由于 $p>q+n$ 时 $R^pQ_X(W)=0$，$p>q+1+n$ 时 $R^pQ_X(V)=0$。故 (V$_{q+1}$) 成立。

由于函子 $\Gamma(U,-)$ 的右同调维数 $\le\dim U\le\dim X$，可见所有 $\mathcal{O}_X$-模满足 (V$_{\dim X}$)。因此对任意拟凝聚 $\mathcal{O}_X$-模 $V$，当 $p>\dim X+n$ 时 $R^pQ_X(V)=0$。
\end{proof}
\begin{annote}
$\mathcal Q$（coherator）是 $\mathcal D_X$-模层到凝聚 $\mathcal D_X$-模的右伴随，$\mathcal Q(\mathcal M)=\varinjlim\mathcal H om_{\mathcal D_X}(\mathcal E,\mathcal M)$，$\mathcal E$ 取遍凝聚 $\mathcal D_X$-模。它修补了“$\mathcal D_X$ 非凝聚导致的层论不灵”，使 $\mathcal D_X$-模可以直接在 $D^+(\mathcal D_X)$ 里做内射消解。Stacks 与 nLab 均无独立条目，本条只作原书内部的记号澄清。[来源：原书 §4, p90–96；译者按]
\end{annote}



\section{射影空间上的 $D$-模}

设 $X=P^n$ 是 $k$ 上 $n$ 维射影空间。记 $Y=k^{n+1}$，$Y^*=Y-\{0\}$。设 $\pi:Y^*\to X$ 是自然投射
$$\pi(x_0,x_1,\dots,x_n)=[x_0,x_1,\dots,x_n].$$
令 $U_0=\{[x_0,x_1,\dots,x_n]\in P^n\mid x_0\neq 0\}$。则 $U_0$ 是 $P^n$ 中同构于 $k^n$ 的开集，同构 $c:k^n\to U_0$ 由 $c(x_1,x_2,\dots,x_n)=[1,x_1,x_2,\dots,x_n]$ 给出。此外 $\pi^{-1}(U_0)=\{(y_0,y_1,y_2,\dots,y_n)\mid y_0\neq 0\}$。因此由
$$\varphi(t,z_1,z_2,\dots,z_n)=(t,tz_1,tz_2,\dots,tz_n)$$
定义的映射 $\varphi:k^*\times k^n\to\pi^{-1}(U_0)$ 是 $k^*\times k^n$ 与 $\pi^{-1}(U_0)$ 的同构，且满足 $\pi(\varphi(t,z_1,\dots,z_n))=[1,z_1,\dots,z_n]$。因此如下交换图交换：
$$\begin{tikzcd}
k^*\times k^n \arrow[r,"\varphi"] \arrow[d,"1\times c"' ] & \pi^{-1}(U_0) \arrow[d,"\pi"]\\
k^*\times U_0 \arrow[r,"\operatorname{pr}_2"'] & U_0
\end{tikzcd}$$
显然 $\alpha=\varphi\circ(1\times c)^{-1}:k^*\times U_0\to\pi^{-1}(U_0)$ 是由
$$\alpha(t,[1,x_1,x_2,\dots,x_n])=(t,tx_1,tx_2,\dots,tx_n)$$
给出的同构，它在开集 $U_0$ 上平凡化纤维化 $\pi:Y^*\to X$。由于 $\operatorname{GL}(n+1,k)$ 在 $P^n$ 上传递作用，这证明如下断言。

\begin{lemma}
态射 $\pi:Y^*\to X$ 是纤维同构于 $k^*$ 的局部平凡纤维化。
\end{lemma}

显然 $U_0$ 和 $\pi^{-1}(U_0)$ 是仿射簇。因此若 $V$ 是拟凝聚 $\mathcal{D}_X$-模，有
$$\Gamma\bigl(\pi^{-1}(U_0),\pi^+(V)\bigr)=R\bigl(\pi^{-1}(U_0)\bigr)\otimes_{R(U_0)}\Gamma(U_0,V).$$
另一方面，同构 $\alpha:k^*\times U_0\to\pi^{-1}(U_0)$ 诱导同构 $\alpha^*:R(\pi^{-1}(U_0))\to R(k^*)\otimes_kR(U_0)$。在此同构下，$R(U_0)$-模作用对应于第二个因子中的乘法。因此此同构诱导 $\Gamma(\pi^{-1}(U_0),\pi^+(V))$ 与 $R(k^*)\otimes_k\Gamma(U_0,V)$ 的同构。

再次利用 $\operatorname{GL}(n+1,k)$ 作用的传递性，立即得到如下推论。

\begin{lemma}
逆像函子 $\pi^+:\mathcal{M}_{qc}(\mathcal{D}_X)\to\mathcal{M}_{qc}(\mathcal{D}_{Y^*})$ 是正合的。
\end{lemma}

因此若
$$0\longrightarrow V_1\longrightarrow V_2\longrightarrow V_3\longrightarrow 0$$
是拟凝聚 $\mathcal{D}_X$-模的正合列，有拟凝聚 $\mathcal{D}_{Y^*}$-模的正合列
$$0\longrightarrow\pi^+(V_1)\longrightarrow\pi^+(V_2)\longrightarrow\pi^+(V_3)\longrightarrow 0.$$
此外若以 $j:Y^*\to Y$ 记自然浸入，有拟凝聚 $\mathcal{D}_Y$-模的正合列
$$0\longrightarrow j_+(\pi^+(V_1))\longrightarrow j_+(\pi^+(V_2))\longrightarrow j_+(\pi^+(V_3))\longrightarrow K\longrightarrow 0,$$
且 $K$ 支集在 $\{0\}$。由于 $Y$ 是仿射空间，这意味着
$$0\to\Gamma(Y,j_+(\pi^+(V_1)))\to\Gamma(Y,j_+(\pi^+(V_2)))\to\Gamma(Y,j_+(\pi^+(V_3)))\to\Gamma(Y,K)\to 0$$
是 $D(n+1)$-模的正合列。对任意 $\mathcal{D}_{Y^*}$-模 $W$，由于 $j_+(W)=j_\cdot(W)$，有 $\Gamma(Y,j_+(W))=\Gamma(Y^*,W)$。故断定
$$0\to\Gamma(Y^*,\pi^+(V_1))\to\Gamma(Y^*,\pi^+(V_2))\to\Gamma(Y^*,\pi^+(V_3))\to\Gamma(Y,K)\to 0$$
是 $D(n+1)$-模的正合列。

令 $E=\sum_{i=0}^ny_i\dfrac{\partial}{\partial y_i}$ 为 $Y$ 上 Euler 算子。微分算子 $E$ 是 $Y$ 上向量场。若 $y=(y_0,y_1,\dots,y_n)$ 是 $Y^*$ 中一点，$E(y)$ 是曲线 $t\longmapsto(ty_0,ty_1,\dots,ty_n)$ 在 $y$ 处的切向量。故有如下结果。

\begin{lemma}
对任意 $y\in Y^*$，Euler 算子 $E$ 在 $y$ 处的值属于微分 $\pi:Y^*\to X$ 的核。
\end{lemma}

在同构 $\alpha:k^*\times U_0\to\pi^{-1}(U_0)$ 下，Euler 算子对应于 $k^*\times U_0$ 上的微分算子 $t\partial_t$，其中 $t$ 是 $k^*$ 上自然坐标。显然 $R(k^*)$ 是环 $k((t))$，即多项式环 $k[t]$ 关于乘法系 $t^n$（$n\in\mathbb{Z}_+$）的局部化。因此在 $\alpha$ 给出的同构下，
$$\Gamma\bigl(\pi^{-1}(U_0),\pi^+(V)\bigr)\cong k((t))\otimes_k\Gamma(U_0,V).$$
因此 $\Gamma(\pi^{-1}(U_0),\pi^+(V))$ 的每个截面被 $\prod_{p\in S}(E-p)$ 零化，其中 $S$ 是依赖于该截面的 $\mathbb{Z}$ 的有限子集。

\begin{lemma}
设 $V$ 是拟凝聚 $\mathcal{D}_X$-模。则 $\Gamma(Y^*,\pi^+(V))$ 是特征值属于 $\mathbb{Z}$ 的 $E$-特征空间的直和。
\end{lemma}

\begin{proof}
由于 $U_0$ 在 $\operatorname{GL}(n+1,k)$ 作用下的有限多个平移覆盖 $X$，$\pi^{-1}(U_0)$ 的有限多个平移覆盖 $Y^*$。由上述论证，对 $V$ 的任意整体截面 $v$，存在有限子集 $S\subset\mathbb{Z}$ 使 $\prod_{p\in S}(E-p)v=0$。
\end{proof}

因此所考虑的正合列在 $E$ 作用下分裂为对应于 $E$ 不同特征值的无穷族正合列。特别地若以 $\Gamma(Y^*,\pi^+(V))_{(p)}$ 记 $\Gamma(Y^*,\pi^+(V))$ 中特征值 $p\in\mathbb{Z}$ 的 $E$-特征空间，则对任意 $p\in\mathbb{Z}$，
$$0\to\Gamma(Y^*,\pi^+(V_1))_{(p)}\to\Gamma(Y^*,\pi^+(V_2))_{(p)}\to\Gamma(Y^*,\pi^+(V_3))_{(p)}\to\Gamma(Y,K)_{(p)}\to 0.$$

另一方面，有自然线性映射 $\gamma_U:\Gamma(U,V)\to\Gamma(\pi^{-1}(U),\pi^+(V))$，它与限制相容，即对 $X$ 中任意两个开集 $V\subset U$ 图
$$\begin{tikzcd}
\Gamma(U,V) \arrow[r,"\gamma_U"] \arrow[d,"r_{V,U}"'] & \Gamma(\pi^{-1}(U),\pi^+(V)) \arrow[d,"r_{\pi^+(V),\pi^{-1}(U)}"]\\
\Gamma(V,V) \arrow[r,"\gamma_V"'] & \Gamma(\pi^{-1}(V),\pi^+(V))
\end{tikzcd}$$
交换。显然在上述等同下 $\gamma_{U_0}$ 对应于映射 $v\longmapsto 1\otimes v|_{U_0}$（从 $\Gamma(U_0,V)$ 到 $k((t))\otimes_k\Gamma(U_0,V)$）。因此它是 $\Gamma(U_0,V)$ 到 $\Gamma(\pi^{-1}(U_0),\pi^+(V))_{(0)}$ 的同构。故可把 $\gamma=\gamma_X$ 看作从 $\Gamma(X,V)$ 到 $\Gamma(Y^*,\pi^+(V))_{(0)}$ 的线性映射。

显然 $\gamma(v)=0$ 蕴含 $\gamma(v)|_{U_0}=0$，从而 $v|_{U_0}=0$。故再次利用 $\operatorname{GL}(n+1,k)$ 在 $X$ 上作用的传递性得 $v=0$。由此可知 $\gamma:\Gamma(X,V)\to\Gamma(Y^*,\pi^+(V))_{(0)}$ 是单射。

\begin{lemma}
设 $V$ 是拟凝聚 $\mathcal{D}_X$-模。则 $\gamma:\Gamma(X,V)\to\Gamma(Y^*,\pi^+(V))_{(0)}$ 是同构。
\end{lemma}

\begin{proof}
设 $s$ 是 $\Gamma(Y^*,\pi^+(V))_{(0)}$ 的整体截面。则由上述讨论，其在 $\pi^{-1}(U_0)$ 上的限制等于某截面 $v_0\in V(U_0)$ 的 $\gamma_{U_0}(v_0)$。由于 $U_0$ 在 $\operatorname{GL}(n+1,k)$ 作用下的有限多个平移覆盖 $X$，可见存在 $X$ 的覆盖 $U_0,\dots,U_m$ 和 $V$ 在这些开集上的截面 $v_0,\dots,v_m$，使得 $\gamma_{U_k}(v_k)=s|_{\pi^{-1}(U_k)}$（$0\le k\le m$）。由于 $\gamma_{U_k}$ 是单射且 $U_k$ 和 $\pi^{-1}(U_k)$ 是仿射簇，由局部化得 $\gamma_{U_k\cap U_l}$ 是单射。故对每对 $0\le k<l\le m$ 有 $v_k|_{U_k\cap U_l}=v_l|_{U_k\cap U_l}$。由此可知存在 $V$ 的整体截面 $v$ 使得对所有 $0\le k\le m$ 有 $v|_{U_k}=v_k$。因此
$$\gamma(v)|_{\pi^{-1}(U_k)}=\gamma_{U_k}(v|_{U_k})=\gamma_{U_k}(v_k)=s|_{\pi^{-1}(U_k)}$$
对 $0\le k\le m$ 成立，$\gamma(v)=s$。这证明 $\gamma$ 的满射性。
\end{proof}

以下结果由直接计算得出。

\begin{lemma}
对任意 $0\le i\le n$，
\begin{enumerate}[label=(\roman*)]
  \item $[E,y_i]=y_i$；
  \item $\bigl[E,\dfrac{\partial}{\partial y_i}\bigr]=-\dfrac{\partial}{\partial y_i}$。
\end{enumerate}
\end{lemma}

\begin{lemma}
设 $Y=k^{n+1}$，$V$ 是支集在 $\{0\}$ 的拟凝聚 $\mathcal{D}_Y$-模。则 $\Gamma(Y,V)$ 是特征值属于 $\{-(n+k)\mid k\in\mathbb{N}\}$ 的 $E$-特征空间的直和。
\end{lemma}

\begin{proof}
设 $V_0=\Gamma(Y,V)$。由 I.13.7 知 $V_0$ 由子空间 $W$ 生成，其中 $W$ 由被 $y_i$（$0\le i\le n$）零化的所有整体截面构成。取 $v\in W$，则 $v=[\partial_i,y_i]v=-y_i\partial_iv$，故 $Ev=-(n+1)v$。此外 $V_0=\bigoplus_{I\in\mathbb{Z}_+^{n+1}}\partial^I(W)$。由前面的结果，对任意 $w=\partial^Iv$（$v\in W$）有
$$Ew=E\partial^Iv=\partial^IEv-|I|\partial^Iv=\bigl(-(n+1)-|I|\bigr)\partial^Iv=-(n+1+|I|)w,$$
即 $V_0$ 是 $E$ 的属于特征值 $\{-(n+k)\mid k\in\mathbb{N}\}$ 的特征空间的直和。
\end{proof}

特别地，由于 $K$ 支集在 $\{0\}$，有 $\Gamma(Y,K)_{(0)}=\{0\}$，且
$$0\longrightarrow\Gamma(X,V_1)\longrightarrow\Gamma(X,V_2)\longrightarrow\Gamma(X,V_3)\longrightarrow 0$$
正合。因此 $\Gamma(X,-)$ 在 $\mathcal{M}_{qc}(\mathcal{D}_X)$ 上是正合函子。

\begin{theorem}
设 $V$ 是 $n$ 维射影空间 $X=P^n$ 上拟凝聚 $\mathcal{D}_X$-模。则对 $p\in\mathbb{N}$ 有 $H^p(X,V)=0$。
\end{theorem}

\begin{proof}
$V$ 有由入射拟凝聚 $\mathcal{D}_X$-模构成的右分解 $I^\cdot$。入射拟凝聚 $\mathcal{D}_X$-模是 $\Gamma$-无环的。因此 $H^p(X,V)=H^p(\Gamma(X,I^\cdot))=0$（$p>0$），因为 $\Gamma$ 在拟凝聚 $\mathcal{D}_X$-模上正合。
\end{proof}

因此若以 $D_X=\Gamma(X,\mathcal{D}_X)$ 记 $X$ 上整体微分算子，$\Gamma:\mathcal{M}_{qc}(\mathcal{D}_X)\to\mathcal{M}(D_X)$ 是正合函子。还可定义函子 $\Delta:\mathcal{M}(D_X)\to\mathcal{M}(D_X)$ 为 $\Delta(V)=D_X\otimes_{D_X}V$，则 $\Delta$ 是右正合的。设 $V\in\mathcal{M}(D_X)$，则存在正合列 $D_X^{(I)}\to D_X^{(J)}\to V\to 0$，应用 $\Delta$ 后得到 $D_X$-模的正合列 $D_X^{(I)}\to D_X^{(J)}\to\Delta(V)\to 0$。因此 $\Delta(V)$ 是拟凝聚 $\mathcal{D}_X$-模。

函子 $\Gamma\circ\Delta$ 是 $\mathcal{M}(D_X)$ 到自身的右正合函子。此外对任意 $V\in\mathcal{M}(D_X)$ 存在自然态射 $\lambda_V:V\to\Gamma(X,\Delta(V))$。显然 $\lambda$ 是恒等函子到 $\Gamma\circ\Delta$ 的自然变换。

\begin{lemma}
自然变换 $\lambda$ 是 $\mathcal{M}(D_X)$ 上恒等函子到函子 $\Gamma\circ\Delta$ 的同构。
\end{lemma}

\begin{proof}
显然对任意自由 $D_X$-模 $F$，$\lambda_F:F\to\Gamma(X,\Delta(F))$ 是同构。因此取 $D_X$-模的正合列 $D_X^{(I)}\to D_X^{(J)}\to V\to 0$，得到如下 $D_X$-模的交换图
$$\begin{tikzcd}
D_X^{(I)} \arrow[r] \arrow[d,"\lambda_{D_X^{(I)}}"'] & D_X^{(J)} \arrow[r] \arrow[d,"\lambda_{D_X^{(J)}}"'] & V \arrow[r] \arrow[d,"\lambda_V"] & 0\\
D_X^{(I)} \arrow[r] & D_X^{(J)} \arrow[r] & \Gamma(X,\Delta(V)) \arrow[r] & 0
\end{tikzcd}$$
其行正合且前两个竖直箭头是同构，因此 $\lambda_V$ 是同构。
\end{proof}

\begin{lemma}
设 $V$ 是拟凝聚 $\mathcal{D}_X$-模。若 $\Gamma(X,V)=0$，则 $V=0$。
\end{lemma}

\begin{proof}
设 $\Gamma(X,V)=0$。由前面的结果这意味着 $\Gamma(Y^*,\pi^+(V))_{(0)}=0$。断言实际上 $\Gamma(Y^*,\pi^+(V))=0$。反设不然，则存在 $s\in\mathbb{Z}-\{0\}$ 使 $\Gamma(Y^*,\pi^+(V))_{(s)}\neq 0$。若 $s>0$，取 $0\neq v\in\Gamma(Y^*,\pi^+(V))_{(s)}$，由前面的引理 (ii) 对 $0\le i\le n$ 有 $E\partial_iv=(s-1)\partial_iv$。显然所有 $\partial_iv=0$（$0\le i\le n$）不可能，因为那会蕴含 $Ev=0$。因此 $\Gamma(Y^*,\pi^+(V))_{(s-1)}\neq 0$，由 $s$ 上的递降归纳得矛盾。故 $s$ 必为负。此时若 $0\neq v\in\Gamma(Y^*,\pi^+(V))_{(s)}$，由前面的引理 (i) 对 $0\le i\le n$ 有 $Ey_iv=(s+1)y_iv$。若 $y_iv=0$，截面 $v$ 的支集含于 $Y^*$ 与第 $i$ 个坐标超平面的交中。由于所有坐标超平面与 $Y^*$ 的交为空，至少对某 $0\le i\le n$ 有 $y_iv\neq 0$。这意味着 $\Gamma(Y^*,\pi^+(V))_{(s+1)}\neq 0$，由 $s$ 上的归纳再次得矛盾。

由此可知 $\Gamma(Y,j_+(\pi^+(V)))=\Gamma(Y^*,\pi^+(V))=0$，且由于 $Y$ 是仿射簇，$j_+(\pi^+(V))=0$。这进而蕴含 $\pi^+(V)=0$，$V=0$。
\end{proof}

对任意拟凝聚 $\mathcal{D}_X$-模 $V$，存在自然态射 $\mu_V:\Delta(\Gamma(X,V))\to V$。显然 $\mu$ 是函子 $\Delta\circ\Gamma$ 到 $\mathcal{M}_{qc}(\mathcal{D}_X)$ 上恒等函子的自然变换。

\begin{lemma}
自然变换 $\mu$ 是函子 $\Delta\circ\Gamma$ 到 $\mathcal{M}_{qc}(\mathcal{D}_X)$ 上恒等函子的同构。
\end{lemma}

\begin{proof}
考虑拟凝聚 $\mathcal{D}_X$-模的正合列
$$0\longrightarrow K\longrightarrow\Delta(\Gamma(X,V))\longrightarrow V\longrightarrow C\longrightarrow 0.$$
由于 $\Gamma(X,-)$ 正合，得到 $D_X$-模的正合列
$$0\to\Gamma(X,K)\to\Gamma(X,\Delta(\Gamma(X,V)))\to\Gamma(X,V)\to\Gamma(X,C)\to 0.$$
由前面的引理中间箭头是同构，故 $\Gamma(X,K)=\Gamma(X,C)=0$。由此最终断定 $K=C=0$。
\end{proof}

这立即蕴含如下结果。

\begin{theorem}
函子 $\Gamma(X,-)$ 是范畴 $\mathcal{M}_{qc}(\mathcal{D}_X)$ 与 $\mathcal{M}(D_X)$ 的等价。其逆为 $\Delta$。
\end{theorem}

\begin{corollary}
$X=P^n$ 上任意拟凝聚 $\mathcal{D}_X$-模由其整体截面生成。
\end{corollary}

现在把这些结果推广到光滑仿射簇与射影空间的乘积。设 $X=P^n$，$Y$ 是光滑仿射簇，$\pi:X\times Y\to X$ 是自然投影。则 $\pi$ 是仿射态射。事实上若 $U\subset X$ 是开仿射子簇，$\pi^{-1}(U)=U\times Y$ 是 $X\times Y$ 的开仿射子簇。因此对任意拟凝聚 $\mathcal{O}_X$-模 $V$，
$$H^p(X\times Y,V)=H^p(X,\pi_*(V)),\quad p\in\mathbb{Z}.$$

若 $V$ 是 $\mathcal{D}_{X\times Y}$-模，$U\subset X$ 是开仿射子集，则
$$\pi_*(V)(U)=V(\pi^{-1}(U))=V(U\times Y)=\Gamma(U\times Y,V)$$
是 $D(U\times Y)$-模。由于 $D(U\times Y)=D(U)\otimes_kD(Y)$，$\pi_*(V)(U)$ 有由映射 $T\longmapsto T\otimes1$ 诱导的自然 $D(U)$-模结构。因此 $\pi_*(V)$ 有自然的 $\mathcal{D}_X$-模结构，此结构与 $\mathcal{O}_X$-模结构相容。由于 $\pi_*$ 保持拟凝聚性，若 $V$ 是拟凝聚 $\mathcal{D}_{X\times Y}$-模，直像 $\pi_*(V)$ 是拟凝聚 $\mathcal{D}_X$-模。故
$$H^p(X\times Y,V)=H^p(X,\pi_*(V))=0,\quad p>0,$$
且函子 $\Gamma(X\times Y,-)$ 在 $\mathcal{M}_{qc}(\mathcal{D}_{X\times Y})$ 上正合。另一方面若 $0=\Gamma(X\times Y,V)=\Gamma(X,\pi_*(V))$，由前面的结果 $\pi_*(V)=0$。若 $U\subset X$ 是开仿射子集，$0=\pi_*(V)(U)=V(U\times Y)$。由于 $U\times Y$ 是仿射簇且 $V$ 拟凝聚，可知 $V|_{U\times Y}=0$。由于 $U$ 任意，这意味着 $V=0$。

这证明前一定理的如下推广。

\begin{theorem}
设 $X=P^n$，$Y$ 是光滑仿射簇。则函子 $\Gamma(X\times Y,-)$ 是范畴 $\mathcal{M}_{qc}(\mathcal{D}_{X\times Y})$ 与 $\mathcal{M}(D_{X\times Y})$ 的等价。其逆为 $\Delta$。
\end{theorem}

特别地，拟凝聚 $\mathcal{D}_{X\times Y}$-模由其整体截面生成。
\chapter{D-模的直像与逆像}

在本章中，$\varphi:X\to Y$ 总表示光滑代数簇之间的态射。我们以
$\mathcal{M}_L(\mathcal{D}_X)$（相应地 $\mathcal{M}_R(\mathcal{D}_X)$）记左（相应地右）
$\mathcal{D}_X$-模层的范畴，并略去下标时以 $D_X=\Gamma(X,\mathcal{D}_X)$ 记其整体微分算子环。

\section{双模 $\mathcal{D}_{X\to Y}$}

设 $X$ 是代数簇，$V$ 与 $W$ 是两个 $\mathcal{O}_X$-模。$V$ 到 $W$ 的 $k$-线性态射 $T$ 称为\emph{阶 $\le n$ 的微分态射}，如果对任意开集 $U$ 及任意 $(n+1)$-元组正则函数
$f_0,f_1,\dots,f_n\in R(U)$，在 $U$ 上有
$$[\dots[[T,f_0],f_1],\dots,f_n]=0.$$
以 $\Diff(V,W)$ 记 $V$ 到 $W$ 的所有微分态射构成的空间。同时规定：当 $p<0$ 时
$F_p\Diff(V,W)=0$，当 $p$ 为正时 $F_p\Diff(V,W)$ 为阶 $\le p$ 的微分态射构成的子空间。
显然 $F\Diff(V,W)$ 是 $\Diff(V,W)$ 穷尽的滤过。这个概念推广了 $X$ 上微分算子的概念：
若 $V=W=\mathcal{O}_X$，则 $\mathcal{O}_X$ 的微分自同态恰是 $X$ 上的微分算子。

类似地，可以定义 $V$ 到 $W$ 的微分态射层 $\mathcal{D}iff(V,W)$。

\begin{lemma}\label{lem:4.1.1}
设 $T$、$S$ 分别是阶 $\le n$、$\le m$ 的微分态射。则 $T\circ S$ 是阶 $\le n+m$ 的微分自同态。
\end{lemma}

\begin{proof}
对 $n+m$ 作归纳。若 $n=m=0$，则 $T$、$S$ 是 $\mathcal{O}_X$-模的态射，故 $T\circ S$ 也是
$\mathcal{O}_X$-模的态射，从而是阶 $\le 0$ 的微分态射。

现设 $n+m>0$。则
$$[T\circ S,f]=TSf-fTS=T[S,f]+[T,f]S,$$
且 $[T,f]$、$[S,f]$ 分别是阶 $\le n-1$ 与 $\le m-1$ 的微分态射。由归纳假设，这个微分态射的阶
$\le n+m-1$。因此 $T\circ S$ 的阶 $\le n+m$。
\end{proof}

因此，$\mathcal{O}_X$-模 $V$ 的所有微分自同态构成一个滤过环，而局部微分自同态构成一个滤过环层。

设 $X$、$Y$ 是光滑簇，$\varphi:X\to Y$ 是簇的态射。令
$\mathcal{D}_{X\to Y}=\varphi^*(\mathcal{D}_Y)$。则它既是 $\mathcal{O}_X$-模，又通过右乘成为右
$\varphi^{-1}\mathcal{D}_Y$-模。设 $\mathcal{C}$ 是 $\mathcal{O}_X$-模 $\mathcal{D}_{X\to Y}$ 的所有局部微分自同态、同时也是 $\varphi^{-1}\mathcal{D}_Y$-自同态的层。下面描述环层 $\mathcal{C}$ 的结构。

首先注意到 $\mathcal{O}_X$ 自然是 $F_0\mathcal{C}$ 的子环。

切层 $\mathcal{T}_Y$ 是 $\mathcal{D}_Y$ 的 $\mathcal{O}_Y$-子模。以 $\mathcal{J}_Y$ 记 $\mathcal{D}_Y$ 中由
$\mathcal{T}_Y$ 生成的左理想层。设 $y\in Y$，则存在 $y$ 的仿射开邻域 $U$ 及其上的坐标系
$(f_1,\dots,f_m;D_1,\dots,D_m)$，使得 $(D^I;I\in\mathbb{Z}_+^m)$ 是自由 $\mathcal{O}_U$-模 $\mathcal{D}_U$
关于左乘的一组基。左理想层 $\mathcal{J}_Y|_U$ 由 $(D^I;I\in\mathbb{Z}_+^m,\ |I|>0)$ 张成。因此有
$\mathcal{D}_Y=\mathcal{O}_Y\oplus\mathcal{J}_Y$。作为 $\mathcal{O}_X$-模，这导致直和分解
$$\mathcal{D}_{X\to Y}=\mathcal{O}_X\oplus\varphi^*(\mathcal{J}_Y).$$
以 $\alpha:\mathcal{D}_{X\to Y}\to\mathcal{O}_X$ 记相应的投影。

设 $S\in\mathcal{C}(U)$。我们断言 $\varphi^*(\mathcal{J}_Y)|_U$ 是 $S$-不变的。通过限制，可设 $U$ 是仿射的且
$\varphi(U)$ 包含在带有坐标系 $(f_1,\dots,f_m;D_1,\dots,D_m)$ 的“小”仿射开集 $V$ 中。于是
$\varphi^*(\mathcal{J}_Y)|_U$ 的整体截面是以 $(1\otimes D^I;I\in\mathbb{Z}_+^m,\ |I|>0)$ 为基的自由
$R(U)$-模。由于 $S$ 是右 $\mathcal{D}_V$-模 $\mathcal{D}_{X\to Y}(U)$ 的自同态，对任意 $g_I\in R(U)$ 有
$$S\Bigl(\sum_{|I|>0}g_I\otimes D^I\Bigr)=\sum_{|I|>0}S(g_I\otimes 1)D^I\in\varphi^*(\mathcal{J}_Y).$$
因此 $\varphi^*(\mathcal{J}_Y)|_U$ 在 $S$ 下不变。

由此可知 $\varphi^*(\mathcal{J}_Y)$ 是 $\mathcal{D}_{X\to Y}$ 的 $\mathcal{C}$-子模。这意味着商
$\mathcal{D}_{X\to Y}/\varphi^*(\mathcal{J}_Y)$ 是 $\mathcal{C}$-模。由上述讨论，自然单射
$\mathcal{O}_X\to\mathcal{D}_{X\to Y}$ 与该商映射的复合是 $\mathcal{O}_X$-模的同构。这个同构反过来定义了环层态射
$\gamma:\mathcal{C}\to\mathcal{E}nd_k(\mathcal{O}_X)$，由
$$\gamma_U(S)(f)=\alpha(S(f\otimes 1))$$
给出，其中 $f\in\mathcal{O}_X(U)$；显然 $\gamma$ 在 $\mathcal{O}_X$ 上为恒等。于是 $\gamma$ 把 $\mathcal{C}$ 映入
$\mathcal{O}_X$ 的微分自同态，即 $X$ 上的局部微分算子。因此构造出了滤过环层同态
$\gamma:\mathcal{C}\to\mathcal{D}_X$。本节的主要结果如下。

\begin{theorem}\label{thm:4.1.2}
映射 $\gamma:\mathcal{C}\to\mathcal{D}_X$ 是滤过环层的同构。
\end{theorem}

先证 $\gamma$ 是单射。从一个特殊情形开始。

\begin{lemma}\label{lem:4.1.3}
设 $U\subset X$ 是 $X$ 中的开集，$S$ 是 $\mathcal{C}(U)$ 中阶 $\le 0$ 的元素。则 $S=f\in\mathcal{O}_X(U)$。
\end{lemma}

\begin{proof}
由于 $S$ 的阶 $\le 0$，它实际上是 $\mathcal{O}_U$-模 $\mathcal{D}_{X\to Y}|_U$ 的自同态。只需证明 $S$ 在 $U$ 的一个开覆盖上的限制由函数给出即可。因此可设 $U$ 是仿射的且 $\varphi(U)$ 包含在 $Y$ 中带坐标系
$(f_1,\dots,f_m;D_1,\dots,D_m)$ 的“小”仿射开集 $V$ 内。更确切地说，由 2.10 与 3.5，可设 $Y$ 上存在坐标系
$(f_1,\dots,f_m;D_1,\dots,D_m)$ 使得
\begin{enumerate}[label=(\roman*)]
  \item $[D_i,D_j]=0$，$1\le i\le m$；
  \item $D^I$（$I\in\mathbb{Z}_+^m$）是自由 $R(V)$-模 $\mathcal{D}_V$（关于左乘）的一组基。
\end{enumerate}
于是在此情形下 $\mathcal{D}_{X\to Y}(U)=R(U)\otimes_{R(V)}\mathcal{D}_V$ 是以 $1\otimes D^I$（$I\in\mathbb{Z}_+^m$）为基的自由 $R(U)$-模。由于 $S$ 与 $R(X)$ 元素的左乘及 $\mathcal{D}_Y$ 元素的右乘都可交换，它完全由其在 $1\otimes 1$ 上的值决定。设
$S(1\otimes 1)=\sum a_I\otimes D^I$，其中 $a_I\in R(X)$。则
$$\sum a_I\otimes D^I f_j=S(1\otimes 1)f_j=S(1\otimes f_j)=S(f_j\circ\varphi\otimes 1)
=\sum a_I(f_j\circ\varphi)\otimes D^I=\sum a_I\otimes f_jD^I$$
对任意 $1\le j\le m$ 成立。这意味着
$$\sum a_I\otimes[D^I,f_j]=0$$
对任意 $1\le j\le m$ 成立。设 $I=(i_1,\dots,i_m)$。显然若 $i_j=0$ 则 $[D^I,f_j]=0$；若 $i_j>0$，令
$I'=(i_1,\dots,i_{j-1},i_j-1,i_{j+1},\dots,i_m)$，则 $[D^I,f_j]=i_jD^{I'}$。这立即使对所有 $I\ne 0$ 有 $a_I=0$。因此 $S=a_0$。
\end{proof}

现在证明 $\gamma_U$ 的单射性。设 $S\in\mathcal{C}(U)$ 且 $\gamma_U(S)=0$。对 $S$ 的阶 $p$ 作归纳证明 $S=0$。若 $p=0$，由引理 \ref{lem:4.1.3} 有 $S=f\in\mathcal{O}_X(U)$ 且 $\gamma_U(S)=\gamma_U(f)=f$，故 $S=0$。设结论对所有阶 $\le p-1$（$p>1$）的 $T\in\mathcal{C}(V)$ 及所有开集 $V\subset X$ 成立。若 $S$ 的阶 $\le p$，则对任意 $g\in\mathcal{O}_X(V)$、$V\subset U$，有
$\gamma([S,g])=[\gamma(S),g]=0$，而 $[S,g]$ 的阶 $\le p-1$，由归纳假设 $[S,g]=0$。这意味着 $S$ 的阶 $\le 0$，由证明的第一部分得 $S=0$。

这表明 $\gamma:\mathcal{C}\to\mathcal{D}_X$ 是单射。还需证明 $\gamma$ 是满射。

先证明所有局部向量场都在 $\gamma$ 的像中。以 $\mathcal{T}_X$、$\mathcal{T}_Y$ 分别记 $X$、$Y$ 的切层。由 2.12，$\mathcal{O}_Y$-模 $\mathcal{T}_Y$ 是有限秩局部自由的，因此其逆像 $\varphi^*(\mathcal{T}_Y)$ 是有限秩局部自由 $\mathcal{O}_X$-模。于是 $\varphi^*(\mathcal{T}_Y)$ 在开集 $U\subset X$ 上的截面完全由它在所有 $x\in U$ 的几何纤维中的像决定：
$$T_x(\varphi^*(\mathcal{T}_Y))=\mathcal{O}_{X,x}/\mathfrak{m}_{X,x}\otimes_{\mathcal{O}_{Y,\varphi(x)}}\mathcal{T}_{Y,\varphi(x)}=T_{\varphi(x)}(Y).$$
设 $\Psi\in\operatorname{Hom}_{\mathcal{O}_X}(\mathcal{T}_X,\varphi^*(\mathcal{T}_Y))$。则对任意 $x\in X$，这个态射决定几何纤维之间的线性映射
$$T_x(\Psi):T_x(X)=T_x(\mathcal{T}_X)\longrightarrow T_x(\varphi^*(\mathcal{T}_Y))=T_{\varphi(x)}(Y).$$

\begin{lemma}\label{lem:4.1.4}
存在唯一的 $\mathcal{O}_X$-模态射 $\Phi:\mathcal{T}_X\to\varphi^*(\mathcal{T}_Y)$，使得对任意 $x\in X$，由它在几何纤维上诱导的线性映射 $T_x(\Phi)$ 等于切线性映射
$T_x(\varphi):T_x(X)\to T_{\varphi(x)}(Y)$。
\end{lemma}

\begin{proof}
唯一性由上注明显。设 $U$ 是 $X$ 中开集，$T$ 是 $U$ 上向量场。设 $\mathcal{U}=(U_i;1\le i\le n)$ 是 $U$ 的一个开覆盖，其元素的像包含在 $Y$ 的“小”开集内，即 $\varphi(U_i)\subset V_i$，且在 $V_i$ 上存在坐标系
$(f_1,\dots,f_m;D_1,\dots,D_m)$。此时对任意 $x\in U_i$ 的切向量 $\xi\in T_x(X)$，有
$T_x(\varphi)(\xi)=\sum_{j=1}^m a_jD_j(\varphi(x))$。此外，对任意 $1\le j\le m$，
$$\bigl(T_x(\varphi)(\xi)\bigr)(df_j(\varphi(x)))=\sum_{k=1}^m a_kD_k(\varphi(x))(df_j(\varphi(x)))
=\sum_{k=1}^m a_kD_k(f_j)(\varphi(x))=a_j.$$
因此
$$a_j=\xi(T_x^*(\varphi)(df_j(\varphi(x))))=\xi(d(f_j\circ\varphi)(x)).$$
这最终给出
$$T_x(\varphi)(\xi)=\sum_{j=1}^m\xi(d(f_j\circ\varphi)(x))D_j(\varphi(x)).$$
函数 $x\mapsto\bigl((T|_{U_i})(f_j\circ\varphi)\bigr)(x)$（$1\le j\le m$）在 $U_i$ 上正则。于是
$$\sum_{j=1}^m(T|_{U_i})(f_j\circ\varphi)\otimes D_j$$
是 $\varphi^*(\mathcal{T}_Y)$ 在 $U_i$ 上的一个截面。此外，由上述公式，对任意 $x\in U_i$，这个截面在 $\varphi^*(\mathcal{T}_Y)$ 于 $\varphi(x)$ 的几何纤维 $T_{\varphi(x)}(Y)$ 中的像等于
$$\sum_{j=1}^m T(f_j\circ\varphi)(x)D_j(\varphi(x))=\sum_{j=1}^m T(x)(d(f_j\circ\varphi)(x))D_j(\varphi(x))=T_x(\varphi)(T(x)).$$
因此在 $U_i$ 上构造了一个截面，其在 $x\in U_i$ 的几何纤维 $T_{\varphi(x)}(Y)$ 中的像为 $T_x(\varphi)(T(x))$。由于截面完全由其在几何纤维中的像决定，这个截面唯一。可将这些截面在 $\mathcal{U}$ 的所有元素上粘合，得到 $U$ 上的截面 $\Phi(T)$。显然映射 $\Phi:T\mapsto\Phi(T)$ 具有所需性质。
\end{proof}

现在注意到局部地 $\mathcal{T}_Y$ 是 $\mathcal{D}_Y$ 的 $\mathcal{O}_Y$-模直和因子，故自然态射
$\varphi^*(\mathcal{T}_Y)\to\mathcal{D}_{X\to Y}$ 是单射，可将 $\varphi^*(\mathcal{T}_Y)$ 与 $\mathcal{D}_{X\to Y}$ 的子模等同。因此可把 $\Phi$ 看作 $\mathcal{T}_X$ 到 $\mathcal{D}_{X\to Y}$ 的态射。

\begin{lemma}\label{lem:4.1.5}
设 $U\subset X$ 是开集，$T$ 是 $U$ 上向量场。则存在唯一元素 $\delta_U(T)\in\mathcal{C}(U)$ 使得
$$\delta_U(T)|_V(g\otimes 1)=T(g)\otimes 1+g\Phi(T)|_V$$
对任意 $g\in\mathcal{O}_X(V)$ 及任意开集 $V\subset U$ 成立。
\end{lemma}

\begin{proof}
假设存在 $\mathcal{C}(U)$ 的元素 $S$ 满足
$$S|_V(g\otimes 1)=T(g)\otimes 1+g\Phi(T)|_V$$
对任意 $g\in\mathcal{O}_X(V)$ 及任意开集 $V\subset U$ 成立。设 $x\in U$，则存在 $\varphi(x)$ 在 $Y$ 中的仿射开邻域 $W$ 及 $x$ 的仿射开邻域 $V\subset U$ 使 $\varphi(V)\subset W$。于是 $\mathcal{D}_{X\to Y}(V)=R(V)\otimes_{R(W)}\mathcal{D}_W$。对任意 $g\in R(V)$、$Q\in\mathcal{D}_W$，由于 $S|_V$ 是右 $\mathcal{D}_W$-模自同态，有 $S|_V(g\otimes Q)=S(g\otimes 1)Q$。由此可知这样的 $S|_V$ 是唯一的，故 $S$ 由上述性质唯一确定。

由唯一性，为证存在性，只需对 $U$ 换成其开覆盖的元素来证明。因此一开始就可设 $U$ 是仿射开集且 $\varphi(U)$ 包含在 $Y$ 中带坐标系 $(f_1,\dots,f_m;D_1,\dots,D_m)$ 的“小”仿射开集 $W$ 内。如上，此时 $\mathcal{D}_{X\to Y}(V)=R(V)\otimes_{R(W)}\mathcal{D}_W$。首先定义双线性映射 $R(V)\times\mathcal{D}_W\to R(V)\otimes_{R(W)}\mathcal{D}_W$ 为
$$(g,S)\longmapsto T(g)\otimes S+g\Phi(T)|_V S.$$
现取 $h\in R(W)$，其与 $\varphi|_V$ 的复合是 $V$ 上正则函数。我们断言
$$\Phi(T)|_V h-(h\circ\varphi)\Phi(T)|_V=T(h\circ\varphi)\otimes 1.$$
利用引理 \ref{lem:4.1.4} 证明中的计算，有
$$\Phi(T)|_V h-(h\circ\varphi)\Phi(T)|_V
=\sum_{j=1}^m\bigl(T(f_j\circ\varphi)\otimes D_jh-(h\circ\varphi)T(f_j\circ\varphi)\otimes D_j\bigr)$$
$$=\sum_{j=1}^m T(f_j\circ\varphi)\otimes[D_j,h]
=\sum_{j=1}^m T(f_j\circ\varphi)\otimes D_j(h)
=\sum_{j=1}^m T(f_j\circ\varphi)(D_j(h)\circ\varphi)\otimes 1,$$
即这个表达式是一个函数。它在 $x\in V$ 处的值为
$$\sum_{j=1}^m T(f_j\circ\varphi)(x)D_j(h)(\varphi(x))
=\sum_{j=1}^m T(x)(d(f_j\circ\varphi)(x))D_j(\varphi(x))(dh(\varphi(x)))$$
$$=\bigl(T_x(\varphi)(T(x))\bigr)(dh(\varphi(x)))=T(x)(d(h\circ\varphi)(x))=T(h\circ\varphi)(x),$$
这证明了上述关系。因此对任意 $h\in R(Y)$ 有
\begin{align*}
T(g(h\circ\varphi))\otimes S+g(h\circ\varphi)\Phi(T)S
&=T(g)\otimes hS+gT(h\circ\varphi)\otimes S+g(h\circ\varphi)\Phi(T)S\\
&=T(g)\otimes hS+g\Phi(T)hS,
\end{align*}
故该映射可分解通过 $\mathcal{D}_{X\to Y}(V)$，从而定义了 $\delta_V(T)\in\mathcal{E}nd_k(\mathcal{D}_{X\to Y}(V))$。若 $V_f$（$f\in R(V)$）是 $V$ 中主开集，则
$\mathcal{D}_{X\to Y}(V_f)=R(V_f)\otimes_{R(W)}\mathcal{D}_W=R(V)_f\otimes_{R(W)}\mathcal{D}_W$，因此由局部化，$\delta_V(T)$“延拓”为 $\mathcal{D}_{X\to Y}(V_f)$ 的自同态，它本质上由同一公式给出，因而等于 $\delta_{V_f}(T)$。这意味着 $\delta_V(T)$ 定义了 $\mathcal{E}nd_k(\mathcal{D}_{X\to Y})$ 的元素。显然对任意 $g,h\in R(V)$，
$$[\delta_V(T),g](h\otimes 1)=\delta_V(T)(gh\otimes 1)-g\delta_V(T)(h\otimes 1)
=T(gh)\otimes 1-gT(h)\otimes 1=T(g)(h\otimes 1),$$
故 $\delta_V(T)$ 是微分自同态。此外它还与 $\mathcal{D}_W$ 的右作用可交换，故 $\delta_V(T)\in\mathcal{C}(V)$。存在性得证。
\end{proof}

另一方面，由于 $\varphi^*(\mathcal{T}_Y)\subset\varphi^*(\mathcal{J}_Y)=\ker\alpha$，对任意 $g\in\mathcal{O}_X$ 有
$$\gamma(\delta(T))(g)=\alpha(\delta(T)(g\otimes 1))=\alpha(T(g)\otimes 1)=T(g),$$
即 $(\gamma\circ\delta)(T)=T$。这说明 $\gamma$ 的像包含 $X$ 上所有向量场，而由 3.7.(ii)，这些向量场与 $\mathcal{O}_X$ 生成环层 $\mathcal{D}_X$，故 $\gamma$ 是满射。定理 \ref{thm:4.1.2} 得证。

因此可把 $\mathcal{D}_{X\to Y}$ 视为双模层，带左 $\mathcal{D}_X$-作用与右 $\varphi^{-1}(\mathcal{D}_Y)$-作用。

现在讨论 $X=k^n$、$Y=k^m$ 的情形。此时在 I.12 中构造了 $\mathcal{D}_{X\to Y}$ 上由
$$\frac{\partial}{\partial x_i}(P\otimes S)=\frac{\partial P}{\partial x_i}\otimes S
+\sum_{j=1}^m P\frac{\partial(y_j\circ\varphi)}{\partial x_i}\otimes\frac{\partial}{\partial y_j}S$$
给出的左 $\mathcal{D}_X$-模结构，其中 $P\in R(X)$、$S\in\mathcal{D}_Y$。显然这个左 $\mathcal{D}_X$-作用与刚才一般定义的作用一致。因此双模层 $\mathcal{D}_{X\to Y}$ 是 I 中多项式映射情形构造的自然推广。

\begin{remark}\label{rem:4.1.6}
设 $U$ 是光滑代数簇 $X$ 的开子集，考虑自然包含 $i:U\to X$。则
$i^{-1}(\mathcal{D}_X)=\mathcal{D}_X|_U=\mathcal{D}_U$。因此 $\mathcal{D}_{U\to X}$ 作为右 $i^{-1}(\mathcal{D}_U)$-模同构于 $\mathcal{D}_U$。显然 $\mathcal{D}_{U\to X}$ 通过左乘具有左 $\mathcal{D}_U$-模结构，此作用与右 $i^{-1}(\mathcal{D}_X)$-作用可交换。于是得到一个自然态射 $\mathcal{D}_U\to\mathcal{C}$。由 $\gamma:\mathcal{C}\to\mathcal{D}_U$ 的定义知，对任意开集 $V\subset U$ 与向量场 $T\in\mathcal{D}_U(V)$，
$$\gamma_V(T)(f)=\alpha(T(f))=\alpha([T,f]+fT)=T(f).$$
故自然包含与 $\gamma$ 的复合在局部向量场上是同构。由于 $\mathcal{D}_U$ 局部由向量场生成，这说明该复合是 $\mathcal{D}_U$ 的同构。因此 $\mathcal{D}_U$ 在双模 $\mathcal{D}_{U\to X}$ 中的作用就是 $\mathcal{D}_U$ 的自然作用。
\end{remark}

设 $X$、$Y$、$Z$ 是光滑簇，$\varphi:X\to Y$、$\psi:Y\to Z$ 是簇的态射。则
$$\mathcal{D}_{X\to Z}=(\psi\circ\varphi)^*(\mathcal{D}_Z)=\varphi^*(\psi^*(\mathcal{D}_Z))
=\varphi^*(\mathcal{D}_{Y\to Z})$$
$$=\mathcal{O}_X\otimes_{\varphi^{-1}(\mathcal{O}_Y)}\varphi^{-1}(\mathcal{D}_{Y\to Z})
=\mathcal{O}_X\otimes_{\varphi^{-1}(\mathcal{O}_Y)}\bigl(\varphi^{-1}(\mathcal{D}_Y)
\otimes_{\varphi^{-1}(\mathcal{D}_Y)}\varphi^{-1}(\mathcal{D}_{Y\to Z})\bigr)$$
$$=\bigl(\mathcal{O}_X\otimes_{\varphi^{-1}(\mathcal{O}_Y)}\varphi^{-1}(\mathcal{D}_Y)\bigr)
\otimes_{\varphi^{-1}(\mathcal{D}_Y)}\varphi^{-1}(\mathcal{D}_{Y\to Z})
=\mathcal{D}_{X\to Y}\otimes_{\varphi^{-1}(\mathcal{D}_Y)}\varphi^{-1}(\mathcal{D}_{Y\to Z}).$$
这个同构显然与 $\mathcal{D}_{X\to Z}$ 和 $\mathcal{D}_{X\to Y}$ 中由左乘给出的 $\mathcal{O}_X$-模结构相容。对于由右乘给出的右
$(\psi\circ\varphi)^{-1}(\mathcal{D}_Z)=\varphi^{-1}(\psi^{-1}(\mathcal{D}_Z))$-模结构，以及右
$\varphi^{-1}(\mathcal{D}_{Y\to Z})$-模结构，同样成立。我们断言左 $\mathcal{D}_X$-模作用也相容。

设 $x\in X$，$T$ 是 $x$ 的仿射开邻域 $U$ 上的向量场。则
$$(\mathcal{D}_{X\to Y}\otimes_{\varphi^{-1}(\mathcal{D}_Y)}\varphi^{-1}(\mathcal{D}_{Y\to Z}))_x
=\mathcal{D}_{X\to Y,x}\otimes_{\mathcal{D}_{Y,\varphi(x)}}\mathcal{D}_{Y\to Z,\varphi(x)}.$$
设 $f$ 是 $x$ 处正则函数的芽。则在 $\mathcal{D}_{X\to Y}$ 中
$$T(f\otimes 1)=T(f)\otimes 1+\sum_i g_i\otimes T_i,$$
其中 $g_i\in\mathcal{O}_{X,x}$、$T_i\in\mathcal{J}_{Y,\varphi(x)}$。这意味着 $T$ 在茎
$\mathcal{D}_{X\to Y,x}\otimes_{\mathcal{D}_{Y,\varphi(x)}}\mathcal{D}_{Y\to Z,\varphi(x)}$ 的元素
$(f\otimes 1)\otimes(1\otimes 1)$ 上的作用为
$$T((f\otimes 1)\otimes(1\otimes 1))=T(f\otimes 1)\otimes(1\otimes 1)$$
$$=(T(f)\otimes 1)\otimes(1\otimes 1)+\sum_i(g_i\otimes T_i)\otimes(1\otimes 1)
=(T(f)\otimes 1)\otimes(1\otimes 1)+\sum_i(g_i\otimes 1)\otimes T_i(1\otimes 1).$$
前面已指出 $\psi^*(\mathcal{J}_Z)_{\varphi(x)}$ 是 $\mathcal{D}_{Y,\varphi(x)}$-不变的。此外，由作用的构造，$\varphi(x)$ 处任一局部向量场 $S$ 把 $1\otimes 1$ 映入 $\psi^*(\mathcal{J}_Z)_{\varphi(x)}$。这说明对所有 $i$ 有 $T_i(1\otimes 1)\in\psi^*(\mathcal{J}_Z)_{\varphi(x)}$。因此在上述同构下，$T((f\otimes 1)\otimes(1\otimes 1))$ 映为
$T(f)\otimes 1+\sum_j h_j\otimes S_j$，其中 $S_j\in\mathcal{J}_{Z,\psi(\varphi(x))}$。这说明 $T$ 的两个作用一致。因此建立了如下结果。

\begin{proposition}\label{prop:4.1.7}
设 $X$、$Y$、$Z$ 是光滑代数簇，$\varphi:X\to Y$、$\psi:Y\to Z$ 是代数簇的态射。则作为双模层有
$$\mathcal{D}_{X\to Z}=\mathcal{D}_{X\to Y}\otimes_{\varphi^{-1}(\mathcal{D}_Y)}\varphi^{-1}(\mathcal{D}_{Y\to Z}).$$
\end{proposition}
\begin{annote}
$f_+,f^+$ 与 $\boxtimes$、$\otimes^L$、$R\mathcal H om$、$R\Gamma_Z$ 合起来是 $\mathcal D$-模的六种运算。$f_+$ 有两套等价定义：经双边模 $\mathcal D_{Y\leftarrow X}=\mathcal D_Y\otimes_{f^{-1}\mathcal D_X}f^{-1}\mathcal D_X$ 的张量，或经相对 $\mathcal D_{X\to Y}$ 的转移；$f^+$ 是 $f_+$ 的右伴随、只能作导出函子。整套形式对应 nLab 的 six operations 公理，其中基变换为 Beck–Chevalley：$(p_2)_!\circ(p_1)^*\cong g^*\circ f_!$。[来源：nLab 条目 six+operations（Last revised 2026-09-30）；原书 §1–§2, p97–108]
\end{annote}



\section{仿射簇的逆像与直像}

设 $X$、$Y$ 是两个光滑仿射簇，$\varphi:X\to Y$ 是代数簇的态射。则
$D_{X\to Y}=\Gamma(X,\mathcal{D}_{X\to Y})$ 具有自然的（左 $D_X$、右 $D_Y$）-双模结构。定义以下两个右正合函子：

\begin{enumerate}[label=(\roman*)]
  \item 逆像函子 $\varphi^+:\mathcal{M}_L(D_Y)\to\mathcal{M}_L(D_X)$，对任意左 $D_Y$-模 $M$ 令
    $$\varphi^+(M)=D_{X\to Y}\otimes_{D_Y}M;$$
  \item 直像函子 $\varphi_+:\mathcal{M}_R(D_X)\to\mathcal{M}_R(D_Y)$，对任意右 $D_X$-模 $N$ 令
    $$\varphi_+(N)=N\otimes_{D_X}D_{X\to Y}.$$
\end{enumerate}
显然这些定义推广了 I.12 中仿射空间情形的定义。

设 $Z$ 是另一个光滑仿射簇，$\psi:Y\to Z$ 是簇的态射。则由前述结果，$D_{X\to Y}\otimes_{D_Y}D_{Y\to Z}$ 到 $D_{X\to Z}$ 有自然态射，且与左 $D_X$-模和右 $D_Z$-模作用相容。

\begin{proposition}\label{prop:4.2.1}
$D_{X\to Z}=D_{X\to Y}\otimes_{D_Y}D_{Y\to Z}$。
\end{proposition}

\begin{proof}
由前述结果，
$$D_{X\to Z}=\Gamma(X,\mathcal{D}_{X\to Z})=\Gamma(X,(\psi\circ\varphi)^*(\mathcal{D}_Z))
=\Gamma(X,\varphi^*(\psi^*(\mathcal{D}_Z)))$$
$$=\Gamma(X,\varphi^*(\mathcal{D}_{Y\to Z}))=R(X)\otimes_{R(Y)}\Gamma(Y,\mathcal{D}_{Y\to Z})
=R(X)\otimes_{R(Y)}D_{Y\to Z}$$
$$=R(X)\otimes_{R(Y)}(D_Y\otimes_{D_Y}D_{Y\to Z})
=(R(X)\otimes_{R(Y)}D_Y)\otimes_{D_Y}D_{Y\to Z}=D_{X\to Y}\otimes_{D_Y}D_{Y\to Z}.$$
\end{proof}

这立即蕴含：对任意左 $D_Z$-模 $M$，
$$(\varphi^+\circ\psi^+)(M)=\varphi^+(\psi^+(M))=D_{X\to Y}\otimes_{D_Y}(D_{Y\to Z}\otimes_{D_Z}M)$$
$$=(D_{X\to Y}\otimes_{D_Y}D_{Y\to Z})\otimes_{D_Z}M=D_{X\to Z}\otimes_{D_Z}M=(\psi\circ\varphi)^+(M);$$
对任意右 $D_X$-模 $N$，
$$(\psi_+\circ\varphi_+)(N)=\psi_+(\varphi_+(N))=(N\otimes_{D_X}D_{X\to Y})\otimes_{D_Y}D_{Y\to Z}$$
$$=N\otimes_{D_X}(D_{X\to Y}\otimes_{D_Y}D_{Y\to Z})=N\otimes_{D_X}D_{X\to Z}=(\psi\circ\varphi)_+(N).$$
因此建立了如下结果。

\begin{theorem}\label{thm:4.2.2}
设 $X$、$Y$、$Z$ 是光滑仿射簇，$\varphi:X\to Y$、$\psi:Y\to Z$ 是簇的态射。则
\begin{enumerate}[label=(\roman*)]
  \item $\varphi^+\circ\psi^+=(\psi\circ\varphi)^+$；
  \item $\psi_+\circ\varphi_+=(\psi\circ\varphi)_+$。
\end{enumerate}
\end{theorem}

现研究逆像 $\varphi^+$ 与直像 $\varphi_+$ 的左导出函子。先从以下结果开始。

\begin{lemma}\label{lem:4.2.3}
设 $V$ 是仿射簇 $X$ 上有限秩局部自由 $\mathcal{O}_X$-模。则 $\Gamma(X,V)$ 是投射 $R(X)$-模。
\end{lemma}

\begin{proof}
首先注意 $\Gamma(X,V)$ 是有限生成 $R(X)$-模。因此由 I.6.1 的类似结论，对任意 $R(X)$-模 $N$ 及任意 $x\in X$ 有
$$\operatorname{Ext}^1_{R(X)}(\Gamma(X,V),N)_x=\operatorname{Ext}^1_{\mathcal{O}_{X,x}}(V_x,N_x).$$
由假设，每个 $V_x$ 是自由 $\mathcal{O}_{X,x}$-模。故对任意 $x$（原文误写为 $x\in R(X)$）有
$\operatorname{Ext}^1_{R(X)}(\Gamma(X,V),N)_x=0$，这意味着
$\operatorname{Ext}^1_{R(X)}(\Gamma(X,V),N)=0$，即 $\Gamma(X,V)$ 是投射 $R(X)$-模。
\end{proof}

\begin{corollary}\label{cor:4.2.4}
设 $X$ 是光滑仿射簇。则 $D_X$ 关于左乘是投射 $R(X)$-模。
\end{corollary}

\begin{proof}
设 $p\in\mathbb{N}$。则有 $\mathcal{O}_X$-模的正合列
$$0\longrightarrow F_{p-1}D_X\longrightarrow F_pD_X\longrightarrow \operatorname{Gr}_pD_X\longrightarrow 0.$$
由于 $X$ 仿射，这给出 $R(X)$-模的正合列
$$0\longrightarrow\Gamma(X,F_{p-1}D_X)\longrightarrow\Gamma(X,F_pD_X)
\longrightarrow\Gamma(X,\operatorname{Gr}_pD_X)\longrightarrow 0.$$
由前述，$\operatorname{Gr}_pD_X$ 是有限秩局部自由 $\mathcal{O}_X$-模。因此由引理 \ref{lem:4.2.3}，$\Gamma(X,\operatorname{Gr}_pD_X)$ 是投射 $R(X)$-模。上述正合列分裂，故
$$\Gamma(X,F_pD_X)=\Gamma(X,F_{p-1}D_X)\oplus\Gamma(X,\operatorname{Gr}_pD_X).$$
由归纳得
$$\Gamma(X,F_pD_X)=\bigoplus_{i=0}^p\Gamma(X,\operatorname{Gr}_iD_X),$$
以及作为 $R(X)$-模
$$D_X=\Gamma(X,\mathcal{D}_X)=\bigoplus_{i=0}^{\infty}\Gamma(X,\operatorname{Gr}_iD_X).$$
因此 $D_X$ 是投射 $R(X)$-模的直和，从而是投射的。
\end{proof}

\begin{lemma}\label{lem:4.2.5}
设 $\varphi:X\to Y$ 是光滑仿射簇之间的态射。则
$D_{X\to Y}=R(X)\otimes_{R(Y)}D_Y$ 是投射 $R(X)$-模。
\end{lemma}

\begin{proof}
由推论 \ref{cor:4.2.4}，$D_Y$ 是投射 $R(Y)$-模，故它是某个自由 $R(Y)$-模的直和因子。于是
$D_{X\to Y}=R(X)\otimes_{R(Y)}D_Y$ 是自由 $R(X)$-模的直和因子，从而是投射 $R(X)$-模。
\end{proof}

\begin{proposition}\label{prop:4.2.6}
设 $X$、$Y$、$Z$ 是光滑仿射簇，$\varphi:X\to Y$、$\psi:Y\to Z$ 是仿射簇的态射。则对 $j\in\mathbb{N}$ 有
$$\operatorname{Tor}^{D_Y}_j(D_{X\to Y},D_{Y\to Z})=0.$$
\end{proposition}

\begin{proof}
设 $M$ 是左 $D_Y$-模，$F^\cdot$ 是它由自由 $D_Y$-模构成的左分解。由推论 \ref{cor:4.2.4}，$D_Y$ 关于左乘是投射 $R(Y)$-模，故也可把 $F^\cdot$ 看作由投射 $R(Y)$-模构成的分解。这意味着
\begin{align*}
\operatorname{Tor}^{R(Y)}_j(R(X),M)&=H^{-j}(R(X)\otimes_{R(Y)}F^\cdot)\\
&=H^{-j}\bigl((R(X)\otimes_{R(Y)}D_Y)\otimes_{D_Y}F^\cdot\bigr)
=\operatorname{Tor}^{D_Y}_j(D_{X\to Y},M).
\end{align*}
由于由引理 \ref{lem:4.2.5}，$D_{Y\to Z}$ 是投射 $R(Y)$-模，故对 $j\in\mathbb{N}$ 有
$\operatorname{Tor}^{R(Y)}_j(R(X),D_{Y\to Z})=0$，结论随之成立。
\end{proof}

下一个结果把 I.12.2 与 I.12.6 推广到光滑仿射簇。

\begin{theorem}\label{thm:4.2.7}
设 $X$、$Y$、$Z$ 是三个光滑仿射簇，$\varphi:X\to Y$、$\psi:Y\to Z$ 是簇的态射。则
\begin{enumerate}[label=(\roman*)]
  \item 对任意左 $D_Z$-模 $M$，存在谱序列，其 $E_2$ 项为
    $E_2^{pq}=L_p\varphi^+(L_q\psi^+(M))$，收敛到 $L_{p+q}(\psi\circ\varphi)^+(M)$；
  \item 对任意右 $D_X$-模 $N$，存在谱序列，其 $E_2$ 项为
    $E_2^{pq}=L_p\psi_+(L_q\varphi_+(N))$，收敛到 $L_{p+q}(\psi\circ\varphi)_+(N)$。
\end{enumerate}
\end{theorem}

\begin{proof}
两个结论都可由 Grothendieck 谱序列得出。
\begin{enumerate}[label=(\roman*)]
  \item 设 $P$ 是投射左 $D_Z$-模，则它是某个自由 $D_Z$-模的直和因子。因此 $\psi^+(P)$ 是某个
    $D_{Y\to Z}^{(I)}$ 的直和因子。由命题 \ref{prop:4.2.6}，这意味着 $\psi^+(P)$ 是 $\varphi^+$-无循环的。
  \item 设 $Q$ 是投射右 $D_X$-模，则它是某个自由 $D_X$-模的直和因子。因此 $\varphi_+(Q)$ 是某个
    $D_{X\to Y}^{(J)}$ 的直和因子。再次应用命题 \ref{prop:4.2.6}，得 $\varphi_+(Q)$ 是 $\psi_+$-无循环的。
\end{enumerate}
\end{proof}

\section{逆像函子}

设 $X$、$Y$ 是两个光滑代数簇，$\varphi:X\to Y$ 是代数簇的态射。定义函子
$$\varphi^+(\mathcal{V})=\mathcal{D}_{X\to Y}\otimes_{\varphi^{-1}(\mathcal{D}_Y)}\varphi^{-1}(\mathcal{V})$$
从左 $\mathcal{D}_Y$-模范畴 $\mathcal{M}_L(\mathcal{D}_Y)$ 到左 $\mathcal{D}_X$-模范畴 $\mathcal{M}_L(\mathcal{D}_X)$。这个函子称为\emph{逆像函子}。由于 $\varphi^{-1}$ 正合而 $\mathcal{D}_{X\to Y}\otimes_{\varphi^{-1}(\mathcal{D}_Y)}-$ 右正合，逆像函子是右正合的。

\begin{remark}\label{rem:4.3.1}
设 $U$ 是 $X$ 的开子集。由注记 \ref{rem:4.1.6} 知 $\mathcal{D}_{U\to X}=\mathcal{D}_U$。故立刻得到
$i^+(\mathcal{V})=i^{-1}(\mathcal{V})=\mathcal{V}|_U$，即 $i$ 的逆像就是限制到 $U$ 的函子
$\mathcal{V}\mapsto\mathcal{V}|_U$。特别地，它是正合函子。
\end{remark}

\begin{lemma}\label{lem:4.3.2}
平坦左 $\mathcal{D}_Y$-模的类对 $\varphi^+$ 是左适合的（left adapted）。
\end{lemma}

\begin{proof}
平坦左 $\mathcal{D}_Y$-模的类对直和封闭，且每个左 $\mathcal{D}_Y$-模都是某个平坦模的商。

设 $F$ 是平坦左 $\mathcal{D}_Y$-模。则对任意 $y\in Y$，$F_y$ 是平坦左 $\mathcal{D}_{Y,y}$-模。因此对任意 $x\in X$，
$\varphi^{-1}(F)_x=F_{\varphi(x)}$ 在 $\varphi^{-1}(\mathcal{D}_Y)_x=\mathcal{D}_{Y,\varphi(x)}$ 上平坦。故
$\varphi^{-1}(F)$ 是平坦左 $\varphi^{-1}(\mathcal{D}_Y)$-模。

因此，若 $F^\cdot$ 是上有界的平坦左 $\mathcal{D}_Y$-模的无循环复形，则 $\varphi^{-1}(F^\cdot)$ 是上有界的平坦左
$\varphi^{-1}(\mathcal{D}_Y)$-模的无循环复形。于是 $\varphi^+(F^\cdot)$ 是上有界无循环复形。

这表明所有平坦左 $\mathcal{D}_Y$-模的类对 $\varphi^+$ 是左适合的。
\end{proof}

因此可定义左导出函子
$$L\varphi^+:D^-(\mathcal{M}_L(\mathcal{D}_Y))\longrightarrow D^-(\mathcal{M}_L(\mathcal{D}_X)).$$

\begin{lemma}\label{lem:4.3.3}
函子 $\varphi^+$ 的左同调维数有限。
\end{lemma}

\begin{proof}
设 $\mathcal{V}$ 是左 $\mathcal{D}_Y$-模，$F^\cdot$ 是它由平坦左 $\mathcal{D}_Y$-模构成的左分解。则由引理 \ref{lem:4.3.2}，
$$L_p\varphi^+(\mathcal{V})=H^p(\varphi^+(F^\cdot))
=H^p\bigl(\mathcal{D}_{X\to Y}\otimes_{\varphi^{-1}(\mathcal{D}_Y)}\varphi^{-1}(F^\cdot)\bigr).$$
设 $x\in X$，则 $F_{\varphi(x)}^\cdot$ 是 $\mathcal{V}_{\varphi(x)}$ 的 $\mathcal{D}_{Y,\varphi(x)}$-平坦分解。故
$$(L_p\varphi^+(\mathcal{V}))_x
=H^p\bigl(\mathcal{D}_{X\to Y}\otimes_{\varphi^{-1}(\mathcal{D}_Y)}\varphi^{-1}(F^\cdot)\bigr)_x
=H^p\bigl(\mathcal{D}_{X\to Y,x}\otimes_{\varphi^{-1}(\mathcal{D}_Y)_x}\varphi^{-1}(F^\cdot)_x\bigr)$$
$$=H^p\bigl(\mathcal{D}_{X\to Y,x}\otimes_{\mathcal{D}_{Y,\varphi(x)}}F_{\varphi(x)}^\cdot\bigr)
=\operatorname{Tor}^{-\,p}_{\mathcal{D}_{Y,\varphi(x)}}(\mathcal{D}_{X\to Y,x},\mathcal{V}_{\varphi(x)}).$$
由于环 $\mathcal{D}_{Y,\varphi(x)}$ 的同调维数 $\le 2\dim Y$，可知当 $p<-2\dim Y$ 时
$(L_p\varphi^+(\mathcal{V}))_x=0$。又因 $x\in X$ 任意，故当 $p<-2\dim Y$ 时 $L_p\varphi^+=0$。
\end{proof}

因此，左导出函子 $L\varphi^+$ 可延拓为无界复形导出范畴之间的导出函子
$$L\varphi^+:D(\mathcal{M}_L(\mathcal{D}_Y))\longrightarrow D(\mathcal{M}_L(\mathcal{D}_X)).$$

\begin{lemma}\label{lem:4.3.4}
设 $\mathcal{V}^\cdot$ 是左 $\mathcal{D}_Y$-模的复形。则
$$L\varphi^+(\mathcal{V}^\cdot)=\mathcal{D}_{X\to Y}\overset{L}{\otimes}_{\varphi^{-1}(\mathcal{D}_Y)}\varphi^{-1}(\mathcal{V}^\cdot).$$
\end{lemma}

\begin{proof}
设 $F$ 是平坦左 $\mathcal{D}_Y$-模。如引理 \ref{lem:4.3.2} 证明中所述，$\varphi^{-1}(F)$ 是平坦左
$\varphi^{-1}(\mathcal{D}_Y)$-模，因而对函子 $\mathcal{D}_{X\to Y}\otimes_{\varphi^{-1}(\mathcal{D}_Y)}-$ 是无循环的。
\end{proof}

设 $\mathcal{V}$ 是左 $\mathcal{D}_Y$-模。则
$$\varphi^+(\mathcal{V})=\mathcal{D}_{X\to Y}\otimes_{\varphi^{-1}(\mathcal{D}_Y)}\varphi^{-1}(\mathcal{V})
=\bigl(\mathcal{O}_X\otimes_{\varphi^{-1}(\mathcal{O}_Y)}\varphi^{-1}(\mathcal{D}_Y)\bigr)
\otimes_{\varphi^{-1}(\mathcal{D}_Y)}\varphi^{-1}(\mathcal{V})$$
$$=\mathcal{O}_X\otimes_{\varphi^{-1}(\mathcal{O}_Y)}
\bigl(\varphi^{-1}(\mathcal{D}_Y)\otimes_{\varphi^{-1}(\mathcal{D}_Y)}\varphi^{-1}(\mathcal{V})\bigr)
=\mathcal{O}_X\otimes_{\varphi^{-1}(\mathcal{O}_Y)}\varphi^{-1}(\mathcal{V})=\varphi^*(\mathcal{V}),$$
即若忘掉 $\mathcal{D}_X$-模结构，则 D-模逆像 $\varphi^+(\mathcal{V})$ 等于 O-模逆像 $\varphi^*(\mathcal{V})$。现证明这对导出函子仍成立。先需要以下结果。

\begin{lemma}\label{lem:4.3.5}
设 $F$ 是平坦左 $\mathcal{D}_Y$-模。则 $F$ 是平坦 $\mathcal{O}_Y$-模。
\end{lemma}

\begin{proof}
设 $y\in Y$，只需证 $F_y$ 是平坦 $\mathcal{O}_{Y,y}$-模。由假设，$F_y$ 是平坦左 $\mathcal{D}_{Y,y}$-模。设 $W$ 是右
$\mathcal{O}_{Y,y}$-模。则
$$W\otimes_{\mathcal{O}_{Y,y}}F_y=W\otimes_{\mathcal{O}_{Y,y}}(\mathcal{D}_{Y,y}\otimes_{\mathcal{D}_{Y,y}}F_y)
=(W\otimes_{\mathcal{O}_{Y,y}}\mathcal{D}_{Y,y})\otimes_{\mathcal{D}_{Y,y}}F_y.$$
由于 $\mathcal{D}_{Y,y}$ 关于左乘是自由 $\mathcal{O}_{Y,y}$-模，故 $W\mapsto W\otimes_{\mathcal{O}_{Y,y}}\mathcal{D}_{Y,y}$ 是正合函子。于是
$W\mapsto(W\otimes_{\mathcal{O}_{Y,y}}\mathcal{D}_{Y,y})\otimes_{\mathcal{D}_{Y,y}}F_y$ 也是正合的。这立即使
$W\mapsto W\otimes_{\mathcal{O}_{Y,y}}F_y$ 正合，即 $F_y$ 是 $\mathcal{O}_{Y,y}$-平坦的。
\end{proof}

这导致以下结果。

\begin{theorem}\label{thm:4.3.6}
若竖直箭头表示遗忘函子，则如下函子图
$$\begin{tikzcd}
D(\mathcal{M}_L(\mathcal{D}_Y)) \arrow[r,"L\varphi^+"] \arrow[d,"\operatorname{For}"'] &
D(\mathcal{M}_L(\mathcal{D}_X)) \arrow[d,"\operatorname{For}"]\\
D(\mathcal{M}(\mathcal{O}_Y)) \arrow[r,"L\varphi^*"'] & D(\mathcal{M}(\mathcal{O}_X))
\end{tikzcd}$$
在同构意义下交换。
\end{theorem}

\begin{proof}
这可由引理 \ref{lem:4.3.4} 以及关于函子复合的导出函子定理，分别应用于遗忘函子与 $\varphi^+$ 的复合、$\varphi^*$ 与遗忘函子的复合而得出。
\end{proof}

这个结果结合前述结论，有以下直接推论。

\begin{theorem}\label{thm:4.3.7}
设 $\mathcal{V}^\cdot$ 是左 $\mathcal{D}_Y$-模的复形，且对所有 $p\in\mathbb{Z}$，$H^p(\mathcal{V}^\cdot)$ 是拟凝聚的。则对所有 $p\in\mathbb{Z}$，上同调模 $H^p(L\varphi^+(\mathcal{V}^\cdot))$ 是拟凝聚左 $\mathcal{D}_X$-模。
\end{theorem}

现设 $X$、$Y$ 是光滑仿射簇。一般地，对任意左 $\mathcal{D}_Y$-模 $\mathcal{V}$，有自然态射
$$\mathcal{D}_{X\to Y}\otimes_{\mathcal{D}_Y}\Gamma(Y,\mathcal{V})\longrightarrow\Gamma(X,\varphi^+(\mathcal{V})).$$
下一个结果蕴含：对拟凝聚的 $\mathcal{V}$，这个态射是同构。

\begin{proposition}\label{prop:4.3.8}
设 $X$、$Y$ 是光滑仿射簇，$\varphi:X\to Y$ 是代数簇的态射。则对任意拟凝聚左 $\mathcal{D}_Y$-模 $\mathcal{V}$，有
$$\Gamma(X,L_p\varphi^+(\mathcal{V}))=\operatorname{Tor}^{-\,p}_{\mathcal{D}_Y}(D_{X\to Y},\Gamma(Y,\mathcal{V}))$$
对 $p\in\mathbb{Z}$ 成立。
\end{proposition}

\begin{proof}
首先注意若 $\mathcal{V}=\mathcal{D}_Y$，则 $\varphi^+(\mathcal{D}_Y)=\mathcal{D}_{X\to Y}$，上述态射显然是同构。现考虑任意拟凝聚左 $\mathcal{D}_Y$-模 $\mathcal{V}$。设 $F^\cdot$ 是 $\Gamma(Y,\mathcal{V})$ 的自由左分解。则它的局部化 $\mathcal{F}^\cdot=\Delta(F^\cdot)$ 是 $\mathcal{V}$ 的自由分解，即 $\mathcal{V}$ 的 $\mathcal{D}_Y$-平坦分解。因此
$$L_p\varphi^+(\mathcal{V})=H^p(L\varphi^+(\mathcal{D}(\mathcal{V})))=H^p(\varphi^+(\mathcal{F}^\cdot))$$
对 $p\in\mathbb{Z}$ 成立。由于 $\varphi^+(\mathcal{F}^\cdot)$ 是拟凝聚 $\mathcal{D}_X$-模的复形，
$$\Gamma(X,L_p\varphi^+(\mathcal{V}))=\Gamma(X,H^p(\varphi^+(\mathcal{F}^\cdot)))=H^p(\Gamma(X,\varphi^+(\mathcal{F}^\cdot))).$$
最后，由于 $\Gamma(X,-)$ 与 $\varphi^+$ 都与直和可交换，由证明的第一部分得
$$\Gamma(X,L_p\varphi^+(\mathcal{V}))=H^p(D_{X\to Y}\otimes_{\mathcal{D}_Y}F^\cdot)
=\operatorname{Tor}^{-\,p}_{\mathcal{D}_Y}(D_{X\to Y},\Gamma(Y,\mathcal{V}))$$
对 $p\in\mathbb{Z}$ 成立。
\end{proof}

因此，在光滑仿射簇的态射与拟凝聚 D-模的情形下，我们恢复了 §2 中的旧定义。

设 $X$、$Y$、$Z$ 是三个光滑代数簇，$\varphi:X\to Y$、$\psi:Y\to Z$ 是两个簇的态射。则对任意左 $\mathcal{D}_Z$-模 $\mathcal{V}$，有
$$(\varphi^+\circ\psi^+)(\mathcal{V})=\mathcal{D}_{X\to Y}\otimes_{\varphi^{-1}(\mathcal{D}_Y)}\varphi^{-1}(\psi^+(\mathcal{V}))$$
$$=\mathcal{D}_{X\to Y}\otimes_{\varphi^{-1}(\mathcal{D}_Y)}\varphi^{-1}
(\mathcal{D}_{Y\to Z}\otimes_{\psi^{-1}(\mathcal{D}_Z)}\psi^{-1}(\mathcal{V}))$$
$$=\mathcal{D}_{X\to Y}\otimes_{\varphi^{-1}(\mathcal{D}_Y)}
\bigl(\varphi^{-1}(\mathcal{D}_{Y\to Z})\otimes_{\varphi^{-1}(\psi^{-1}(\mathcal{D}_Z))}\varphi^{-1}(\psi^{-1}(\mathcal{V}))\bigr)$$
$$=\bigl(\mathcal{D}_{X\to Y}\otimes_{\varphi^{-1}(\mathcal{D}_Y)}\varphi^{-1}(\mathcal{D}_{Y\to Z})\bigr)
\otimes_{(\psi\circ\varphi)^{-1}(\mathcal{D}_Z)}(\psi\circ\varphi)^{-1}(\mathcal{V}).$$
故由前述，
$$(\varphi^+\circ\psi^+)(\mathcal{V})=\mathcal{D}_{X\to Z}
\otimes_{(\psi\circ\varphi)^{-1}(\mathcal{D}_Z)}(\psi\circ\varphi)^{-1}(\mathcal{V})=(\psi\circ\varphi)^+(\mathcal{V}).$$
因此我们证明了以下结果。

\begin{lemma}\label{lem:4.3.9}
$\varphi^+\circ\psi^+=(\psi\circ\varphi)^+$。
\end{lemma}

下一个结果把它推广到导出范畴。

\begin{theorem}\label{thm:4.3.10}
正合函子 $L\varphi^+\circ L\psi^+$ 与 $L(\psi\circ\varphi)^+$（从 $D(\mathcal{M}_L(\mathcal{D}_Z))$ 到
$D(\mathcal{M}_L(\mathcal{D}_X))$）是同构的。
\end{theorem}

\begin{proof}
由上述讨论及注记 \ref{rem:4.3.1}，只需验证对任意平坦左 $\mathcal{D}_Z$-模 $F$，$\psi^+(F)$ 是 $\varphi^+$-无循环的。由引理 \ref{lem:4.3.5}，只需证 $\psi^*(F)$ 是 $\varphi^*$-无循环的。这由引理 \ref{lem:4.3.4} 及前述讨论得出。
\end{proof}

显然有
$$\varphi^+(\mathcal{D}_Y)=\mathcal{D}_{X\to Y}\otimes_{\varphi^{-1}(\mathcal{D}_Y)}\varphi^{-1}(\mathcal{D}_Y)=\mathcal{D}_{X\to Y}.$$
由于 $\mathcal{D}_Y$ 是平坦左 $\mathcal{D}_Y$-模，由注记 \ref{rem:4.3.1} 得
$$L\varphi^+(\mathcal{D}(\mathcal{D}_Y))=\mathcal{D}(\mathcal{D}_{X\to Y}).$$
故由引理 \ref{lem:4.3.9} 得
$$\mathcal{D}(\mathcal{D}_{X\to Z})=L(\psi\circ\varphi)^+(\mathcal{D}(\mathcal{D}_Z))
=L\varphi^+(L\psi^+(\mathcal{D}(\mathcal{D}_Z)))$$
$$=L\varphi^+(\mathcal{D}(\mathcal{D}_{Y\to Z}))
=\mathcal{D}_{X\to Y}\overset{L}{\otimes}_{\varphi^{-1}(\mathcal{D}_Y)}\mathcal{D}(\varphi^{-1}(\mathcal{D}_{Y\to Z})).$$

\begin{corollary}\label{cor:4.3.11}
有
$$\mathcal{D}(\mathcal{D}_{X\to Z})=\mathcal{D}(\mathcal{D}_{X\to Y})
\overset{L}{\otimes}_{\varphi^{-1}(\mathcal{D}_Y)}\mathcal{D}(\varphi^{-1}(\mathcal{D}_{Y\to Z})).$$
\end{corollary}

\section{投影公式}

设 $X$、$Y$ 是两个拓扑空间，$\varphi:X\to Y$ 是连续映射。设 $R$ 是 $Y$ 上的环层，则 $\varphi^{-1}(R)$ 是 $X$ 上的环层。

设 $A$ 是 $X$ 上右 $\varphi^{-1}(R)$-模，$B$ 是 $Y$ 上左 $R$-模。则可考虑 $Y$ 上的阿贝尔群层
$\varphi_\bullet(A\otimes_{\varphi^{-1}(R)}\varphi^{-1}(B))$ 与 $\varphi_\bullet(A)\otimes_R B$。第一个由
$$V\longmapsto\varphi_\bullet(A\otimes_{\varphi^{-1}(R)}\varphi^{-1}(B))(V)
=(A\otimes_{\varphi^{-1}(R)}\varphi^{-1}(B))(\varphi^{-1}(V))$$
给出，第二个是预层
$$V\longmapsto\varphi_\bullet(A)(V)\otimes_{R(V)}B(V)=A(\varphi^{-1}(V))\otimes_{R(V)}B(V)$$
相伴的层。

由于 $A\otimes_{\varphi^{-1}(R)}\varphi^{-1}(B)$ 是 $X$ 上预层
$$U\longmapsto A(U)\otimes_{\varphi^{-1}(R)(U)}\varphi^{-1}(B)(U)$$
相伴的层，对任意开集 $V\subset Y$，有从
$$A(\varphi^{-1}(V))\otimes_{R(V)}B(V)$$
到
$$A(\varphi^{-1}(V))\otimes_{\varphi^{-1}(R)(\varphi^{-1}(V))}\varphi^{-1}(B)(\varphi^{-1}(V))$$
再到群
$$(A\otimes_{\varphi^{-1}(R)}\varphi^{-1}(B))(\varphi^{-1}(V))$$
的自然态射。因此有自然态射
$$\varphi_\bullet(A)\otimes_R B\longrightarrow\varphi_\bullet(A\otimes_{\varphi^{-1}(R)}\varphi^{-1}(B)).$$
这显然是双函子的态射。因此它诱导有上界复形同伦范畴之间相应双函子的态射：若 $A^\cdot$ 是上有界右 $\varphi^{-1}(R)$-模的复形，$B^\cdot$ 是上有界左 $R$-模的复形，则有典范态射
$$\varphi_\bullet(A^\cdot)\otimes_R B^\cdot\longrightarrow
\varphi_\bullet(A^\cdot\otimes_{\varphi^{-1}(R)}\varphi^{-1}(B^\cdot)).$$
现设 $\varphi_\bullet$ 具有有限右同调维数。则有典范态射
$$\varphi_\bullet(A^\cdot\otimes_{\varphi^{-1}(R)}\varphi^{-1}(B^\cdot))
\longrightarrow R\varphi_\bullet(A^\cdot\otimes_{\varphi^{-1}(R)}\varphi^{-1}(B^\cdot)).$$
设 $A^\cdot$ 是 $\varphi_\bullet$-无循环的且 $B^\cdot$ 是 $R$-平坦的。则由于 $\varphi^{-1}(B^\cdot)$ 是 $\varphi^{-1}(R)$-平坦的，有典范态射
$$\varphi_\bullet(A^\cdot)\otimes_R B^\cdot\longrightarrow
R\varphi_\bullet(A^\cdot\otimes_{\varphi^{-1}(R)}\varphi^{-1}(B^\cdot))
=R\varphi_\bullet(A^\cdot\overset{L}{\otimes}_{\varphi^{-1}(R)}\varphi^{-1}(B^\cdot)).$$
最后，由于任意上有界右 $\varphi^{-1}(R)$-模复形拟同构于某个上有界 $\varphi_\bullet$-无循环复形，任意上有界左 $R$-模复形拟同构于某个上有界 $R$-平坦复形，得到自然态射
$$R\varphi_\bullet(A^\cdot)\overset{L}{\otimes}_R B^\cdot\longrightarrow
R\varphi_\bullet(A^\cdot\overset{L}{\otimes}_{\varphi^{-1}(R)}\varphi^{-1}(B^\cdot)).$$
因此在导出范畴之间得到了相应双函子的态射。

在某些条件下，这个自然态射是同构。相应的陈述通常称为\emph{投影公式}。例如，设 $X$、$Y$ 是代数簇，$\varphi:X\to Y$ 是代数簇的态射。此时 $\varphi_\bullet$ 具有有限右同调维数。

\begin{proposition}\label{prop:4.4.1}
设 $X$、$Y$ 是两个代数簇，$\varphi:X\to Y$ 是代数簇的态射。则
$$R\varphi_\bullet(\mathcal{V}^\cdot)\overset{L}{\otimes}_{\mathcal{O}_Y}\mathcal{W}^\cdot
=R\varphi_\bullet(\mathcal{V}^\cdot\overset{L}{\otimes}_{\varphi^{-1}(\mathcal{O}_Y)}\varphi^{-1}(\mathcal{W}^\cdot))$$
对任意 $D^-(\mathcal{M}(\varphi^{-1}(\mathcal{O}_Y)))$ 中的 $\mathcal{V}^\cdot$ 与
$D^-(\mathcal{M}_{qc}(\mathcal{O}_Y))$ 中的 $\mathcal{W}^\cdot$ 成立。
\end{proposition}

\begin{proof}
只需证明典范态射在上同调群上诱导同构。由于 $\varphi_\bullet$ 的右同调维数有限，通过截断点可设 $\mathcal{W}^\cdot$ 是有界复形。此外，由于所有拟凝聚 $\mathcal{O}_Y$-模构成有界导出范畴的一个生成类，可设 $\mathcal{W}^\cdot=\mathcal{D}(\mathcal{W})$，其中 $\mathcal{W}$ 是拟凝聚 $\mathcal{O}_Y$-模。

由于该陈述关于 $Y$ 是局部的，还可设 $Y$ 是仿射簇。此时可用自由分解 $F^\cdot$ 代替 $\mathcal{W}$。由于 $\varphi_\bullet$ 与直和可交换，显然此时自然态射是同构。
\end{proof}

类似地可证明以下陈述。

\begin{proposition}\label{prop:4.4.2}
设 $X$、$Y$ 是两个代数簇，$\varphi:X\to Y$ 是代数簇的态射。则
$$R\varphi_\bullet(\mathcal{V}^\cdot)\overset{L}{\otimes}_{\mathcal{D}_Y}\mathcal{W}^\cdot
=R\varphi_\bullet(\mathcal{V}^\cdot\overset{L}{\otimes}_{\varphi^{-1}(\mathcal{D}_Y)}\varphi^{-1}(\mathcal{W}^\cdot))$$
对任意 $D^-(\mathcal{M}_R(\varphi^{-1}(\mathcal{D}_Y)))$ 中的 $\mathcal{V}^\cdot$ 与
$D^-(\mathcal{M}^L_{qc}(\mathcal{D}_Y))$ 中的 $\mathcal{W}^\cdot$ 成立。
\end{proposition}

设 $Y$ 是光滑簇。则 $\mathcal{O}_{Y,y}$ 的同调维数 $\le\dim Y$，$\mathcal{D}_{Y,y}$ 的同调维数 $\le 2\dim Y$。因此用标准的截断点论证，可建立前两个结果的下述变体。

\begin{proposition}\label{prop:4.4.3}
设 $X$、$Y$ 是两个光滑代数簇，$\varphi:X\to Y$ 是代数簇的态射。则
$$R\varphi_\bullet(\mathcal{V}^\cdot)\overset{L}{\otimes}_{\mathcal{O}_Y}\mathcal{W}^\cdot
=R\varphi_\bullet(\mathcal{V}^\cdot\overset{L}{\otimes}_{\varphi^{-1}(\mathcal{O}_Y)}\varphi^{-1}(\mathcal{W}^\cdot))$$
对任意 $D(\mathcal{M}(\varphi^{-1}(\mathcal{O}_Y)))$ 中的 $\mathcal{V}^\cdot$ 与
$D^b(\mathcal{M}_{qc}(\mathcal{O}_Y))$ 中的 $\mathcal{W}^\cdot$ 成立。
\end{proposition}

\begin{proposition}\label{prop:4.4.4}
设 $X$、$Y$ 是两个光滑代数簇，$\varphi:X\to Y$ 是代数簇的态射。则
$$R\varphi_\bullet(\mathcal{V}^\cdot)\overset{L}{\otimes}_{\mathcal{D}_Y}\mathcal{W}^\cdot
=R\varphi_\bullet(\mathcal{V}^\cdot\overset{L}{\otimes}_{\varphi^{-1}(\mathcal{D}_Y)}\varphi^{-1}(\mathcal{W}^\cdot))$$
对任意 $D(\mathcal{M}_R(\varphi^{-1}(\mathcal{D}_Y)))$ 中的 $\mathcal{V}^\cdot$ 与
$D^b(\mathcal{M}^L_{qc}(\mathcal{D}_Y))$ 中的 $\mathcal{W}^\cdot$ 成立。
\end{proposition}

特别地，我们需要以下形式的投影公式。设 $\varphi:X\to Y$ 与 $\psi:Y\to Z$ 是光滑代数簇的态射。

\begin{lemma}\label{lem:4.4.5}
设 $\mathcal{V}^\cdot$ 是右 $\varphi^{-1}(\mathcal{D}_Y)$-模的复形。则在
$D(\mathcal{M}_R(\psi^{-1}(\mathcal{D}_Z)))$ 中有右 $\psi^{-1}(\mathcal{D}_Z)$-模复形的自然同构
$$R\varphi_\bullet(\mathcal{V}^\cdot)\overset{L}{\otimes}_{\mathcal{D}_Y}\mathcal{D}_{Y\to Z}
=R\varphi_\bullet(\mathcal{V}^\cdot\overset{L}{\otimes}_{\varphi^{-1}(\mathcal{D}_Y)}
\varphi^{-1}(\mathcal{D}_{Y\to Z})).$$
\end{lemma}

\begin{proof}
可把双模 $\mathcal{D}_{Y\to Z}$ 看作环层 $\mathcal{D}_Y\otimes_k\psi^{-1}(\mathcal{D}_Z^{opp})$ 上的模层，其中
$\mathcal{D}_Z^{opp}$ 是 $\mathcal{D}_Z$ 的对立环层。显然，平坦的
$\mathcal{D}_Y\otimes_k\psi^{-1}(\mathcal{D}_Z^{opp})$-模作为 $\mathcal{D}_Y$-模是平坦的。因此从
$D(\mathcal{M}_R(\varphi^{-1}(\mathcal{D}_Y)))$ 到 $D(\mathcal{M}_R(\psi^{-1}(\mathcal{D}_Z)))$ 的函子典范态射
$$R\varphi_\bullet(\mathcal{V}^\cdot)\overset{L}{\otimes}_{\mathcal{D}_Y}\mathcal{D}_{Y\to Z}
\longrightarrow R\varphi_\bullet(\mathcal{V}^\cdot\overset{L}{\otimes}_{\varphi^{-1}(\mathcal{D}_Y)}
\varphi^{-1}(\mathcal{D}_{Y\to Z}))$$
若忘掉 $\psi^{-1}(\mathcal{D}_Z^{opp})$-模作用，即诱导出命题 \ref{prop:4.4.4} 的典范态射。另一方面，由命题 \ref{prop:4.4.4}，这个态射是同构。
\end{proof}

\section{直像函子}

设 $X$、$Y$ 是两个光滑代数簇，$\varphi:X\to Y$ 是代数簇的态射。由于对 $x\in X$，$\mathcal{D}_{X,x}$ 的同调维数 $\le 2\dim X$，函子 $U\mapsto U\otimes_{\mathcal{D}_X}\mathcal{D}_{X\to Y}$ 具有有限左同调维数。因此有函子
$\mathcal{V}^\cdot\mapsto\mathcal{V}^\cdot\overset{L}{\otimes}_{\mathcal{D}_X}\mathcal{D}_{X\to Y}$，介于
$D(\mathcal{M}_R(\mathcal{D}_X))$ 与 $D(\mathcal{M}_R(\varphi^{-1}(\mathcal{D}_Y)))$ 之间。另一方面，$\varphi_\bullet$ 也具有有限同调维数，故有函子
$\mathcal{V}^\cdot\mapsto R\varphi_\bullet(\mathcal{V}^\cdot\overset{L}{\otimes}_{\mathcal{D}_X}\mathcal{D}_{X\to Y})$，从
$D(\mathcal{M}_R(\mathcal{D}_X))$ 到 $D(\mathcal{M}_R(\mathcal{D}_Y))$。我们称这个函子为直像函子
$\varphi_+:D(\mathcal{M}_R(\mathcal{D}_X))\to D(\mathcal{M}_R(\mathcal{D}_Y))$，特别地
$$\varphi_+(\mathcal{V}^\cdot)=R\varphi_\bullet(\mathcal{V}^\cdot\overset{L}{\otimes}_{\mathcal{D}_X}\mathcal{D}_{X\to Y})$$
对任意 $D(\mathcal{M}_R(\mathcal{D}_X))$ 中的 $\mathcal{V}^\cdot$ 成立。

\begin{lemma}\label{lem:4.5.1}
正合函子 $\varphi_+:D(\mathcal{M}_R(\mathcal{D}_X))\to D(\mathcal{M}_R(\mathcal{D}_Y))$ 具有有限振幅（amplitude）。
\end{lemma}

\begin{proof}
由于 $U\mapsto U\otimes_{\mathcal{D}_X}\mathcal{D}_{X\to Y}$ 具有有限左同调维数，函子
$\mathcal{V}^\cdot\mapsto\mathcal{V}^\cdot\overset{L}{\otimes}_{\mathcal{D}_X}\mathcal{D}_{X\to Y}$ 具有有限振幅。类似地，$R\varphi_\bullet$ 具有有限振幅。这意味着它们的复合具有有限振幅。
\end{proof}

\begin{remark}\label{rem:4.5.2}
设 $U$ 是 $X$ 中开集，$i:U\to X$ 是自然包含。由注记 \ref{rem:4.1.6} 有 $\mathcal{D}_{U\to X}=\mathcal{D}_U$。故此时对任意右 $\mathcal{D}_U$-模复形 $\mathcal{V}^\cdot$ 有
$$i_+(\mathcal{V}^\cdot)=Ri_\bullet(\mathcal{V}^\cdot).$$
特别地，在开包含的情形下，直像函子等于层直像函子 $i_\bullet$ 的右导出函子。
\end{remark}

设 $Z$ 是另一个光滑簇，$\psi:Y\to Z$ 是代数簇的态射。

\begin{theorem}\label{thm:4.5.3}
正合函子 $\psi_+\circ\varphi_+$ 与 $(\psi\circ\varphi)_+$（从 $D(\mathcal{M}_R(\mathcal{D}_X))$ 到
$D(\mathcal{M}_R(\mathcal{D}_Z))$）是同构的。
\end{theorem}

\begin{proof}
设 $\mathcal{V}^\cdot$ 是 $D(\mathcal{M}_R(\mathcal{D}_X))$ 中的复形。则
$$(\psi_+\circ\varphi_+)(\mathcal{V}^\cdot)=\psi_+(\varphi_+(\mathcal{V}^\cdot))
=R\psi_\bullet(\varphi_+(\mathcal{V}^\cdot)\overset{L}{\otimes}_{\mathcal{D}_Y}\mathcal{D}_{Y\to Z})$$
$$=R\psi_\bullet\bigl(R\varphi_\bullet(\mathcal{V}^\cdot\overset{L}{\otimes}_{\mathcal{D}_X}\mathcal{D}_{X\to Y})
\overset{L}{\otimes}_{\mathcal{D}_Y}\mathcal{D}_{Y\to Z}\bigr).$$
由投影公式，
$$R\varphi_\bullet(\mathcal{V}^\cdot\overset{L}{\otimes}_{\mathcal{D}_X}\mathcal{D}_{X\to Y})
\overset{L}{\otimes}_{\mathcal{D}_Y}\mathcal{D}_{Y\to Z}
=R\varphi_\bullet\bigl((\mathcal{V}^\cdot\overset{L}{\otimes}_{\mathcal{D}_X}\mathcal{D}_{X\to Y})
\overset{L}{\otimes}_{\varphi^{-1}(\mathcal{D}_Y)}\varphi^{-1}(\mathcal{D}_{Y\to Z})\bigr).$$
故由前述，
$$(\psi_+\circ\varphi_+)(\mathcal{V}^\cdot)
=R\psi_\bullet\bigl(R\varphi_\bullet(\mathcal{V}^\cdot\overset{L}{\otimes}_{\mathcal{D}_X}
(\mathcal{D}_{X\to Y}\overset{L}{\otimes}_{\varphi^{-1}(\mathcal{D}_Y)}\varphi^{-1}(\mathcal{D}_{Y\to Z})))\bigr)$$
$$=R(\psi\circ\varphi)_\bullet(\mathcal{V}^\cdot\overset{L}{\otimes}_{\mathcal{D}_X}\mathcal{D}_{X\to Z})
=(\psi\circ\varphi)_+(\mathcal{V}^\cdot).$$
\end{proof}

设 $\varphi:X\to Y$ 是光滑代数簇的态射。设 $V$ 是 $Y$ 中开集，$U=\varphi^{-1}(V)$。设 $\mathcal{V}^\cdot$ 是右 $\mathcal{D}_X$-模的复形。则
$$\varphi_+(\mathcal{V}^\cdot)|_V=R\varphi_\bullet(\mathcal{V}^\cdot\overset{L}{\otimes}_{\mathcal{D}_X}\mathcal{D}_{X\to Y})|_V
=R(\varphi|_U)_\bullet\bigl((\mathcal{V}^\cdot\overset{L}{\otimes}_{\mathcal{D}_X}\mathcal{D}_{X\to Y})|_U\bigr)$$
$$=R(\varphi|_U)_\bullet(\mathcal{V}^\cdot|_U\overset{L}{\otimes}_{\mathcal{D}_U}\mathcal{D}_{U\to V})
=(\varphi|_U)_+(\mathcal{V}^\cdot|_U),$$
即直像函子关于目标簇是局部的。

现设 $X$、$Y$ 是光滑仿射簇。则对任意拟凝聚右 $\mathcal{D}_X$-模 $\mathcal{V}$，设 $F^\cdot$ 是 $\Gamma(X,\mathcal{V})$ 的自由分解。则 $\mathcal{F}^\cdot=\Delta(F^\cdot)$ 是 $\mathcal{V}$ 的自由分解。因此
$$\varphi_+(\mathcal{D}(\mathcal{V}))=R\varphi_\bullet(\mathcal{D}(\mathcal{V})\overset{L}{\otimes}_{\mathcal{D}_X}\mathcal{D}_{X\to Y})
=R\varphi_\bullet(F^\cdot\otimes_{\mathcal{D}_X}\mathcal{D}_{X\to Y}).$$
显然对任意 $p\in\mathbb{Z}$，$F^p\otimes_{\mathcal{D}_X}\mathcal{D}_{X\to Y}$ 是 $\mathcal{D}_{X\to Y}$ 若干拷贝的直和，故它是拟凝聚左 $\mathcal{D}_X$-模。特别地，由前述它对 $\varphi_\bullet$ 是无循环的。这意味着
$$\varphi_+(\mathcal{D}(\mathcal{V}))=\varphi_\bullet(F^\cdot\otimes_{\mathcal{D}_X}\mathcal{D}_{X\to Y}).$$

\begin{lemma}\label{lem:4.5.4}
若 $\varphi:X\to Y$ 是光滑仿射簇的态射，则
\begin{enumerate}[label=(\roman*)]
  \item $H^p(\varphi_+(\mathcal{D}(\mathcal{D}_X)))=0$ 对所有 $p\ne 0$；
  \item $H^0(\varphi_+(\mathcal{D}(\mathcal{D}_X)))=\varphi_\bullet(\mathcal{D}_{X\to Y})$ 是拟凝聚右 $\mathcal{D}_Y$-模。
\end{enumerate}
\end{lemma}

\begin{proof}
显然
$$\varphi_+(\mathcal{D}(\mathcal{D}_X))=\varphi_\bullet(\mathcal{D}(\mathcal{D}_X)\otimes_{\mathcal{D}_X}\mathcal{D}_{X\to Y})
=\mathcal{D}(\varphi_\bullet(\mathcal{D}_{X\to Y})).$$
这蕴含 (i)。只需检验 $H^0(\varphi_+(\mathcal{D}(\mathcal{D}_X)))=\varphi_\bullet(\mathcal{D}_{X\to Y})$ 是拟凝聚的。首先注意
$$\Gamma(Y,\varphi_\bullet(\mathcal{D}_{X\to Y}))=\Gamma(X,\mathcal{D}_{X\to Y})=R(X)\otimes_{R(Y)}D_Y,$$
因为 $\mathcal{D}_{X\to Y}$ 是拟凝聚 $\mathcal{D}_X$-模。设 $g\in R(Y)$、$f=g\circ\varphi$。则 $\varphi^{-1}(Y_g)=X_f$，且由于 $\varphi_+$ 是局部的，
$$H^0(\varphi_+(\mathcal{D}(\mathcal{D}_X)))|_{Y_g}=H^0(\varphi_+(\mathcal{D}(\mathcal{D}_{X_f}))),$$
以及
$$\varphi_\bullet(\mathcal{D}_{X\to Y})(Y_g)=H^0(\varphi_+(\mathcal{D}(\mathcal{D}_X)))(Y_g)
=H^0(\varphi_+(\mathcal{D}(\mathcal{D}_{X_f})))(Y_g)$$
$$=\Gamma(Y_g,(\varphi|_{X_f})_\bullet(\mathcal{D}_{X_f\to Y_g}))
=R(X_f)\otimes_{R(Y_g)}D_{Y_g}=R(X)_f\otimes_{R(Y_g)}D_{Y_g}$$
$$=\bigl(R(X)\otimes_{R(Y)}R(Y_g)\bigr)\otimes_{R(Y_g)}D_{Y_g}
=R(X)\otimes_{R(Y)}D_{Y_g}=\bigl(R(X)\otimes_{R(Y)}D_Y\bigr)_g=(\mathcal{D}_{X\to Y})_g,$$
其中局部化是关于第二个因子中的右乘。这说明 $\varphi_\bullet(\mathcal{D}_{X\to Y})$ 是 $\mathcal{D}_{X\to Y}$ 作为右 $\mathcal{D}_Y$-模的局部化，故它是拟凝聚的。
\end{proof}

因此，$\varphi_\bullet(F^\cdot\otimes_{\mathcal{D}_X}\mathcal{D}_{X\to Y})$ 中的右 $\mathcal{D}_Y$-模是拟凝聚的。这意味着该复形的上同调群也是拟凝聚右 $\mathcal{D}_Y$-模。于是 $H^p(\varphi_+(\mathcal{D}(\mathcal{V})))$ 是拟凝聚右 $\mathcal{D}_Y$-模。此外，
$$\Gamma(Y,H^p(\varphi_+(\mathcal{D}(\mathcal{V}))))=H^p(R\Gamma(\varphi_+(\mathcal{D}(\mathcal{V}))))$$
$$=H^p(R\Gamma(\varphi_\bullet(F^\cdot\otimes_{\mathcal{D}_X}\mathcal{D}_{X\to Y})))
=H^p(R\Gamma(F^\cdot\otimes_{\mathcal{D}_X}\mathcal{D}_{X\to Y}))$$
$$=H^p(F^\cdot\otimes_{\mathcal{D}_X}\mathcal{D}_{X\to Y})
=H^p(\Gamma(X,\mathcal{V})\otimes_{\mathcal{D}_X}\mathcal{D}_{X\to Y})
=\operatorname{Tor}^{-\,p}_{\mathcal{D}_X}(\Gamma(X,\mathcal{V}),\mathcal{D}_{X\to Y}).$$
因此我们证明了以下结果。

\begin{proposition}\label{prop:4.5.5}
设 $\varphi:X\to Y$ 是光滑仿射簇的态射。设 $\mathcal{V}$ 是拟凝聚右 $\mathcal{D}_X$-模。则：
\begin{enumerate}[label=(\roman*)]
  \item 对 $p\in\mathbb{Z}$，$H^p(\varphi_+(\mathcal{D}(\mathcal{V})))$ 是拟凝聚右 $\mathcal{D}_Y$-模；
  \item 对 $p>0$，$H^p(\varphi_+(\mathcal{D}(\mathcal{V})))=0$；
  \item 对 $p<0$，有
    $$\Gamma(Y,H^p(\varphi_+(\mathcal{D}(\mathcal{V}))))
    =\operatorname{Tor}^{-\,p}_{\mathcal{D}_X}(\Gamma(X,\mathcal{V}),\mathcal{D}_{X\to Y}).$$
\end{enumerate}
\end{proposition}

因此，在光滑仿射簇上拟凝聚右 D-模的情形下，我们的定义与前述旧定义一致。

设 $\varphi:X\to Y$ 是仿射态射。于是对 $Y$ 中仿射开集 $V$，集合 $U=\varphi^{-1}(V)$ 也是 $X$ 中仿射开子集。设 $\mathcal{V}$ 是拟凝聚右 $\mathcal{D}_X$-模。则
$$\varphi_+(\mathcal{D}(\mathcal{V}))|_V=(\varphi|_U)_+(\mathcal{D}(\mathcal{V}|_U)).$$
因此前述结果有以下推论。

\begin{corollary}\label{cor:4.5.6}
设 $\varphi:X\to Y$ 是光滑簇之间的仿射态射。设 $\mathcal{V}$ 是拟凝聚右 $\mathcal{D}_X$-模。则：
\begin{enumerate}[label=(\roman*)]
  \item 对 $p\in\mathbb{Z}$，$H^p(\varphi_+(\mathcal{D}(\mathcal{V})))$ 是拟凝聚右 $\mathcal{D}_Y$-模；
  \item 对 $p>0$，$H^p(\varphi_+(\mathcal{D}(\mathcal{V})))=0$。
\end{enumerate}
\end{corollary}

第一个陈述在一般情形下也成立。先考虑上述结果的一个特殊情形。

\begin{corollary}\label{cor:4.5.7}
设 $U$ 是 $X$ 中仿射开集，$i:U\to X$ 是自然包含。设 $\mathcal{V}$ 是拟凝聚右 $\mathcal{D}_U$-模。则对任意拟凝聚右 $\mathcal{D}_U$-模 $\mathcal{V}$ 与 $p\ne 0$ 有 $H^p(i_+(\mathcal{D}(\mathcal{V})))=0$。此外，函子 $H^0\circ i_+\circ\mathcal{D}=i_\bullet$ 从拟凝聚右 $\mathcal{D}_U$-模范畴到右 $\mathcal{D}_X$-模范畴是正合的。
\end{corollary}

现在可以证明上述结果的如下推广。

\begin{theorem}\label{thm:4.5.8}
设 $\varphi:X\to Y$ 是光滑代数簇的态射。设 $\mathcal{V}^\cdot$ 是右 $\mathcal{D}_X$-模的复形，且对所有 $p\in\mathbb{Z}$，$H^p(\mathcal{V}^\cdot)$ 是拟凝聚右 $\mathcal{D}_X$-模。则对所有 $p\in\mathbb{Z}$，$H^p(\varphi_+(\mathcal{V}^\cdot))$ 是拟凝聚右 $\mathcal{D}_Y$-模。
\end{theorem}

\begin{proof}
设 $U$ 是 $X$ 中仿射开子集，$i:U\to X$ 是自然浸入。设 $\mathcal{W}$ 是拟凝聚右 $\mathcal{D}_U$-模。则由推论 \ref{cor:4.5.7}，$i_+(\mathcal{D}(\mathcal{W}))=Ri_\bullet(\mathcal{D}(\mathcal{W}))=\mathcal{D}(i_\bullet(\mathcal{W}))$ 是具有拟凝聚上同调的复形。又因 $\varphi\circ i:U\to Y$ 是仿射态射，由定理 \ref{thm:4.5.3} 有
$\varphi_+(i_+(\mathcal{D}(\mathcal{W})))=(\varphi\circ i)_+(\mathcal{D}(\mathcal{W}))$，而由推论 \ref{cor:4.5.6}，后者是具有拟凝聚上同调的复形。由于 $i_\bullet(\mathcal{W})$ 形式的模生成 $D^b_{qc}(\mathcal{M}_R(\mathcal{D}_X))$，可知该陈述对有界复形成立。最终结论由标准截断点论证得出。
\end{proof}

\section{浸入的直像}

设 $X$ 是光滑代数簇，$Y$ 是其光滑子簇。令 $n=\dim X$、$m=\dim Y$。设 $i:Y\to X$ 是典范浸入。考虑 $X$ 上（相应地 $Y$ 上）右 D-模的范畴 $\mathcal{M}_R(\mathcal{D}_X)$（相应地 $\mathcal{M}_R(\mathcal{D}_Y)$），以及相应的有界导出范畴 $D^b(\mathcal{M}_R(\mathcal{D}_Y))$、$D^b(\mathcal{M}_R(\mathcal{D}_X))$。则有直像函子
$i_+:D^b(\mathcal{M}_R(\mathcal{D}_Y))\to D^b(\mathcal{M}_R(\mathcal{D}_X))$。

\begin{proposition}\label{prop:4.6.1}
模 $\mathcal{D}_{Y\to X}$ 是局部自由 $\mathcal{D}_Y$-模。
\end{proposition}

为证明这个结果，先构造一个适应当前情形的局部坐标系。

\begin{lemma}\label{lem:4.6.2}
设 $Y$ 是光滑 $n$ 维簇 $X$ 的光滑 $m$ 维子簇。设 $y\in Y$。则存在 $y$ 在 $X$ 中的仿射开邻域 $U$ 及 $U$ 上的坐标系
$(f_1,f_2,\dots,f_n;D_1,D_2,\dots,D_n)$，具有如下性质：
\begin{enumerate}[label=(\roman*)]
  \item $U\cap Y$ 是 $U$ 的闭子簇；
  \item $R(U)$ 中所有在 $U\cap Y$ 上消没的函数构成的理想 $I(U\cap Y)$ 由
    $f_{m+1},f_{m+2},\dots,f_n$ 生成；
  \item 向量场 $D_1,D_2,\dots,D_m$ 把理想 $I(U\cap Y)$ 映入自身。
\end{enumerate}
\end{lemma}

\begin{proof}
这个证明是 2.9 证明的一个小变体。由于陈述是局部的，显然可设 $Y$ 在 $X$ 中闭且 $X$ 是某个 $k^p$ 的闭子簇。以 $I_X$、$I_Y$ 分别记 $A=k[X_1,X_2,\dots,X_p]$ 中所有在 $X$、$Y$ 上消没的多项式构成的理想。显然 $I_X\subset I_Y$。由于
$\dim_k T_y(X)=\dim X=n$、$\dim_k T_y(Y)=\dim Y=m$，由 2.6 可找到多项式
$P_{m+1},P_{m+2},\dots,P_p\in I_Y$ 使得：
\begin{enumerate}[label=(\alph*)]
  \item $P_{n+1},P_{n+2},\dots,P_p$ 在 $I_X$ 中；
  \item 矩阵 $[(\partial_iP_j)(y)]$ 的秩为 $p-m$。
\end{enumerate}
这意味着该矩阵的秩在 $y\in k^p$ 的某个邻域 $V$ 上等于 $p-m$，且
$$T_y(X)=\{(\xi_1,\dots,\xi_p)\in k^p\mid
\sum_{i=1}^p\xi_i(\partial_iP_j)(y)=0,\ n+1\le j\le p\},$$
$$T_y(Y)=\{(\xi_1,\dots,\xi_p)\in k^p\mid
\sum_{i=1}^p\xi_i(\partial_iP_j)(y)=0,\ m+1\le j\le p\}.$$
现在如 2.9 的证明，可找到 $g\in A$ 使得 $(I_X)_g$ 在 $A_g$ 中由 $P_{n+1},\dots,P_p$ 生成，$(I_Y)_g$ 由 $P_{m+1},\dots,P_p$ 生成。

可找到多项式 $P_1,P_2,\dots,P_m\in A$ 使得矩阵 $[(\partial_iP_j)(y);1\le i,j\le p]$ 非奇异。因此必要时改变 $g$，还可设它在 $k^p$ 的主开集 $V'$ 上非奇异。以 $Q$ 记此矩阵的逆。则 $Q$ 的矩阵系数在 $A_g$ 中。因此在 $V'$ 上可定义微分算子
$\delta_i=\sum_{j=1}^pQ_{ij}\partial_j$（$1\le i\le n$）。显然对任意 $1\le j\le p$ 满足
$$\delta_i(P_j)=\sum_{k=1}^pQ_{ik}\partial_kP_j=\delta_{ij}.$$
由于任意 $f\in(I_X)_g$ 可表示为 $f=\sum_{j=n+1}^ph_jP_j$（$h_j\in A_g$），有
$$\delta_i(f)=\delta_i\Bigl(\sum_{j=n+1}^ph_jP_j\Bigr)
=\sum_{j=n+1}^p\bigl(\delta_i(h_j)P_j+h_j\delta_i(P_j)\bigr)
=\sum_{j=n+1}^p\delta_i(h_j)P_j\in(I_X)_g,$$
即 $(I_X)_g$ 在 $\delta_i$（$1\le i\le n$）作用下不变。设 $U=X\cap V'$。由于 $R(U)=A_g/(I_X)_g$，这说明 $\delta_i$（$1\le i\le n$）诱导出 $U=X\cap V'$ 上的局部向量场，记为 $D_i$（$1\le i\le n$）。

显然 $I(U\cap Y)=(I_Y)_g/(I_X)_g$，故它由函数 $f_i=P_i|_U$（$m+1\le i\le n$）生成。这证明了 (ii)。

由类似计算还可见 $(I_Y)_g$ 在 $\delta_i$（$1\le i\le m$）作用下不变。因此 $D_1,D_2,\dots,D_m$ 把理想 $I(U\cap Y)$ 映入自身。显然 $D_i(f_j)=\delta_i(P_j)=\delta_{ij}$，故 $(f_1,\dots,f_n;D_1,\dots,D_n)$ 是 $U$ 上的坐标系。
\end{proof}

现在证明命题 \ref{prop:4.6.1}。先设 $X$ 在如下意义下是“小”的。存在 $X$ 上的坐标系 $(f_1,\dots,f_n;D_1,\dots,D_n)$ 使得：
\begin{enumerate}[label=(\roman*)]
  \item 对所有 $I,J\in\mathbb{Z}_+^n$，$D^I\circ D^J=D^{I+J}$；
  \item $(D^I;I\in\mathbb{Z}_+^n)$ 是自由 $R(X)$-模 $\mathcal{D}_X$ 关于左乘（相应地右乘）的一组基；
  \item $R(X)$ 中所有在 $Y$ 上消没的函数构成的理想 $I(Y)$ 由 $f_{m+1},\dots,f_n$ 生成；
  \item 向量场 $D_1,D_2,\dots,D_m$ 把理想 $I(Y)$ 映入自身。
\end{enumerate}
由 3.5 与 8，任意 $y\in Y$ 有这样的“小”邻域 $U$。此外，由 (iv)，$D_1,D_2,\dots,D_m$ 诱导出 $Y$ 上向量场 $T_1,T_2,\dots,T_m$。若以 $g_i$ 记 $f_i$ 在 $Y$ 上的限制（$1\le i\le m$），可见 $(g_1,\dots,g_m;T_1,\dots,T_m)$ 是 $Y$ 上的坐标系。必要时缩小 $X$，由 3.5 还可设
\begin{enumerate}[label=(\roman*),resume]
  \item $(T^I;I\in\mathbb{Z}_+^m)$ 是自由 $R(Y)$-模 $\mathcal{D}_Y$ 关于左乘（相应地右乘）的一组基。
\end{enumerate}
在这些条件下有如下结果。

\begin{lemma}\label{lem:4.6.3}
$(D^I;I\in\{0\}\times\mathbb{Z}_+^{n-m})$ 是自由左 $\mathcal{D}_Y$-模 $\mathcal{D}_{Y\to X}$ 的一组基。
\end{lemma}

\begin{proof}
由 (ii) 可见正合列
$$0\longrightarrow I(Y)\longrightarrow R(X)\longrightarrow R(Y)\longrightarrow 0$$
导致（$R(X)$，右 $\mathcal{D}_X$）-双模的正合列
$$0\longrightarrow I(Y)\otimes_{R(X)}\mathcal{D}_X\longrightarrow\mathcal{D}_X
\longrightarrow\mathcal{D}_{Y\to X}\longrightarrow 0.$$
由于 $D_i$（$1\le i\le m$）保持 $I(Y)$ 不变，左乘 $D_i$ 把 $I(Y)\otimes_{R(X)}\mathcal{D}_X$ 映入自身，并诱导出 $R(Y)$-模 $\mathcal{D}_{Y\to X}$ 的微分自同态，且与 $\mathcal{D}_X$ 的右作用可交换。此外它把 $g\otimes1\in\mathcal{D}_{Y\to X}$ 映为 $T_i(g)\otimes1+g\otimes D_i$，故它等于 $T_i$ 在 $\mathcal{D}_{Y\to X}$ 上的左作用。因此 $T_i$（$1\le i\le m$）在 $\mathcal{D}_{Y\to X}$ 上的左作用为
$$T_i(g\otimes S)=T_i(g)\otimes S+g\otimes D_iS$$
对 $g\in R(Y)$、$S\in\mathcal{D}_X$ 成立。这意味着对任意
$I\in\mathbb{Z}_+^m\times\{0\}\subset\mathbb{Z}_+^n$ 与 $J\in\{0\}\times\mathbb{Z}_+^{n-m}\subset\mathbb{Z}_+^n$ 有
$T^I(1\otimes D^J)=1\otimes D^{I+J}$。特别地，$(D^I;I\in\{0\}\times\mathbb{Z}_+^{n-m})$ 生成左 $\mathcal{D}_Y$-模 $\mathcal{D}_{Y\to X}$。由于由 (ii)，$(D^I;I\in\mathbb{Z}_+^n)$ 是自由 $R(Y)$-模 $\mathcal{D}_{Y\to X}$ 的一组基，由 (v) 可知 $(D^I;I\in\{0\}\times\mathbb{Z}_+^{n-m})$ 是自由左 $\mathcal{D}_Y$-模 $\mathcal{D}_{Y\to X}$ 的一组基。
\end{proof}

因此，由命题 \ref{prop:4.6.1}，对 $D^b(\mathcal{M}_R(\mathcal{D}_Y))$ 中任意有界复形有
$$i_+(\mathcal{V}^\cdot)=Ri_\bullet(\mathcal{V}^\cdot\overset{L}{\otimes}_{\mathcal{D}_Y}\mathcal{D}_{Y\to X})
=Ri_\bullet(\mathcal{V}^\cdot\otimes_{\mathcal{D}_Y}\mathcal{D}_{Y\to X}).$$
下一个结果表明 $i_+$ 对浸入是右导出函子。

\begin{theorem}\label{thm:4.6.4}
设 $Y$ 是光滑簇 $X$ 的光滑子簇，$i:Y\to X$ 是典范浸入。则：
\begin{enumerate}[label=(\roman*)]
  \item 函子 $H^0\circ i_+\circ\mathcal{D}:\mathcal{M}_R(\mathcal{D}_Y)\to\mathcal{M}_R(\mathcal{D}_X)$，由
    $$\mathcal{V}\longmapsto H^0(i_+(\mathcal{D}(\mathcal{V})))
    =i_\bullet(\mathcal{V}\otimes_{\mathcal{D}_Y}\mathcal{D}_{Y\to X})$$
    给出，是左正合的；
  \item 函子 $i_+$ 是 $H^0\circ i_+\circ\mathcal{D}:\mathcal{M}_R(\mathcal{D}_Y)\to\mathcal{M}_R(\mathcal{D}_X)$ 的右导出函子；
  \item $H^0(i_+(\mathcal{D}(\mathcal{V})))$ 的支集等于 $\operatorname{supp}(\mathcal{V})$ 在 $X$ 中的闭包。
\end{enumerate}
\end{theorem}

\begin{proof}
\begin{enumerate}[label=(\roman*)]
  \item 这个断言是显然的，因为
    $$H^0(i_+(\mathcal{D}(\mathcal{V})))=H^0(Ri_\bullet(\mathcal{D}(\mathcal{V})\otimes_{\mathcal{D}_Y}\mathcal{D}_{Y\to X}))$$
    $$=H^0(Ri_\bullet(\mathcal{D}(\mathcal{V}\otimes_{\mathcal{D}_Y}\mathcal{D}_{Y\to X})))
    =i_\bullet(\mathcal{V}\otimes_{\mathcal{D}_Y}\mathcal{D}_{Y\to X}),$$
    而由命题 \ref{prop:4.6.1}，函子 $\mathcal{V}\mapsto\mathcal{V}\otimes_{\mathcal{D}_Y}\mathcal{D}_{Y\to X}$ 是正合的，$i_\bullet$ 是左正合的。
  \item 设 $I$ 是内射右 $\mathcal{D}_Y$-模。存在 $Y$ 的开覆盖 $\{U_j;1\le j\le p\}$ 使得
    $\mathcal{D}_{Y\to X}|_{U_j}$ 是自由 $\mathcal{D}_{U_j}$-模。这意味着
    $(I\otimes_{\mathcal{D}_Y}\mathcal{D}_{Y\to X})|_{U_j}$ 是 $I|_{U_j}$ 无穷多个拷贝的直和。由于
    $I|_{U_j}$ 也是内射的，它是松的（flabby），因此 $(I\otimes_{\mathcal{D}_Y}\mathcal{D}_{Y\to X})|_{U_j}$ 是松的。由此可知正合函子
    $\mathcal{V}\mapsto\mathcal{V}\otimes_{\mathcal{D}_Y}\mathcal{D}_{Y\to X}$ 把 $\mathcal{M}_R(\mathcal{D}_Y)$ 中的内射对象映为松层，即映为对 $i_\bullet$ 无循环的层。这意味着相应导出函子的复合
    $\mathcal{V}^\cdot\mapsto Ri_\bullet(\mathcal{V}^\cdot\otimes_{\mathcal{D}_Y}\mathcal{D}_{Y\to X})$ 是左正合函子
    $\mathcal{V}\mapsto i_\bullet(\mathcal{V}\otimes_{\mathcal{D}_Y}\mathcal{D}_{Y\to X})$ 的右导出函子。
  \item 设 $U$ 是 $X$ 中开集。则
    $$H^0(i_+(\mathcal{D}(\mathcal{V})))(U)=i_\bullet(\mathcal{V}\otimes_{\mathcal{D}_Y}\mathcal{D}_{Y\to X})(U)
    =(\mathcal{V}\otimes_{\mathcal{D}_Y}\mathcal{D}_{Y\to X})(U\cap Y).$$
    因此 $H^0(i_+(\mathcal{D}(\mathcal{V})))$ 在 $\operatorname{supp}(\mathcal{V})$ 的补集的任一开集上的限制等于 0。故
    $$\operatorname{supp}(H^0(i_+(\mathcal{D}(\mathcal{V}))))\subset\operatorname{supp}(\mathcal{V}).$$
    设 $y\in\operatorname{supp}(\mathcal{V})$。先假设 $y$ 不在 $\operatorname{supp}(H^0(i_+(\mathcal{D}(\mathcal{V}))))$ 中。则存在引理 \ref{lem:4.6.2} 所述的 $y$ 在 $X$ 中的“小”仿射邻域 $U$，使得 $H^0(i_+(\mathcal{D}(\mathcal{V})))(U)=0$。由引理 \ref{lem:4.6.3}，$(D^I;I\in\{0\}\times\mathbb{Z}_+^{n-m})$ 是自由 $\mathcal{D}_Y$-模 $\mathcal{D}_{Y\to X}$ 的一组基。这说明 $(\mathcal{V}\otimes_{\mathcal{D}_Y}\mathcal{D}_{Y\to X})|_{U\cap Y}$ 是 $\mathcal{V}|_{U\cap Y}$ 无穷多个拷贝的直和。故 $H^0(i_+(\mathcal{D}(\mathcal{V})))(U)=0$ 将蕴含 $\mathcal{V}|_{U\cap Y}=0$，这是不可能的。因此 $y\in\operatorname{supp}(H^0(i_+(\mathcal{D}(\mathcal{V}))))$。于是
    $$\operatorname{supp}(\mathcal{V})\subset\operatorname{supp}(H^0(i_+(\mathcal{D}(\mathcal{V})))).$$
    这证明了等式。
\end{enumerate}
\end{proof}

结合前述结果，得到以下推论。

\begin{corollary}\label{cor:4.6.5}
设 $i:Y\to X$ 是仿射浸入。则函子
$H^0\circ i_+\circ\mathcal{D}:\mathcal{M}^R_{qc}(\mathcal{D}_Y)\to\mathcal{M}^R_{qc}(\mathcal{D}_X)$ 是正合的。
\end{corollary}

若 $Y$ 是 $X$ 的闭光滑子簇，$i_\bullet$ 是正合函子。因此有如下结果。

\begin{corollary}\label{cor:4.6.6}
设 $Y$ 是 $X$ 的闭光滑子簇。则函子
$H^0\circ i_+\circ\mathcal{D}:\mathcal{M}_R(\mathcal{D}_Y)\to\mathcal{M}_R(\mathcal{D}_X)$ 是正合的。
\end{corollary}

\begin{proof}
由于 $i_\bullet$ 是正合函子，
$$\mathcal{V}\longmapsto H^0(i_+(\mathcal{D}(\mathcal{V})))=i_\bullet(\mathcal{V}\otimes_{\mathcal{D}_Y}\mathcal{D}_{Y\to X})$$
是从 $\mathcal{M}_R(\mathcal{D}_Y)$ 到 $\mathcal{M}_R(\mathcal{D}_X)$ 的正合函子。
\end{proof}

借由记号上的滥用，此时也把函子 $H^0\circ i_+\circ\mathcal{D}$ 记为 $i_+$。因此有
$$i_+(\mathcal{V})=i_\bullet(\mathcal{V}\otimes_{\mathcal{D}_Y}\mathcal{D}_{Y\to X})$$
对任意 $\mathcal{M}_R(\mathcal{D}_Y)$ 中的 $\mathcal{V}$ 成立，且对任意 $D^b(\mathcal{M}_R(\mathcal{D}_X))$ 中的 $\mathcal{V}^\cdot$ 有
$(i_+(\mathcal{V}^\cdot))^p=i_+(\mathcal{V}^p)$。

\begin{proposition}\label{prop:4.6.7}
设 $Y$ 是 $X$ 的闭光滑子簇。则对任意凝聚右 $\mathcal{D}_Y$-模 $\mathcal{V}$，右 $\mathcal{D}_X$-模 $i_+(\mathcal{V})$ 也是凝聚的。
\end{proposition}

\begin{proof}
只需证明对 $X$ 中任意仿射开集 $U$，$i_+(\mathcal{V})|_U$ 是凝聚 $\mathcal{D}_U$-模。因此用 $U$ 代替 $X$、$Y\cap U$ 代替 $Y$，可设 $X$ 是光滑仿射簇。此时 $i_+(\mathcal{V})$ 是右 $\mathcal{D}_X$-模 $\Gamma(Y,\mathcal{V})\otimes_{\mathcal{D}_Y}\mathcal{D}_{Y\to X}$ 的局部化。由前述，只需证 $\Gamma(Y,\mathcal{V})\otimes_{\mathcal{D}_Y}\mathcal{D}_{Y\to X}$ 是有限生成右 $\mathcal{D}_X$-模。由于 $Y$ 是 $X$ 的闭子簇，$R(Y)$ 是 $R(X)$ 的商，且 $\mathcal{D}_{Y\to X}=R(Y)\otimes_{R(X)}\mathcal{D}_X$ 是 $\mathcal{D}_X$ 作为右 $\mathcal{D}_X$-模的商。因此它由元素 $1\otimes1$ 生成。由假设，$\Gamma(Y,\mathcal{V})$ 是有限生成 $\mathcal{D}_Y$-模。设 $v_1,\dots,v_n$ 是它的一组生成元。则 $v_1\otimes1\otimes1,\dots,v_n\otimes1\otimes1$ 生成
$\Gamma(Y,\mathcal{V})\otimes_{\mathcal{D}_Y}\mathcal{D}_{Y\to X}$ 作为右 $\mathcal{D}_X$-模。
\end{proof}

特别地，若 $\mathcal{V}$ 是凝聚右 $\mathcal{D}_Y$-模，可比较 $\mathcal{V}$ 与 $i_+(\mathcal{V})$ 的特征簇。设 $x\in Y$。则 $T_x(Y)\subset T_x(X)$，故有自然投影 $p_x:T_x^*(X)\to T_x^*(Y)$。

设 $Y$ 是代数簇 $X$ 的闭子集。定义 $\dim_YX=\sup_{x\in Z}(\dim_xX)$（原文如此，应为 $x\in Y$）。

\begin{theorem}\label{thm:4.6.8}
设 $Y$ 是 $X$ 的闭光滑子簇。则对任意凝聚右 $\mathcal{D}_Y$-模 $\mathcal{V}$ 有
$$\operatorname{Ch}(i_+(\mathcal{V}))=\{(x,\omega)\in T^*(X)\mid(x,p_x(\omega))\in\operatorname{Ch}(\mathcal{V})\}.$$
特别地，
$$\dim_{p_X^{-1}(x)}\operatorname{Ch}(i_+(\mathcal{V}))
=\dim_{p_Y^{-1}(x)}\operatorname{Ch}(\mathcal{V})+\operatorname{codim}_xY.$$
\end{theorem}

\begin{proof}
由于 $i_+(\mathcal{V})$ 的支集在 $Y$ 中，只需证明每个 $y\in Y$ 有 $X$ 中开邻域 $U$ 使得
$$\operatorname{Ch}(i_+(\mathcal{V})|_U)=\{(x,\omega)\in T^*(U)\mid(x,p_x(\omega))\in\operatorname{Ch}(\mathcal{V}|_{Y\cap U})\}.$$
因此可设该邻域满足引理 \ref{lem:4.6.2} 的条件。用此邻域代替 $X$，可设 $X$ 是“小”的。此时由引理 \ref{lem:4.6.3}，$\mathcal{D}_{Y\to X}$ 是自由左 $\mathcal{D}_Y$-模。设 $V=\Gamma(Y,\mathcal{V})$。于是 $i_+(\mathcal{V})$ 是右 $\mathcal{D}_X$-模 $V\otimes_{\mathcal{D}_Y}\mathcal{D}_{Y\to X}$ 的局部化，且作为向量空间
$$\Gamma(X,i_+(\mathcal{V}))=V\otimes_{\mathcal{D}_Y}\mathcal{D}_{Y\to X}
=\bigoplus_{I\in\{0\}\times\mathbb{Z}_+^{n-m}}V\otimes_kD^I.$$
设 $FV$ 是右 $\mathcal{D}_Y$-模 $V$ 的好滤过。则由 I.3.1，$\operatorname{Gr}V$ 是有限生成 $\operatorname{Gr}\mathcal{D}_Y$-模。可通过下式定义 $\Gamma(X,i_+(\mathcal{V}))$ 的滤过：
$$F_p\Gamma(X,i_+(\mathcal{V}))=\bigoplus_{I\in\{0\}\times\mathbb{Z}_+^{n-m},\ s+|I|\le p}F_sV\otimes D^I.$$
显然这是穷尽的递增滤过，且对充分负的 $p\in\mathbb{Z}$ 有 $F_p\Gamma(X,i_+(\mathcal{V}))=0$。

断言这是右 $\mathcal{D}_X$-模滤过。设 $v\otimes D^I\in\Gamma(X,i_+(\mathcal{V}))$，
$v\in F_sV$，$I\in\{0\}\times\mathbb{Z}_+^{n-m}$。断言对任意 $f\in R(X)$ 有
$$(v\otimes D^I)f\in F_{s+|I|}\Gamma(X,i_+(\mathcal{V})).$$
对 $|I|$ 作归纳。若 $|I|=0$，则
$$(v\otimes1)f=vf|_Y\otimes1\in F_sV\otimes1.$$
设 $|I|>0$。则可找到 $m+1\le j\le n$ 与 $I'\in\{0\}\times\mathbb{Z}_+^{n-m}$、$|I'|=|I|-1$，使得 $D^I=D^{I'}D_j$。故
$$(v\otimes D^I)f=(v\otimes D^{I'})D_jf=(v\otimes D^{I'})D_j(f)+(v\otimes D^{I'})fD_j.$$
由归纳假设，$(v\otimes D^{I'})D_j(f)$ 与 $(v\otimes D^{I'})f$ 都在 $F_{s+|I'|}\Gamma(X,i_+(\mathcal{V}))$ 中。故
$$(v\otimes D^{I'})fD_j\in F_{s+|I'|}\Gamma(X,i_+(\mathcal{V}))D_j
\subset F_{s+|I|}\Gamma(X,i_+(\mathcal{V})).$$
这证明了我们的断言，即 $F\Gamma(X,i_+(\mathcal{V}))$ 是由 $R(X)$-子模构成的滤过。

设 $1\le j\le n$。断言有
$$(v\otimes D^I)D_j\in F_{s+|I|+1}\Gamma(X,i_+(\mathcal{V})).$$
当 $m+1\le j\le n$ 时显然。若 $1\le j\le m$，则
$$(v\otimes D^I)D_j=vT_j\otimes D^I\in F_{s+|I|+1}\Gamma(X,i_+(\mathcal{V})),$$
因为 $vT_j\in F_{s+1}V$。由前述，这说明 $F\Gamma(X,i_+(\mathcal{V}))$ 是 $\mathcal{D}_X$-模滤过。

显然
$$\operatorname{Gr}\Gamma(X,i_+(\mathcal{V}))=\operatorname{Gr}V\otimes_kk[\xi_{m+1},\dots,\xi_n],$$
且 $\operatorname{Gr}_p\Gamma(X,i_+(\mathcal{V}))$ 由 $v\otimes\xi^I$（$v\in\operatorname{Gr}_sV$，$I\in\{0\}\times\mathbb{Z}_+^{n-m}$，$|I|=p-s$）张成。又有
$\operatorname{Gr}\mathcal{D}_X=R(X)[\xi_1,\dots,\xi_n]=R(X)[\xi_1,\dots,\xi_m]\otimes_kk[\xi_{m+1},\dots,\xi_n]$。$\operatorname{Gr}\mathcal{D}_X$ 在 $\operatorname{Gr}\Gamma(X,i_+(\mathcal{V}))$ 上的作用如下：$f\in R(X)$ 乘以第一因子中限制 $f|_Y$；$\xi_j$（$1\le j\le m$）在第一因子中相乘；$\xi_j$（$m+1\le j\le n$）在第二因子中相乘。

因此，若 $v_1,\dots,v_k$ 是 $\operatorname{Gr}V$ 作为 $\operatorname{Gr}\mathcal{D}_Y$-模的生成元，则 $v_1\otimes1,\dots,v_k\otimes1$ 生成 $\operatorname{Gr}\Gamma(X,i_+(\mathcal{V}))$ 作为 $\operatorname{Gr}\mathcal{D}_X$-模。于是 $\operatorname{Gr}\Gamma(X,i_+(\mathcal{V}))$ 是有限生成 $\operatorname{Gr}\mathcal{D}_X$-模。故由 I.3.1，$F\Gamma(X,i_+(\mathcal{V}))$ 是好滤过。

设 $A\subset R(Y)[\xi_1,\dots,\xi_m]$ 是 $\operatorname{Gr}V$ 在 $\operatorname{Gr}\mathcal{D}_Y$ 中的零化子。限制映射 $f\mapsto f|_Y$ 定义了 $R(X)$ 到 $R(Y)$ 的满同态。它诱导 $R(X)[\xi_1,\dots,\xi_m]$ 到 $R(Y)[\xi_1,\dots,\xi_m]$ 的满同态 $\varphi$。令 $B=\varphi^{-1}(A)$。则 $B\otimes_kk[\xi_{m+1},\dots,\xi_n]$ 零化 $\operatorname{Gr}\Gamma(X,i_+(\mathcal{V}))$。设 $S\otimes\xi^J$ 是 $\operatorname{Gr}\Gamma(X,i_+(\mathcal{V}))$ 零化子中的齐次元素。则对 $v\in\operatorname{Gr}V$ 有
$(S\otimes\xi^J)(v\otimes1)=Sv\otimes\xi^J=0$，即 $Sv=0$，$S\in B$。这说明 $\operatorname{Gr}\Gamma(X,i_+(\mathcal{V}))$ 的零化子等于 $B\otimes_kk[\xi_{m+1},\dots,\xi_n]$。

我们把 $R(T^*(Y))$ 与 $R(Y)[\xi_1,\dots,\xi_m]$ 的等同，对应于把 $T^*(Y)$ 与 $Y\times k^m$ 通过
$(y,\eta)\mapsto(y,(\eta(T_1(y)),\dots,\eta(T_m(y))))$ 的等同。类似地，$R(T^*(X))$ 与 $R(X)[\xi_1,\dots,\xi_n]$ 的等同对应于把 $T^*(X)$ 与 $X\times k^n$ 通过 $(x,\omega)\mapsto(x,(\omega(D_1(x)),\dots,\omega(D_n(x))))$ 的等同。在这些等同下，特征簇 $\operatorname{Ch}(\mathcal{V})$ 对应于 $Y\times k^m$ 中 $A$ 的零点集构成的子簇，而 $\operatorname{Ch}(i_+(\mathcal{V}))$ 对应于 $X\times k^n$ 中
$B\otimes_kk[\xi_{m+1},\dots,\xi_n]$ 的零点集构成的子簇。可通过映射
$(y,\xi_1,\dots,\xi_m)\mapsto(y,\xi_1,\dots,\xi_m,0,\dots,0)$ 把 $Y\times k^m$ 嵌入 $X\times k^n$。这个嵌入对应于 $R(X)[\xi_1,\dots,\xi_n]$ 到 $R(Y)[\xi_1,\dots,\xi_m]$ 的自然投影。在此嵌入下有等同 $\operatorname{Ch}(i_+(\mathcal{V}))=\operatorname{Ch}(\mathcal{V})\times k^{n-m}$。这立即蕴含
$$\dim\operatorname{Ch}(i_+(\mathcal{V}))=\dim\operatorname{Ch}(\mathcal{V})+n-m.$$
此外，点 $(x,\omega)$ 属于 $\operatorname{Ch}(i_+(\mathcal{V}))$ 当且仅当
$(x,(\omega(D_1(x)),\dots,\omega(D_m(x))))$ 对应于 $\operatorname{Ch}(\mathcal{V})$ 中的点，故 $x\in Y$。由于对 $x\in Y$、$1\le i\le m$ 有 $\omega(D_i(x))=p_x(\omega)(T_i(x))$，可见 $(x,\omega)\in\operatorname{Ch}(i_+(\mathcal{V}))$ 等价于 $x\in Y$ 且 $(x,p_x(\omega))\in\operatorname{Ch}(\mathcal{V})$。
\end{proof}

设 $\mathcal{V}$ 是 $X$ 上非零凝聚 $\mathcal{D}_X$-模。对任意 $x\in X$，记
$$\operatorname{holdef}_x(\mathcal{V})=\dim_{p_X^{-1}(x)}\operatorname{Ch}(\mathcal{V})-\dim_xX.$$
称这个数为 $\mathcal{V}$ 在 $x$ 处的\emph{完整亏格（holonomic defect）}。显然对任意 $x\notin\operatorname{supp}(\mathcal{V})$，$\mathcal{V}$ 的完整亏格等于 $-\dim_xX$。

\begin{corollary}\label{cor:4.6.9}
设 $Y$ 是 $X$ 的闭光滑子簇。则对任意非零凝聚右 $\mathcal{D}_Y$-模 $\mathcal{V}$ 与任意 $x\in Y$ 有
$$\operatorname{holdef}_x(\mathcal{V})=\operatorname{holdef}_x(i_+(\mathcal{V})).$$
\end{corollary}

\section{Bernstein 不等式}

设 $X$ 是光滑簇。下一个结果对 D-模理论具有根本重要性。它推广了前述结论。

\begin{theorem}[Bernstein 不等式]\label{thm:4.7.1}
设 $\mathcal{V}$ 是凝聚 $\mathcal{D}_X$-模。则对任意 $x\in\operatorname{supp}(\mathcal{V})$ 有
$$\dim_{p_X^{-1}(x)}\operatorname{Ch}(\mathcal{V})\ge\dim_xX.$$
\end{theorem}

\begin{proof}
设 $\mathcal{V}\ne0$。则可取一点 $x\in\operatorname{supp}(\mathcal{V})$ 及 $x$ 的连通开仿射邻域 $U$，使得
$p_X^{-1}(U)$ 只与 $\operatorname{Ch}(\mathcal{V})$ 中和 $p_X^{-1}(x)$ 相交的不可约分支相交。因此
$\dim U=\dim_xX$，且 $\mathcal{V}|_U$ 是非零凝聚 $\mathcal{D}_U$-模，满足
$\dim\operatorname{Ch}(\mathcal{V}|_U)=\dim_{p_X^{-1}(x)}\operatorname{Ch}(\mathcal{V})$。
显然只需对 $\mathcal{V}|_U$ 证明该陈述。此时可设 $U$ 是某个 $k^n$（$n\in\mathbb{Z}_+$）的闭光滑子簇。设
$i:U\to k^n$ 是相应的闭浸入。由前述，
$$\dim_{p_X^{-1}(x)}\operatorname{Ch}(\mathcal{V})-\dim_xX
=\dim_{p_U^{-1}(x)}\operatorname{Ch}(\mathcal{V}|_U)-\dim_xU=\operatorname{holdef}_x(\mathcal{V}|_U)$$
$$=\operatorname{holdef}_x(i_+(\mathcal{V}|_U))=\dim\operatorname{Ch}(i_+(\mathcal{V}|_U))-n.$$
而由前述可知 $\dim\operatorname{Ch}(i_+(\mathcal{V}|_U))\ge n$。
\end{proof}
\begin{annote}
Bernstein 不等式 $\dim_{\pi^{-1}(x)}\mathrm{Ch}(\mathcal V)\ge\dim_x X$ 的证明是“化到局部”：取 $x$ 的仿射邻域并令 $f$ 为到该邻域的光滑投影，用 $i_+$ 把问题搬到特征簇上，再用完整亏格 $\mathrm{holdef}$ 在 $i_+$ 下的不变性把两项维数对齐。讲义此处的 $\mathrm{holdef}$ 链条与原文一致；注意第二步的 $\dim\mathrm{Ch}(i_+(\mathcal V|_U))\ge n$ 在原书里写作“$\mathrm{Ch}(i_+(\mathcal V|_U))\ge n$”，量纲不严，译文照原文。[来源：原书 §7, p119–124；nLab 条目 holonomic+D-module（2014-02-16）]
\end{annote}



\section{闭浸入与 Kashiwara 定理}

设 $X$ 是光滑簇，$Y$ 是其闭光滑子簇。设 $i:Y\to X$ 是 $Y$ 到 $X$ 的自然包含。我们在前述中证明了直像函子
$$i_+(\mathcal{V})=i_\bullet(\mathcal{V}\otimes_{\mathcal{D}_Y}\mathcal{D}_{Y\to X})$$
是从 $\mathcal{M}_R(\mathcal{D}_Y)$ 到 $\mathcal{M}_R(\mathcal{D}_X)$ 的正合函子。

\begin{proposition}\label{prop:4.8.1}
函子 $i_+:\mathcal{M}_R(\mathcal{D}_Y)\to\mathcal{M}_R(\mathcal{D}_X)$ 有右伴随
$i^!:\mathcal{M}_R(\mathcal{D}_X)\to\mathcal{M}_R(\mathcal{D}_Y)$。
\end{proposition}

\begin{proof}
显然，由于 $i:Y\to X$ 是闭浸入，直像函子
$i_\bullet:\mathcal{M}_R(i^{-1}(\mathcal{D}_X))\to\mathcal{M}_R(\mathcal{D}_X)$ 有右伴随
$i^{-1}:\mathcal{M}_R(\mathcal{D}_X)\to\mathcal{M}_R(i^{-1}(\mathcal{D}_X))$。因此对
$\mathcal{M}_R(\mathcal{D}_Y)$ 与 $\mathcal{M}_R(\mathcal{D}_X)$ 中的 $\mathcal{V}$、$\mathcal{W}$ 有
$$\operatorname{Hom}_{\mathcal{D}_X}(i_+(\mathcal{V}),\mathcal{W})
=\operatorname{Hom}_{\mathcal{D}_X}(i_\bullet(\mathcal{V}\otimes_{\mathcal{D}_Y}\mathcal{D}_{Y\to X}),\mathcal{W})$$
$$=\operatorname{Hom}_{i^{-1}(\mathcal{D}_X)}(\mathcal{V}\otimes_{\mathcal{D}_Y}\mathcal{D}_{Y\to X},i^{-1}(\mathcal{W})).$$
现在利用张量积的性质，得
$$\operatorname{Hom}_{\mathcal{D}_X}(i_+(\mathcal{V}),\mathcal{W})
=\operatorname{Hom}_{\mathcal{D}_Y}(\mathcal{V},
\operatorname{Hom}_{i^{-1}(\mathcal{D}_X)}(\mathcal{D}_{Y\to X},i^{-1}(\mathcal{W}))).$$
因此函子
$$i^!(\mathcal{W})=\operatorname{Hom}_{i^{-1}(\mathcal{D}_X)}(\mathcal{D}_{Y\to X},i^{-1}(\mathcal{W}))$$
是 $i_+$ 的右伴随。
\end{proof}

显然，右伴随 $i^!:\mathcal{M}_R(\mathcal{D}_X)\to\mathcal{M}_R(\mathcal{D}_Y)$ 是左正合函子。此外，由于 $i_+$ 正合，$i^!$ 把内射模映为内射模。

现在想找 $i^!$ 的另一种描述。显然
$$i^!(\mathcal{W})=\operatorname{Hom}_{i^{-1}(\mathcal{D}_X)}(\mathcal{D}_{Y\to X},i^{-1}(\mathcal{W}))$$
$$=\operatorname{Hom}_{i^{-1}(\mathcal{D}_X)}(\mathcal{O}_Y\otimes_{i^{-1}(\mathcal{O}_X)}i^{-1}(\mathcal{D}_X),i^{-1}(\mathcal{W}))
=\operatorname{Hom}_{i^{-1}(\mathcal{O}_X)}(\mathcal{O}_Y,i^{-1}(\mathcal{W})).$$
设 $\mathcal{J}$ 是 $\mathcal{O}_X$ 中由沿 $Y$ 消没的函数构成的理想层。则有自然正合列
$$0\longrightarrow i^{-1}(\mathcal{J})\longrightarrow i^{-1}(\mathcal{O}_X)\longrightarrow\mathcal{O}_Y\longrightarrow0,$$
它导致正合列
$$0\longrightarrow\operatorname{Hom}_{i^{-1}(\mathcal{O}_X)}(\mathcal{O}_Y,i^{-1}(\mathcal{W}))
\longrightarrow\operatorname{Hom}_{i^{-1}(\mathcal{O}_X)}(i^{-1}(\mathcal{O}_X),i^{-1}(\mathcal{W}))$$
$$\longrightarrow\operatorname{Hom}_{i^{-1}(\mathcal{O}_X)}(i^{-1}(\mathcal{J}),i^{-1}(\mathcal{W})).$$
映射 $T\mapsto T(1)$ 把 $\operatorname{Hom}_{i^{-1}(\mathcal{O}_X)}(i^{-1}(\mathcal{O}_X),i^{-1}(\mathcal{W}))$ 与
$i^{-1}(\mathcal{W})$ 等同。在此同构下，$\operatorname{Hom}_{i^{-1}(\mathcal{O}_X)}(\mathcal{O}_Y,i^{-1}(\mathcal{W}))$ 对应于 $i^{-1}(\mathcal{W})$ 中所有被 $i^{-1}(\mathcal{J})$ 的所有元素零化的截面构成的子层 $\mathcal{W}_0$。

设 $Z$ 是 $Y$ 的闭光滑子簇，$j:Z\to Y$ 是自然包含。则由前述 $(i\circ j)_+=i_+\circ j_+$。因此由伴随函子的唯一性得出以下事实。

\begin{lemma}\label{lem:4.8.2}
有 $(i\circ j)^!\cong j^!\circ i^!$。
\end{lemma}

最后，函子 $i^!$ 保持拟凝聚性，即有如下结果。

\begin{proposition}\label{prop:4.8.3}
设 $\mathcal{V}$ 是拟凝聚右 $\mathcal{D}_X$-模。则 $i^!(\mathcal{V})$ 是拟凝聚右 $\mathcal{D}_Y$-模。
\end{proposition}

\begin{proof}
为证明这个结果，首先注意这个结果是局部的，且只需考虑 $\mathcal{O}_Y$-模结构。因此可设 $X$ 是连通仿射簇，带整体坐标系 $(f_1,\dots,f_n;D_1,\dots,D_n)$，使得 $Y$ 是 $(f_{m+1},\dots,f_n)$ 的零点集。此时集合 $Y_k$（$m\le k\le n$）是 $(f_{k+1},\dots,f_n)$ 的公共零点集，是包含 $Y$ 的 $X$ 的闭光滑子簇。由于 $Y_m=Y$ 且 $Y_m\subset Y_{m+1}$ 余维为 1，由引理 \ref{lem:4.8.2} 可把证明归结为 $Y$ 在 $X$ 中余维为 1 的情形。因此可设 $Y$ 是函数 $f=f_n$ 的零点集。于是由于此情形下 $\mathcal{J}$ 由 $f$ 生成，可考虑正合列
$$0\longrightarrow\mathcal{O}_X\xrightarrow{\ f\ }\mathcal{O}_X\longrightarrow\mathcal{O}_X/\mathcal{J}\longrightarrow0,$$
其中第一个态射是乘以 $f$。把它限制到 $Y$，得到正合列
$$0\longrightarrow i^{-1}(\mathcal{O}_X)\xrightarrow{\ f\ }i^{-1}(\mathcal{O}_X)\longrightarrow\mathcal{O}_Y\longrightarrow0.$$
因此复形
$$0\longrightarrow i^{-1}(\mathcal{O}_X)\xrightarrow{\ f\ }i^{-1}(\mathcal{O}_X)\longrightarrow0$$
是 $\mathcal{O}_Y$ 关于 $i^{-1}(\mathcal{O}_X)$-模的自由分解。把这个分解在 $i^{-1}(\mathcal{D}_X)$ 上与 $i^{-1}(\mathcal{V})$ 作张量积，得到复形
$$0\longrightarrow i^{-1}(\mathcal{V})\xrightarrow{\ f\ }i^{-1}(\mathcal{V})\longrightarrow0,$$
它表示 $Li^*(\mathcal{D}(\mathcal{V}))$。由此可知，作为 $\mathcal{O}_Y$-模有 $i^!(\mathcal{V})=L^{-1}i^*(\mathcal{V})$。由前述，这意味着对拟凝聚 $\mathcal{D}_X$-模 $\mathcal{V}$，$\mathcal{O}_Y$-模 $i^!(\mathcal{V})$ 也是拟凝聚的。
\end{proof}

现在研究伴随态射 $\operatorname{id}\to i^!\circ i_+$ 与 $i_+\circ i^!\to\operatorname{id}$。

由于 $\mathcal{D}_{Y\to X}=\mathcal{O}_Y\otimes_{i^{-1}(\mathcal{O}_X)}i^{-1}(\mathcal{D}_X)$，它由 $1\otimes1$ 决定一个典范截面。因此有由
$T\mapsto T(1\otimes1)$ 给出的典范态射 $\mathcal{D}_Y\to\mathcal{D}_{Y\to X}$。由前述 $\mathcal{D}_{Y\to X}$ 作为 $\mathcal{D}_Y$-模的局部描述可知，这个态射是单射，且有正合列
$$0\longrightarrow\mathcal{D}_Y\longrightarrow\mathcal{D}_{Y\to X}\longrightarrow Q\longrightarrow0$$
其中 $Q$ 是局部自由 $\mathcal{D}_Y$-模。此外，$\mathcal{D}_Y$ 的像包含在 $\mathcal{D}_{Y\to X}$ 中被 $i^{-1}(\mathcal{J})$ 的所有元素零化的截面构成的子层中。

设 $\mathcal{V}$ 是右 $\mathcal{D}_Y$-模。则在 $\mathcal{D}_Y$ 上与典范态射 $\mathcal{D}_Y\to\mathcal{D}_{Y\to X}$ 作张量积，得到自然单射
$$\mathcal{V}\longrightarrow\mathcal{V}\otimes_{\mathcal{D}_Y}\mathcal{D}_{Y\to X}=i^{-1}(i_+(\mathcal{V})),$$
其像包含在 $i^!(i_+(\mathcal{V}))=(i_+(\mathcal{V}))_0$ 中。因此有典范单射 $\alpha_{\mathcal{V}}:\mathcal{V}\to i^!(i_+(\mathcal{V}))$。显然这个态射正是伴随态射之一。

设 $y\in Y$。则可找到 $y$ 的连通仿射邻域，带坐标系 $(f_1,\dots,f_n;D_1,\dots,D_n)$，使得 $Y$ 是 $f_{m+1},\dots,f_n$ 的零点集。此时如前述，$\mathcal{D}_{Y\to X}|_U$ 是以 $(D^I;I\in\{0\}\times\mathbb{Z}_+^{n-m})$ 为基的自由 $\mathcal{D}_{Y\cap U}$-模。令
$$E=\sum_{j=m+1}^nf_jD_j.$$
以下结果由直接计算得出。

\begin{lemma}\label{lem:4.8.4}
\begin{enumerate}[label=(\roman*)]
  \item $[E,f_j]=f_j$，$m+1\le j\le n$；
  \item $[E,D^I]=-|I|D^I$，$I\in\{0\}\times\mathbb{Z}_+^{n-m}$。
\end{enumerate}
\end{lemma}

设 $\mathcal{V}$ 是右 $\mathcal{D}_Y$-模。则
$$i_+(\mathcal{V})_y=\mathcal{V}_y\otimes_{\mathcal{D}_{Y,y}}\mathcal{D}_{Y\to X,y}$$
且 $\mathcal{D}_{Y\to X,y}$ 是自由左 $\mathcal{D}_{Y,y}$-模。此外，$D^I$（$I\in\{0\}\times\mathbb{Z}_+^{n-m}$）的像构成自由 $\mathcal{D}_{Y,y}$-模 $\mathcal{D}_{Y\to X,y}$ 的一组基。故有
$$i_+(\mathcal{V})_y=\bigoplus_{I\in\{0\}\times\mathbb{Z}_+^{n-m}}\mathcal{V}_yD^I.$$
设 $v\in i_+(\mathcal{V})_y$。则有唯一分解 $v=\sum_Iv_ID^I$，$v_I\in\mathcal{V}_y$。定义
$$\operatorname{ord}(v)=\max\{|I|\mid v_I\ne0\}.$$
由于 $\sum_Iv_IE=\sum_jv_If_jD_j=0$，有
$$vE=\sum_Iv_ID^IE=\sum_Iv_I[D^I,E]=\sum_I|I|v_ID^I.$$
因此有如下结果。

\begin{lemma}\label{lem:4.8.5}
设 $v\in i_+(\mathcal{V})_y$。则
\begin{enumerate}[label=(\roman*)]
  \item $\operatorname{ord}(vf_j)\le\operatorname{ord}(v)-1$，$m+1\le j\le n$；
  \item $\operatorname{ord}\bigl(v(E-\operatorname{ord}(v))\bigr)\le\operatorname{ord}(v)-1$。
\end{enumerate}
\end{lemma}

\begin{proof}
\begin{enumerate}[label=(\roman*)]
  \item 设 $v=\sum_Iv_ID^I$，$v_I\in\mathcal{V}_y$。则
    $$vf_j=\sum_Iv_ID^If_j=\sum_Iv_I[D^I,f_j].$$
    若 $I=(i_1,\dots,i_j,\dots,i_n)$ 且 $I'=(i_1,\dots,i_j-1,\dots,i_n)$，可见
    $[D^I,f_j]=i_jD^{I'}$。这立即使 (i) 成立。
  \item 由上述公式立即得出。
\end{enumerate}
\end{proof}

现在可以分析第一个伴随态射。

\begin{lemma}\label{lem:4.8.6}
态射 $\alpha_{\mathcal{V}}$ 诱导 $\mathcal{V}$ 到 $i^!(i_+(\mathcal{V}))$ 的同构。
\end{lemma}

\begin{proof}
只需证明 $\alpha_{\mathcal{V}}$ 是满射，或等价地，$\alpha_{\mathcal{V},y}:\mathcal{V}_y\to i^!(i_+(\mathcal{V}))_y$ 是满射。设 $v$ 是 $i^!(i_+(\mathcal{V}))_y$ 中截面的芽。则对 $m+1\le j\le n$ 有 $vf_j=0$。故 $vE=0$，由上述计算得 $\operatorname{ord}(v)=0$，即 $v\in\mathcal{V}_y$。
\end{proof}

因此，伴随态射 $\operatorname{id}\to i^!\circ i_+$ 是函子的同构。

现在研究另一个伴随态射 $i_+\circ i^!\to\operatorname{id}$。设 $\mathcal{W}$ 是 $\mathcal{D}_X$-模，$\beta_{\mathcal{W}}:i_+(i^!(\mathcal{W}))\to\mathcal{W}$ 是相应自然态射。

\begin{lemma}\label{lem:4.8.7}
态射 $\beta_{\mathcal{W}}:i_+(i^!(\mathcal{W}))\to\mathcal{W}$ 是单射。
\end{lemma}

\begin{proof}
设 $y\in Y$。须证 $\beta_{\mathcal{W},y}:i_+(i^!(\mathcal{W}))_y\to\mathcal{W}_y$ 是单射。设 $\beta_{\mathcal{W},y}$ 不是单射。取 $\beta_{\mathcal{W},y}$ 的核中元素 $v\ne0$。若 $p=\operatorname{ord}(v)>0$，则存在 $m+1\le j\le n$ 使 $vf_j\ne0$。由引理 \ref{lem:4.8.5}(i)，$\operatorname{ord}(vf_j)\le p-1$。故由对 $\operatorname{ord}(v)$ 的递减归纳，得出 $\beta_{\mathcal{W},y}$ 的核中存在 $\operatorname{ord}(v)=0$ 的元素 $v\ne0$。但这显然不可能，因为 $i^!(\mathcal{W})_y$ 是 $\mathcal{W}_y$ 的子空间。
\end{proof}

因此可把 $i_+(i^!(\mathcal{W}))$ 看作 $\mathcal{W}$ 的子模。由于它支集在 $Y$ 中，它是 $\mathcal{W}$ 中所有支集在 $Y$ 中的局部截面构成的模 $\Gamma_Y(\mathcal{W})$ 的子模。

归纳构造一列递减的理想层 $(\mathcal{J}^p;p\in\mathbb{Z}_+)$：
$$\mathcal{J}^0=\mathcal{O}_X,\qquad \mathcal{J}^p\text{ 是 }\mathcal{J}^{p-1}\otimes_{\mathcal{O}_X}\mathcal{J}\to\mathcal{O}_X\text{ 的像},\quad p\in\mathbb{N}.$$
设 $\mathcal{V}$ 是 $\mathcal{O}_X$-模。则可定义 $\mathcal{V}[Y,p]$ 为所有被 $\mathcal{J}^p$ 的所有元素零化的截面构成的子层。于是 $\mathcal{V}[Y,p]$ 是 $\mathcal{V}$ 的 $\mathcal{O}_X$-子模。显然 $\mathcal{V}[Y,p]$ 的截面支集在 $Y$ 中，因此它们是 $\mathcal{V}$ 中所有支集在 $Y$ 中的局部截面构成的层 $\Gamma_Y(\mathcal{V})$ 的子模。显然对所有 $p\in\mathbb{Z}_+$ 有 $\mathcal{V}[Y,p]\subset\mathcal{V}[Y,p+1]$。令 $\Gamma_{[Y]}(\mathcal{V})$ 为所有子层 $\mathcal{V}[Y,p]$（$p\in\mathbb{Z}_+$）的并。则 $\Gamma_{[Y]}(\mathcal{V})$ 是 $\Gamma_Y(\mathcal{V})$ 的 $\mathcal{O}_X$-子模。

\begin{lemma}\label{lem:4.8.8}
设 $\mathcal{W}$ 是右 $\mathcal{D}_X$-模。则
$$\Gamma_{[Y]}(\mathcal{W})=i_+(i^!(\mathcal{W})).$$
\end{lemma}

\begin{proof}
设 $y\in Y$、$v\in i_+(i^!(\mathcal{W}))_y$。断言对任意 $p\in\mathbb{N}$ 有
$$\{v\in i_+(i^!(\mathcal{W}))_y\mid\operatorname{ord}(v)<p\}=\mathcal{W}[Y,p]_y.$$
显然 $\operatorname{ord}(v)=0$ 当且仅当 $v\in i^!(\mathcal{W})_y=\mathcal{W}[Y,1]_y$。这证明了 $p=1$ 的情形。设对某个 $p\ge1$ 成立。

先证明
$$\{v\in i_+(i^!(\mathcal{W}))_y\mid\operatorname{ord}(v)<p+1\}\subset\mathcal{W}[Y,p+1]_y.$$
若 $\operatorname{ord}(v)=p+1$，由引理 \ref{lem:4.8.5}(i)，对任意 $m+1\le j\le n$ 有 $\operatorname{ord}(vf_j)\le p$。故由归纳假设，对任意 $m+1\le j\le n$ 有 $vf_j\in\mathcal{W}[Y,p]_y$。由此得出 $v\in\mathcal{W}[Y,p+1]_y$。

反之，若 $v\in\mathcal{W}[Y,p+1]_y$，则对每个 $m+1\le j\le n$ 有 $vf_j\in\mathcal{W}[Y,p]_y$。由归纳假设，$vf_j\in i_+(i^!(\mathcal{W}))_y$ 且 $\operatorname{ord}(vf_j)<p$。因此
$vE=\sum_{j=m+1}^nvf_jD^j\in i_+(i^!(\mathcal{W}))_y$ 且 $\operatorname{ord}(vE)\le p$。此外，由引理 \ref{lem:4.8.5}(ii)，
$$v(E-p)f_j=vEf_j-pvf_j=v[E,f_j]+vf_jE-pvf_j=vf_j(E-(p-1))$$
是阶 $<p-1$ 的元素。故对任意 $m+1\le j\le n$ 有 $v(E-p)f_j\in\mathcal{W}[Y,p-1]_y$。因此 $v(E-p)\in\mathcal{W}[Y,p]_y$，由归纳假设 $v(E-p)\in i_+(i^!(\mathcal{W}))_y$ 且 $\operatorname{ord}(v(E-p))<p$。故 $pv=vE-v(E-p)\in i_+(i^!(\mathcal{W}))_y$ 且阶 $\le p$。由于 $p>0$，得 $v\in i_+(i^!(\mathcal{W}))_y$ 且 $\operatorname{ord}(v)<p+1$。
\end{proof}

因此 $\Gamma_{[Y]}(\mathcal{W})$ 是 $\mathcal{W}$ 的 $\mathcal{D}_X$-子模。此外，由于 $i_+$ 是正合函子而 $i^!$ 是左正合的，立即得到以下推论。

\begin{proposition}\label{prop:4.8.9}
函子 $\Gamma_{[Y]}:\mathcal{M}(\mathcal{D}_X)\to\mathcal{M}(\mathcal{D}_X)$ 是左正合的。
\end{proposition}

最后有如下结果，它给出了两个伴随态射的描述。

\begin{theorem}\label{thm:4.8.10}
设 $Y$ 是光滑簇 $X$ 的闭光滑子簇，$i:Y\to X$ 是自然包含。则
$i^!\circ i_+\cong\operatorname{id}$ 且 $i_+\circ i^!\cong\Gamma_{[Y]}$。
\end{theorem}

另一方面，有如下简单事实。

\begin{lemma}\label{lem:4.8.11}
设 $\mathcal{V}$ 是拟凝聚 $\mathcal{O}_X$-模。则 $\Gamma_{[Y]}(\mathcal{V})=\Gamma_Y(\mathcal{V})$。
\end{lemma}

\begin{proof}
可设 $X$ 是仿射簇且 $Y\ne X$。设 $v\in\Gamma(X,\mathcal{V})$ 是支集在 $Y$ 中的截面。设 $g\in R(X)$ 是在 $Y$ 上消没的非零函数，$U$ 是 $g$ 对应的主开集。则 $v|_U=0$。又 $\Gamma(U,\mathcal{V})=\Gamma(X,\mathcal{V})_g$，故对充分大的 $p\in\mathbb{Z}_+$ 有 $g^pv=0$。
\end{proof}

因此，若把最后结果应用于拟凝聚 $\mathcal{D}_X$-模，得到以下结果。

\begin{proposition}\label{prop:4.8.12}
若 $\mathcal{V}$ 是拟凝聚右 $\mathcal{D}_X$-模，则 $\Gamma_Y(\mathcal{V})=i_+(i^!(\mathcal{V}))$。
\end{proposition}

特别地，若 $\mathcal{V}$ 是支集包含在 $Y$ 中的拟凝聚右 $\mathcal{D}_X$-模，则 $i_+(i^!(\mathcal{V}))=\mathcal{V}$。这证明了以下属于 Kashiwara 的结果。

以 $\mathcal{M}^R_{qc,Y}(\mathcal{D}_X)$（相应地 $\mathcal{M}^R_{coh,Y}(\mathcal{D}_X)$，$\operatorname{Hol}_Y(\mathcal{D}_X)$）记
$\mathcal{M}^R_{qc}(\mathcal{D}_X)$（相应地 $\mathcal{M}^R_{coh}(\mathcal{D}_X)$，$\operatorname{Hol}(\mathcal{D}_X)$）中由支集在 $Y$ 中的模构成的全子范畴。

\begin{theorem}[Kashiwara 定理]\label{thm:4.8.13}
\begin{enumerate}[label=(\roman*)]
  \item 函子 $i_+:\mathcal{M}^R_{qc}(\mathcal{D}_Y)\to\mathcal{M}^R_{qc}(\mathcal{D}_X)$ 是
    $\mathcal{M}^R_{qc}(\mathcal{D}_Y)$ 与 $\mathcal{M}^R_{qc,Y}(\mathcal{D}_X)$ 之间范畴的等价。函子 $i^!$ 是 $i_+$ 的拟逆。
  \item 这些等价诱导出 $\mathcal{M}^R_{coh}(\mathcal{D}_Y)$ 与 $\mathcal{M}^R_{coh,Y}(\mathcal{D}_X)$ 之间保持完整亏格的等价。
\end{enumerate}
\end{theorem}

\begin{proof}
(i) 立即由命题 \ref{prop:4.8.3}、定理 \ref{thm:4.8.10} 与命题 \ref{prop:4.8.12} 得出。
(ii) 由前述结果得出。
\end{proof}

\section{D-模的局部上同调}

设 $Y$ 是光滑簇 $X$ 的闭光滑子簇。以 $i:Y\to X$ 记自然浸入。令 $U=X-Y$，$j:U\to X$ 是相应的开浸入。则对任意下方有界的右 $\mathcal{D}_X$-模复形 $\mathcal{V}^\cdot$，在
$D^+(\mathcal{M}_R(\mathcal{D}_X))$ 中有一个可区别三角
$$\begin{tikzcd}
& Rj_\bullet(\mathcal{V}^\cdot|_U) \arrow[dl,"+1"'] & \\
R\Gamma_Y(\mathcal{V}^\cdot) \arrow[rr] && \mathcal{V}^\cdot \arrow[lu]
\end{tikzcd}$$
且这个三角关于 $\mathcal{V}^\cdot$ 是函子性的（例如见 [3]）。

\begin{lemma}\label{lem:4.9.1}
函子 $\Gamma_Y$ 的右同调维数 $\le\dim X+1$。
\end{lemma}

\begin{proof}
由 $\mathcal{D}(\mathcal{V})$ 的上述可区别三角附带的长正合序列可知，存在从
$R^pj_\bullet(\mathcal{V})$ 到 $H^{p+1}_Y(\mathcal{V})$ 的自然满射。由于 $j_\bullet$ 的右同调维数 $\le\dim X$，结论成立。
\end{proof}

因此存在右导出函子
$R\Gamma_Y:D(\mathcal{M}_R(\mathcal{D}_X))\to D(\mathcal{M}_R(\mathcal{D}_X))$。由于
$D^+(\mathcal{M}^R_{qc}(\mathcal{D}_X))$ 等价于 $D^+(\mathcal{M}_R(\mathcal{D}_X))$ 中由具有拟凝聚上同调的复形构成的全子范畴，且内射拟凝聚 $\mathcal{D}_X$-模是松的、松层是 $\Gamma_Y$-无循环的，可见 $R\Gamma_Y$ 诱导 $D^+(\mathcal{M}^R_{qc}(\mathcal{D}_X))$ 到自身的一个正合函子，同构于
$\Gamma_Y:\mathcal{M}^R_{qc}(\mathcal{D}_X)\to\mathcal{M}^R_{qc}(\mathcal{D}_X)$ 的右导出函子。

如前，考虑伴随函子对
$i_+:\mathcal{M}^R_{qc}(\mathcal{D}_Y)\to\mathcal{M}^R_{qc}(\mathcal{D}_X)$ 与
$i^!:\mathcal{M}^R_{qc}(\mathcal{D}_X)\to\mathcal{M}^R_{qc}(\mathcal{D}_Y)$。如前所述，$i_+$ 正合而 $i^!$ 左正合。因此可考虑其导出函子
$Ri^!:D^+(\mathcal{M}^R_{qc}(\mathcal{D}_X))\to D^+(\mathcal{M}^R_{qc}(\mathcal{D}_Y))$。由前述立即得出
$$R\Gamma_Y\cong i_+\circ Ri^!$$
作为从 $D^+(\mathcal{M}^R_{qc}(\mathcal{D}_X))$ 到自身的函子。特别地，若 $\mathcal{V}$ 是拟凝聚右 $\mathcal{D}_X$-模，则对任意 $p\in\mathbb{Z}_+$ 有
$$H_Y^p(\mathcal{V})=i_+(R^pi^!(\mathcal{V})).$$
这意味着对任意 $p\in\mathbb{Z}_+$ 有
$$i^!(H^p_Y(\mathcal{V}))=i^!(i_+(R^pi^!(\mathcal{V})))=R^pi^!(\mathcal{V}),$$
这结合前述证明了以下结果。

\begin{lemma}\label{lem:4.9.2}
函子 $i^!:\mathcal{M}^R_{qc}(\mathcal{D}_X)\to\mathcal{M}^R_{qc}(\mathcal{D}_Y)$ 的右同调维数 $\le\dim X+1$。
\end{lemma}

因此存在右导出函子
$Ri^!:D(\mathcal{M}^R_{qc}(\mathcal{D}_X))\to D(\mathcal{M}^R_{qc}(\mathcal{D}_Y))$。由前述讨论得到以下结果。

\begin{theorem}\label{thm:4.9.3}
函子 $R\Gamma_Y$ 与 $i_+\circ Ri^!$ 作为从三角范畴 $D(\mathcal{M}^R_{qc}(\mathcal{D}_X))$ 到自身的正合函子是同构的。
\end{theorem}

\begin{lemma}\label{lem:4.9.4}
设 $\mathcal{V}$ 是支集在 $Y$ 中的右 $\mathcal{D}_X$-模。则 $\mathcal{V}$ 是 $\Gamma_Y$-无循环的。
\end{lemma}

\begin{proof}
如前述，由可区别三角
$$\begin{tikzcd}
& Rj_\bullet(\mathcal{D}(\mathcal{V}|_U)) \arrow[dl,"+1"'] & \\
R\Gamma_Y(\mathcal{D}(\mathcal{V})) \arrow[rr] && \mathcal{D}(\mathcal{V}) \arrow[lu]
\end{tikzcd}$$
附带的长正合序列可得，对 $p\in\mathbb{Z}_+$，$H^{p+1}_Y(\mathcal{V})$ 是
$R^pj_\bullet(\mathcal{V}|_U)$ 的商。由于 $\mathcal{V}|_U=0$，这意味着对 $p\ge1$ 有 $H^p_Y(\mathcal{V})=0$。
\end{proof}

由于对拟凝聚右 $\mathcal{D}_X$-模 $\mathcal{V}$ 有 $R^pi^!(\mathcal{V})=i^!(H_Y^p(\mathcal{V}))$（$p\in\mathbb{Z}_+$），这立即蕴含以下结果。

\begin{corollary}\label{cor:4.9.5}
设 $\mathcal{V}$ 是支集在 $Y$ 中的拟凝聚右 $\mathcal{D}_X$-模。则 $\mathcal{V}$ 是 $i^!$-无循环的。
\end{corollary}

如前，以 $\mathcal{M}^R_{qc,Y}(\mathcal{D}_X)$ 记 $\mathcal{M}^R_{qc}(\mathcal{D}_X)$ 中由支集包含在 $Y$ 中的模构成的全子范畴。这个范畴是 $\mathcal{M}^R_{qc}(\mathcal{D}_X)$ 的厚子范畴。以
$D_Y^*(\mathcal{M}^R_{qc}(\mathcal{D}_X))$ 记 $D^*(\mathcal{M}^R_{qc}(\mathcal{D}_X))$ 中由满足所有 $p\in\mathbb{Z}$ 时 $H^p(\mathcal{V}^\cdot)$ 支集在 $Y$ 中的复形 $\mathcal{V}^\cdot$ 构成的全子范畴。由前述，这是 $D^*(\mathcal{M}^R_{qc}(\mathcal{D}_X))$ 的三角子范畴。显然对任意拟凝聚右 $\mathcal{D}_X$-模复形 $\mathcal{V}^\cdot$，$R\Gamma_Y(\mathcal{V}^\cdot)$ 在
$D_Y(\mathcal{M}^R_{qc}(\mathcal{D}_X))$ 中。

现考虑 $R\Gamma_Y$ 到 $D(\mathcal{M}^R_{qc}(\mathcal{D}_X))$ 上恒等函子的自然变换。它在 $D_Y(\mathcal{M}^R_{qc}(\mathcal{D}_X))$ 上的限制诱导
$R\Gamma_Y:D_Y(\mathcal{M}^R_{qc}(\mathcal{D}_X))\to D_Y(\mathcal{M}^R_{qc}(\mathcal{D}_X))$ 到恒等函子的自然变换 $\alpha$。

\begin{proposition}\label{prop:4.9.6}
自然变换 $\alpha:R\Gamma_Y:D_Y(\mathcal{M}^R_{qc}(\mathcal{D}_X))\to D_Y(\mathcal{M}^R_{qc}(\mathcal{D}_X))$ 到恒等函子是函子的同构。
\end{proposition}

\begin{proof}
先设 $\mathcal{V}^\cdot\in D_Y(\mathcal{M}^R_{qc}(\mathcal{D}_X))$ 是下方有界复形。则由可区别三角
$$\begin{tikzcd}
& Rj_\bullet(\mathcal{V}^\cdot|_U) \arrow[dl,"+1"'] & \\
R\Gamma_Y(\mathcal{V}^\cdot) \arrow[rr] && \mathcal{V}^\cdot \arrow[lu]
\end{tikzcd}$$
可见该陈述等价于 $Rj_\bullet(\mathcal{V}^\cdot|_U)=0$。但这是显然的，因为对所有 $p\in\mathbb{Z}$ 有
$H^p(\mathcal{V}^\cdot|_U)=H^p(\mathcal{V}^\cdot)|_U=0$。一般情形由前述得出。
\end{proof}

这个结果结合 Kashiwara 定理，有以下直接推论。

\begin{theorem}\label{thm:4.9.7}
函子 $i_+:D(\mathcal{M}^R_{qc}(\mathcal{D}_Y))\to D_Y(\mathcal{M}^R_{qc}(\mathcal{D}_X))$ 是三角范畴的等价。拟逆是函子
$Ri^!:D_Y(\mathcal{M}^R_{qc}(\mathcal{D}_X))\to D(\mathcal{M}^R_{qc}(\mathcal{D}_Y))$。
\end{theorem}

\begin{proof}
由定理 \ref{thm:4.9.3} 与命题 \ref{prop:4.9.6} 可知 $i_+\circ Ri^!$ 同构于 $D_Y(\mathcal{M}^R_{qc}(\mathcal{D}_X))$ 上的恒等函子。另一方面，由于
$i_+(\mathcal{W})$（$\mathcal{W}\in\mathcal{M}^R_{qc}(\mathcal{D}_Y)$）形式的模由推论 \ref{cor:4.9.5} 是 $i^!$-无循环的，也得出 $Ri^!\circ i_+$ 同构于恒等函子。
\end{proof}

特别地，$D_Y(\mathcal{M}^R_{qc}(\mathcal{D}_X))$ 中每个复形 $\mathcal{V}^\cdot$ 都同构于由支集在 $Y$ 中的 $\mathcal{D}_X$-模构成的 $i_+(Ri^!(\mathcal{V}^\cdot))$。

现考虑 $Y$ 的闭光滑子簇 $Z$ 与自然浸入 $i_1:Z\to Y$。则有以下结果。

\begin{theorem}\label{thm:4.9.8}
$R(i\circ i_1)^!\cong Ri_1^!\circ Ri^!$ 作为从 $D(\mathcal{M}^R_{qc}(\mathcal{D}_X))$ 到
$D(\mathcal{M}^R_{qc}(\mathcal{D}_Z))$ 的函子。
\end{theorem}

\begin{proof}
由前述知 $(i\circ i_1)^!=i_1^!\circ i^!$。由于 $i^!$ 是正合函子 $i_+$ 的右伴随，它把内射拟凝聚右
$\mathcal{D}_X$-模映为内射拟凝聚右 $\mathcal{D}_Y$-模。
\end{proof}

现在想证明对 $\Gamma_Y:\mathcal{M}^R_{qc}(\mathcal{D}_X)\to\mathcal{M}^R_{qc}(\mathcal{D}_X)$ 与
$i^!:\mathcal{M}^R_{qc}(\mathcal{D}_X)\to\mathcal{M}^R_{qc}(\mathcal{D}_Y)$ 的右同调维数的更精确估计。

\begin{theorem}\label{thm:4.9.9}
\begin{enumerate}[label=(\roman*)]
  \item 函子 $\Gamma_Y:\mathcal{M}^R_{qc}(\mathcal{D}_X)\to\mathcal{M}^R_{qc}(\mathcal{D}_X)$ 的右同调维数 $\le\dim X-\dim Y$；
  \item 函子 $i^!:\mathcal{M}^R_{qc}(\mathcal{D}_X)\to\mathcal{M}^R_{qc}(\mathcal{D}_Y)$ 的右同调维数 $\le\dim X-\dim Y$。
\end{enumerate}
\end{theorem}

\begin{proof}
由于对任意拟凝聚右 $\mathcal{D}_X$-模 $\mathcal{V}$ 与 $p\in\mathbb{Z}_+$ 有
$R^pi^!(\mathcal{V})=i^!(H_Y^p(\mathcal{V}))$，这些结果等价。

此外，由于层的内射性是局部性质，第一个断言显然是局部的。因此为证明它，可设 $X$ 是仿射的，带坐标系 $(f_1,\dots,f_n;D_1,\dots,D_n)$，且 $Y$ 是 $f_{m+1},\dots,f_n$ 的公共零点集。

对 $n-m$ 作归纳。先考虑 $m=n-1$ 的情形。此时
$U=\{x\in X\mid f_n(x)\ne0\}$，即 $X$ 中主开集。故 $U$ 仿射且 $j:U\to X$ 是仿射态射。这意味着对任意拟凝聚右 $\mathcal{D}_X$-模 $\mathcal{V}$ 与 $p\ge1$ 有
$R^pj_\bullet(\mathcal{V}|_U)=0$。由于对 $p\in\mathbb{Z}_+$，$H^{p+1}_Y(\mathcal{V})$ 是
$R^pj_\bullet(\mathcal{V}|_U)$ 的商，可见对 $p\ge2$ 有 $H_Y^p(\mathcal{V})=0$。因此 (i)、(ii) 在此情形成立。

现设 $m<n-1$。则可考虑光滑子簇 $Z=\{x\in X\mid f_n(x)=0\}$。设 $i_1:Y\to Z$、$i_2:Z\to X$ 是典范浸入。由归纳假设，$i_1^!$、$i_2^!$ 具有性质 (ii)，故由定理 \ref{thm:4.9.8}，它们的复合 $i$ 也具有性质 (ii)。
\end{proof}

最后指出基变换的一个特殊情形。

\begin{lemma}\label{lem:4.9.10}
设 $Y$ 是 $X$ 的闭光滑子簇，$U=X-Y$，$i:Y\to X$、$j:U\to X$ 是自然浸入。则在
$D(\mathcal{M}_{qc}(\mathcal{D}_U))$ 上 $Ri^!\circ Rj_\bullet=0$。
\end{lemma}

\begin{proof}
设 $\mathcal{W}^\cdot$ 是下方有界的拟凝聚 $\mathcal{D}_U$-模复形。则存在 $j_\bullet$-无循环复形 $\mathcal{J}^\cdot$ 与拟同构
$\mathcal{W}^\cdot\to\mathcal{J}^\cdot$。故有交换图
$$\begin{tikzcd}
Rj_\bullet(\mathcal{W}^\cdot) \arrow[r] \arrow[d,"\sim"'] &
Rj_\bullet(Rj_\bullet(\mathcal{W}^\cdot)|_U) \arrow[d,"\sim"]\\
j_\bullet(\mathcal{J}^\cdot) \arrow[r,"="'] & j_\bullet(j_\bullet(\mathcal{J}^\cdot)|_U)
\end{tikzcd}$$
其中竖线是同构，下方水平箭头是恒等。这说明上方水平箭头也是同构。由可区别三角
$$\begin{tikzcd}
& Rj_\bullet(Rj_\bullet(\mathcal{W}^\cdot)|_U) \arrow[dl,"+1"'] & \\
R\Gamma_Y(Rj_\bullet(\mathcal{W}^\cdot)) \arrow[rr] && Rj_\bullet(\mathcal{W}^\cdot) \arrow[lu]
\end{tikzcd}$$
得出 $R\Gamma_Y(Rj_\bullet(\mathcal{W}^\cdot))=0$。于是由前述，在
$D^+(\mathcal{M}^R_{qc}(\mathcal{D}_U))$ 上 $i_+\circ Ri^!\circ Rj_\bullet=0$。再由定理 \ref{thm:4.9.7}，
$$0=Ri^!\circ i_+\circ Ri^!\circ Rj_\bullet\cong Ri^!\circ Rj_\bullet$$
在 $D^+(\mathcal{M}^R_{qc}(\mathcal{D}_U))$ 上。由标准截断点论证得出在 $D(\mathcal{M}^R_{qc}(\mathcal{D}_U))$ 上同样成立。
\end{proof}
\begin{annote}
$R\Gamma_Z$ 是“支集落在 $Z$”的截面的导出函子；在 D-模语境里它是 $R\mathcal H om_{\mathcal D_X}(\mathcal D_{X,\infty Z},-)$ 的另一种写法，$Z$ 闭、$\mathcal D_{X,\infty Z}$ 为 $Z$ 上无穷阶极点的算子层。最容易混淆的是它与 $\mathcal D_X$-模的 $\mathcal H^0_Z$（代数局部上同调）之间的换位：习题式的等式 $R\Gamma_Z(\mathcal M)\cong R\mathcal H om(\mathcal O_{X,\infty Z},\mathcal M)$ 只在 $\mathcal M$ 拟凝聚时成立。Stacks 的对应条目是 dualizing complex 一章里 torsion 与 complete 的对偶关系，可作形式对照。[来源：Stacks Project Tag 0A6Y（Lemma 47.12.3, torsion vs. complete）；原书 §9, p127–130]
\end{annote}



\section{基变换}

设 $X$、$Y$、$S$ 是三个代数簇，$\varphi:X\to S$、$\psi:Y\to S$ 是簇的两个态射。则可考虑如下代数簇的交换图
$$\begin{tikzcd}
X\times Y \arrow[r,"p_1"] \arrow[d,"p_2"' ] & X \arrow[d,"\varphi"]\\
Y \arrow[r,"\psi"'] & S
\end{tikzcd}$$
其中 $p_1$、$p_2$ 分别是到第一、第二因子的投影。设 $\Delta_S$ 是 $S\times S$ 中的对角线。由于 $S$ 是簇，$\Delta_S$ 是 $S\times S$ 的闭子簇。令
$$X\times_SY=\{(x,y)\in X\times Y\mid\varphi(x)=\psi(y)\}$$
$$=\{(x,y)\in X\times Y\mid(\varphi\circ p_1)(x,y)=(\psi\circ p_2)(x,y)\}
=((\varphi\circ p_1)\times(\psi\circ p_2))^{-1}(\Delta_S),$$
则这是 $X\times Y$ 的闭子簇。称 $X\times_SY$ 为 $X$ 与 $Y$ 在 $S$ 上的\emph{纤维积}。投影 $p_1$、$p_2$ 诱导纤维积 $X\times_SY$ 到 $X$、$Y$ 的态射，使得交换图
$$\begin{tikzcd}
X\times_SY \arrow[r,"\psi'"] \arrow[d,"\varphi'"' ] & X \arrow[d,"\varphi"]\\
Y \arrow[r,"\psi"'] & S
\end{tikzcd}$$
交换。纤维积具有以下泛性质。设 $T$ 是代数簇，$\alpha:T\to X$、$\beta:T\to Y$ 是使得交换图
$$\begin{tikzcd}
T \arrow[r,"\alpha"] \arrow[d,"\beta"' ] & X \arrow[d,"\varphi"]\\
Y \arrow[r,"\psi"'] & S
\end{tikzcd}$$
交换的簇的态射。则由积的泛性质得到一个态射 $\gamma:T\to X\times Y$ 使得 $\alpha=p_1\circ\gamma$、$\beta=p_2\circ\gamma$。因此 $\gamma(T)\subset X\times_SY$，$\gamma$ 诱导态射 $\delta:T\to X\times_SY$ 使得 $\alpha=\psi'\circ\delta$、$\beta=\varphi'\circ\delta$。

设 $Z$ 是另一个代数簇，$\chi:Z\to Y$。则有如下交换图
$$\begin{tikzcd}
(X\times_SY)\times_YZ \arrow[r,"\chi'"] \arrow[d,"\varphi''"'] &
X\times_SY \arrow[r,"\psi'"] \arrow[d,"\varphi'"'] & X \arrow[d,"\varphi"]\\
Z \arrow[r,"\chi"'] & Y \arrow[r,"\psi"'] & S
\end{tikzcd}$$
由纤维积的泛性质，得到态射 $\omega:(X\times_SY)\times_YZ\to X\times_SZ$。

\begin{lemma}\label{lem:4.10.1}
态射 $\omega:(X\times_SY)\times_YZ\to X\times_SZ$ 是同构。
\end{lemma}

\begin{proof}
由上述讨论，
$$(X\times_SY)\times_YZ=\{(x,y,z)\in X\times Y\times Z\mid\varphi'(x,y)=\chi(z),\ (x,y)\in X\times_SY\}$$
$$=\{(x,y,z)\in X\times Y\times Z\mid y=\chi(z),\ \varphi(x)=\psi(y)\}$$
且 $\omega(x,y,z)=(x,z)$。因此 $\omega$ 的像等于
$$X\times_SZ=\{(x,z)\in X\times Z\mid\varphi(x)=\psi(\chi(z))\},$$
其逆映射由 $(p_1,\chi\circ p_2,p_2)$ 在 $X\times_SY$ 上的限制给出。
\end{proof}

现考虑两种特殊情形。设 $Y$ 是 $S$ 的闭子簇，$i:Y\to S$ 是自然浸入。则
$$X\times_SY=\{(x,y)\in X\times Y\mid\varphi(x)=y\}\subset X\times S$$
等于 $\varphi$ 的图 $\Gamma_\varphi$ 与 $X\times Y$ 的交。因此 $\psi':X\times_SY\to X$ 是到 $X$ 的闭子簇 $\varphi^{-1}(Y)$ 上的同构。故纤维积图如下
$$\begin{tikzcd}
\varphi^{-1}(Y) \arrow[r,"i'"] \arrow[d,"\varphi'"' ] & X \arrow[d,"\varphi"]\\
Y \arrow[r,"i"'] & S
\end{tikzcd}$$
其中 $\varphi'$ 是 $\varphi$ 的限制，$i'$ 是自然包含。

设 $Y=Z\times S$，$\psi:Y\to S$ 是到第二个变量的投影。则可把 $Y$ 看作 $\alpha:Z\to\{pt\}$ 与 $\beta:S\to\{pt\}$ 的纤维积。故由引理 \ref{lem:4.10.1}，
$$(Z\times S)\times_SX=Z\times X,$$
且纤维积图如下
$$\begin{tikzcd}
Z\times X \arrow[r,"p_2"] \arrow[d,"\operatorname{id}\times\varphi"'] & X \arrow[d,"\varphi"]\\
Z\times S \arrow[r,"\psi"'] & S
\end{tikzcd}$$

若 $\psi:Y\to S$ 是任意态射，可用图分解 $\psi=\alpha\circ\beta$，其中 $\beta$ 是 $Y$ 到图 $\Gamma_\varphi\subset Y\times S$ 的同构，$\alpha:Y\times S\to S$ 是到第二因子的投影。故由引理 \ref{lem:4.10.1}，每个纤维积都可看作分两步得到，每步都是上述特殊类型之一。

本节的主要结果如下定理。

\begin{theorem}\label{thm:4.10.2}
设 $X$、$Y$、$S$ 是光滑代数簇，$\varphi:X\to S$、$\psi:Y\to S$ 是代数簇的态射，且纤维积 $X\times_SY$ 是光滑代数簇。则交换图
$$\begin{tikzcd}
X\times_SY \arrow[r,"\psi'"] \arrow[d,"\varphi'"' ] & X \arrow[d,"\varphi"]\\
Y \arrow[r,"\psi"'] & S
\end{tikzcd}$$
决定了从 $D(\mathcal{M}_{qc}(\mathcal{D}_X))$ 到 $D(\mathcal{M}_{qc}(\mathcal{D}_Y))$ 的函子的同构
$$\psi^!\circ\varphi_+\cong\varphi'_+\circ\psi'^!.$$
\end{theorem}

\begin{proof}
如前述，纤维积的构造总可分两步：第一步一个态射是投影，第二步是闭浸入。在积的情形下，纤维积的光滑性是自动的；在闭浸入的情形下，这是一个附加条件。

先考虑后一种情形。设 $i:Y\to S$ 是闭浸入，且 $\varphi^{-1}(Y)$ 是 $X$ 的光滑闭子簇。设 $i'$ 是 $\varphi^{-1}(Y)$ 到 $X$ 的自然浸入。令 $U=S-Y$，$j:U\to S$ 是自然浸入。又有 $X-\varphi^{-1}(Y)=\varphi^{-1}(U)$。则有如下交换图
$$\begin{tikzcd}
\varphi^{-1}(Y) \arrow[r,"i'"] & X & \arrow[l,"j'"'] \varphi^{-1}(U) \arrow[d,"\varphi''"]\\
Y \arrow[r,"i"'] \arrow[u,"\varphi'"'] & S & \arrow[l,"j"'] U
\end{tikzcd}$$
其中 $\varphi'$、$\varphi''$ 是 $\varphi$ 的相应限制。

先设 $\mathcal{V}^\cdot$ 是下方有界的拟凝聚 $\mathcal{D}_X$-模复形。考虑 $D(\mathcal{M}_{qc}(\mathcal{D}_X))$ 中的可区别三角
$$\begin{tikzcd}
& Rj'_\bullet(\mathcal{V}^\cdot|_{\varphi^{-1}(U)}) \arrow[dl,"+1"'] & \\
R\Gamma_{\varphi^{-1}(Y)}(\mathcal{V}^\cdot) \arrow[rr] && \mathcal{V}^\cdot \arrow[lu]
\end{tikzcd}$$
应用导出函子 $\varphi_+$，得到 $D(\mathcal{M}_{qc}(\mathcal{D}_S))$ 中的可区别三角
$$\begin{tikzcd}
& \varphi_+(Rj'_\bullet(\mathcal{V}^\cdot|_{\varphi^{-1}(U)})) \arrow[dl,"+1"'] & \\
\varphi_+(R\Gamma_{\varphi^{-1}(Y)}(\mathcal{V}^\cdot)) \arrow[rr] && \varphi_+(\mathcal{V}^\cdot) \arrow[lu]
\end{tikzcd}$$
再应用导出函子 $Ri^!$，得到 $D(\mathcal{M}_{qc}(\mathcal{D}_Y))$ 中的可区别三角
$$\begin{tikzcd}
& Ri^!(\varphi_+(Rj'_\bullet(\mathcal{V}^\cdot|_{\varphi^{-1}(U)}))) \arrow[dl,"+1"'] & \\
Ri^!(\varphi_+(R\Gamma_{\varphi^{-1}(Y)}(\mathcal{V}^\cdot))) \arrow[rr] && Ri^!(\varphi_+(\mathcal{V}^\cdot)) \arrow[lu]
\end{tikzcd}$$
由于
$$\varphi_+(Rj'_\bullet(\mathcal{V}^\cdot|_{\varphi^{-1}(U)}))
=(\varphi\circ j')_+(\mathcal{V}^\cdot|_{\varphi^{-1}(U)})
=(j\circ\varphi'')_+(\mathcal{V}^\cdot|_{\varphi^{-1}(U)})$$
$$=Rj_\bullet(\varphi''_+(\mathcal{V}^\cdot|_{\varphi^{-1}(U)}))=Rj_\bullet(\varphi_+(\mathcal{V}^\cdot)|_U),$$
可见
$$Ri^!(\varphi_+(Rj'_\bullet(\mathcal{V}^\cdot|_{\varphi^{-1}(U)})))=Ri^!(Rj_\bullet(\varphi_+(\mathcal{V}^\cdot)|_U)).$$
由引理 \ref{lem:4.9.10}，这等于 0。因此自然态射
$$Ri^!(\varphi_+(R\Gamma_{\varphi^{-1}(Y)}(\mathcal{V}^\cdot)))\longrightarrow Ri^!(\varphi_+(\mathcal{V}^\cdot))$$
是关于 $\mathcal{V}^\cdot$ 函子性的同构。由于所涉函子都具有有限同调维数，由截断点论证可知这个自然态射对任意拟凝聚 $\mathcal{D}_X$-模复形 $\mathcal{V}^\cdot$ 都是同构。

另一方面，由前述有
$$Ri^!(\varphi_+(R\Gamma_{\varphi^{-1}(Y)}(\mathcal{V}^\cdot)))
=Ri^!(\varphi_+(i'_+(Ri'^!(\mathcal{V}^\cdot))))
=Ri^!((\varphi\circ i')_+(Ri'^!(\mathcal{V}^\cdot)))$$
$$=Ri^!((i\circ\varphi')_+(Ri'^!(\mathcal{V}^\cdot)))
=Ri^!(i_+(\varphi'_+(Ri'^!(\mathcal{V}^\cdot))))=\varphi'_+(Ri'^!(\mathcal{V}^\cdot)).$$
因此在这个特殊情形下结论得证。

现考虑投影的情形。若 $p:Z\times S\to S$ 是到第二因子的投影，如前得到纤维积交换图
$$\begin{tikzcd}
Z\times X \arrow[r,"p_2"] \arrow[d,"\operatorname{id}_Z\times\varphi"'] & X \arrow[d,"\varphi"]\\
Z\times S \arrow[r,"p"'] & S
\end{tikzcd}$$
于是由前述有
$$(p^!\circ\varphi_+)(\mathcal{V}^\cdot)=Lp_+(\varphi_+(\mathcal{V}^\cdot))[-\dim S]
=\mathcal{O}_Z\boxtimes\varphi_+(\mathcal{V}^\cdot)[-\dim S]$$
$$=(\operatorname{id}_Z\times\varphi)_+\bigl((\mathcal{O}_Z\boxtimes\mathcal{V}^\cdot)[-\dim S]\bigr)
=(\operatorname{id}_Z\times\varphi)_+\bigl(p_2^+(\mathcal{V}^\cdot)[-\dim S]\bigr)$$
$$=(\operatorname{id}_Z\times\varphi)_+(p_2^!(\mathcal{V}^\cdot))
=((\operatorname{id}_Z\times\varphi)_+\circ p_2^!)(\mathcal{V}^\cdot),$$
这在此情形下建立了基变换。
\end{proof}

这使我们能够推广引理 \ref{lem:4.9.10}。

\begin{corollary}\label{cor:4.10.3}
设 $X$、$Y$、$Z$ 是光滑代数簇，$\varphi:X\to Z$、$\psi:Y\to Z$ 是代数簇的态射，使得
$\varphi(X)\cap\psi(Y)=\varnothing$。则在 $D(\mathcal{M}_{qc}(\mathcal{D}_X))$ 上 $\psi^!\circ\varphi_+=0$。
\end{corollary}

\chapter{完整 D-模}

设 $X$ 是光滑代数簇。以 $T^*(X)$ 记其余切丛，$p_X:T^*(X)\to X$ 记相应的自然投影。对任意凝聚
$\mathcal{D}_X$-模 $\mathcal{V}$，其特征簇 $\operatorname{Ch}(\mathcal{V})$ 是 $T^*(X)$ 的闭子簇。在映射 $p_X$ 下，特征簇
$\operatorname{Ch}(\mathcal{V})$ 投射到 $\mathcal{V}$ 的支集 $\operatorname{supp}(\mathcal{V})$ 上。我们已证明：对支集
$\operatorname{supp}(\mathcal{V})$ 中任意一点 $x$，有
$\dim_{p_X^{-1}(x)}\operatorname{Ch}(\mathcal{V})\ge\dim_xX$。因此，称凝聚 $\mathcal{D}_X$-模是\emph{完整（holonomic）的}，若对任意 $x\in\operatorname{supp}(\mathcal{V})$ 有
$\dim_{p_X^{-1}(x)}\operatorname{Ch}(\mathcal{V})=\dim_xX$。这推广了前述定义。

以 $\operatorname{Hol}(\mathcal{D}_X)$ 记 $\mathcal{M}_{coh}(\mathcal{D}_X)$ 中由所有完整 $\mathcal{D}_X$-模构成的全子范畴。

\section{完整 D-模}

\begin{theorem}\label{thm:5.1.1}
设 $X$ 是光滑簇。设
$$0\longrightarrow\mathcal{V}_1\longrightarrow\mathcal{V}_2\longrightarrow\mathcal{V}_3\longrightarrow0$$
是凝聚 $\mathcal{D}_X$-模的短正合列。则：
\begin{enumerate}[label=(\roman*)]
  \item 若 $\mathcal{V}_2$ 是完整模，则 $\mathcal{V}_1$ 与 $\mathcal{V}_3$ 也是完整的；
  \item 若 $\mathcal{V}_1$ 与 $\mathcal{V}_3$ 是完整的，则 $\mathcal{V}_2$ 是完整的。
\end{enumerate}
\end{theorem}

\begin{proof}
若 $\mathcal{V}_1$ 或 $\mathcal{V}_3$ 为零，断言显然成立。若它们非零，由前述可见
$\operatorname{Ch}(\mathcal{V}_2)=\operatorname{Ch}(\mathcal{V}_1)\cup\operatorname{Ch}(\mathcal{V}_3)$。这立即蕴含
$$\dim_{p_X^{-1}(x)}\operatorname{Ch}(\mathcal{V}_2)
=\max\bigl(\dim_{p_X^{-1}(x)}\operatorname{Ch}(\mathcal{V}_1),
\dim_{p_X^{-1}(x)}\operatorname{Ch}(\mathcal{V}_3)\bigr).$$
此外，由前述知 $\operatorname{supp}(\mathcal{V}_2)=\operatorname{supp}(\mathcal{V}_1)\cup\operatorname{supp}(\mathcal{V}_3)$。因此断言立即成立。
\end{proof}

因此，$\mathcal{M}_{coh}(\mathcal{D}_X)$ 的全子范畴 $\operatorname{Hol}(\mathcal{D}_X)$ 是 Abel 的且是厚的。

\begin{lemma}\label{lem:5.1.2}
设 $X$ 是光滑簇，$Y$ 是其闭光滑子簇。设 $i:Y\to X$ 是自然包含。则对任意 $\mathcal{D}_Y$-模 $\mathcal{V}$，以下陈述等价：
\begin{enumerate}[label=(\roman*)]
  \item 模 $\mathcal{V}$ 是完整的；
  \item 模 $i_+(\mathcal{V})$ 是完整的。
\end{enumerate}
\end{lemma}

\begin{proof}
显然，由 Kashiwara 定理，$\mathcal{V}$ 为零当且仅当 $i_+(\mathcal{V})$ 为零。此外，由前述，对任意 $x\in\operatorname{supp}(\mathcal{V})$ 有
$$\operatorname{holdef}_x(\mathcal{V})=\operatorname{holdef}_x(i_+(\mathcal{V})).$$
故对任意 $x\in\operatorname{supp}(\mathcal{V})$，$\dim_{p_Y^{-1}(x)}\operatorname{Ch}(\mathcal{V})=\dim_xY$ 当且仅当
$\dim_{p_X^{-1}(x)}\operatorname{Ch}(i_+(\mathcal{V}))=\dim_xX$。
\end{proof}

下一个结果推广了前述结论。

\begin{theorem}\label{thm:5.1.3}
每个完整 $\mathcal{D}_X$-模都是有限长度的。
\end{theorem}

\begin{proof}
显然，完整模在开集上的限制仍是完整的。由于由 III.2.8，为有限长度是局部性质，只需对光滑仿射簇证明该断言。此时可设 $X$ 是仿射空间 $k^n$ 的光滑闭子簇。设 $i:X\to k^n$ 是自然包含。则由引理 \ref{lem:5.1.2}，对完整模 $\mathcal{V}$，模 $i_+(\mathcal{V})$ 是完整的。由 I.8.1 可知 $i_+(\mathcal{V})$ 是有限长度的。显然 $i_+(\mathcal{V})$ 的任意子模或商模的支集都在 $X$ 中。故 $i_+(\mathcal{V})$ 作为所有 D-模范畴中对象的长度，等于它作为所有支集在 $X$ 中的模构成的子范畴中对象的长度。由 Kashiwara 定理，这意味着 $\mathcal{V}$ 是有限长度的。
\end{proof}

\section{联络（Connections）}

设 $X$ 是特征零的代数闭域 $k$ 上的光滑代数簇，$\mathcal{D}_X$ 是 $X$ 上微分算子层。本节描述一些非常简单的完整 $\mathcal{D}_X$-模。

若 $\mathcal{D}_X$-模 $\mathcal{V}$ 作为 $\mathcal{O}_X$-模是凝聚的，则称之为\emph{联络}。

\begin{theorem}\label{thm:5.2.1}
设 $X$ 是连通光滑代数簇，$\mathcal{V}$ 是非零凝聚 $\mathcal{D}_X$-模。则以下条件等价：
\begin{enumerate}[label=(\roman*)]
  \item $\mathcal{V}$ 是联络；
  \item $\mathcal{V}$ 的特征簇 $\operatorname{Ch}(\mathcal{V})$ 包含在余切丛 $T^*(X)$ 的零截面中；
  \item $\mathcal{V}$ 的特征簇 $\operatorname{Ch}(\mathcal{V})$ 就是余切丛 $T^*(X)$ 的零截面；
  \item $\mathcal{V}$ 是有限秩局部自由 $\mathcal{O}_X$-模。
\end{enumerate}
\end{theorem}

\begin{proof}
显然 (iv)$\Rightarrow$(i) 与 (iii)$\Rightarrow$(ii)。

\textbf{(iv)$\Rightarrow$(iii)} 由于 $\mathcal{V}$ 是有限秩局部自由 $\mathcal{O}_X$-模，几何纤维
$T_x(\mathcal{V})=\mathcal{V}_x/\mathfrak{m}_x\mathcal{V}_x$ 的维数是局部常值。因为 $X$ 连通，这意味着它是非零常值，从而
$\operatorname{supp}(\mathcal{V})=X$。可在 $\mathcal{V}$ 上定义滤过 $F\mathcal{V}$：当 $p<0$ 时 $F_p\mathcal{V}=0$，当 $p\ge0$ 时 $F_p\mathcal{V}=\mathcal{V}$。这显然是 $\mathcal{V}$ 上的好滤过。分次模 $\operatorname{Gr}\mathcal{V}$ 除
$\operatorname{Gr}_0\mathcal{V}=\mathcal{V}$ 外所有齐次分量都为 0。因此 $\operatorname{Gr}\mathcal{V}$ 的零化子包含
$\bigoplus_{p=1}^{\infty}\operatorname{Gr}_p\mathcal{D}_X$，从而 $\mathcal{V}$ 的特征簇包含在 $T^*(X)$ 的零截面中。由于由 III.3.9 它投射到 $\operatorname{supp}(\mathcal{V})$ 上，得出 $\operatorname{Ch}(\mathcal{V})$ 等于 $T^*(X)$ 的零截面。

\textbf{(ii)$\Rightarrow$(i)} 该陈述是局部的，故由 II.2.10，可设 $X$ 是仿射的且有整体坐标系
$(f_1,f_2,\dots,f_n;D_1,D_2,\dots,D_n)$，使得 $(D^I;I\in\mathbb{Z}_+^n)$ 构成 $X$ 上微分算子自由 $R(X)$-模的一组基。设 $F\mathcal{V}$ 是 $\mathcal{V}$ 的好滤过，$J$ 是 $\Gamma(X,\operatorname{Gr}\mathcal{V})$ 在 $R(T^*(X))$ 中的零化子。则 $J$ 在 $T^*(X)$ 中的零点集包含在 $T^*(X)$ 的零截面中。由 Hilbert 零点定理，$J$ 的根包含由 $D_1,D_2,\dots,D_n$ 的主符号生成的理想。这意味着存在 $m\in\mathbb{Z}_+$ 使得 $D_1^m,D_2^m,\dots,D_n^m$ 的主符号零化 $\Gamma(X,\operatorname{Gr}\mathcal{V})$。此外，任意阶 $q\ge nm$ 的微分算子的 $q$ 次主符号零化 $\Gamma(X,\operatorname{Gr}\mathcal{V})$。由于 $\Gamma(X,\operatorname{Gr}\mathcal{V})$ 是有限生成 $\operatorname{Gr}\mathcal{D}_X$-模，这意味着对充分大的 $p\in\mathbb{Z}_+$ 有 $\operatorname{Gr}_p\mathcal{V}=0$。因此对充分大的 $p\in\mathbb{Z}_+$ 有 $F_p\mathcal{V}=\mathcal{V}$，即 $\mathcal{V}$ 是凝聚 $\mathcal{O}_X$-模。

\textbf{(i)$\Rightarrow$(iv)} 由于该陈述是局部的，可设 $X$ 是仿射且“小”的。设 $x\in X$。由于 $\mathcal{V}$ 是凝聚 $\mathcal{O}_X$-模，几何纤维
$T_x(\mathcal{V})=\mathcal{V}_x/\mathfrak{m}_x\mathcal{V}_x=\mathcal{O}_{X,x}/\mathfrak{m}_x\otimes_{\mathcal{O}_{X,x}}\mathcal{V}_x$
是 $k$ 上有限维向量空间。设 $s_1,s_2,\dots,s_q$ 是 $\mathcal{V}$ 的整体截面族，其像 $s_1(x),s_2(x),\dots,s_q(x)$ 在 $T_x(\mathcal{V})$ 中构成该向量空间的一组基。这些截面定义了自由 $\mathcal{O}_X$-模 $\mathcal{O}_X^q$ 到 $\mathcal{V}$ 的自然态射。记其像为 $\mathcal{U}$。则有凝聚 $\mathcal{O}_X$-模的正合列
$$\mathcal{O}_X^q\longrightarrow\mathcal{V}\longrightarrow\mathcal{V}/\mathcal{U}\longrightarrow0,$$
它导致正合列
$$T_x(\mathcal{O}_X^q)\longrightarrow T_x(\mathcal{V})\longrightarrow T_x(\mathcal{V}/\mathcal{U})\longrightarrow0.$$
由于由构造第一个箭头是满射，得出 $T_x(\mathcal{V}/\mathcal{U})=0$，由 Nakayama 引理（I.2.1）得 $(\mathcal{V}/\mathcal{U})_x=0$。因此 $x$ 不在凝聚 $\mathcal{O}_X$-模 $\mathcal{V}/\mathcal{U}$ 的支集中。缩小 $X$ 可设 $\mathcal{V}/\mathcal{U}=0$，即 $\mathcal{V}$ 由 $s_1,s_2,\dots,s_q$ 生成。此时对 $X$ 上任意整体向量场 $T$ 有
$$Ts_i=\sum_{j=1}^q a^i_js_j$$
其中 $a^i_j\in R(X)$。欲证 $s_1,s_2,\dots,s_q$ 是自由 $\mathcal{O}_{X,x}$-模 $\mathcal{V}_x$ 的一组基。设 $\sum_i g_is_i=0$，其中 $g_1,\dots,g_q\in\mathcal{O}_{X,x}$。断言这意味着对任意 $p\in\mathbb{Z}_+$ 有 $g_1,\dots,g_q\in\mathfrak{m}_x^p$。

由于 $s_1(x),\dots,s_q(x)$ 线性无关，$\sum_i g_i(x)s_i(x)=0$ 蕴含
$g_1(x)=\dots=g_q(x)=0$，从而 $g_1,\dots,g_q\in\mathfrak{m}_x$。因此该陈述对 $p=1$ 成立。设它对 $p-1$ 成立。对 $X$ 上任意整体向量场 $T$ 有
$$0=T\sum_{i=1}^qg_is_i=\sum_{i=1}^q(T(g_i)s_i+g_iTs_i)$$
$$=\sum_{i=1}^qT(g_i)s_i+\sum_{i,j=1}^q a_j^ig_is_j
=\sum_{i=1}^q\Bigl(T(g_i)+\sum_{j=1}^qa_i^jg_j\Bigr)s_i.$$
故由归纳假设，$T(g_i)+\sum_{j=1}^qa_i^jg_j\in\mathfrak{m}_x^{p-1}$ 且 $g_i\in\mathfrak{m}_x^{p-1}$（$1\le i\le q$）。这意味着 $T(g_i)\in\mathfrak{m}_x^{p-1}$（$1\le i\le q$）。特别地，$D_j(g_i)\in\mathfrak{m}_x^{p-1}$（$1\le i\le q$，$1\le j\le n$）。这导致对所有 $|I|<p$ 的 $I\in\mathbb{Z}_+^n$ 有 $D^Ig_i\in\mathfrak{m}_x$，由 II.2.15 得出 $g_i\in\mathfrak{m}_x^p$（$1\le i\le q$）。

因此由对 $p$ 的归纳，对所有 $p\in\mathbb{Z}_+$ 有 $g_i\in\mathfrak{m}_x^p$（$1\le i\le q$）。故 $g_1=\dots=g_q=0$。于是 $s_1,\dots,s_q$ 是自由 $\mathcal{O}_{X,x}$-模 $\mathcal{V}_x$ 的一组基。

因此可考虑自然短正合列
$$0\longrightarrow\mathcal{K}\longrightarrow\mathcal{O}_X^q\longrightarrow\mathcal{V}\longrightarrow0,$$
其中 $\mathcal{K}$ 是自然态射 $\mathcal{O}_X^q\to\mathcal{V}$ 的核。显然 $\mathcal{K}$ 是凝聚 $\mathcal{O}_X$-模，由上述结果 $\mathcal{K}_x=0$。于是 $x\notin\operatorname{supp}(\mathcal{K})$，必要时缩小 $X$ 可设 $\mathcal{K}=0$。这意味着 $s_1,\dots,s_q$ 是自由 $\mathcal{O}_X$-模 $\mathcal{V}$ 的一组基。
\end{proof}

这个结果有以下明显推论。

\begin{corollary}\label{cor:5.2.2}
联络是完整 $\mathcal{D}_X$-模。
\end{corollary}

上述证明表明，若 $X$ 足够“小”且 $\mathcal{V}$ 是 $X$ 上联络，可找到自由 $\mathcal{O}_X$-模 $\mathcal{V}$ 的一组基 $s_1,\dots,s_q$，把 $\mathcal{V}$ 与 $\mathcal{O}_X^q$ 等同。设 $\Psi:\mathcal{O}_X^q\to\mathcal{V}$ 是由
$\Psi(f_1,\dots,f_q)=\sum_i f_is_i$ 给出的相应 $\mathcal{O}_X$-模同构。则存在 $q\times q$ 矩阵 $A_i=(A^i_{kj})$（$1\le i\le n$），其元素在 $R(X)$ 中，使得
$$D_is_j=\sum_{k=1}^qA^i_{kj}s_k.$$
于是
$$(D_i\circ\Psi)(f_1,\dots,f_q)=D_i\Bigl(\sum_{j=1}^qf_js_j\Bigr)=\sum_{j=1}^q\bigl(D_i(f_j)s_j+f_jD_is_j\bigr)$$
$$=\sum_{j=1}^qD_i(f_j)s_j+\sum_{j,k=1}^qf_jA^i_{kj}s_k
=\sum_{j=1}^q\Bigl(D_i(f_j)+\sum_{k=1}^qf_kA^i_{jk}\Bigr)s_j.$$
故
$$\Psi^{-1}\circ D_i\circ\Psi=D_i+A_i$$
对任意 $1\le i\le n$ 成立。此外，由于在 $\mathcal{V}$ 上 $[D_i,D_j]=0$，有
$$0=[D_i+A_i,D_j+A_j]=[D_i,A_j]+[A_i,D_j]+[A_i,A_j]=D_i(A_j)-D_j(A_i)+[A_i,A_j],$$
即
$$D_j(A_i)-D_i(A_j)=[A_i,A_j]$$
对所有 $1\le i,j\le n$ 成立。因此，联络局部地对应于可积联络（integrable connections）的经典概念。
\begin{annote}
$\mathcal O_X$-模配平坦联络 $\nabla$ 等价于局部系统（局部常值层）；取 de Rham 复形得到 $\mathrm{dR}:D^b_{rh}(X)\cong D^b_c(X^{an},\mathbb C)$（Mebkhout–Kashiwara–Beilinson–Bernstein）。这把“完整模的分类”翻译成“构造可层/反常层的分类”，是第五章所有结论的收束点。[来源：nLab 条目 flat+connection（Last revised 2023-06-04）、Riemann-Hilbert+correspondence（2024-08-09）；原书 §2, p133–138]
\end{annote}



\section{完整性与直像}

本节证明完整性在直像下保持。

先考虑光滑仿射簇态射的情形。设 $X$、$Y$ 是两个光滑仿射簇，$\varphi:X\to Y$ 是态射。以 $D_X$、$D_Y$ 分别记 $X$、$Y$ 上微分算子环，以 $\mathcal{M}_R(D_X)$、$\mathcal{M}_R(D_Y)$ 记相应的右 D-模范畴。则可考虑直像函子
$\varphi_+:\mathcal{M}_R(D_X)\to\mathcal{M}_R(D_Y)$ 及其左导出函子 $L^p\varphi_+$。

\begin{theorem}\label{thm:5.3.1}
设 $\mathcal{V}$ 是完整右 $D_X$-模。则 $L^p\varphi_+(\mathcal{V})$ 是完整右 $D_Y$-模。
\end{theorem}

\begin{proof}
在仿射空间 $X=k^n$、$Y=k^m$ 的情形下，这个结果已在 I.13.5 中证明。现在把定理的证明归结到该情形。显然可把 $X$、$Y$ 作为闭代数集嵌入仿射空间 $k^n$、$k^m$。设
$i_X:X\to k^n$、$i_Y:Y\to k^m$。由 Kashiwara 定理，直像函子
$i_{X,+}:\mathcal{M}^R_{qc}(D_X)\to\mathcal{M}^R_{qc}(D_{k^n})$ 与
$i_{Y,+}:\mathcal{M}^R_{qc}(D_Y)\to\mathcal{M}^R_{qc}(D_{k^m})$ 是正合的。借由记号上的滥用，用相同字母表示 $\mathcal{M}_R(D_X)$（相应地 $\mathcal{M}_R(D_Y)$）与 $\mathcal{M}_R(D^{(n)})$（相应地 $\mathcal{M}_R(D^{(m)})$）之间的相应直像函子。于是对 $X$ 有正合函子的交换图
$$\begin{tikzcd}
\mathcal{M}_R(D_X) \arrow[r,"i_{X,+}"] \arrow[d,"\operatorname{loc}"' ] &
\mathcal{M}_R(D^{(n)}) \arrow[d,"\operatorname{loc}"]\\
\mathcal{M}^R_{qc}(D_X) \arrow[r,"i_{X,+}"] & \mathcal{M}^R_{qc}(D_{k^n})
\end{tikzcd}$$
其中竖直箭头是局部化函子，$Y$ 有类似的图。由于竖直箭头是范畴的等价，而下箭头由 Kashiwara 定理是完全忠实的，得出上水平箭头也是完全忠实的，并建立 $\mathcal{M}_R(D_X)$ 与 $\mathcal{M}_R(D^{(n)})$ 中由支集在 $X$ 中的模构成的全子范畴的等价。

设 $\Phi:k^n\to k^m$ 是诱导 $\varphi:X\to Y$ 的多项式映射。则有态射的交换图
$$\begin{tikzcd}
X \arrow[r,"i_X"] \arrow[d,"\varphi"'] & k^n \arrow[d,"\Phi"]\\
Y \arrow[r,"i_Y"'] & k^m
\end{tikzcd}$$
由 IV.2.7 知，对任意 $p\in-\mathbb{Z}_+$ 有
$$L^p\Phi_+\circ i_{X,+}=L^p(\Phi_+\circ i_{X,+})=L^p(i_{Y,+}\circ\varphi_+)=i_{Y,+}\circ L^p\varphi_+.$$
设 $\mathcal{V}$ 是 $X$ 上完整右 $D_X$-模。由于由引理 \ref{lem:5.1.2}，函子 $i_{X,+}$ 把完整模映为完整模，模 $i_{X,+}(\mathcal{V})$ 是完整右 $D^{(n)}$-模。故由 I.13.5，$L^p\Phi_+(i_{X,+}(\mathcal{V}))$ 对所有 $p\in-\mathbb{Z}_+$ 是完整 $D^{(n)}$-模。这意味着 $i_{Y,+}(L^p\varphi_+(\mathcal{V}))$ 是完整 $D^{(n)}$-模。再次使用引理 \ref{lem:5.1.2}，得出对所有 $p\in-\mathbb{Z}_+$，$L^p\varphi_+(\mathcal{V})$ 是完整 $D_Y$-模。
\end{proof}

对定理 \ref{thm:5.3.1} 局部化可见：对任意完整右 $D_X$-模 $\mathcal{V}$，对任意 $p\in\mathbb{Z}$，右 $D_Y$-模 $H^p(\varphi_+(\mathcal{D}(\mathcal{V})))$ 是完整的。

现考虑一般情形。若 $\mathcal{D}_X$-模的复形 $\mathcal{V}^\cdot$ 满足所有 $p\in\mathbb{Z}$ 时 $H^p(\mathcal{V}^\cdot)$ 是完整 $\mathcal{D}_X$-模，则称它是\emph{完整复形}。以 $D^b(\mathcal{M}(\mathcal{D}_X))$ 记 $\mathcal{D}_X$-模的有界导出范畴。由于范畴 $\operatorname{Hol}(\mathcal{D}_X)$ 是 $\mathcal{D}_X$-模范畴的厚 Abel 子范畴，$D^b(\mathcal{M}(\mathcal{D}_X))$ 中由所有完整复形构成的全子范畴
$D^b_{hol}(\mathcal{M}(\mathcal{D}_X))$ 是三角子范畴。下一个结果是 [2] 的一个特例。

\begin{proposition}\label{prop:5.3.2}
完整 $\mathcal{D}_X$-模构成三角范畴 $D^b_{hol}(\mathcal{M}(\mathcal{D}_X))$ 的一个生成类。
\end{proposition}

事实上，可在 $D^b_{hol}(\mathcal{M}(\mathcal{D}_X))$ 中找到一个更小的生成类。设 $U\subset X$ 是仿射开集，$i:U\to X$ 是自然包含。设 $\mathcal{V}$ 是 $U$ 上完整模。断言 $i_+(\mathcal{V})$ 是 $X$ 上完整模。事实上，由于完整性是局部性质，只需证明对任意仿射开集 $V\subset X$，限制 $i_\bullet(\mathcal{V})|_V$ 是完整的。设 $j:U\cap V\to V$ 是自然包含。则
$i_\bullet(\mathcal{V})|_V=j_\bullet(\mathcal{V}|_{U\cap V})$，且由于 $U\cap V$ 仿射，由讨论的第一部分，$j_\bullet(\mathcal{V}|_{U\cap V})$ 是完整模。这证明了 $i_\bullet(\mathcal{V})$ 是完整模。

\begin{proposition}\label{prop:5.3.3}
设 $\mathcal{U}=(U_1,U_2,\dots,U_n)$ 是 $X$ 的仿射开覆盖，$i_j:U_j\to X$ 是自然包含。则模族 $\mathcal{G}$（由 $i_{j,\bullet}(\mathcal{V})$ 组成，其中 $\mathcal{V}$ 是 $U_j$ 上任意完整模且 $1\le j\le n$）构成
$D^b_{hol}(\mathcal{M}(\mathcal{D}_X))$ 的一个生成族。
\end{proposition}

\begin{proof}
设 $\mathcal{V}$ 是 $X$ 上完整模，$C^\cdot(\mathcal{U},\mathcal{V})$ 是其 Čech 分解。则 $\mathcal{V}$ 拟同构于
$C^\cdot(\mathcal{U},\mathcal{V})$。由相关结果，$\mathcal{G}$ 是 $D^b_{hol}(\mathcal{M}(\mathcal{D}_X))$ 的一个生成类。
\end{proof}

这最终使我们能够证明 I.13.5 的如下推广。

\begin{theorem}\label{thm:5.3.4}
设 $X$、$Y$ 是光滑代数簇，$\varphi:X\to Y$ 是簇的态射。若 $\mathcal{V}^\cdot$ 是有界完整 $\mathcal{D}_X$-模复形，则
$\varphi_+(\mathcal{V}^\cdot)$ 是有界完整 $\mathcal{D}_Y$-模复形。
\end{theorem}

\begin{proof}
由于完整性是局部性质，可设 $Y$ 是光滑仿射簇。

设 $U$ 是 $X$ 中仿射开集，$i:U\to X$ 是自然浸入，$\mathcal{V}$ 是 $U$ 上完整模。则 $i_+(\mathcal{V})$ 是完整模。由于由 IV.5.7，$i_+(\mathcal{D}(\mathcal{V}))=\mathcal{D}(i_\bullet(\mathcal{V}))$，可见
$$\varphi_+(\mathcal{D}(i_\bullet(\mathcal{V})))=\varphi_+(i_+(\mathcal{D}(\mathcal{V})))
=(\varphi\circ i)_+(\mathcal{D}(\mathcal{V})).$$
由于 $\varphi\circ i:U\to Y$ 是仿射簇之间的态射，由定理 \ref{thm:5.3.1} 的局部化，对 $p\in\mathbb{Z}$，
$H^p(\varphi_+(\mathcal{D}(i_\bullet(\mathcal{V}))))$ 是 $Y$ 上完整模。因此，由命题 \ref{prop:5.3.3}，结论成立。
\end{proof}

\section{不可约完整模的分类}

现给出不可约完整 $\mathcal{D}_X$-模的分类。它基于以下结果。

\begin{lemma}\label{lem:5.4.1}
设 $U$ 是 $X$ 的开子集，$i:U\to X$ 是自然浸入，$\mathcal{V}$ 是不可约完整 $\mathcal{D}_U$-模。则
\begin{enumerate}[label=(\roman*)]
  \item $i_\bullet(\mathcal{V})$ 包含唯一的不可约 $\mathcal{D}_X$-子模 $\mathcal{W}$；
  \item $\mathcal{W}|_U=\mathcal{V}$。
\end{enumerate}
\end{lemma}

\begin{proof}
由定理 \ref{thm:5.3.4}，$i_\bullet(\mathcal{V})$ 是完整 $\mathcal{D}_X$-模。因此由定理 \ref{thm:5.1.3}，它是有限长度的。设 $\mathcal{W}$ 是 $i_\bullet(\mathcal{V})$ 的不可约 $\mathcal{D}_X$-子模。由于 $i_\bullet$ 是限制到 $U$ 的函子的右伴随，$\mathcal{W}$ 的限制 $\mathcal{W}|_U$ 非零，因此等于 $\mathcal{V}$。这意味着 $i_\bullet(\mathcal{V})$ 的任意两个不可约 $\mathcal{D}_X$-子模的交非零，即 $\mathcal{W}$ 是 $i_\bullet(\mathcal{V})$ 的唯一不可约 $\mathcal{D}_X$-子模。
\end{proof}

这个结果蕴含不可约完整模的如下延拓结果。

\begin{corollary}\label{cor:5.4.2}
设 $U$ 是 $X$ 的开子集，$\mathcal{V}$ 是不可约完整 $\mathcal{D}_U$-模。则存在不可约完整 $\mathcal{D}_X$-模 $\mathcal{W}$，使得 $\mathcal{W}|_U$ 同构于 $\mathcal{V}$。此外，$\mathcal{W}$ 在同构意义下唯一，且 $\operatorname{supp}(\mathcal{W})$ 是 $\operatorname{supp}(\mathcal{V})$ 在 $X$ 中的闭包。
\end{corollary}

\begin{proof}
存在性部分立即由引理 \ref{lem:5.4.1} 得出。设 $\mathcal{W}'$ 是另一个不可约完整 $\mathcal{D}_X$-模，使得 $\mathcal{W}'|_U$ 同构于 $\mathcal{V}$。由于 $i_\bullet$ 是限制到 $U$ 的右伴随，存在自然态射
$\alpha:\mathcal{W}'\to i_\bullet(\mathcal{V})$，使得 $\alpha|_U$ 是 $\mathcal{W}'|_U$ 到 $\mathcal{V}$ 的同构。由于 $\mathcal{W}'$ 不可约，$\alpha$ 的核为零，其像是 $i_\bullet(\mathcal{V})$ 的不可约 $\mathcal{D}_X$-子模。由引理 \ref{lem:5.4.1}，它必等于 $\mathcal{W}$。

显然 $\operatorname{supp}(\mathcal{W})\cap U=\operatorname{supp}(\mathcal{V})$。因此 $\mathcal{V}$ 的支集在 $X$ 中的闭包包含在 $\mathcal{W}$ 的支集中。另一方面，$\mathcal{W}$ 的支集包含在 $\operatorname{supp}(i_\bullet(\mathcal{V}))$ 中，后者等于 $\operatorname{supp}(\mathcal{V})$ 在 $X$ 中的闭包。
\end{proof}

设 $V$ 是 $X$ 中连通光滑子簇，$j:V\to X$ 是自然浸入。设 $\tau$ 是 $V$ 上不可约联络。则直像
$$I(V,\tau)=j_+(\tau)$$
是完整 $\mathcal{D}_X$-模，称为附属于 $(V,\tau)$ 的\emph{标准 D-模}。由推论 \ref{cor:5.4.2}，其支集等于 $\bar{V}$。由于 $V$ 在 $X$ 中局部闭，可找到开集 $U\subset X$ 使得 $V$ 是 $U$ 的闭子簇。以 $j_U$ 记 $V$ 到 $U$ 的浸入，$i$ 记 $U$ 到 $X$ 的开浸入。则由前述有
$$I(V,\tau)=j_+(\tau)=i_+((j_U)_+(\tau)).$$
由 Kashiwara 定理，$(j_U)_+(\tau)$ 的直像是支集等于 $V$ 的不可约完整 $\mathcal{D}_U$-模。由引理 \ref{lem:5.4.1}，它延拓为不可约完整 $\mathcal{D}_X$-模，即 $I(V,\tau)$ 的唯一不可约子模。记为 $L(V,\tau)$，称为附属于数据 $(V,\tau)$ 的\emph{不可约模}。由推论 \ref{cor:5.4.2}，$L(V,\tau)$ 的支集等于 $\bar{V}$。因此我们证明了以下结果的第一部分。

\begin{proposition}\label{prop:5.4.3}
设 $I(V,\tau)$ 是附属于 $(V,\tau)$ 的标准 $\mathcal{D}_X$-模。则它包含唯一不可约子模 $L(V,\tau)$。$L(V,\tau)$ 的支集等于 $\bar{V}$，且 $I(V,\tau)/L(V,\tau)$ 的支集包含在 $\bar{V}-V$ 中。
\end{proposition}

\begin{proof}
只需证明 $\mathcal{Q}=I(V,\tau)/L(V,\tau)$ 的支集在 $\bar{V}-V$ 中。考虑短正合列
$$0\longrightarrow L(V,\tau)\longrightarrow I(V,\tau)\longrightarrow\mathcal{Q}\longrightarrow0.$$
令 $U=X-(\bar{V}-V)$。则 $U$ 是 $X$ 的开子集且 $V$ 在 $U$ 中闭。因此，由上述讨论与引理 \ref{lem:5.4.1}，从这个短正合列限制到 $U$ 得出 $\mathcal{Q}|_U=0$。故 $\operatorname{supp}(\mathcal{Q})\subset\bar{V}-V$。
\end{proof}

现在可以分类 $X$ 上不可约完整模。

\begin{theorem}\label{thm:5.4.4}
\begin{enumerate}[label=(\roman*)]
  \item 设 $\mathcal{V}$ 是不可约完整 $\mathcal{D}_X$-模。则存在 $\mathcal{V}$ 的支集中的不可约开光滑仿射子簇 $V$ 及 $V$ 上不可约联络 $\tau$，使得 $\mathcal{V}$ 同构于 $L(V,\tau)$。
  \item 设 $V$、$V'$ 是 $X$ 的两个不可约光滑仿射子簇，$\tau$、$\tau'$ 分别是 $V$、$V'$ 上不可约联络。则 $L(V,\tau)$ 同构于 $L(V',\tau')$ 当且仅当
    \begin{enumerate}[label=(\alph*)]
      \item $\bar{V}=\bar{V}'$；
      \item 存在 $V\cap V'$ 的非空开仿射子簇 $V''$ 使得 $\tau|_{V''}\cong\tau'|_{V''}$。
    \end{enumerate}
\end{enumerate}
\end{theorem}

\begin{proof}
\begin{enumerate}[label=(\roman*)]
  \item 由前述，$\mathcal{V}$ 的支集是 $X$ 的不可约闭子集。故存在 $X$ 的开仿射子集 $U$，使得 $V=\operatorname{supp}(\mathcal{V})\cap U$ 是 $U$ 的闭光滑子簇。显然 $\mathcal{V}|_U$ 是不可约完整 $\mathcal{D}_U$-模。若以 $j_V$ 记 $V$ 到 $U$ 的自然浸入，由 Kashiwara 定理，存在不可约完整 $\mathcal{D}_V$-模 $\mathcal{W}$，使得 $(j_V)_+(\mathcal{W})\cong\mathcal{V}|_U$。此外，$\operatorname{supp}(\mathcal{W})=V$。因此，由前述，存在 $V$ 的开稠密仿射子簇 $V'$ 使得 $\mathcal{W}|_{V'}$ 是联络。故必要时缩小 $U$，可进一步设 $\mathcal{W}$ 是联络。令 $\tau=\mathcal{W}$，则 $\mathcal{V}|_U$ 与 $L(V,\tau)|_U$ 同构。由推论 \ref{cor:5.4.2}，这意味着 $\mathcal{V}$ 与 $L(V,\tau)$ 同构。
  \item 若 $L(V,\tau)$ 同构于 $L(V',\tau')$，则它们的支集相等，(a) 成立。因此可设 $L(V,\tau)$ 与 $L(V',\tau')$ 有公共支集 $S$。于是 $V$ 与 $V'$ 在 $S$ 中开稠密，故 $V\cap V'$ 是 $V$ 与 $V'$ 的非空开仿射子簇。设 $V''$ 是 $V\cap V'$ 的非空开仿射子簇，则 $V''$ 不可约。设 $U$ 是 $X$ 的开子集使得 $V''=S\cap U$。则由推论 \ref{cor:5.4.2}，$L(V,\tau)|_U$ 与 $L(V',\tau')|_U$ 是支集等于 $V''$ 的同构的不可约完整 $\mathcal{D}_U$-模，当且仅当 $L(V,\tau)$ 与 $L(V',\tau')$ 同构。此外，若以 $j$ 记 $V''$ 到 $U$ 的浸入，则 $L(V,\tau)|_U=j_+(\tau|_{V''})$，$L(V',\tau')|_U=j_+(\tau'|_{V''})$。由于 $V''$ 是 $U$ 的光滑闭子簇，由 Kashiwara 定理，$j_+(\tau|_{V''})$ 同构于 $j_+(\tau'|_{V''})$ 当且仅当 $\tau|_{V''}$ 同构于 $\tau'|_{V''}$。
\end{enumerate}
\end{proof}
\begin{annote}
$\dim X=1$ 时完整模就是 $X$ 上有限秩平坦联络的推广；正则奇点情形下局部单值化给出 Deligne 的“正则奇点方程必有解”，即 $D_{rh}(X)$ 的每个对象都来自局部系统。这是全书中唯一能完全算出分类的情形，也是理解“完整亏格”的几何意义的入口。[来源：nLab 条目 Riemann-Hilbert+correspondence（2024-08-09）；原书 §2、§4, p133–142]
\end{annote}



\section{完整模的局部上同调}

设 $Y$ 是 $X$ 的闭光滑子簇，$U=X-Y$。以 $i:Y\to X$、$j:U\to X$ 记典范包含映射。

设 $\mathcal{V}^\cdot$ 是有界完整 $\mathcal{D}_X$-模复形。则 $\mathcal{V}^\cdot$ 到开集 $U$ 的限制 $\mathcal{V}^\cdot|_U$ 显然是 $\mathcal{D}_U$-模的有界完整复形。因此，由前述，$j_+(\mathcal{V}^\cdot|_U)=Rj_\bullet(\mathcal{V}^\cdot|_U)$ 是 $\mathcal{D}_X$-模的有界完整复形。由可区别三角
$$\begin{tikzcd}
& Rj_\bullet(\mathcal{V}^\cdot|_U) \arrow[dl,"+1"'] & \\
R\Gamma_Y(\mathcal{V}^\cdot) \arrow[rr] && \mathcal{V}^\cdot \arrow[lu]
\end{tikzcd}$$
可见 $R\Gamma_Y(\mathcal{V}^\cdot)$ 也是 $\mathcal{D}_X$-模的有界完整复形。

因此我们证明了以下结果。

\begin{lemma}\label{lem:5.5.1}
设 $\mathcal{V}^\cdot$ 是有界完整 $\mathcal{D}_X$-模复形。则 $R\Gamma_Y(\mathcal{V}^\cdot)$ 是有界完整 $\mathcal{D}_X$-模复形。
\end{lemma}

\section{完整性与逆像}

由前述结果已知 $R\Gamma_Y(\mathcal{V}^\cdot)=i_+(Ri^!(\mathcal{V}^\cdot))$。

应用 $Ri^!$ 并使用前述，可见 $Ri^!(\mathcal{V}^\cdot)$ 也是 $\mathcal{D}_Y$-模的有界完整复形。这进而蕴含 $Li_+(\mathcal{V}^\cdot)$ 是 $\mathcal{D}_Y$-模的有界完整复形。

因此我们证明了以下结果。

\begin{lemma}\label{lem:5.6.1}
设 $\mathcal{V}^\cdot$ 是有界完整 $\mathcal{D}_X$-模复形。则 $Li_+(\mathcal{V}^\cdot)$ 是有界完整 $\mathcal{D}_Y$-模复形。
\end{lemma}

另一方面，若 $p:X\times Y\to Y$ 是投影映射，则 $p_+(\mathcal{V}^\cdot)=\mathcal{O}_X\boxtimes\mathcal{V}^\cdot$。因此，$p_+:\mathcal{M}_{qc}(\mathcal{D}_Y)\to\mathcal{M}_{qc}(\mathcal{D}_{X\times Y})$ 是正合函子。此外有如下结果。

\begin{lemma}\label{lem:5.6.2}
设 $\mathcal{V}^\cdot$ 是有界完整 $\mathcal{D}_Y$-模复形。则 $p_+(\mathcal{V}^\cdot)$ 是有界完整 $\mathcal{D}_{X\times Y}$-模复形。
\end{lemma}

把这些结果与图构造相结合并使用前述，得到以下推广前述结论的结果。

\begin{theorem}\label{thm:5.6.3}
设 $f:X\to Y$ 是光滑代数簇的态射。则逆像函子 $Lf^+$ 把 $\mathcal{D}_Y$-模的有界完整复形映为 $\mathcal{D}_X$-模的有界完整复形。
\end{theorem}


\backmatter
\chapter*{参考文献}
\addcontentsline{toc}{chapter}{参考文献}

\begin{thebibliography}{9}

\bibitem{har} Robin Hartshorne, \emph{Algebraic Geometry}, Graduate Texts in Mathematics 52, Springer--Verlag, New York, 1977.

\bibitem{mil1} Dragan Miličić, \emph{Lectures on Derived Categories}, 未出版手稿.

\bibitem{mil2} Dragan Miličić, \emph{Lectures on Sheaf Theory}, 未出版手稿.

\bibitem{mil0} Dragan Miličić, \emph{Lectures on Algebraic Theory of $D$-Modules}, University of Utah, 1999 年修订版（本译本底本）.

\end{thebibliography}

\section*{校注所用的外部来源}
\addcontentsline{toc}{section}{校注所用的外部来源}

书中小字段校注的依据分三类，逐条在注文末尾以“来源”标出。以下列出所使用的在线资料及其核对日期（2026 年 11 月 19 日）。

\textbf{Stacks Project}（\url{https://stacks.math.columbia.edu}）：
\begin{itemize}[leftmargin=2.2em,itemsep=0.15em]
  \item Tag 00IN：Lemma 10.51.2（Artin--Rees 引理）；
  \item Tag 00IP：Lemma 10.51.4（Krull 交定理）；
  \item Tag 00KE：Definition 10.60.2（环的 Krull 维数）；Tag 0055：Definition 5.10.1（拓扑空间的 Krull 维数）；
  \item Tag 00LA：Definition 10.63.1（结合素）；
  \item Tag 07BK：Definition 10.131.2（Kähler 微分模）；
  \item Tag 01BD（拟凝聚模）、Tag 01BU（凝聚模）；
  \item Tag 0A6Y：Lemma 47.12.3（局部化中 torsion 与 complete 的对偶，仅作形式对照）。
\end{itemize}

\textbf{nLab}（\url{https://ncatlab.org/nlab/show/}），各条目的最后修订日期一并列出：
\begin{itemize}[leftmargin=2.2em,itemsep=0.15em]
  \item \emph{D-module}（2024-07-27）；\emph{holonomic D-module}（2014-02-16）；
  \item \emph{six operations}（2026-09-30）；\emph{flat connection}（2023-06-04）；
  \item \emph{Riemann--Hilbert correspondence}（2024-08-09）；\emph{local system}（2025-04-10）；
  \item \emph{de Rham complex}（2026-09-26）；\emph{flag variety}（2026-09-28）。
\end{itemize}
其中 \emph{characteristic variety}（2017-08-29）与 \emph{coherent D-module}（2021-10-01）两个条目无实质陈述，故未用于任何断言，相应的校注只依据原书页码。


\end{document}
